% !TeX program = xelatex
% ================================================================
%  《收获与播种》(Récoltes et Semailles) 中文译本
%  卷一 · 前奏 —— 主题的呈现:或四个乐章的前奏
%  原书 p.23-122, 第 1-4 章:
%    1. 权作序言   2. 巡游一部作品(或孩子与母亲)
%    3. 一封信     4. 引言
%  --------------------------------------------------------------
%  五卷本之一。共用导言区 common.tex, 封面 cover-art.tex。
%  构建: xelatex ×4 (目录页码需 4 遍才收敛)
% ================================================================
\PassOptionsToPackage{AutoFallBack}{xeCJK}
\documentclass[
    PaperSize=a4paper,
    HeadingFont=hei,
    Biblatex=false,
    BookTitleCN={收获与播种},
    Translator={数译},
    LogoFile=logo.pdf
]{mathtranslation}

% ================================================================
%  《收获与播种》中文译本 —— 各卷共用导言区
%  由 vol-prelude / vol-part1 / vol-part2 / vol-part3 / vol-part4
%  以及全书合订 main-all 共同 \input。基类 mathtranslation.cls v3.2 不改。
%  --------------------------------------------------------------
%  样式均在此改写, 以贴近原书:
%    · 章标题仿原书 "5. Travail et découverte"(块状左对齐黑体)
%    · 节/小节: 黑体左对齐, 编号后带句点
%    · 脚注: 人工编号照录原书——前奏与第一部分页内重新编号,
%      第二至四部各自连续编号并带(*)角标; 第一部分另有 \fnoteP
%      处理 "(n)" 式连续注号
%    · 边码: 原打字稿页码 \origpage{N}(前奏作 A1、P1、L1、罗马数字等)
%    · ⋄ 记号: 段首补写标记 \pmark 与上标形式 \dmark
% ================================================================

% 题名含空格, 不能走类选项(会被 kvoptions 吃掉), 在此赋值
\renewcommand{\BookTitleEN}{Récoltes et Semailles}
\renewcommand{\OriginalAuthor}{Alexandre Grothendieck}
\renewcommand{\TranslationDate}{\today}

% --------------------------------------------------------------
%  罕用字后备(如"祐"): Fandol 不载时退到 SimSun
% --------------------------------------------------------------
\setCJKfallbackfamilyfont{\CJKrmdefault}{{SimSun}}

% --------------------------------------------------------------
%  hyperref 元信息(各卷在 % ================================================================
%  《收获与播种》中文译本 —— 各卷共用导言区
%  由 vol-prelude / vol-part1 / vol-part2 / vol-part3 / vol-part4
%  以及全书合订 main-all 共同 \input。基类 mathtranslation.cls v3.2 不改。
%  --------------------------------------------------------------
%  样式均在此改写, 以贴近原书:
%    · 章标题仿原书 "5. Travail et découverte"(块状左对齐黑体)
%    · 节/小节: 黑体左对齐, 编号后带句点
%    · 脚注: 人工编号照录原书——前奏与第一部分页内重新编号,
%      第二至四部各自连续编号并带(*)角标; 第一部分另有 \fnoteP
%      处理 "(n)" 式连续注号
%    · 边码: 原打字稿页码 \origpage{N}(前奏作 A1、P1、L1、罗马数字等)
%    · ⋄ 记号: 段首补写标记 \pmark 与上标形式 \dmark
% ================================================================

% 题名含空格, 不能走类选项(会被 kvoptions 吃掉), 在此赋值
\renewcommand{\BookTitleEN}{Récoltes et Semailles}
\renewcommand{\OriginalAuthor}{Alexandre Grothendieck}
\renewcommand{\TranslationDate}{\today}

% --------------------------------------------------------------
%  罕用字后备(如"祐"): Fandol 不载时退到 SimSun
% --------------------------------------------------------------
\setCJKfallbackfamilyfont{\CJKrmdefault}{{SimSun}}

% --------------------------------------------------------------
%  hyperref 元信息(各卷在 % ================================================================
%  《收获与播种》中文译本 —— 各卷共用导言区
%  由 vol-prelude / vol-part1 / vol-part2 / vol-part3 / vol-part4
%  以及全书合订 main-all 共同 \input。基类 mathtranslation.cls v3.2 不改。
%  --------------------------------------------------------------
%  样式均在此改写, 以贴近原书:
%    · 章标题仿原书 "5. Travail et découverte"(块状左对齐黑体)
%    · 节/小节: 黑体左对齐, 编号后带句点
%    · 脚注: 人工编号照录原书——前奏与第一部分页内重新编号,
%      第二至四部各自连续编号并带(*)角标; 第一部分另有 \fnoteP
%      处理 "(n)" 式连续注号
%    · 边码: 原打字稿页码 \origpage{N}(前奏作 A1、P1、L1、罗马数字等)
%    · ⋄ 记号: 段首补写标记 \pmark 与上标形式 \dmark
% ================================================================

% 题名含空格, 不能走类选项(会被 kvoptions 吃掉), 在此赋值
\renewcommand{\BookTitleEN}{Récoltes et Semailles}
\renewcommand{\OriginalAuthor}{Alexandre Grothendieck}
\renewcommand{\TranslationDate}{\today}

% --------------------------------------------------------------
%  罕用字后备(如"祐"): Fandol 不载时退到 SimSun
% --------------------------------------------------------------
\setCJKfallbackfamilyfont{\CJKrmdefault}{{SimSun}}

% --------------------------------------------------------------
%  hyperref 元信息(各卷在 \input{common} 之后可用 \hypersetup 覆写)
% --------------------------------------------------------------
\hypersetup{
  pdftitle={收获与播种},
  pdfauthor={Alexandre Grothendieck / 数译 译},
  pdfsubject={Récoltes et Semailles — Réflexions et témoignage sur un passé de mathématicien}
}

% --------------------------------------------------------------
%  标题样式: 仿原书 —— 黑体左对齐, 编号后带句点
% --------------------------------------------------------------
\titleformat{\section}[block]
  {\normalfont\Large\headfont}{\thesection.}{0.5em}{}
\titlespacing*{\section}{0pt}{3.2ex plus 1.2ex minus .2ex}{1.6ex plus .2ex}
\titleformat{\subsection}[block]
  {\normalfont\large\headfont}{\thesubsection.}{0.5em}{}
\titlespacing*{\subsection}{0pt}{2.6ex plus 1ex minus .2ex}{1.2ex plus .2ex}
\titleformat{\subsubsection}[block]
  {\normalfont\normalsize\headfont}{\thesubsubsection.}{0.5em}{}
\titlespacing*{\subsubsection}{0pt}{2.2ex plus .8ex minus .2ex}{1ex plus .2ex}
\titleformat{name=\subsubsection,numberless}[block]
  {\normalfont\normalsize\headfont}{}{0em}{}
\setcounter{secnumdepth}{3}
\setcounter{tocdepth}{3}
% 为含长外文书名与行内公式的段落提供必要的断行余量
\setlength{\emergencystretch}{3em}

% 章标题样式: 原书章标题与节同级排印("13. A) HERITAGE..."),
% 故章标题取块状左对齐黑体, 字号介于节与部之间
\titleformat{\chapter}[block]
  {\normalfont\LARGE\headfont}{\thechapter.}{0.5em}{}
\titlespacing*{\chapter}{0pt}{3.6ex plus 1.4ex minus .2ex}{2ex plus .2ex}
% 目录中章号照原书用阿拉伯数字("13. A) ..." 而非 "第十三章")
\renewcommand{\CTEXthechapter}{\thechapter\hspace{.3em}}

% --------------------------------------------------------------
%  目录观感: 条目与引导点之间、引导点与页码之间各留 0.5em
%  内核 \@dottedtocline 让 \leaders 从标题末字起排, 第一个点几乎
%  贴着末字, 读起来像多标了一个句点; 页码前也没有余量。两处各补
%  0.5em。
% --------------------------------------------------------------
\makeatletter
\patchcmd{\@dottedtocline}{\nobreak\leaders}{\nobreak\hspace{0.5em}\leaders}%
  {\typeout{MT: TOC 引导点前已留 0.5em}}%
  {\typeout{MT-WARN: TOC 引导点前间距补丁未生效}}
\patchcmd{\@dottedtocline}{\hfill\nobreak\hb@xt@\@pnumwidth}%
  {\hfill\nobreak\hspace{0.5em}\hb@xt@\@pnumwidth}%
  {\typeout{MT: TOC 页码前已留 0.5em}}%
  {\typeout{MT-WARN: TOC 页码前间距补丁未生效}}
\makeatother

% --------------------------------------------------------------
%  正文半角标点后的间距
%  本书体例用半角 , ; : ? ! 与 ASCII 括号, 由 xeCJK 按 HalfRight
%  (右半留白) 处理。放宽到 0.45em, 与全角标点右半留白大致齐平。
% --------------------------------------------------------------
\xeCJKsetup{CJKecglue={\hskip 0.45em plus 0.06em minus 0.09em}}

% --------------------------------------------------------------
%  页眉页脚: 偶数页眉右对齐斜体节题; 页码页脚外侧
% --------------------------------------------------------------
\fancyhf{}
\fancyhead[RE]{\small\itshape\rightmark}
\fancyfoot[LE,RO]{\small\thepage}
\renewcommand{\headrulewidth}{0pt}
\renewcommand{\sectionmark}[1]{\markright{\thesection.\ #1}}
\fancypagestyle{plain}{%
  \fancyhf{}%
  \fancyfoot[LE,RO]{\small \thepage}%
  \renewcommand{\headrulewidth}{0pt}%
}

% --------------------------------------------------------------
%  原书版式记号宏
% --------------------------------------------------------------
% 边码: 原打字稿页码(G. 手稿系统里的自引页码), 排在页边
\usepackage{marginnote}
\newcommand{\origpage}[1]{\marginnote{\normalfont\scriptsize p.~#1}}

% ⋄ 记号: 段首/行内正常大小; 上标形式为作者后补文字的位置标记
\DeclareRobustCommand{\pmark}{\ensuremath{\diamond}\hspace{0.15em}}
\DeclareRobustCommand{\dmark}{\textsuperscript{\ensuremath{\diamond}}}

% 脚注人工编号: \fnote{795}{(*)}{文本} —— 编号与(*)角标均照录原书
\newcommand{\fnote}[3]{%
  \begingroup
  \renewcommand{\thefootnote}{#1#2}%
  \footnote{#3}%
  \endgroup
}

% 第一部分的脚注: 页内序号作上标(#1), 注文行首印 "序号(连续注号)",
% 正文在序号之后另行括注同一连续注号 —— \fnoteP{1}{(3)}{注文}
% 原书有三条注无连续注号而改标日期(ch07 注3 = 30 Septembre,
% ch09 注11/12 = 8 août), 第二参数留空, 日期写进注文之首。
\newcommand{\fnoteP}[3]{%
  \begingroup
  \renewcommand{\thefootnote}{#1}\footnotemark
  \renewcommand{\thefootnote}{#1#2}\footnotetext{#3}%
  \endgroup
}

% 正文呼号与注文行首标号不同(原书照录)
\newcommand{\fnotevariant}[4]{%
  \begingroup
  \renewcommand{\thefootnote}{#1#2}\footnotemark
  \renewcommand{\thefootnote}{#1#3}\footnotetext{#4}%
  \endgroup
}

% 大注块标题: \notehead{171³}{斜体评论} —— 正文随后另起
\newcommand{\notehead}[2]{%
  \par\bigskip
  \noindent{\headfont 注~#1}\ifblank{#2}{}{\hspace{0.5em}{\itshape #2}}%
  \par\nobreak\smallskip\nobreak
}

% ∗ ∗ ∗ 分隔
\newcommand{\asterisms}{%
  \par\vspace{1.2ex}%
  \noindent\hspace*{2em}$\ast$\hspace{1.2em}$\ast$\hspace{1.2em}$\ast$%
  \par\vspace{1.2ex}%
}

% 无编号子目的目录条目
\newcommand{\tocsubsub}[1]{%
  \phantomsection
  \addcontentsline{toc}{subsubsection}{#1}%
}

% 部扉: 仿原书部题页
\newcommand{\partdivider}[3]{%
  \clearpage
  \thispagestyle{empty}%
  \phantomsection
  \addcontentsline{toc}{part}{#1\quad #3}%
  \vspace*{6.5cm}
  \begin{center}
    {\fontsize{16pt}{26pt}\selectfont\headfont #1\par}
    \vspace{0.9cm}
    {\fontsize{22pt}{34pt}\selectfont\headfont #2\par}
  \end{center}
  \clearpage
}

% 献辞页: 仿原书部题后的题献, 独占一页且不入目录
\newcommand{\dedication}[1]{%
  \thispagestyle{empty}%
  \vspace*{4cm}%
  \begin{center}#1\end{center}%
  \clearpage
}

% --------------------------------------------------------------
%  各卷"译者说明"的共用体例段(各卷的范围段自写, 后接本段)
% --------------------------------------------------------------
\newcommand{\CommonNoteBody}{%
为便于与原书及学界通行的打字稿页码对勘, 本册保留:

\begin{itemize}
  \item 原书全部章节编号(章~1--18, 节如~5.2、18.2, 子目如~18.1.1)与原书小节标题;
  \item 原书脚注的编号与角标——
        \begin{itemize}
          \item 前奏(第~1--4 章)与第一部的注在各页内重新编号, 正文中的上标即该
                页序号; 第一部紧随其后的括注(如“(3)”)是该部分的连续注号, 与第~12
                章“第一部分的注”里的“注~3”一一对应, 注文行首照原书印作“1(3)”式。
                其中三条注原书未标连续注号而改标日期(ch7 注~3 = 30 Septembre、
                ch9 注~11/12 = 8 août), 本册照录其日期;
          \item 第二部第~13--16~章各自重新编号, 分别为~1--41、1--86、1--132、1--29;
                第三部第~17~章自~1~起至~12, 第~18~章自~1~重新起算并延续至~794;
                末段自~\textbf{795(*)}~续至~1053(*)。各段与原书一一对应;
        \end{itemize}
  \item 页边码“p.~1”“p.~A3”“p.~1066”式小注——这是作者引用自己打字稿原稿时
        所用的页码, 与本排印本页码不同。前奏四章的边码各有前缀: 第~1 章作
        \texttt{A1…A6}, 第~2 章作 \texttt{P1…P65}, 第~3 章作 \texttt{L1…L56},
        第~4 章作罗马数字 \texttt{I…XXII}, 自第一部起作纯数字 \texttt{1…171};
  \item 作者用以标记段落的~$\diamond$~记号, 以及标示后来补写文字的上标
        ~\textsuperscript{$\diamond$}~记号;
  \item 原书大小注(“注~44'”“注~46$^1$”“注~171³”等)的编号与层级。
\end{itemize}

第~12 章“第一部分的注”把第~5 至第~11 章各页脚注汇为一编, 并给每条注配上节题
(如“12.19. 耶稣与十二使徒”), 因此部分注文与正文脚注重复——这是原书的体例, 非
重出之误。这一部分属随笔式的自省文字, 同一句法语在正文脚注与章末汇编注中
各有措辞, 本册有意保留各自的译笔, 不强求统一。

除上述标记外, 译文力求贴近作者驳杂、跳跃而辛辣的文风; 数学名词的译法见各卷
术语表(如~\emph{topos}~不译, \emph{faisceaux pervers}~作“反常层”)。文中的日期
括注(如“(5~月~30~日)”“(1984年3月补写)”)为作者历次补写时的纪日, 一律照录。文中
另有一些对注号、页码的引用(如“n°~1695”“打字稿第~227~页”)沿袭作者打字稿的编号
系统, 与本排印本不尽一致, 亦无法在本书中落实, 悉照原文录出。

全套五卷(另有全书合订一册)的分工是:

\begin{itemize}
  \item \textbf{卷一·前奏}——主题的呈现: 或四个乐章的前奏(原书 p.23--122);
  \item \textbf{卷二·第一部}——愚妄与更新(原书 p.123--254);
  \item \textbf{卷三·第二部}——葬礼(I): 或中国皇帝的新衣(原书 p.255--412);
  \item \textbf{卷四·第三部}——葬礼(II): 或阴阳之钥(原书 p.413--632);
  \item \textbf{卷五·第四部}——葬礼(III): 或四种运算(原书 p.633--929)。
\end{itemize}
}
 之后可用 \hypersetup 覆写)
% --------------------------------------------------------------
\hypersetup{
  pdftitle={收获与播种},
  pdfauthor={Alexandre Grothendieck / 数译 译},
  pdfsubject={Récoltes et Semailles — Réflexions et témoignage sur un passé de mathématicien}
}

% --------------------------------------------------------------
%  标题样式: 仿原书 —— 黑体左对齐, 编号后带句点
% --------------------------------------------------------------
\titleformat{\section}[block]
  {\normalfont\Large\headfont}{\thesection.}{0.5em}{}
\titlespacing*{\section}{0pt}{3.2ex plus 1.2ex minus .2ex}{1.6ex plus .2ex}
\titleformat{\subsection}[block]
  {\normalfont\large\headfont}{\thesubsection.}{0.5em}{}
\titlespacing*{\subsection}{0pt}{2.6ex plus 1ex minus .2ex}{1.2ex plus .2ex}
\titleformat{\subsubsection}[block]
  {\normalfont\normalsize\headfont}{\thesubsubsection.}{0.5em}{}
\titlespacing*{\subsubsection}{0pt}{2.2ex plus .8ex minus .2ex}{1ex plus .2ex}
\titleformat{name=\subsubsection,numberless}[block]
  {\normalfont\normalsize\headfont}{}{0em}{}
\setcounter{secnumdepth}{3}
\setcounter{tocdepth}{3}
% 为含长外文书名与行内公式的段落提供必要的断行余量
\setlength{\emergencystretch}{3em}

% 章标题样式: 原书章标题与节同级排印("13. A) HERITAGE..."),
% 故章标题取块状左对齐黑体, 字号介于节与部之间
\titleformat{\chapter}[block]
  {\normalfont\LARGE\headfont}{\thechapter.}{0.5em}{}
\titlespacing*{\chapter}{0pt}{3.6ex plus 1.4ex minus .2ex}{2ex plus .2ex}
% 目录中章号照原书用阿拉伯数字("13. A) ..." 而非 "第十三章")
\renewcommand{\CTEXthechapter}{\thechapter\hspace{.3em}}

% --------------------------------------------------------------
%  目录观感: 条目与引导点之间、引导点与页码之间各留 0.5em
%  内核 \@dottedtocline 让 \leaders 从标题末字起排, 第一个点几乎
%  贴着末字, 读起来像多标了一个句点; 页码前也没有余量。两处各补
%  0.5em。
% --------------------------------------------------------------
\makeatletter
\patchcmd{\@dottedtocline}{\nobreak\leaders}{\nobreak\hspace{0.5em}\leaders}%
  {\typeout{MT: TOC 引导点前已留 0.5em}}%
  {\typeout{MT-WARN: TOC 引导点前间距补丁未生效}}
\patchcmd{\@dottedtocline}{\hfill\nobreak\hb@xt@\@pnumwidth}%
  {\hfill\nobreak\hspace{0.5em}\hb@xt@\@pnumwidth}%
  {\typeout{MT: TOC 页码前已留 0.5em}}%
  {\typeout{MT-WARN: TOC 页码前间距补丁未生效}}
\makeatother

% --------------------------------------------------------------
%  正文半角标点后的间距
%  本书体例用半角 , ; : ? ! 与 ASCII 括号, 由 xeCJK 按 HalfRight
%  (右半留白) 处理。放宽到 0.45em, 与全角标点右半留白大致齐平。
% --------------------------------------------------------------
\xeCJKsetup{CJKecglue={\hskip 0.45em plus 0.06em minus 0.09em}}

% --------------------------------------------------------------
%  页眉页脚: 偶数页眉右对齐斜体节题; 页码页脚外侧
% --------------------------------------------------------------
\fancyhf{}
\fancyhead[RE]{\small\itshape\rightmark}
\fancyfoot[LE,RO]{\small\thepage}
\renewcommand{\headrulewidth}{0pt}
\renewcommand{\sectionmark}[1]{\markright{\thesection.\ #1}}
\fancypagestyle{plain}{%
  \fancyhf{}%
  \fancyfoot[LE,RO]{\small \thepage}%
  \renewcommand{\headrulewidth}{0pt}%
}

% --------------------------------------------------------------
%  原书版式记号宏
% --------------------------------------------------------------
% 边码: 原打字稿页码(G. 手稿系统里的自引页码), 排在页边
\usepackage{marginnote}
\newcommand{\origpage}[1]{\marginnote{\normalfont\scriptsize p.~#1}}

% ⋄ 记号: 段首/行内正常大小; 上标形式为作者后补文字的位置标记
\DeclareRobustCommand{\pmark}{\ensuremath{\diamond}\hspace{0.15em}}
\DeclareRobustCommand{\dmark}{\textsuperscript{\ensuremath{\diamond}}}

% 脚注人工编号: \fnote{795}{(*)}{文本} —— 编号与(*)角标均照录原书
\newcommand{\fnote}[3]{%
  \begingroup
  \renewcommand{\thefootnote}{#1#2}%
  \footnote{#3}%
  \endgroup
}

% 第一部分的脚注: 页内序号作上标(#1), 注文行首印 "序号(连续注号)",
% 正文在序号之后另行括注同一连续注号 —— \fnoteP{1}{(3)}{注文}
% 原书有三条注无连续注号而改标日期(ch07 注3 = 30 Septembre,
% ch09 注11/12 = 8 août), 第二参数留空, 日期写进注文之首。
\newcommand{\fnoteP}[3]{%
  \begingroup
  \renewcommand{\thefootnote}{#1}\footnotemark
  \renewcommand{\thefootnote}{#1#2}\footnotetext{#3}%
  \endgroup
}

% 正文呼号与注文行首标号不同(原书照录)
\newcommand{\fnotevariant}[4]{%
  \begingroup
  \renewcommand{\thefootnote}{#1#2}\footnotemark
  \renewcommand{\thefootnote}{#1#3}\footnotetext{#4}%
  \endgroup
}

% 大注块标题: \notehead{171³}{斜体评论} —— 正文随后另起
\newcommand{\notehead}[2]{%
  \par\bigskip
  \noindent{\headfont 注~#1}\ifblank{#2}{}{\hspace{0.5em}{\itshape #2}}%
  \par\nobreak\smallskip\nobreak
}

% ∗ ∗ ∗ 分隔
\newcommand{\asterisms}{%
  \par\vspace{1.2ex}%
  \noindent\hspace*{2em}$\ast$\hspace{1.2em}$\ast$\hspace{1.2em}$\ast$%
  \par\vspace{1.2ex}%
}

% 无编号子目的目录条目
\newcommand{\tocsubsub}[1]{%
  \phantomsection
  \addcontentsline{toc}{subsubsection}{#1}%
}

% 部扉: 仿原书部题页
\newcommand{\partdivider}[3]{%
  \clearpage
  \thispagestyle{empty}%
  \phantomsection
  \addcontentsline{toc}{part}{#1\quad #3}%
  \vspace*{6.5cm}
  \begin{center}
    {\fontsize{16pt}{26pt}\selectfont\headfont #1\par}
    \vspace{0.9cm}
    {\fontsize{22pt}{34pt}\selectfont\headfont #2\par}
  \end{center}
  \clearpage
}

% 献辞页: 仿原书部题后的题献, 独占一页且不入目录
\newcommand{\dedication}[1]{%
  \thispagestyle{empty}%
  \vspace*{4cm}%
  \begin{center}#1\end{center}%
  \clearpage
}

% --------------------------------------------------------------
%  各卷"译者说明"的共用体例段(各卷的范围段自写, 后接本段)
% --------------------------------------------------------------
\newcommand{\CommonNoteBody}{%
为便于与原书及学界通行的打字稿页码对勘, 本册保留:

\begin{itemize}
  \item 原书全部章节编号(章~1--18, 节如~5.2、18.2, 子目如~18.1.1)与原书小节标题;
  \item 原书脚注的编号与角标——
        \begin{itemize}
          \item 前奏(第~1--4 章)与第一部的注在各页内重新编号, 正文中的上标即该
                页序号; 第一部紧随其后的括注(如“(3)”)是该部分的连续注号, 与第~12
                章“第一部分的注”里的“注~3”一一对应, 注文行首照原书印作“1(3)”式。
                其中三条注原书未标连续注号而改标日期(ch7 注~3 = 30 Septembre、
                ch9 注~11/12 = 8 août), 本册照录其日期;
          \item 第二部第~13--16~章各自重新编号, 分别为~1--41、1--86、1--132、1--29;
                第三部第~17~章自~1~起至~12, 第~18~章自~1~重新起算并延续至~794;
                末段自~\textbf{795(*)}~续至~1053(*)。各段与原书一一对应;
        \end{itemize}
  \item 页边码“p.~1”“p.~A3”“p.~1066”式小注——这是作者引用自己打字稿原稿时
        所用的页码, 与本排印本页码不同。前奏四章的边码各有前缀: 第~1 章作
        \texttt{A1…A6}, 第~2 章作 \texttt{P1…P65}, 第~3 章作 \texttt{L1…L56},
        第~4 章作罗马数字 \texttt{I…XXII}, 自第一部起作纯数字 \texttt{1…171};
  \item 作者用以标记段落的~$\diamond$~记号, 以及标示后来补写文字的上标
        ~\textsuperscript{$\diamond$}~记号;
  \item 原书大小注(“注~44'”“注~46$^1$”“注~171³”等)的编号与层级。
\end{itemize}

第~12 章“第一部分的注”把第~5 至第~11 章各页脚注汇为一编, 并给每条注配上节题
(如“12.19. 耶稣与十二使徒”), 因此部分注文与正文脚注重复——这是原书的体例, 非
重出之误。这一部分属随笔式的自省文字, 同一句法语在正文脚注与章末汇编注中
各有措辞, 本册有意保留各自的译笔, 不强求统一。

除上述标记外, 译文力求贴近作者驳杂、跳跃而辛辣的文风; 数学名词的译法见各卷
术语表(如~\emph{topos}~不译, \emph{faisceaux pervers}~作“反常层”)。文中的日期
括注(如“(5~月~30~日)”“(1984年3月补写)”)为作者历次补写时的纪日, 一律照录。文中
另有一些对注号、页码的引用(如“n°~1695”“打字稿第~227~页”)沿袭作者打字稿的编号
系统, 与本排印本不尽一致, 亦无法在本书中落实, 悉照原文录出。

全套五卷(另有全书合订一册)的分工是:

\begin{itemize}
  \item \textbf{卷一·前奏}——主题的呈现: 或四个乐章的前奏(原书 p.23--122);
  \item \textbf{卷二·第一部}——愚妄与更新(原书 p.123--254);
  \item \textbf{卷三·第二部}——葬礼(I): 或中国皇帝的新衣(原书 p.255--412);
  \item \textbf{卷四·第三部}——葬礼(II): 或阴阳之钥(原书 p.413--632);
  \item \textbf{卷五·第四部}——葬礼(III): 或四种运算(原书 p.633--929)。
\end{itemize}
}
 之后可用 \hypersetup 覆写)
% --------------------------------------------------------------
\hypersetup{
  pdftitle={收获与播种},
  pdfauthor={Alexandre Grothendieck / 数译 译},
  pdfsubject={Récoltes et Semailles — Réflexions et témoignage sur un passé de mathématicien}
}

% --------------------------------------------------------------
%  标题样式: 仿原书 —— 黑体左对齐, 编号后带句点
% --------------------------------------------------------------
\titleformat{\section}[block]
  {\normalfont\Large\headfont}{\thesection.}{0.5em}{}
\titlespacing*{\section}{0pt}{3.2ex plus 1.2ex minus .2ex}{1.6ex plus .2ex}
\titleformat{\subsection}[block]
  {\normalfont\large\headfont}{\thesubsection.}{0.5em}{}
\titlespacing*{\subsection}{0pt}{2.6ex plus 1ex minus .2ex}{1.2ex plus .2ex}
\titleformat{\subsubsection}[block]
  {\normalfont\normalsize\headfont}{\thesubsubsection.}{0.5em}{}
\titlespacing*{\subsubsection}{0pt}{2.2ex plus .8ex minus .2ex}{1ex plus .2ex}
\titleformat{name=\subsubsection,numberless}[block]
  {\normalfont\normalsize\headfont}{}{0em}{}
\setcounter{secnumdepth}{3}
\setcounter{tocdepth}{3}
% 为含长外文书名与行内公式的段落提供必要的断行余量
\setlength{\emergencystretch}{3em}

% 章标题样式: 原书章标题与节同级排印("13. A) HERITAGE..."),
% 故章标题取块状左对齐黑体, 字号介于节与部之间
\titleformat{\chapter}[block]
  {\normalfont\LARGE\headfont}{\thechapter.}{0.5em}{}
\titlespacing*{\chapter}{0pt}{3.6ex plus 1.4ex minus .2ex}{2ex plus .2ex}
% 目录中章号照原书用阿拉伯数字("13. A) ..." 而非 "第十三章")
\renewcommand{\CTEXthechapter}{\thechapter\hspace{.3em}}

% --------------------------------------------------------------
%  目录观感: 条目与引导点之间、引导点与页码之间各留 0.5em
%  内核 \@dottedtocline 让 \leaders 从标题末字起排, 第一个点几乎
%  贴着末字, 读起来像多标了一个句点; 页码前也没有余量。两处各补
%  0.5em。
% --------------------------------------------------------------
\makeatletter
\patchcmd{\@dottedtocline}{\nobreak\leaders}{\nobreak\hspace{0.5em}\leaders}%
  {\typeout{MT: TOC 引导点前已留 0.5em}}%
  {\typeout{MT-WARN: TOC 引导点前间距补丁未生效}}
\patchcmd{\@dottedtocline}{\hfill\nobreak\hb@xt@\@pnumwidth}%
  {\hfill\nobreak\hspace{0.5em}\hb@xt@\@pnumwidth}%
  {\typeout{MT: TOC 页码前已留 0.5em}}%
  {\typeout{MT-WARN: TOC 页码前间距补丁未生效}}
\makeatother

% --------------------------------------------------------------
%  正文半角标点后的间距
%  本书体例用半角 , ; : ? ! 与 ASCII 括号, 由 xeCJK 按 HalfRight
%  (右半留白) 处理。放宽到 0.45em, 与全角标点右半留白大致齐平。
% --------------------------------------------------------------
\xeCJKsetup{CJKecglue={\hskip 0.45em plus 0.06em minus 0.09em}}

% --------------------------------------------------------------
%  页眉页脚: 偶数页眉右对齐斜体节题; 页码页脚外侧
% --------------------------------------------------------------
\fancyhf{}
\fancyhead[RE]{\small\itshape\rightmark}
\fancyfoot[LE,RO]{\small\thepage}
\renewcommand{\headrulewidth}{0pt}
\renewcommand{\sectionmark}[1]{\markright{\thesection.\ #1}}
\fancypagestyle{plain}{%
  \fancyhf{}%
  \fancyfoot[LE,RO]{\small \thepage}%
  \renewcommand{\headrulewidth}{0pt}%
}

% --------------------------------------------------------------
%  原书版式记号宏
% --------------------------------------------------------------
% 边码: 原打字稿页码(G. 手稿系统里的自引页码), 排在页边
\usepackage{marginnote}
\newcommand{\origpage}[1]{\marginnote{\normalfont\scriptsize p.~#1}}

% ⋄ 记号: 段首/行内正常大小; 上标形式为作者后补文字的位置标记
\DeclareRobustCommand{\pmark}{\ensuremath{\diamond}\hspace{0.15em}}
\DeclareRobustCommand{\dmark}{\textsuperscript{\ensuremath{\diamond}}}

% 脚注人工编号: \fnote{795}{(*)}{文本} —— 编号与(*)角标均照录原书
\newcommand{\fnote}[3]{%
  \begingroup
  \renewcommand{\thefootnote}{#1#2}%
  \footnote{#3}%
  \endgroup
}

% 第一部分的脚注: 页内序号作上标(#1), 注文行首印 "序号(连续注号)",
% 正文在序号之后另行括注同一连续注号 —— \fnoteP{1}{(3)}{注文}
% 原书有三条注无连续注号而改标日期(ch07 注3 = 30 Septembre,
% ch09 注11/12 = 8 août), 第二参数留空, 日期写进注文之首。
\newcommand{\fnoteP}[3]{%
  \begingroup
  \renewcommand{\thefootnote}{#1}\footnotemark
  \renewcommand{\thefootnote}{#1#2}\footnotetext{#3}%
  \endgroup
}

% 正文呼号与注文行首标号不同(原书照录)
\newcommand{\fnotevariant}[4]{%
  \begingroup
  \renewcommand{\thefootnote}{#1#2}\footnotemark
  \renewcommand{\thefootnote}{#1#3}\footnotetext{#4}%
  \endgroup
}

% 大注块标题: \notehead{171³}{斜体评论} —— 正文随后另起
\newcommand{\notehead}[2]{%
  \par\bigskip
  \noindent{\headfont 注~#1}\ifblank{#2}{}{\hspace{0.5em}{\itshape #2}}%
  \par\nobreak\smallskip\nobreak
}

% ∗ ∗ ∗ 分隔
\newcommand{\asterisms}{%
  \par\vspace{1.2ex}%
  \noindent\hspace*{2em}$\ast$\hspace{1.2em}$\ast$\hspace{1.2em}$\ast$%
  \par\vspace{1.2ex}%
}

% 无编号子目的目录条目
\newcommand{\tocsubsub}[1]{%
  \phantomsection
  \addcontentsline{toc}{subsubsection}{#1}%
}

% 部扉: 仿原书部题页
\newcommand{\partdivider}[3]{%
  \clearpage
  \thispagestyle{empty}%
  \phantomsection
  \addcontentsline{toc}{part}{#1\quad #3}%
  \vspace*{6.5cm}
  \begin{center}
    {\fontsize{16pt}{26pt}\selectfont\headfont #1\par}
    \vspace{0.9cm}
    {\fontsize{22pt}{34pt}\selectfont\headfont #2\par}
  \end{center}
  \clearpage
}

% 献辞页: 仿原书部题后的题献, 独占一页且不入目录
\newcommand{\dedication}[1]{%
  \thispagestyle{empty}%
  \vspace*{4cm}%
  \begin{center}#1\end{center}%
  \clearpage
}

% --------------------------------------------------------------
%  各卷"译者说明"的共用体例段(各卷的范围段自写, 后接本段)
% --------------------------------------------------------------
\newcommand{\CommonNoteBody}{%
为便于与原书及学界通行的打字稿页码对勘, 本册保留:

\begin{itemize}
  \item 原书全部章节编号(章~1--18, 节如~5.2、18.2, 子目如~18.1.1)与原书小节标题;
  \item 原书脚注的编号与角标——
        \begin{itemize}
          \item 前奏(第~1--4 章)与第一部的注在各页内重新编号, 正文中的上标即该
                页序号; 第一部紧随其后的括注(如“(3)”)是该部分的连续注号, 与第~12
                章“第一部分的注”里的“注~3”一一对应, 注文行首照原书印作“1(3)”式。
                其中三条注原书未标连续注号而改标日期(ch7 注~3 = 30 Septembre、
                ch9 注~11/12 = 8 août), 本册照录其日期;
          \item 第二部第~13--16~章各自重新编号, 分别为~1--41、1--86、1--132、1--29;
                第三部第~17~章自~1~起至~12, 第~18~章自~1~重新起算并延续至~794;
                末段自~\textbf{795(*)}~续至~1053(*)。各段与原书一一对应;
        \end{itemize}
  \item 页边码“p.~1”“p.~A3”“p.~1066”式小注——这是作者引用自己打字稿原稿时
        所用的页码, 与本排印本页码不同。前奏四章的边码各有前缀: 第~1 章作
        \texttt{A1…A6}, 第~2 章作 \texttt{P1…P65}, 第~3 章作 \texttt{L1…L56},
        第~4 章作罗马数字 \texttt{I…XXII}, 自第一部起作纯数字 \texttt{1…171};
  \item 作者用以标记段落的~$\diamond$~记号, 以及标示后来补写文字的上标
        ~\textsuperscript{$\diamond$}~记号;
  \item 原书大小注(“注~44'”“注~46$^1$”“注~171³”等)的编号与层级。
\end{itemize}

第~12 章“第一部分的注”把第~5 至第~11 章各页脚注汇为一编, 并给每条注配上节题
(如“12.19. 耶稣与十二使徒”), 因此部分注文与正文脚注重复——这是原书的体例, 非
重出之误。这一部分属随笔式的自省文字, 同一句法语在正文脚注与章末汇编注中
各有措辞, 本册有意保留各自的译笔, 不强求统一。

除上述标记外, 译文力求贴近作者驳杂、跳跃而辛辣的文风; 数学名词的译法见各卷
术语表(如~\emph{topos}~不译, \emph{faisceaux pervers}~作“反常层”)。文中的日期
括注(如“(5~月~30~日)”“(1984年3月补写)”)为作者历次补写时的纪日, 一律照录。文中
另有一些对注号、页码的引用(如“n°~1695”“打字稿第~227~页”)沿袭作者打字稿的编号
系统, 与本排印本不尽一致, 亦无法在本书中落实, 悉照原文录出。

全套五卷(另有全书合订一册)的分工是:

\begin{itemize}
  \item \textbf{卷一·前奏}——主题的呈现: 或四个乐章的前奏(原书 p.23--122);
  \item \textbf{卷二·第一部}——愚妄与更新(原书 p.123--254);
  \item \textbf{卷三·第二部}——葬礼(I): 或中国皇帝的新衣(原书 p.255--412);
  \item \textbf{卷四·第三部}——葬礼(II): 或阴阳之钥(原书 p.413--632);
  \item \textbf{卷五·第四部}——葬礼(III): 或四种运算(原书 p.633--929)。
\end{itemize}
}


% --------------------------------------------------------------
%  封面参数
% --------------------------------------------------------------
% ================================================================
%  《收获与播种》统一封面 —— TikZ 原生矢量绘制(不依赖任何外部图片)
%  母题: 暮色中的麦田 —— 犁沟向地平线的灭点收束, 前景麦穗剪影,
%        一位播种者把一把种子撒向天空, 种子沿弧线落入田里。
%  用法(在 \begin{document} 之前, 依次):
%      % ================================================================
%  《收获与播种》中文译本 —— 各卷共用导言区
%  由 vol-prelude / vol-part1 / vol-part2 / vol-part3 / vol-part4
%  以及全书合订 main-all 共同 \input。基类 mathtranslation.cls v3.2 不改。
%  --------------------------------------------------------------
%  样式均在此改写, 以贴近原书:
%    · 章标题仿原书 "5. Travail et découverte"(块状左对齐黑体)
%    · 节/小节: 黑体左对齐, 编号后带句点
%    · 脚注: 人工编号照录原书——前奏与第一部分页内重新编号,
%      第二至四部各自连续编号并带(*)角标; 第一部分另有 \fnoteP
%      处理 "(n)" 式连续注号
%    · 边码: 原打字稿页码 \origpage{N}(前奏作 A1、P1、L1、罗马数字等)
%    · ⋄ 记号: 段首补写标记 \pmark 与上标形式 \dmark
% ================================================================

% 题名含空格, 不能走类选项(会被 kvoptions 吃掉), 在此赋值
\renewcommand{\BookTitleEN}{Récoltes et Semailles}
\renewcommand{\OriginalAuthor}{Alexandre Grothendieck}
\renewcommand{\TranslationDate}{\today}

% --------------------------------------------------------------
%  罕用字后备(如"祐"): Fandol 不载时退到 SimSun
% --------------------------------------------------------------
\setCJKfallbackfamilyfont{\CJKrmdefault}{{SimSun}}

% --------------------------------------------------------------
%  hyperref 元信息(各卷在 % ================================================================
%  《收获与播种》中文译本 —— 各卷共用导言区
%  由 vol-prelude / vol-part1 / vol-part2 / vol-part3 / vol-part4
%  以及全书合订 main-all 共同 \input。基类 mathtranslation.cls v3.2 不改。
%  --------------------------------------------------------------
%  样式均在此改写, 以贴近原书:
%    · 章标题仿原书 "5. Travail et découverte"(块状左对齐黑体)
%    · 节/小节: 黑体左对齐, 编号后带句点
%    · 脚注: 人工编号照录原书——前奏与第一部分页内重新编号,
%      第二至四部各自连续编号并带(*)角标; 第一部分另有 \fnoteP
%      处理 "(n)" 式连续注号
%    · 边码: 原打字稿页码 \origpage{N}(前奏作 A1、P1、L1、罗马数字等)
%    · ⋄ 记号: 段首补写标记 \pmark 与上标形式 \dmark
% ================================================================

% 题名含空格, 不能走类选项(会被 kvoptions 吃掉), 在此赋值
\renewcommand{\BookTitleEN}{Récoltes et Semailles}
\renewcommand{\OriginalAuthor}{Alexandre Grothendieck}
\renewcommand{\TranslationDate}{\today}

% --------------------------------------------------------------
%  罕用字后备(如"祐"): Fandol 不载时退到 SimSun
% --------------------------------------------------------------
\setCJKfallbackfamilyfont{\CJKrmdefault}{{SimSun}}

% --------------------------------------------------------------
%  hyperref 元信息(各卷在 \input{common} 之后可用 \hypersetup 覆写)
% --------------------------------------------------------------
\hypersetup{
  pdftitle={收获与播种},
  pdfauthor={Alexandre Grothendieck / 数译 译},
  pdfsubject={Récoltes et Semailles — Réflexions et témoignage sur un passé de mathématicien}
}

% --------------------------------------------------------------
%  标题样式: 仿原书 —— 黑体左对齐, 编号后带句点
% --------------------------------------------------------------
\titleformat{\section}[block]
  {\normalfont\Large\headfont}{\thesection.}{0.5em}{}
\titlespacing*{\section}{0pt}{3.2ex plus 1.2ex minus .2ex}{1.6ex plus .2ex}
\titleformat{\subsection}[block]
  {\normalfont\large\headfont}{\thesubsection.}{0.5em}{}
\titlespacing*{\subsection}{0pt}{2.6ex plus 1ex minus .2ex}{1.2ex plus .2ex}
\titleformat{\subsubsection}[block]
  {\normalfont\normalsize\headfont}{\thesubsubsection.}{0.5em}{}
\titlespacing*{\subsubsection}{0pt}{2.2ex plus .8ex minus .2ex}{1ex plus .2ex}
\titleformat{name=\subsubsection,numberless}[block]
  {\normalfont\normalsize\headfont}{}{0em}{}
\setcounter{secnumdepth}{3}
\setcounter{tocdepth}{3}
% 为含长外文书名与行内公式的段落提供必要的断行余量
\setlength{\emergencystretch}{3em}

% 章标题样式: 原书章标题与节同级排印("13. A) HERITAGE..."),
% 故章标题取块状左对齐黑体, 字号介于节与部之间
\titleformat{\chapter}[block]
  {\normalfont\LARGE\headfont}{\thechapter.}{0.5em}{}
\titlespacing*{\chapter}{0pt}{3.6ex plus 1.4ex minus .2ex}{2ex plus .2ex}
% 目录中章号照原书用阿拉伯数字("13. A) ..." 而非 "第十三章")
\renewcommand{\CTEXthechapter}{\thechapter\hspace{.3em}}

% --------------------------------------------------------------
%  目录观感: 条目与引导点之间、引导点与页码之间各留 0.5em
%  内核 \@dottedtocline 让 \leaders 从标题末字起排, 第一个点几乎
%  贴着末字, 读起来像多标了一个句点; 页码前也没有余量。两处各补
%  0.5em。
% --------------------------------------------------------------
\makeatletter
\patchcmd{\@dottedtocline}{\nobreak\leaders}{\nobreak\hspace{0.5em}\leaders}%
  {\typeout{MT: TOC 引导点前已留 0.5em}}%
  {\typeout{MT-WARN: TOC 引导点前间距补丁未生效}}
\patchcmd{\@dottedtocline}{\hfill\nobreak\hb@xt@\@pnumwidth}%
  {\hfill\nobreak\hspace{0.5em}\hb@xt@\@pnumwidth}%
  {\typeout{MT: TOC 页码前已留 0.5em}}%
  {\typeout{MT-WARN: TOC 页码前间距补丁未生效}}
\makeatother

% --------------------------------------------------------------
%  正文半角标点后的间距
%  本书体例用半角 , ; : ? ! 与 ASCII 括号, 由 xeCJK 按 HalfRight
%  (右半留白) 处理。放宽到 0.45em, 与全角标点右半留白大致齐平。
% --------------------------------------------------------------
\xeCJKsetup{CJKecglue={\hskip 0.45em plus 0.06em minus 0.09em}}

% --------------------------------------------------------------
%  页眉页脚: 偶数页眉右对齐斜体节题; 页码页脚外侧
% --------------------------------------------------------------
\fancyhf{}
\fancyhead[RE]{\small\itshape\rightmark}
\fancyfoot[LE,RO]{\small\thepage}
\renewcommand{\headrulewidth}{0pt}
\renewcommand{\sectionmark}[1]{\markright{\thesection.\ #1}}
\fancypagestyle{plain}{%
  \fancyhf{}%
  \fancyfoot[LE,RO]{\small \thepage}%
  \renewcommand{\headrulewidth}{0pt}%
}

% --------------------------------------------------------------
%  原书版式记号宏
% --------------------------------------------------------------
% 边码: 原打字稿页码(G. 手稿系统里的自引页码), 排在页边
\usepackage{marginnote}
\newcommand{\origpage}[1]{\marginnote{\normalfont\scriptsize p.~#1}}

% ⋄ 记号: 段首/行内正常大小; 上标形式为作者后补文字的位置标记
\DeclareRobustCommand{\pmark}{\ensuremath{\diamond}\hspace{0.15em}}
\DeclareRobustCommand{\dmark}{\textsuperscript{\ensuremath{\diamond}}}

% 脚注人工编号: \fnote{795}{(*)}{文本} —— 编号与(*)角标均照录原书
\newcommand{\fnote}[3]{%
  \begingroup
  \renewcommand{\thefootnote}{#1#2}%
  \footnote{#3}%
  \endgroup
}

% 第一部分的脚注: 页内序号作上标(#1), 注文行首印 "序号(连续注号)",
% 正文在序号之后另行括注同一连续注号 —— \fnoteP{1}{(3)}{注文}
% 原书有三条注无连续注号而改标日期(ch07 注3 = 30 Septembre,
% ch09 注11/12 = 8 août), 第二参数留空, 日期写进注文之首。
\newcommand{\fnoteP}[3]{%
  \begingroup
  \renewcommand{\thefootnote}{#1}\footnotemark
  \renewcommand{\thefootnote}{#1#2}\footnotetext{#3}%
  \endgroup
}

% 正文呼号与注文行首标号不同(原书照录)
\newcommand{\fnotevariant}[4]{%
  \begingroup
  \renewcommand{\thefootnote}{#1#2}\footnotemark
  \renewcommand{\thefootnote}{#1#3}\footnotetext{#4}%
  \endgroup
}

% 大注块标题: \notehead{171³}{斜体评论} —— 正文随后另起
\newcommand{\notehead}[2]{%
  \par\bigskip
  \noindent{\headfont 注~#1}\ifblank{#2}{}{\hspace{0.5em}{\itshape #2}}%
  \par\nobreak\smallskip\nobreak
}

% ∗ ∗ ∗ 分隔
\newcommand{\asterisms}{%
  \par\vspace{1.2ex}%
  \noindent\hspace*{2em}$\ast$\hspace{1.2em}$\ast$\hspace{1.2em}$\ast$%
  \par\vspace{1.2ex}%
}

% 无编号子目的目录条目
\newcommand{\tocsubsub}[1]{%
  \phantomsection
  \addcontentsline{toc}{subsubsection}{#1}%
}

% 部扉: 仿原书部题页
\newcommand{\partdivider}[3]{%
  \clearpage
  \thispagestyle{empty}%
  \phantomsection
  \addcontentsline{toc}{part}{#1\quad #3}%
  \vspace*{6.5cm}
  \begin{center}
    {\fontsize{16pt}{26pt}\selectfont\headfont #1\par}
    \vspace{0.9cm}
    {\fontsize{22pt}{34pt}\selectfont\headfont #2\par}
  \end{center}
  \clearpage
}

% 献辞页: 仿原书部题后的题献, 独占一页且不入目录
\newcommand{\dedication}[1]{%
  \thispagestyle{empty}%
  \vspace*{4cm}%
  \begin{center}#1\end{center}%
  \clearpage
}

% --------------------------------------------------------------
%  各卷"译者说明"的共用体例段(各卷的范围段自写, 后接本段)
% --------------------------------------------------------------
\newcommand{\CommonNoteBody}{%
为便于与原书及学界通行的打字稿页码对勘, 本册保留:

\begin{itemize}
  \item 原书全部章节编号(章~1--18, 节如~5.2、18.2, 子目如~18.1.1)与原书小节标题;
  \item 原书脚注的编号与角标——
        \begin{itemize}
          \item 前奏(第~1--4 章)与第一部的注在各页内重新编号, 正文中的上标即该
                页序号; 第一部紧随其后的括注(如“(3)”)是该部分的连续注号, 与第~12
                章“第一部分的注”里的“注~3”一一对应, 注文行首照原书印作“1(3)”式。
                其中三条注原书未标连续注号而改标日期(ch7 注~3 = 30 Septembre、
                ch9 注~11/12 = 8 août), 本册照录其日期;
          \item 第二部第~13--16~章各自重新编号, 分别为~1--41、1--86、1--132、1--29;
                第三部第~17~章自~1~起至~12, 第~18~章自~1~重新起算并延续至~794;
                末段自~\textbf{795(*)}~续至~1053(*)。各段与原书一一对应;
        \end{itemize}
  \item 页边码“p.~1”“p.~A3”“p.~1066”式小注——这是作者引用自己打字稿原稿时
        所用的页码, 与本排印本页码不同。前奏四章的边码各有前缀: 第~1 章作
        \texttt{A1…A6}, 第~2 章作 \texttt{P1…P65}, 第~3 章作 \texttt{L1…L56},
        第~4 章作罗马数字 \texttt{I…XXII}, 自第一部起作纯数字 \texttt{1…171};
  \item 作者用以标记段落的~$\diamond$~记号, 以及标示后来补写文字的上标
        ~\textsuperscript{$\diamond$}~记号;
  \item 原书大小注(“注~44'”“注~46$^1$”“注~171³”等)的编号与层级。
\end{itemize}

第~12 章“第一部分的注”把第~5 至第~11 章各页脚注汇为一编, 并给每条注配上节题
(如“12.19. 耶稣与十二使徒”), 因此部分注文与正文脚注重复——这是原书的体例, 非
重出之误。这一部分属随笔式的自省文字, 同一句法语在正文脚注与章末汇编注中
各有措辞, 本册有意保留各自的译笔, 不强求统一。

除上述标记外, 译文力求贴近作者驳杂、跳跃而辛辣的文风; 数学名词的译法见各卷
术语表(如~\emph{topos}~不译, \emph{faisceaux pervers}~作“反常层”)。文中的日期
括注(如“(5~月~30~日)”“(1984年3月补写)”)为作者历次补写时的纪日, 一律照录。文中
另有一些对注号、页码的引用(如“n°~1695”“打字稿第~227~页”)沿袭作者打字稿的编号
系统, 与本排印本不尽一致, 亦无法在本书中落实, 悉照原文录出。

全套五卷(另有全书合订一册)的分工是:

\begin{itemize}
  \item \textbf{卷一·前奏}——主题的呈现: 或四个乐章的前奏(原书 p.23--122);
  \item \textbf{卷二·第一部}——愚妄与更新(原书 p.123--254);
  \item \textbf{卷三·第二部}——葬礼(I): 或中国皇帝的新衣(原书 p.255--412);
  \item \textbf{卷四·第三部}——葬礼(II): 或阴阳之钥(原书 p.413--632);
  \item \textbf{卷五·第四部}——葬礼(III): 或四种运算(原书 p.633--929)。
\end{itemize}
}
 之后可用 \hypersetup 覆写)
% --------------------------------------------------------------
\hypersetup{
  pdftitle={收获与播种},
  pdfauthor={Alexandre Grothendieck / 数译 译},
  pdfsubject={Récoltes et Semailles — Réflexions et témoignage sur un passé de mathématicien}
}

% --------------------------------------------------------------
%  标题样式: 仿原书 —— 黑体左对齐, 编号后带句点
% --------------------------------------------------------------
\titleformat{\section}[block]
  {\normalfont\Large\headfont}{\thesection.}{0.5em}{}
\titlespacing*{\section}{0pt}{3.2ex plus 1.2ex minus .2ex}{1.6ex plus .2ex}
\titleformat{\subsection}[block]
  {\normalfont\large\headfont}{\thesubsection.}{0.5em}{}
\titlespacing*{\subsection}{0pt}{2.6ex plus 1ex minus .2ex}{1.2ex plus .2ex}
\titleformat{\subsubsection}[block]
  {\normalfont\normalsize\headfont}{\thesubsubsection.}{0.5em}{}
\titlespacing*{\subsubsection}{0pt}{2.2ex plus .8ex minus .2ex}{1ex plus .2ex}
\titleformat{name=\subsubsection,numberless}[block]
  {\normalfont\normalsize\headfont}{}{0em}{}
\setcounter{secnumdepth}{3}
\setcounter{tocdepth}{3}
% 为含长外文书名与行内公式的段落提供必要的断行余量
\setlength{\emergencystretch}{3em}

% 章标题样式: 原书章标题与节同级排印("13. A) HERITAGE..."),
% 故章标题取块状左对齐黑体, 字号介于节与部之间
\titleformat{\chapter}[block]
  {\normalfont\LARGE\headfont}{\thechapter.}{0.5em}{}
\titlespacing*{\chapter}{0pt}{3.6ex plus 1.4ex minus .2ex}{2ex plus .2ex}
% 目录中章号照原书用阿拉伯数字("13. A) ..." 而非 "第十三章")
\renewcommand{\CTEXthechapter}{\thechapter\hspace{.3em}}

% --------------------------------------------------------------
%  目录观感: 条目与引导点之间、引导点与页码之间各留 0.5em
%  内核 \@dottedtocline 让 \leaders 从标题末字起排, 第一个点几乎
%  贴着末字, 读起来像多标了一个句点; 页码前也没有余量。两处各补
%  0.5em。
% --------------------------------------------------------------
\makeatletter
\patchcmd{\@dottedtocline}{\nobreak\leaders}{\nobreak\hspace{0.5em}\leaders}%
  {\typeout{MT: TOC 引导点前已留 0.5em}}%
  {\typeout{MT-WARN: TOC 引导点前间距补丁未生效}}
\patchcmd{\@dottedtocline}{\hfill\nobreak\hb@xt@\@pnumwidth}%
  {\hfill\nobreak\hspace{0.5em}\hb@xt@\@pnumwidth}%
  {\typeout{MT: TOC 页码前已留 0.5em}}%
  {\typeout{MT-WARN: TOC 页码前间距补丁未生效}}
\makeatother

% --------------------------------------------------------------
%  正文半角标点后的间距
%  本书体例用半角 , ; : ? ! 与 ASCII 括号, 由 xeCJK 按 HalfRight
%  (右半留白) 处理。放宽到 0.45em, 与全角标点右半留白大致齐平。
% --------------------------------------------------------------
\xeCJKsetup{CJKecglue={\hskip 0.45em plus 0.06em minus 0.09em}}

% --------------------------------------------------------------
%  页眉页脚: 偶数页眉右对齐斜体节题; 页码页脚外侧
% --------------------------------------------------------------
\fancyhf{}
\fancyhead[RE]{\small\itshape\rightmark}
\fancyfoot[LE,RO]{\small\thepage}
\renewcommand{\headrulewidth}{0pt}
\renewcommand{\sectionmark}[1]{\markright{\thesection.\ #1}}
\fancypagestyle{plain}{%
  \fancyhf{}%
  \fancyfoot[LE,RO]{\small \thepage}%
  \renewcommand{\headrulewidth}{0pt}%
}

% --------------------------------------------------------------
%  原书版式记号宏
% --------------------------------------------------------------
% 边码: 原打字稿页码(G. 手稿系统里的自引页码), 排在页边
\usepackage{marginnote}
\newcommand{\origpage}[1]{\marginnote{\normalfont\scriptsize p.~#1}}

% ⋄ 记号: 段首/行内正常大小; 上标形式为作者后补文字的位置标记
\DeclareRobustCommand{\pmark}{\ensuremath{\diamond}\hspace{0.15em}}
\DeclareRobustCommand{\dmark}{\textsuperscript{\ensuremath{\diamond}}}

% 脚注人工编号: \fnote{795}{(*)}{文本} —— 编号与(*)角标均照录原书
\newcommand{\fnote}[3]{%
  \begingroup
  \renewcommand{\thefootnote}{#1#2}%
  \footnote{#3}%
  \endgroup
}

% 第一部分的脚注: 页内序号作上标(#1), 注文行首印 "序号(连续注号)",
% 正文在序号之后另行括注同一连续注号 —— \fnoteP{1}{(3)}{注文}
% 原书有三条注无连续注号而改标日期(ch07 注3 = 30 Septembre,
% ch09 注11/12 = 8 août), 第二参数留空, 日期写进注文之首。
\newcommand{\fnoteP}[3]{%
  \begingroup
  \renewcommand{\thefootnote}{#1}\footnotemark
  \renewcommand{\thefootnote}{#1#2}\footnotetext{#3}%
  \endgroup
}

% 正文呼号与注文行首标号不同(原书照录)
\newcommand{\fnotevariant}[4]{%
  \begingroup
  \renewcommand{\thefootnote}{#1#2}\footnotemark
  \renewcommand{\thefootnote}{#1#3}\footnotetext{#4}%
  \endgroup
}

% 大注块标题: \notehead{171³}{斜体评论} —— 正文随后另起
\newcommand{\notehead}[2]{%
  \par\bigskip
  \noindent{\headfont 注~#1}\ifblank{#2}{}{\hspace{0.5em}{\itshape #2}}%
  \par\nobreak\smallskip\nobreak
}

% ∗ ∗ ∗ 分隔
\newcommand{\asterisms}{%
  \par\vspace{1.2ex}%
  \noindent\hspace*{2em}$\ast$\hspace{1.2em}$\ast$\hspace{1.2em}$\ast$%
  \par\vspace{1.2ex}%
}

% 无编号子目的目录条目
\newcommand{\tocsubsub}[1]{%
  \phantomsection
  \addcontentsline{toc}{subsubsection}{#1}%
}

% 部扉: 仿原书部题页
\newcommand{\partdivider}[3]{%
  \clearpage
  \thispagestyle{empty}%
  \phantomsection
  \addcontentsline{toc}{part}{#1\quad #3}%
  \vspace*{6.5cm}
  \begin{center}
    {\fontsize{16pt}{26pt}\selectfont\headfont #1\par}
    \vspace{0.9cm}
    {\fontsize{22pt}{34pt}\selectfont\headfont #2\par}
  \end{center}
  \clearpage
}

% 献辞页: 仿原书部题后的题献, 独占一页且不入目录
\newcommand{\dedication}[1]{%
  \thispagestyle{empty}%
  \vspace*{4cm}%
  \begin{center}#1\end{center}%
  \clearpage
}

% --------------------------------------------------------------
%  各卷"译者说明"的共用体例段(各卷的范围段自写, 后接本段)
% --------------------------------------------------------------
\newcommand{\CommonNoteBody}{%
为便于与原书及学界通行的打字稿页码对勘, 本册保留:

\begin{itemize}
  \item 原书全部章节编号(章~1--18, 节如~5.2、18.2, 子目如~18.1.1)与原书小节标题;
  \item 原书脚注的编号与角标——
        \begin{itemize}
          \item 前奏(第~1--4 章)与第一部的注在各页内重新编号, 正文中的上标即该
                页序号; 第一部紧随其后的括注(如“(3)”)是该部分的连续注号, 与第~12
                章“第一部分的注”里的“注~3”一一对应, 注文行首照原书印作“1(3)”式。
                其中三条注原书未标连续注号而改标日期(ch7 注~3 = 30 Septembre、
                ch9 注~11/12 = 8 août), 本册照录其日期;
          \item 第二部第~13--16~章各自重新编号, 分别为~1--41、1--86、1--132、1--29;
                第三部第~17~章自~1~起至~12, 第~18~章自~1~重新起算并延续至~794;
                末段自~\textbf{795(*)}~续至~1053(*)。各段与原书一一对应;
        \end{itemize}
  \item 页边码“p.~1”“p.~A3”“p.~1066”式小注——这是作者引用自己打字稿原稿时
        所用的页码, 与本排印本页码不同。前奏四章的边码各有前缀: 第~1 章作
        \texttt{A1…A6}, 第~2 章作 \texttt{P1…P65}, 第~3 章作 \texttt{L1…L56},
        第~4 章作罗马数字 \texttt{I…XXII}, 自第一部起作纯数字 \texttt{1…171};
  \item 作者用以标记段落的~$\diamond$~记号, 以及标示后来补写文字的上标
        ~\textsuperscript{$\diamond$}~记号;
  \item 原书大小注(“注~44'”“注~46$^1$”“注~171³”等)的编号与层级。
\end{itemize}

第~12 章“第一部分的注”把第~5 至第~11 章各页脚注汇为一编, 并给每条注配上节题
(如“12.19. 耶稣与十二使徒”), 因此部分注文与正文脚注重复——这是原书的体例, 非
重出之误。这一部分属随笔式的自省文字, 同一句法语在正文脚注与章末汇编注中
各有措辞, 本册有意保留各自的译笔, 不强求统一。

除上述标记外, 译文力求贴近作者驳杂、跳跃而辛辣的文风; 数学名词的译法见各卷
术语表(如~\emph{topos}~不译, \emph{faisceaux pervers}~作“反常层”)。文中的日期
括注(如“(5~月~30~日)”“(1984年3月补写)”)为作者历次补写时的纪日, 一律照录。文中
另有一些对注号、页码的引用(如“n°~1695”“打字稿第~227~页”)沿袭作者打字稿的编号
系统, 与本排印本不尽一致, 亦无法在本书中落实, 悉照原文录出。

全套五卷(另有全书合订一册)的分工是:

\begin{itemize}
  \item \textbf{卷一·前奏}——主题的呈现: 或四个乐章的前奏(原书 p.23--122);
  \item \textbf{卷二·第一部}——愚妄与更新(原书 p.123--254);
  \item \textbf{卷三·第二部}——葬礼(I): 或中国皇帝的新衣(原书 p.255--412);
  \item \textbf{卷四·第三部}——葬礼(II): 或阴阳之钥(原书 p.413--632);
  \item \textbf{卷五·第四部}——葬礼(III): 或四种运算(原书 p.633--929)。
\end{itemize}
}
                       % 共用导言区
%      % ================================================================
%  《收获与播种》统一封面 —— TikZ 原生矢量绘制(不依赖任何外部图片)
%  母题: 暮色中的麦田 —— 犁沟向地平线的灭点收束, 前景麦穗剪影,
%        一位播种者把一把种子撒向天空, 种子沿弧线落入田里。
%  用法(在 \begin{document} 之前, 依次):
%      % ================================================================
%  《收获与播种》中文译本 —— 各卷共用导言区
%  由 vol-prelude / vol-part1 / vol-part2 / vol-part3 / vol-part4
%  以及全书合订 main-all 共同 \input。基类 mathtranslation.cls v3.2 不改。
%  --------------------------------------------------------------
%  样式均在此改写, 以贴近原书:
%    · 章标题仿原书 "5. Travail et découverte"(块状左对齐黑体)
%    · 节/小节: 黑体左对齐, 编号后带句点
%    · 脚注: 人工编号照录原书——前奏与第一部分页内重新编号,
%      第二至四部各自连续编号并带(*)角标; 第一部分另有 \fnoteP
%      处理 "(n)" 式连续注号
%    · 边码: 原打字稿页码 \origpage{N}(前奏作 A1、P1、L1、罗马数字等)
%    · ⋄ 记号: 段首补写标记 \pmark 与上标形式 \dmark
% ================================================================

% 题名含空格, 不能走类选项(会被 kvoptions 吃掉), 在此赋值
\renewcommand{\BookTitleEN}{Récoltes et Semailles}
\renewcommand{\OriginalAuthor}{Alexandre Grothendieck}
\renewcommand{\TranslationDate}{\today}

% --------------------------------------------------------------
%  罕用字后备(如"祐"): Fandol 不载时退到 SimSun
% --------------------------------------------------------------
\setCJKfallbackfamilyfont{\CJKrmdefault}{{SimSun}}

% --------------------------------------------------------------
%  hyperref 元信息(各卷在 \input{common} 之后可用 \hypersetup 覆写)
% --------------------------------------------------------------
\hypersetup{
  pdftitle={收获与播种},
  pdfauthor={Alexandre Grothendieck / 数译 译},
  pdfsubject={Récoltes et Semailles — Réflexions et témoignage sur un passé de mathématicien}
}

% --------------------------------------------------------------
%  标题样式: 仿原书 —— 黑体左对齐, 编号后带句点
% --------------------------------------------------------------
\titleformat{\section}[block]
  {\normalfont\Large\headfont}{\thesection.}{0.5em}{}
\titlespacing*{\section}{0pt}{3.2ex plus 1.2ex minus .2ex}{1.6ex plus .2ex}
\titleformat{\subsection}[block]
  {\normalfont\large\headfont}{\thesubsection.}{0.5em}{}
\titlespacing*{\subsection}{0pt}{2.6ex plus 1ex minus .2ex}{1.2ex plus .2ex}
\titleformat{\subsubsection}[block]
  {\normalfont\normalsize\headfont}{\thesubsubsection.}{0.5em}{}
\titlespacing*{\subsubsection}{0pt}{2.2ex plus .8ex minus .2ex}{1ex plus .2ex}
\titleformat{name=\subsubsection,numberless}[block]
  {\normalfont\normalsize\headfont}{}{0em}{}
\setcounter{secnumdepth}{3}
\setcounter{tocdepth}{3}
% 为含长外文书名与行内公式的段落提供必要的断行余量
\setlength{\emergencystretch}{3em}

% 章标题样式: 原书章标题与节同级排印("13. A) HERITAGE..."),
% 故章标题取块状左对齐黑体, 字号介于节与部之间
\titleformat{\chapter}[block]
  {\normalfont\LARGE\headfont}{\thechapter.}{0.5em}{}
\titlespacing*{\chapter}{0pt}{3.6ex plus 1.4ex minus .2ex}{2ex plus .2ex}
% 目录中章号照原书用阿拉伯数字("13. A) ..." 而非 "第十三章")
\renewcommand{\CTEXthechapter}{\thechapter\hspace{.3em}}

% --------------------------------------------------------------
%  目录观感: 条目与引导点之间、引导点与页码之间各留 0.5em
%  内核 \@dottedtocline 让 \leaders 从标题末字起排, 第一个点几乎
%  贴着末字, 读起来像多标了一个句点; 页码前也没有余量。两处各补
%  0.5em。
% --------------------------------------------------------------
\makeatletter
\patchcmd{\@dottedtocline}{\nobreak\leaders}{\nobreak\hspace{0.5em}\leaders}%
  {\typeout{MT: TOC 引导点前已留 0.5em}}%
  {\typeout{MT-WARN: TOC 引导点前间距补丁未生效}}
\patchcmd{\@dottedtocline}{\hfill\nobreak\hb@xt@\@pnumwidth}%
  {\hfill\nobreak\hspace{0.5em}\hb@xt@\@pnumwidth}%
  {\typeout{MT: TOC 页码前已留 0.5em}}%
  {\typeout{MT-WARN: TOC 页码前间距补丁未生效}}
\makeatother

% --------------------------------------------------------------
%  正文半角标点后的间距
%  本书体例用半角 , ; : ? ! 与 ASCII 括号, 由 xeCJK 按 HalfRight
%  (右半留白) 处理。放宽到 0.45em, 与全角标点右半留白大致齐平。
% --------------------------------------------------------------
\xeCJKsetup{CJKecglue={\hskip 0.45em plus 0.06em minus 0.09em}}

% --------------------------------------------------------------
%  页眉页脚: 偶数页眉右对齐斜体节题; 页码页脚外侧
% --------------------------------------------------------------
\fancyhf{}
\fancyhead[RE]{\small\itshape\rightmark}
\fancyfoot[LE,RO]{\small\thepage}
\renewcommand{\headrulewidth}{0pt}
\renewcommand{\sectionmark}[1]{\markright{\thesection.\ #1}}
\fancypagestyle{plain}{%
  \fancyhf{}%
  \fancyfoot[LE,RO]{\small \thepage}%
  \renewcommand{\headrulewidth}{0pt}%
}

% --------------------------------------------------------------
%  原书版式记号宏
% --------------------------------------------------------------
% 边码: 原打字稿页码(G. 手稿系统里的自引页码), 排在页边
\usepackage{marginnote}
\newcommand{\origpage}[1]{\marginnote{\normalfont\scriptsize p.~#1}}

% ⋄ 记号: 段首/行内正常大小; 上标形式为作者后补文字的位置标记
\DeclareRobustCommand{\pmark}{\ensuremath{\diamond}\hspace{0.15em}}
\DeclareRobustCommand{\dmark}{\textsuperscript{\ensuremath{\diamond}}}

% 脚注人工编号: \fnote{795}{(*)}{文本} —— 编号与(*)角标均照录原书
\newcommand{\fnote}[3]{%
  \begingroup
  \renewcommand{\thefootnote}{#1#2}%
  \footnote{#3}%
  \endgroup
}

% 第一部分的脚注: 页内序号作上标(#1), 注文行首印 "序号(连续注号)",
% 正文在序号之后另行括注同一连续注号 —— \fnoteP{1}{(3)}{注文}
% 原书有三条注无连续注号而改标日期(ch07 注3 = 30 Septembre,
% ch09 注11/12 = 8 août), 第二参数留空, 日期写进注文之首。
\newcommand{\fnoteP}[3]{%
  \begingroup
  \renewcommand{\thefootnote}{#1}\footnotemark
  \renewcommand{\thefootnote}{#1#2}\footnotetext{#3}%
  \endgroup
}

% 正文呼号与注文行首标号不同(原书照录)
\newcommand{\fnotevariant}[4]{%
  \begingroup
  \renewcommand{\thefootnote}{#1#2}\footnotemark
  \renewcommand{\thefootnote}{#1#3}\footnotetext{#4}%
  \endgroup
}

% 大注块标题: \notehead{171³}{斜体评论} —— 正文随后另起
\newcommand{\notehead}[2]{%
  \par\bigskip
  \noindent{\headfont 注~#1}\ifblank{#2}{}{\hspace{0.5em}{\itshape #2}}%
  \par\nobreak\smallskip\nobreak
}

% ∗ ∗ ∗ 分隔
\newcommand{\asterisms}{%
  \par\vspace{1.2ex}%
  \noindent\hspace*{2em}$\ast$\hspace{1.2em}$\ast$\hspace{1.2em}$\ast$%
  \par\vspace{1.2ex}%
}

% 无编号子目的目录条目
\newcommand{\tocsubsub}[1]{%
  \phantomsection
  \addcontentsline{toc}{subsubsection}{#1}%
}

% 部扉: 仿原书部题页
\newcommand{\partdivider}[3]{%
  \clearpage
  \thispagestyle{empty}%
  \phantomsection
  \addcontentsline{toc}{part}{#1\quad #3}%
  \vspace*{6.5cm}
  \begin{center}
    {\fontsize{16pt}{26pt}\selectfont\headfont #1\par}
    \vspace{0.9cm}
    {\fontsize{22pt}{34pt}\selectfont\headfont #2\par}
  \end{center}
  \clearpage
}

% 献辞页: 仿原书部题后的题献, 独占一页且不入目录
\newcommand{\dedication}[1]{%
  \thispagestyle{empty}%
  \vspace*{4cm}%
  \begin{center}#1\end{center}%
  \clearpage
}

% --------------------------------------------------------------
%  各卷"译者说明"的共用体例段(各卷的范围段自写, 后接本段)
% --------------------------------------------------------------
\newcommand{\CommonNoteBody}{%
为便于与原书及学界通行的打字稿页码对勘, 本册保留:

\begin{itemize}
  \item 原书全部章节编号(章~1--18, 节如~5.2、18.2, 子目如~18.1.1)与原书小节标题;
  \item 原书脚注的编号与角标——
        \begin{itemize}
          \item 前奏(第~1--4 章)与第一部的注在各页内重新编号, 正文中的上标即该
                页序号; 第一部紧随其后的括注(如“(3)”)是该部分的连续注号, 与第~12
                章“第一部分的注”里的“注~3”一一对应, 注文行首照原书印作“1(3)”式。
                其中三条注原书未标连续注号而改标日期(ch7 注~3 = 30 Septembre、
                ch9 注~11/12 = 8 août), 本册照录其日期;
          \item 第二部第~13--16~章各自重新编号, 分别为~1--41、1--86、1--132、1--29;
                第三部第~17~章自~1~起至~12, 第~18~章自~1~重新起算并延续至~794;
                末段自~\textbf{795(*)}~续至~1053(*)。各段与原书一一对应;
        \end{itemize}
  \item 页边码“p.~1”“p.~A3”“p.~1066”式小注——这是作者引用自己打字稿原稿时
        所用的页码, 与本排印本页码不同。前奏四章的边码各有前缀: 第~1 章作
        \texttt{A1…A6}, 第~2 章作 \texttt{P1…P65}, 第~3 章作 \texttt{L1…L56},
        第~4 章作罗马数字 \texttt{I…XXII}, 自第一部起作纯数字 \texttt{1…171};
  \item 作者用以标记段落的~$\diamond$~记号, 以及标示后来补写文字的上标
        ~\textsuperscript{$\diamond$}~记号;
  \item 原书大小注(“注~44'”“注~46$^1$”“注~171³”等)的编号与层级。
\end{itemize}

第~12 章“第一部分的注”把第~5 至第~11 章各页脚注汇为一编, 并给每条注配上节题
(如“12.19. 耶稣与十二使徒”), 因此部分注文与正文脚注重复——这是原书的体例, 非
重出之误。这一部分属随笔式的自省文字, 同一句法语在正文脚注与章末汇编注中
各有措辞, 本册有意保留各自的译笔, 不强求统一。

除上述标记外, 译文力求贴近作者驳杂、跳跃而辛辣的文风; 数学名词的译法见各卷
术语表(如~\emph{topos}~不译, \emph{faisceaux pervers}~作“反常层”)。文中的日期
括注(如“(5~月~30~日)”“(1984年3月补写)”)为作者历次补写时的纪日, 一律照录。文中
另有一些对注号、页码的引用(如“n°~1695”“打字稿第~227~页”)沿袭作者打字稿的编号
系统, 与本排印本不尽一致, 亦无法在本书中落实, 悉照原文录出。

全套五卷(另有全书合订一册)的分工是:

\begin{itemize}
  \item \textbf{卷一·前奏}——主题的呈现: 或四个乐章的前奏(原书 p.23--122);
  \item \textbf{卷二·第一部}——愚妄与更新(原书 p.123--254);
  \item \textbf{卷三·第二部}——葬礼(I): 或中国皇帝的新衣(原书 p.255--412);
  \item \textbf{卷四·第三部}——葬礼(II): 或阴阳之钥(原书 p.413--632);
  \item \textbf{卷五·第四部}——葬礼(III): 或四种运算(原书 p.633--929)。
\end{itemize}
}
                       % 共用导言区
%      % ================================================================
%  《收获与播种》统一封面 —— TikZ 原生矢量绘制(不依赖任何外部图片)
%  母题: 暮色中的麦田 —— 犁沟向地平线的灭点收束, 前景麦穗剪影,
%        一位播种者把一把种子撒向天空, 种子沿弧线落入田里。
%  用法(在 \begin{document} 之前, 依次):
%      \input{common}                       % 共用导言区
%      \input{cover-art}                    % 本文件
%      \CoverSetup{prelude}{卷一}{前奏}%
%                 {主题的呈现——或四个乐章的前奏}{p.23–122}
%  正文中再 \makecover 即可。配色关键字:
%      prelude / part1 / part2 / part3 / part4 / all
%  五卷共用同一幅画, 只换配色与卷次文字, 以保证"一套书"的观感。
% ================================================================

\RequirePackage{tikz}

% --------------------------------------------------------------
%  配色: 六套(前奏 / 四部 / 全书), 各为同一暮色母题的变调
% --------------------------------------------------------------
\newcommand{\CoverPalette}[1]{%
  \ifstrequal{#1}{prelude}{%
    \definecolor{cvSkyTop}{RGB}{10,26,48}%
    \definecolor{cvSkyMid}{RGB}{46,84,116}%
    \definecolor{cvSkyBot}{RGB}{232,164,92}%
    \definecolor{cvSun}{RGB}{255,224,168}%
    \definecolor{cvFieldTop}{RGB}{150,110,52}%
    \definecolor{cvFieldBot}{RGB}{58,36,18}%
    \definecolor{cvFurrow}{RGB}{40,22,10}%
    \definecolor{cvSil}{RGB}{30,18,10}%
    \definecolor{cvCream}{RGB}{247,240,228}%
  }{%
  \ifstrequal{#1}{part1}{%
    \definecolor{cvSkyTop}{RGB}{22,70,100}%
    \definecolor{cvSkyMid}{RGB}{92,150,168}%
    \definecolor{cvSkyBot}{RGB}{240,200,104}%
    \definecolor{cvSun}{RGB}{255,244,200}%
    \definecolor{cvFieldTop}{RGB}{176,140,58}%
    \definecolor{cvFieldBot}{RGB}{84,54,22}%
    \definecolor{cvFurrow}{RGB}{45,28,12}%
    \definecolor{cvSil}{RGB}{34,22,10}%
    \definecolor{cvCream}{RGB}{250,245,232}%
  }{%
  \ifstrequal{#1}{part2}{%
    \definecolor{cvSkyTop}{RGB}{46,20,30}%
    \definecolor{cvSkyMid}{RGB}{128,52,54}%
    \definecolor{cvSkyBot}{RGB}{226,120,64}%
    \definecolor{cvSun}{RGB}{255,206,150}%
    \definecolor{cvFieldTop}{RGB}{132,84,44}%
    \definecolor{cvFieldBot}{RGB}{48,26,18}%
    \definecolor{cvFurrow}{RGB}{36,16,14}%
    \definecolor{cvSil}{RGB}{26,12,12}%
    \definecolor{cvCream}{RGB}{248,236,224}%
  }{%
  \ifstrequal{#1}{part3}{%
    \definecolor{cvSkyTop}{RGB}{30,20,54}%
    \definecolor{cvSkyMid}{RGB}{92,62,124}%
    \definecolor{cvSkyBot}{RGB}{216,132,140}%
    \definecolor{cvSun}{RGB}{255,214,196}%
    \definecolor{cvFieldTop}{RGB}{128,90,74}%
    \definecolor{cvFieldBot}{RGB}{44,26,30}%
    \definecolor{cvFurrow}{RGB}{34,20,30}%
    \definecolor{cvSil}{RGB}{22,14,22}%
    \definecolor{cvCream}{RGB}{246,238,238}%
  }{%
  \ifstrequal{#1}{part4}{%
    \definecolor{cvSkyTop}{RGB}{12,28,28}%
    \definecolor{cvSkyMid}{RGB}{44,86,72}%
    \definecolor{cvSkyBot}{RGB}{186,168,86}%
    \definecolor{cvSun}{RGB}{240,232,168}%
    \definecolor{cvFieldTop}{RGB}{110,104,50}%
    \definecolor{cvFieldBot}{RGB}{30,34,20}%
    \definecolor{cvFurrow}{RGB}{26,32,20}%
    \definecolor{cvSil}{RGB}{14,20,14}%
    \definecolor{cvCream}{RGB}{240,238,226}%
  }{%
    % 默认: 全书(麦浪金褐)
    \definecolor{cvSkyTop}{RGB}{28,40,32}%
    \definecolor{cvSkyMid}{RGB}{128,110,64}%
    \definecolor{cvSkyBot}{RGB}{232,178,88}%
    \definecolor{cvSun}{RGB}{255,236,180}%
    \definecolor{cvFieldTop}{RGB}{168,124,52}%
    \definecolor{cvFieldBot}{RGB}{62,40,20}%
    \definecolor{cvFurrow}{RGB}{42,26,12}%
    \definecolor{cvSil}{RGB}{28,18,10}%
    \definecolor{cvCream}{RGB}{250,242,226}%
  }}}}}%
}

% --------------------------------------------------------------
%  一键设定:  \CoverSetup{配色}{卷次}{卷名}{副题}{原书页码范围}
% --------------------------------------------------------------
\newcommand{\CoverSetup}[5]{%
  \CoverPalette{#1}%
  \newcommand{\VolNum}{#2}%
  \newcommand{\VolName}{#3}%
  \newcommand{\VolSub}{#4}%
  \newcommand{\VolPages}{#5}%
}

% --------------------------------------------------------------
%  一支麦穗:  \cvWheatEar{x}{底 y}{高}{倾角}
%  茎略弯, 麦粒成对左右斜出, 顶端一芒三针
% --------------------------------------------------------------
\newcommand{\cvWheatEar}[4]{%
  \begin{scope}[shift={(#1,#2)}, rotate=#4]
    \draw[cvSil, line width=0.85mm, line cap=round]
      (0,0) .. controls (0.5,{0.35*#3}) and (0.8,{0.72*#3}) .. ({0.55*#3},#3);
    \foreach \cvK in {0,...,6}{%
      \pgfmathsetmacro{\cvT}{0.50+0.50*\cvK/6}
      \pgfmathsetmacro{\cvGY}{\cvT*#3}
      \pgfmathsetmacro{\cvGX}{0.55*\cvT*\cvT*#3}
      \pgfmathsetmacro{\cvGS}{1.0-0.34*\cvT}
      \begin{scope}[shift={(\cvGX,\cvGY)}, rotate=28]
        \fill[cvSil] (0,{1.9*\cvGS}) ellipse [x radius={0.92*\cvGS}, y radius={2.3*\cvGS}];
      \end{scope}
      \begin{scope}[shift={(\cvGX,\cvGY)}, rotate=-28]
        \fill[cvSil] (0,{-1.9*\cvGS}) ellipse [x radius={0.92*\cvGS}, y radius={2.3*\cvGS}];
      \end{scope}
    }
    \begin{scope}[shift={({0.55*#3},#3)}]
      \fill[cvSil] (0,1.5) ellipse [x radius=0.78, y radius=1.95];
    \end{scope}
    \draw[cvSil, line width=0.26mm, line cap=round]
      ({0.55*#3},{#3+2.3}) -- ({0.55*#3-1.7},{#3+9.5});
    \draw[cvSil, line width=0.26mm, line cap=round]
      ({0.55*#3},{#3+2.3}) -- ({0.55*#3+2.1},{#3+10.5});
    \draw[cvSil, line width=0.26mm, line cap=round]
      ({0.55*#3},{#3+2.3}) -- ({0.55*#3},{#3+11.5});
  \end{scope}
}

% --------------------------------------------------------------
%  封面
% --------------------------------------------------------------
\makeatletter
\renewcommand{\makecover}{%
  \begin{titlepage}
    \thispagestyle{empty}
    \begin{tikzpicture}[remember picture, overlay, x=1mm, y=1mm,
                        shift={(current page.south west)}]
      % ===== 天空 =====
      \shade[top color=cvSkyTop, middle color=cvSkyMid, bottom color=cvSkyBot]
        (0,105) rectangle (210,297);
      % ===== 田野 =====
      \shade[top color=cvFieldTop, bottom color=cvFieldBot]
        (0,0) rectangle (210,105);
      % ===== 落日与光晕 =====
      \begin{scope}
        \clip (0,105) rectangle (210,297);
        \fill[cvSun, opacity=0.14] (152,132) circle (46);
        \fill[cvSun, opacity=0.24] (152,132) circle (30);
      \end{scope}
      \fill[cvSun, opacity=0.95] (152,132) circle (16.5);
      % ===== 犁沟: 向地平线灭点收束 =====
      \begin{scope}
        \clip (0,0) rectangle (210,105);
        \foreach \cvI in {1,...,20}{%
          \pgfmathsetmacro{\cvDX}{3.1*pow(\cvI,1.55)}
          \draw[cvFurrow, opacity=0.30, line width=0.32mm] ({105+\cvDX},0) -- (105,105);
          \draw[cvFurrow, opacity=0.30, line width=0.32mm] ({105-\cvDX},0) -- (105,105);
        }
        \foreach \cvY in {8,18,31,47,68,92}{%
          \draw[cvFurrow, opacity=0.15, line width=0.3mm]
            (0,\cvY) .. controls (60,{\cvY+3.5}) and (150,{\cvY-3}) .. (210,{\cvY+1.5});
        }
      \end{scope}
      % 地平线
      \draw[cvSil, opacity=0.45, line width=0.45mm] (0,105) -- (210,105);
      % ===== 播种者剪影(立于地平线) =====
      \begin{scope}[cvSil]
        \draw[line width=1.9mm, line cap=round, line join=round]
          (55.5,116) -- (50,105)  (55.5,116) -- (60.5,105.5);
        \draw[line width=2.3mm, line cap=round, line join=round]
          (55.5,116) -- (55,123.5);
        \draw[line width=1.6mm, line cap=round, line join=round]
          (55.2,123.2) -- (61.5,126.6) -- (68,129.4);
        \draw[line width=1.5mm, line cap=round, line join=round]
          (55.2,123.2) -- (51.4,119.8) -- (48.6,116.2);
        \fill (55.6,127.3) circle (2.35);
      \end{scope}
      % ===== 种子: 沿贝塞尔弧撒出 =====
      \draw[cvCream, opacity=0.13, line width=0.5mm]
        (68,129.4) .. controls (100,168) and (160,162) .. (196,104);
      % 掌心那一把(尚未撒出)
      \foreach \cvI in {1,...,5}{%
        \pgfmathsetmacro{\cvJX}{1.1*sin(\cvI*53)}
        \pgfmathsetmacro{\cvJY}{1.1*cos(\cvI*83)}
        \fill[cvCream, opacity=0.9] ({68+\cvJX},{129.4+\cvJY}) circle (1.0);
      }
      \foreach \cvI in {1,...,26}{%
        \pgfmathsetmacro{\cvT}{0.035+0.92*\cvI/26}
        \pgfmathsetmacro{\cvU}{1-\cvT}
        \pgfmathsetmacro{\cvSX}{\cvU*\cvU*\cvU*68 + 3*\cvU*\cvU*\cvT*100
                                + 3*\cvU*\cvT*\cvT*160 + \cvT*\cvT*\cvT*196}
        \pgfmathsetmacro{\cvSY}{\cvU*\cvU*\cvU*129.4 + 3*\cvU*\cvU*\cvT*168
                                + 3*\cvU*\cvT*\cvT*162 + \cvT*\cvT*\cvT*104}
        \pgfmathsetmacro{\cvR}{1.55-0.85*\cvT}
        \pgfmathsetmacro{\cvJX}{1.7*sin(\cvI*97)}
        \pgfmathsetmacro{\cvJY}{1.7*cos(\cvI*71)}
        \fill[cvCream, opacity={0.88-0.28*\cvT}] ({\cvSX+\cvJX},{\cvSY+\cvJY}) circle (\cvR);
      }
      % ===== 前景麦穗(两排) =====
      \foreach \cvI in {1,...,11}{%
        \pgfmathsetmacro{\cvWX}{4+19*(\cvI-1)+1.6*sin(\cvI*2.1)}
        \pgfmathsetmacro{\cvWH}{23+7*sin(\cvI*1.7)}
        \pgfmathsetmacro{\cvWL}{3.2*sin(\cvI*2.9)}
        \begin{scope}[opacity=0.45]
          \cvWheatEar{\cvWX}{30}{\cvWH}{\cvWL}
        \end{scope}
      }
      \foreach \cvI in {1,...,10}{%
        \pgfmathsetmacro{\cvWX}{9+20.4*(\cvI-1)+1.5*sin(\cvI*1.3)}
        \pgfmathsetmacro{\cvWH}{41+10*sin(\cvI*2.3)}
        \pgfmathsetmacro{\cvWL}{4.5*sin(\cvI*3.7)}
        \cvWheatEar{\cvWX}{30}{\cvWH}{\cvWL}
      }
      % ===== 压暗一层, 统一色调并保证文字可读 =====
      \fill[black, opacity=0.10] (0,105) rectangle (210,297);
      % ===== 底部信息带 =====
      \fill[cvSil, opacity=0.72] (0,0) rectangle (210,30);
      \draw[cvCream, opacity=0.40, line width=0.4mm] (0,30) -- (210,30);
      % ===== 文字 =====
      \node[anchor=north west, inner sep=0, text=cvCream] at (22,268)
        {\fontsize{38pt}{45pt}\selectfont\headfont \ziju{0.08}\BookTitleCN};
      \node[anchor=north west, inner sep=0, text=cvCream] at (22.6,242)
        {\fontsize{15pt}{20pt}\selectfont\itshape \BookTitleEN};
      \draw[cvCream, opacity=0.55, line width=0.5mm] (22,233) -- (120,233);
      \node[anchor=north west, inner sep=0, text=cvCream] at (22,224)
        {\fontsize{11pt}{15pt}\selectfont 一位数学家对往昔的反思与见证};
      \node[anchor=north west, inner sep=0, text=cvCream, opacity=0.78] at (22.4,208)
        {\fontsize{10pt}{14pt}\selectfont \ziju{0.14}\VolNum};
      \node[anchor=north west, inner sep=0, text=cvCream] at (22,199)
        {\fontsize{30pt}{37pt}\selectfont\headfont \ziju{0.1}\VolName};
      \node[anchor=north west, inner sep=0, text=cvCream, opacity=0.92] at (22,180)
        {\fontsize{13pt}{18pt}\selectfont \VolSub};
      \node[anchor=north west, inner sep=0, text=cvCream] at (22,19)
        {\fontsize{11.5pt}{15pt}\selectfont \OriginalAuthor\quad 著};
      \node[anchor=north west, inner sep=0, text=cvCream, opacity=0.85] at (22,7)
        {\fontsize{9.5pt}{13pt}\selectfont \Translator\ 译\quad ·\quad
         \ifdefempty{\TranslationDate}{\today}{\TranslationDate}};
      \node[anchor=north east, inner sep=0, text=cvCream, opacity=0.85] at (188,13)
        {\fontsize{10pt}{14pt}\selectfont 原书 \VolPages};
    \end{tikzpicture}
  \end{titlepage}%
}
\makeatother
                    % 本文件
%      \CoverSetup{prelude}{卷一}{前奏}%
%                 {主题的呈现——或四个乐章的前奏}{p.23–122}
%  正文中再 \makecover 即可。配色关键字:
%      prelude / part1 / part2 / part3 / part4 / all
%  五卷共用同一幅画, 只换配色与卷次文字, 以保证"一套书"的观感。
% ================================================================

\RequirePackage{tikz}

% --------------------------------------------------------------
%  配色: 六套(前奏 / 四部 / 全书), 各为同一暮色母题的变调
% --------------------------------------------------------------
\newcommand{\CoverPalette}[1]{%
  \ifstrequal{#1}{prelude}{%
    \definecolor{cvSkyTop}{RGB}{10,26,48}%
    \definecolor{cvSkyMid}{RGB}{46,84,116}%
    \definecolor{cvSkyBot}{RGB}{232,164,92}%
    \definecolor{cvSun}{RGB}{255,224,168}%
    \definecolor{cvFieldTop}{RGB}{150,110,52}%
    \definecolor{cvFieldBot}{RGB}{58,36,18}%
    \definecolor{cvFurrow}{RGB}{40,22,10}%
    \definecolor{cvSil}{RGB}{30,18,10}%
    \definecolor{cvCream}{RGB}{247,240,228}%
  }{%
  \ifstrequal{#1}{part1}{%
    \definecolor{cvSkyTop}{RGB}{22,70,100}%
    \definecolor{cvSkyMid}{RGB}{92,150,168}%
    \definecolor{cvSkyBot}{RGB}{240,200,104}%
    \definecolor{cvSun}{RGB}{255,244,200}%
    \definecolor{cvFieldTop}{RGB}{176,140,58}%
    \definecolor{cvFieldBot}{RGB}{84,54,22}%
    \definecolor{cvFurrow}{RGB}{45,28,12}%
    \definecolor{cvSil}{RGB}{34,22,10}%
    \definecolor{cvCream}{RGB}{250,245,232}%
  }{%
  \ifstrequal{#1}{part2}{%
    \definecolor{cvSkyTop}{RGB}{46,20,30}%
    \definecolor{cvSkyMid}{RGB}{128,52,54}%
    \definecolor{cvSkyBot}{RGB}{226,120,64}%
    \definecolor{cvSun}{RGB}{255,206,150}%
    \definecolor{cvFieldTop}{RGB}{132,84,44}%
    \definecolor{cvFieldBot}{RGB}{48,26,18}%
    \definecolor{cvFurrow}{RGB}{36,16,14}%
    \definecolor{cvSil}{RGB}{26,12,12}%
    \definecolor{cvCream}{RGB}{248,236,224}%
  }{%
  \ifstrequal{#1}{part3}{%
    \definecolor{cvSkyTop}{RGB}{30,20,54}%
    \definecolor{cvSkyMid}{RGB}{92,62,124}%
    \definecolor{cvSkyBot}{RGB}{216,132,140}%
    \definecolor{cvSun}{RGB}{255,214,196}%
    \definecolor{cvFieldTop}{RGB}{128,90,74}%
    \definecolor{cvFieldBot}{RGB}{44,26,30}%
    \definecolor{cvFurrow}{RGB}{34,20,30}%
    \definecolor{cvSil}{RGB}{22,14,22}%
    \definecolor{cvCream}{RGB}{246,238,238}%
  }{%
  \ifstrequal{#1}{part4}{%
    \definecolor{cvSkyTop}{RGB}{12,28,28}%
    \definecolor{cvSkyMid}{RGB}{44,86,72}%
    \definecolor{cvSkyBot}{RGB}{186,168,86}%
    \definecolor{cvSun}{RGB}{240,232,168}%
    \definecolor{cvFieldTop}{RGB}{110,104,50}%
    \definecolor{cvFieldBot}{RGB}{30,34,20}%
    \definecolor{cvFurrow}{RGB}{26,32,20}%
    \definecolor{cvSil}{RGB}{14,20,14}%
    \definecolor{cvCream}{RGB}{240,238,226}%
  }{%
    % 默认: 全书(麦浪金褐)
    \definecolor{cvSkyTop}{RGB}{28,40,32}%
    \definecolor{cvSkyMid}{RGB}{128,110,64}%
    \definecolor{cvSkyBot}{RGB}{232,178,88}%
    \definecolor{cvSun}{RGB}{255,236,180}%
    \definecolor{cvFieldTop}{RGB}{168,124,52}%
    \definecolor{cvFieldBot}{RGB}{62,40,20}%
    \definecolor{cvFurrow}{RGB}{42,26,12}%
    \definecolor{cvSil}{RGB}{28,18,10}%
    \definecolor{cvCream}{RGB}{250,242,226}%
  }}}}}%
}

% --------------------------------------------------------------
%  一键设定:  \CoverSetup{配色}{卷次}{卷名}{副题}{原书页码范围}
% --------------------------------------------------------------
\newcommand{\CoverSetup}[5]{%
  \CoverPalette{#1}%
  \newcommand{\VolNum}{#2}%
  \newcommand{\VolName}{#3}%
  \newcommand{\VolSub}{#4}%
  \newcommand{\VolPages}{#5}%
}

% --------------------------------------------------------------
%  一支麦穗:  \cvWheatEar{x}{底 y}{高}{倾角}
%  茎略弯, 麦粒成对左右斜出, 顶端一芒三针
% --------------------------------------------------------------
\newcommand{\cvWheatEar}[4]{%
  \begin{scope}[shift={(#1,#2)}, rotate=#4]
    \draw[cvSil, line width=0.85mm, line cap=round]
      (0,0) .. controls (0.5,{0.35*#3}) and (0.8,{0.72*#3}) .. ({0.55*#3},#3);
    \foreach \cvK in {0,...,6}{%
      \pgfmathsetmacro{\cvT}{0.50+0.50*\cvK/6}
      \pgfmathsetmacro{\cvGY}{\cvT*#3}
      \pgfmathsetmacro{\cvGX}{0.55*\cvT*\cvT*#3}
      \pgfmathsetmacro{\cvGS}{1.0-0.34*\cvT}
      \begin{scope}[shift={(\cvGX,\cvGY)}, rotate=28]
        \fill[cvSil] (0,{1.9*\cvGS}) ellipse [x radius={0.92*\cvGS}, y radius={2.3*\cvGS}];
      \end{scope}
      \begin{scope}[shift={(\cvGX,\cvGY)}, rotate=-28]
        \fill[cvSil] (0,{-1.9*\cvGS}) ellipse [x radius={0.92*\cvGS}, y radius={2.3*\cvGS}];
      \end{scope}
    }
    \begin{scope}[shift={({0.55*#3},#3)}]
      \fill[cvSil] (0,1.5) ellipse [x radius=0.78, y radius=1.95];
    \end{scope}
    \draw[cvSil, line width=0.26mm, line cap=round]
      ({0.55*#3},{#3+2.3}) -- ({0.55*#3-1.7},{#3+9.5});
    \draw[cvSil, line width=0.26mm, line cap=round]
      ({0.55*#3},{#3+2.3}) -- ({0.55*#3+2.1},{#3+10.5});
    \draw[cvSil, line width=0.26mm, line cap=round]
      ({0.55*#3},{#3+2.3}) -- ({0.55*#3},{#3+11.5});
  \end{scope}
}

% --------------------------------------------------------------
%  封面
% --------------------------------------------------------------
\makeatletter
\renewcommand{\makecover}{%
  \begin{titlepage}
    \thispagestyle{empty}
    \begin{tikzpicture}[remember picture, overlay, x=1mm, y=1mm,
                        shift={(current page.south west)}]
      % ===== 天空 =====
      \shade[top color=cvSkyTop, middle color=cvSkyMid, bottom color=cvSkyBot]
        (0,105) rectangle (210,297);
      % ===== 田野 =====
      \shade[top color=cvFieldTop, bottom color=cvFieldBot]
        (0,0) rectangle (210,105);
      % ===== 落日与光晕 =====
      \begin{scope}
        \clip (0,105) rectangle (210,297);
        \fill[cvSun, opacity=0.14] (152,132) circle (46);
        \fill[cvSun, opacity=0.24] (152,132) circle (30);
      \end{scope}
      \fill[cvSun, opacity=0.95] (152,132) circle (16.5);
      % ===== 犁沟: 向地平线灭点收束 =====
      \begin{scope}
        \clip (0,0) rectangle (210,105);
        \foreach \cvI in {1,...,20}{%
          \pgfmathsetmacro{\cvDX}{3.1*pow(\cvI,1.55)}
          \draw[cvFurrow, opacity=0.30, line width=0.32mm] ({105+\cvDX},0) -- (105,105);
          \draw[cvFurrow, opacity=0.30, line width=0.32mm] ({105-\cvDX},0) -- (105,105);
        }
        \foreach \cvY in {8,18,31,47,68,92}{%
          \draw[cvFurrow, opacity=0.15, line width=0.3mm]
            (0,\cvY) .. controls (60,{\cvY+3.5}) and (150,{\cvY-3}) .. (210,{\cvY+1.5});
        }
      \end{scope}
      % 地平线
      \draw[cvSil, opacity=0.45, line width=0.45mm] (0,105) -- (210,105);
      % ===== 播种者剪影(立于地平线) =====
      \begin{scope}[cvSil]
        \draw[line width=1.9mm, line cap=round, line join=round]
          (55.5,116) -- (50,105)  (55.5,116) -- (60.5,105.5);
        \draw[line width=2.3mm, line cap=round, line join=round]
          (55.5,116) -- (55,123.5);
        \draw[line width=1.6mm, line cap=round, line join=round]
          (55.2,123.2) -- (61.5,126.6) -- (68,129.4);
        \draw[line width=1.5mm, line cap=round, line join=round]
          (55.2,123.2) -- (51.4,119.8) -- (48.6,116.2);
        \fill (55.6,127.3) circle (2.35);
      \end{scope}
      % ===== 种子: 沿贝塞尔弧撒出 =====
      \draw[cvCream, opacity=0.13, line width=0.5mm]
        (68,129.4) .. controls (100,168) and (160,162) .. (196,104);
      % 掌心那一把(尚未撒出)
      \foreach \cvI in {1,...,5}{%
        \pgfmathsetmacro{\cvJX}{1.1*sin(\cvI*53)}
        \pgfmathsetmacro{\cvJY}{1.1*cos(\cvI*83)}
        \fill[cvCream, opacity=0.9] ({68+\cvJX},{129.4+\cvJY}) circle (1.0);
      }
      \foreach \cvI in {1,...,26}{%
        \pgfmathsetmacro{\cvT}{0.035+0.92*\cvI/26}
        \pgfmathsetmacro{\cvU}{1-\cvT}
        \pgfmathsetmacro{\cvSX}{\cvU*\cvU*\cvU*68 + 3*\cvU*\cvU*\cvT*100
                                + 3*\cvU*\cvT*\cvT*160 + \cvT*\cvT*\cvT*196}
        \pgfmathsetmacro{\cvSY}{\cvU*\cvU*\cvU*129.4 + 3*\cvU*\cvU*\cvT*168
                                + 3*\cvU*\cvT*\cvT*162 + \cvT*\cvT*\cvT*104}
        \pgfmathsetmacro{\cvR}{1.55-0.85*\cvT}
        \pgfmathsetmacro{\cvJX}{1.7*sin(\cvI*97)}
        \pgfmathsetmacro{\cvJY}{1.7*cos(\cvI*71)}
        \fill[cvCream, opacity={0.88-0.28*\cvT}] ({\cvSX+\cvJX},{\cvSY+\cvJY}) circle (\cvR);
      }
      % ===== 前景麦穗(两排) =====
      \foreach \cvI in {1,...,11}{%
        \pgfmathsetmacro{\cvWX}{4+19*(\cvI-1)+1.6*sin(\cvI*2.1)}
        \pgfmathsetmacro{\cvWH}{23+7*sin(\cvI*1.7)}
        \pgfmathsetmacro{\cvWL}{3.2*sin(\cvI*2.9)}
        \begin{scope}[opacity=0.45]
          \cvWheatEar{\cvWX}{30}{\cvWH}{\cvWL}
        \end{scope}
      }
      \foreach \cvI in {1,...,10}{%
        \pgfmathsetmacro{\cvWX}{9+20.4*(\cvI-1)+1.5*sin(\cvI*1.3)}
        \pgfmathsetmacro{\cvWH}{41+10*sin(\cvI*2.3)}
        \pgfmathsetmacro{\cvWL}{4.5*sin(\cvI*3.7)}
        \cvWheatEar{\cvWX}{30}{\cvWH}{\cvWL}
      }
      % ===== 压暗一层, 统一色调并保证文字可读 =====
      \fill[black, opacity=0.10] (0,105) rectangle (210,297);
      % ===== 底部信息带 =====
      \fill[cvSil, opacity=0.72] (0,0) rectangle (210,30);
      \draw[cvCream, opacity=0.40, line width=0.4mm] (0,30) -- (210,30);
      % ===== 文字 =====
      \node[anchor=north west, inner sep=0, text=cvCream] at (22,268)
        {\fontsize{38pt}{45pt}\selectfont\headfont \ziju{0.08}\BookTitleCN};
      \node[anchor=north west, inner sep=0, text=cvCream] at (22.6,242)
        {\fontsize{15pt}{20pt}\selectfont\itshape \BookTitleEN};
      \draw[cvCream, opacity=0.55, line width=0.5mm] (22,233) -- (120,233);
      \node[anchor=north west, inner sep=0, text=cvCream] at (22,224)
        {\fontsize{11pt}{15pt}\selectfont 一位数学家对往昔的反思与见证};
      \node[anchor=north west, inner sep=0, text=cvCream, opacity=0.78] at (22.4,208)
        {\fontsize{10pt}{14pt}\selectfont \ziju{0.14}\VolNum};
      \node[anchor=north west, inner sep=0, text=cvCream] at (22,199)
        {\fontsize{30pt}{37pt}\selectfont\headfont \ziju{0.1}\VolName};
      \node[anchor=north west, inner sep=0, text=cvCream, opacity=0.92] at (22,180)
        {\fontsize{13pt}{18pt}\selectfont \VolSub};
      \node[anchor=north west, inner sep=0, text=cvCream] at (22,19)
        {\fontsize{11.5pt}{15pt}\selectfont \OriginalAuthor\quad 著};
      \node[anchor=north west, inner sep=0, text=cvCream, opacity=0.85] at (22,7)
        {\fontsize{9.5pt}{13pt}\selectfont \Translator\ 译\quad ·\quad
         \ifdefempty{\TranslationDate}{\today}{\TranslationDate}};
      \node[anchor=north east, inner sep=0, text=cvCream, opacity=0.85] at (188,13)
        {\fontsize{10pt}{14pt}\selectfont 原书 \VolPages};
    \end{tikzpicture}
  \end{titlepage}%
}
\makeatother
                    % 本文件
%      \CoverSetup{prelude}{卷一}{前奏}%
%                 {主题的呈现——或四个乐章的前奏}{p.23–122}
%  正文中再 \makecover 即可。配色关键字:
%      prelude / part1 / part2 / part3 / part4 / all
%  五卷共用同一幅画, 只换配色与卷次文字, 以保证"一套书"的观感。
% ================================================================

\RequirePackage{tikz}

% --------------------------------------------------------------
%  配色: 六套(前奏 / 四部 / 全书), 各为同一暮色母题的变调
% --------------------------------------------------------------
\newcommand{\CoverPalette}[1]{%
  \ifstrequal{#1}{prelude}{%
    \definecolor{cvSkyTop}{RGB}{10,26,48}%
    \definecolor{cvSkyMid}{RGB}{46,84,116}%
    \definecolor{cvSkyBot}{RGB}{232,164,92}%
    \definecolor{cvSun}{RGB}{255,224,168}%
    \definecolor{cvFieldTop}{RGB}{150,110,52}%
    \definecolor{cvFieldBot}{RGB}{58,36,18}%
    \definecolor{cvFurrow}{RGB}{40,22,10}%
    \definecolor{cvSil}{RGB}{30,18,10}%
    \definecolor{cvCream}{RGB}{247,240,228}%
  }{%
  \ifstrequal{#1}{part1}{%
    \definecolor{cvSkyTop}{RGB}{22,70,100}%
    \definecolor{cvSkyMid}{RGB}{92,150,168}%
    \definecolor{cvSkyBot}{RGB}{240,200,104}%
    \definecolor{cvSun}{RGB}{255,244,200}%
    \definecolor{cvFieldTop}{RGB}{176,140,58}%
    \definecolor{cvFieldBot}{RGB}{84,54,22}%
    \definecolor{cvFurrow}{RGB}{45,28,12}%
    \definecolor{cvSil}{RGB}{34,22,10}%
    \definecolor{cvCream}{RGB}{250,245,232}%
  }{%
  \ifstrequal{#1}{part2}{%
    \definecolor{cvSkyTop}{RGB}{46,20,30}%
    \definecolor{cvSkyMid}{RGB}{128,52,54}%
    \definecolor{cvSkyBot}{RGB}{226,120,64}%
    \definecolor{cvSun}{RGB}{255,206,150}%
    \definecolor{cvFieldTop}{RGB}{132,84,44}%
    \definecolor{cvFieldBot}{RGB}{48,26,18}%
    \definecolor{cvFurrow}{RGB}{36,16,14}%
    \definecolor{cvSil}{RGB}{26,12,12}%
    \definecolor{cvCream}{RGB}{248,236,224}%
  }{%
  \ifstrequal{#1}{part3}{%
    \definecolor{cvSkyTop}{RGB}{30,20,54}%
    \definecolor{cvSkyMid}{RGB}{92,62,124}%
    \definecolor{cvSkyBot}{RGB}{216,132,140}%
    \definecolor{cvSun}{RGB}{255,214,196}%
    \definecolor{cvFieldTop}{RGB}{128,90,74}%
    \definecolor{cvFieldBot}{RGB}{44,26,30}%
    \definecolor{cvFurrow}{RGB}{34,20,30}%
    \definecolor{cvSil}{RGB}{22,14,22}%
    \definecolor{cvCream}{RGB}{246,238,238}%
  }{%
  \ifstrequal{#1}{part4}{%
    \definecolor{cvSkyTop}{RGB}{12,28,28}%
    \definecolor{cvSkyMid}{RGB}{44,86,72}%
    \definecolor{cvSkyBot}{RGB}{186,168,86}%
    \definecolor{cvSun}{RGB}{240,232,168}%
    \definecolor{cvFieldTop}{RGB}{110,104,50}%
    \definecolor{cvFieldBot}{RGB}{30,34,20}%
    \definecolor{cvFurrow}{RGB}{26,32,20}%
    \definecolor{cvSil}{RGB}{14,20,14}%
    \definecolor{cvCream}{RGB}{240,238,226}%
  }{%
    % 默认: 全书(麦浪金褐)
    \definecolor{cvSkyTop}{RGB}{28,40,32}%
    \definecolor{cvSkyMid}{RGB}{128,110,64}%
    \definecolor{cvSkyBot}{RGB}{232,178,88}%
    \definecolor{cvSun}{RGB}{255,236,180}%
    \definecolor{cvFieldTop}{RGB}{168,124,52}%
    \definecolor{cvFieldBot}{RGB}{62,40,20}%
    \definecolor{cvFurrow}{RGB}{42,26,12}%
    \definecolor{cvSil}{RGB}{28,18,10}%
    \definecolor{cvCream}{RGB}{250,242,226}%
  }}}}}%
}

% --------------------------------------------------------------
%  一键设定:  \CoverSetup{配色}{卷次}{卷名}{副题}{原书页码范围}
% --------------------------------------------------------------
\newcommand{\CoverSetup}[5]{%
  \CoverPalette{#1}%
  \newcommand{\VolNum}{#2}%
  \newcommand{\VolName}{#3}%
  \newcommand{\VolSub}{#4}%
  \newcommand{\VolPages}{#5}%
}

% --------------------------------------------------------------
%  一支麦穗:  \cvWheatEar{x}{底 y}{高}{倾角}
%  茎略弯, 麦粒成对左右斜出, 顶端一芒三针
% --------------------------------------------------------------
\newcommand{\cvWheatEar}[4]{%
  \begin{scope}[shift={(#1,#2)}, rotate=#4]
    \draw[cvSil, line width=0.85mm, line cap=round]
      (0,0) .. controls (0.5,{0.35*#3}) and (0.8,{0.72*#3}) .. ({0.55*#3},#3);
    \foreach \cvK in {0,...,6}{%
      \pgfmathsetmacro{\cvT}{0.50+0.50*\cvK/6}
      \pgfmathsetmacro{\cvGY}{\cvT*#3}
      \pgfmathsetmacro{\cvGX}{0.55*\cvT*\cvT*#3}
      \pgfmathsetmacro{\cvGS}{1.0-0.34*\cvT}
      \begin{scope}[shift={(\cvGX,\cvGY)}, rotate=28]
        \fill[cvSil] (0,{1.9*\cvGS}) ellipse [x radius={0.92*\cvGS}, y radius={2.3*\cvGS}];
      \end{scope}
      \begin{scope}[shift={(\cvGX,\cvGY)}, rotate=-28]
        \fill[cvSil] (0,{-1.9*\cvGS}) ellipse [x radius={0.92*\cvGS}, y radius={2.3*\cvGS}];
      \end{scope}
    }
    \begin{scope}[shift={({0.55*#3},#3)}]
      \fill[cvSil] (0,1.5) ellipse [x radius=0.78, y radius=1.95];
    \end{scope}
    \draw[cvSil, line width=0.26mm, line cap=round]
      ({0.55*#3},{#3+2.3}) -- ({0.55*#3-1.7},{#3+9.5});
    \draw[cvSil, line width=0.26mm, line cap=round]
      ({0.55*#3},{#3+2.3}) -- ({0.55*#3+2.1},{#3+10.5});
    \draw[cvSil, line width=0.26mm, line cap=round]
      ({0.55*#3},{#3+2.3}) -- ({0.55*#3},{#3+11.5});
  \end{scope}
}

% --------------------------------------------------------------
%  封面
% --------------------------------------------------------------
\makeatletter
\renewcommand{\makecover}{%
  \begin{titlepage}
    \thispagestyle{empty}
    \begin{tikzpicture}[remember picture, overlay, x=1mm, y=1mm,
                        shift={(current page.south west)}]
      % ===== 天空 =====
      \shade[top color=cvSkyTop, middle color=cvSkyMid, bottom color=cvSkyBot]
        (0,105) rectangle (210,297);
      % ===== 田野 =====
      \shade[top color=cvFieldTop, bottom color=cvFieldBot]
        (0,0) rectangle (210,105);
      % ===== 落日与光晕 =====
      \begin{scope}
        \clip (0,105) rectangle (210,297);
        \fill[cvSun, opacity=0.14] (152,132) circle (46);
        \fill[cvSun, opacity=0.24] (152,132) circle (30);
      \end{scope}
      \fill[cvSun, opacity=0.95] (152,132) circle (16.5);
      % ===== 犁沟: 向地平线灭点收束 =====
      \begin{scope}
        \clip (0,0) rectangle (210,105);
        \foreach \cvI in {1,...,20}{%
          \pgfmathsetmacro{\cvDX}{3.1*pow(\cvI,1.55)}
          \draw[cvFurrow, opacity=0.30, line width=0.32mm] ({105+\cvDX},0) -- (105,105);
          \draw[cvFurrow, opacity=0.30, line width=0.32mm] ({105-\cvDX},0) -- (105,105);
        }
        \foreach \cvY in {8,18,31,47,68,92}{%
          \draw[cvFurrow, opacity=0.15, line width=0.3mm]
            (0,\cvY) .. controls (60,{\cvY+3.5}) and (150,{\cvY-3}) .. (210,{\cvY+1.5});
        }
      \end{scope}
      % 地平线
      \draw[cvSil, opacity=0.45, line width=0.45mm] (0,105) -- (210,105);
      % ===== 播种者剪影(立于地平线) =====
      \begin{scope}[cvSil]
        \draw[line width=1.9mm, line cap=round, line join=round]
          (55.5,116) -- (50,105)  (55.5,116) -- (60.5,105.5);
        \draw[line width=2.3mm, line cap=round, line join=round]
          (55.5,116) -- (55,123.5);
        \draw[line width=1.6mm, line cap=round, line join=round]
          (55.2,123.2) -- (61.5,126.6) -- (68,129.4);
        \draw[line width=1.5mm, line cap=round, line join=round]
          (55.2,123.2) -- (51.4,119.8) -- (48.6,116.2);
        \fill (55.6,127.3) circle (2.35);
      \end{scope}
      % ===== 种子: 沿贝塞尔弧撒出 =====
      \draw[cvCream, opacity=0.13, line width=0.5mm]
        (68,129.4) .. controls (100,168) and (160,162) .. (196,104);
      % 掌心那一把(尚未撒出)
      \foreach \cvI in {1,...,5}{%
        \pgfmathsetmacro{\cvJX}{1.1*sin(\cvI*53)}
        \pgfmathsetmacro{\cvJY}{1.1*cos(\cvI*83)}
        \fill[cvCream, opacity=0.9] ({68+\cvJX},{129.4+\cvJY}) circle (1.0);
      }
      \foreach \cvI in {1,...,26}{%
        \pgfmathsetmacro{\cvT}{0.035+0.92*\cvI/26}
        \pgfmathsetmacro{\cvU}{1-\cvT}
        \pgfmathsetmacro{\cvSX}{\cvU*\cvU*\cvU*68 + 3*\cvU*\cvU*\cvT*100
                                + 3*\cvU*\cvT*\cvT*160 + \cvT*\cvT*\cvT*196}
        \pgfmathsetmacro{\cvSY}{\cvU*\cvU*\cvU*129.4 + 3*\cvU*\cvU*\cvT*168
                                + 3*\cvU*\cvT*\cvT*162 + \cvT*\cvT*\cvT*104}
        \pgfmathsetmacro{\cvR}{1.55-0.85*\cvT}
        \pgfmathsetmacro{\cvJX}{1.7*sin(\cvI*97)}
        \pgfmathsetmacro{\cvJY}{1.7*cos(\cvI*71)}
        \fill[cvCream, opacity={0.88-0.28*\cvT}] ({\cvSX+\cvJX},{\cvSY+\cvJY}) circle (\cvR);
      }
      % ===== 前景麦穗(两排) =====
      \foreach \cvI in {1,...,11}{%
        \pgfmathsetmacro{\cvWX}{4+19*(\cvI-1)+1.6*sin(\cvI*2.1)}
        \pgfmathsetmacro{\cvWH}{23+7*sin(\cvI*1.7)}
        \pgfmathsetmacro{\cvWL}{3.2*sin(\cvI*2.9)}
        \begin{scope}[opacity=0.45]
          \cvWheatEar{\cvWX}{30}{\cvWH}{\cvWL}
        \end{scope}
      }
      \foreach \cvI in {1,...,10}{%
        \pgfmathsetmacro{\cvWX}{9+20.4*(\cvI-1)+1.5*sin(\cvI*1.3)}
        \pgfmathsetmacro{\cvWH}{41+10*sin(\cvI*2.3)}
        \pgfmathsetmacro{\cvWL}{4.5*sin(\cvI*3.7)}
        \cvWheatEar{\cvWX}{30}{\cvWH}{\cvWL}
      }
      % ===== 压暗一层, 统一色调并保证文字可读 =====
      \fill[black, opacity=0.10] (0,105) rectangle (210,297);
      % ===== 底部信息带 =====
      \fill[cvSil, opacity=0.72] (0,0) rectangle (210,30);
      \draw[cvCream, opacity=0.40, line width=0.4mm] (0,30) -- (210,30);
      % ===== 文字 =====
      \node[anchor=north west, inner sep=0, text=cvCream] at (22,268)
        {\fontsize{38pt}{45pt}\selectfont\headfont \ziju{0.08}\BookTitleCN};
      \node[anchor=north west, inner sep=0, text=cvCream] at (22.6,242)
        {\fontsize{15pt}{20pt}\selectfont\itshape \BookTitleEN};
      \draw[cvCream, opacity=0.55, line width=0.5mm] (22,233) -- (120,233);
      \node[anchor=north west, inner sep=0, text=cvCream] at (22,224)
        {\fontsize{11pt}{15pt}\selectfont 一位数学家对往昔的反思与见证};
      \node[anchor=north west, inner sep=0, text=cvCream, opacity=0.78] at (22.4,208)
        {\fontsize{10pt}{14pt}\selectfont \ziju{0.14}\VolNum};
      \node[anchor=north west, inner sep=0, text=cvCream] at (22,199)
        {\fontsize{30pt}{37pt}\selectfont\headfont \ziju{0.1}\VolName};
      \node[anchor=north west, inner sep=0, text=cvCream, opacity=0.92] at (22,180)
        {\fontsize{13pt}{18pt}\selectfont \VolSub};
      \node[anchor=north west, inner sep=0, text=cvCream] at (22,19)
        {\fontsize{11.5pt}{15pt}\selectfont \OriginalAuthor\quad 著};
      \node[anchor=north west, inner sep=0, text=cvCream, opacity=0.85] at (22,7)
        {\fontsize{9.5pt}{13pt}\selectfont \Translator\ 译\quad ·\quad
         \ifdefempty{\TranslationDate}{\today}{\TranslationDate}};
      \node[anchor=north east, inner sep=0, text=cvCream, opacity=0.85] at (188,13)
        {\fontsize{10pt}{14pt}\selectfont 原书 \VolPages};
    \end{tikzpicture}
  \end{titlepage}%
}
\makeatother

\CoverSetup{prelude}{卷一}{前奏}{主题的呈现——或四个乐章的前奏}{p.23–122}

\hypersetup{
  pdftitle={收获与播种·卷一·前奏},
  pdfauthor={Alexandre Grothendieck / 数译 译}
}

% --------------------------------------------------------------
\begin{document}

\frontmatter
\makecover

% ---------------------------------------------------------------
%  译者说明
% ---------------------------------------------------------------
\chapter*{译者说明}
\markboth{译者说明}{译者说明}
\addcontentsline{toc}{chapter}{译者说明}

本册是《收获与播种》中文译本的第一卷, 收原书“前奏:主题的呈现——或四个乐章的
前奏”(Présentation des Thèmes ou Prélude en quatre Mouvements), 即原书第~23 至
第~122 页, 含第~1 至第~4~章:

\begin{itemize}
  \item 第~1~章“权作序言”——交代成书缘起与写作方式;
  \item 第~2~章“巡游一部作品(或孩子与母亲)”——逐项清点自己留下的十二个主题;
  \item 第~3~章“一封信”——写给一位同行的一封长信, 兼作全书导引;
  \item 第~4~章“引言”——以四“口气”记述本书的成书过程与内在节奏。
\end{itemize}

\CommonNoteBody

\makecontents

\mainmatter

% ================================================================
%  前奏 (第 1-4 章, 原书 p.23-122)
% ================================================================
\partdivider{前奏}{主题的呈现\\[0.35em]{\normalsize 或四个乐章的前奏}}{主题的呈现——或四个乐章的前奏}

\setcounter{chapter}{0}
\markboth{收获与播种·前奏}{}

% ================================================================
%  第 1 章 权作序言 (原书 p.27-30, 边码 A1-A6)
%  源文: _work/src1/ch01.md
% ================================================================
\chapter{权作序言}

% ---------------- p.27 ----------------

\noindent\pmark 1986年1月30日\par

\origpage{A1} 只差一篇序言没写,就可以把《收获与播种》交给印刷厂了。我发誓,我是怀着世上最好的意愿,想写出点像样的东西来的。这一回,要写点合乎情理的东西。三四页,不能再多,但要写得有分量,用来引荐这头一千多页的庞然大物。要写点能“抓住”厌倦了的读者的东西,让他隐约看到:在这令人生畏的“一千多页”里,也许有些叫他感兴趣的东西(甚至,谁知道呢,与他有关的东西?)。“抓人”不是我的风格,不是。可这一回,我偏要破一次例!那位“疯到敢冒这个险的出版商”(敢出版这头显然无法出版的巨兽),总得勉强收回成本吧。

结果,没成。我尽力了。而且不止一个下午,不像我原先打算的那样三下两下就完事。到明天,整整三个星期了,我一直在写它,纸页越堆越高。写出来的东西,肯定不是人们体面地称之为“序言”的那种。又一次搞砸了,没治!我这把年纪,改不了了——我也不是那块料,不是叫卖的料,也不是让人叫卖的料。即便只为讨人喜欢(讨自己喜欢,也讨朋友喜欢……)。

写出来的,是一种长长的、带评注的“巡游”,穿过我作为数学家的作品。这巡游首先是给“门外汉”的——给那个“对数学一窍不通”的人。也是为我自己,我此前从未有过这样一趟巡游的闲情。一针引一线,我发现自己被引到一些向来留在未言说之中的东西面前,并且把它们说了出来。说来也巧,这些也正是我觉得最本质的东西——在我的工作中,在我的作品里。这些东西没有一点技术性。我这份天真的“传达”是否成功,由你评判——这尝试本身大概也有点疯。我的满足与快乐,将是让你感受到了它们。这些东西,我那些博学的同行里,许多人已经不再感受得到了。也许他们变得太博学、太有名望了。这往往让人失去与单纯而本质之物的接触。

在这趟“巡游一部作品”中,我也稍稍谈到我的一生。也在这里那里,稍稍谈到《收获与播种》讲的是什么。在“巡游”之后的那封“信”(注明去年五月)里,我还会更详细地再谈。这封信原本是写给我从前的学生们,以及数学界里我“昔日的朋友”的。但它同样没有一点技术性。任何读者,只要想通过一份“现场”的叙述,了解究竟哪些前因后果最终促使我写下《收获与播种》,都可以毫无困难地读它。比起“巡游”,它更能让你预先尝到“大世界”——数学的“大世界”——里某种气氛的味道。也(和“巡游”一样)尝到我的表达风格的味道,据说有点特别。还有借这风格表达出来的那种精神——这种精神同样不是人人都欣赏的。

\origpage{A2} 在“巡游”里,以及在《收获与播种》各处,我谈数学工作。这工作我熟悉,是第一手的。我所说的多数东西,对一切创造性工作、一切发现的工作,肯定都是真的。至少对所谓“智力”工作是真的——那种主要“靠脑子”和靠写来进行的工作。这样的工作,其标志是理解的绽放与成熟,即对我们正在探究之物的理解。% ---------------- p.28 ----------------
但举一个处于另一极的例子:爱的激情,同样也是发现的冲动。它为我们打开一种所谓“肉体的”认识,这种认识同样在更新、在舒展、在加深。这两种冲动——不妨说,一种驱动着工作中的数学家,另一种驱动着情人——彼此之近,远过于人们通常的猜度,也远过于人们愿意向自己承认的。我希望《收获与播种》的这些页面,能在你的工作中、在你每日的生活里,帮你感受到这一点。

在巡游中,主要谈的是数学工作本身。相反,对于这工作所处的境况,对于在真正的工作时间之外起作用的那些动机,我几乎一言不发。这有可能给人的形象——给我的形象,或给数学家、给一般“科学家”的形象——添上一层固然讨喜、却已走样的光晕。像是“伟大而高尚的激情”,没有任何修正。总之,合于那位大大的“科学神话”(S 请大写!)的路数。那英雄的、“普罗米修斯式”的神话,文人与学者们争先恐后地掉了进去(如今还在掉)。也许只有史家有时能抵挡这诱人的神话。实情是:在“科学家”的种种动机里——那些有时驱使他毫不吝惜地投入工作的动机——野心与虚荣所起的作用,与在任何别的行当里一样重大,而且几乎同样普遍。其形式或粗或细,因人而异。我绝不认为自己是个例外。读我这份见证,我希望,在这一点上不会留下任何疑问。

\origpage{A3} 也确实,最贪婪的野心也无力发现哪怕一条数学陈述,更无力证明它——正如它(譬如)无力“让人勃起”(就字面意义而言)。无论男女,使人“勃起”的绝不是野心,不是想要出风头、想要展示一种力量的欲望——那力量在此是性方面的——恰恰相反!而是对某种既强有力、又极真实、同时又极细微之物的敏锐知觉。人们可以称之为“美”,而美只是那东西的千张面孔之一。有野心,并不必然妨碍人偶尔感到一个存在、一件事物的美,这我同意。但可以肯定的是:让我们感到这美的,不是野心……

第一个发现并驾驭了火的人,正是和你我一样的人。绝不是人们以“英雄”“半神”之类名目设想的那种人。肯定,像你像我一样,他尝过焦虑的咬啮,也尝过那种虚荣的药膏——它让人忘掉咬啮。但在他“认识”火的那一刻,既没有恐惧,也没有虚荣。英雄神话里的真相就在于此。当神话被用来替我们遮掩事物的另一面——同样真实、同样本质的另一面——它就变得寡淡无味,变成了药膏。

我在《收获与播种》里的用意,是谈这两个方面——谈求知的冲动,也谈恐惧及其虚荣的解药。我相信自己“理解”、至少是认识这冲动及其本性。(也许有一天我会惊奇地发现,自己原来错得多么厉害……)但说到恐惧与虚荣,以及由此派生的对创造力的暗中阻塞,我很清楚自己并未深入这大谜的底。我也不知道,在我余下的岁月里,是否会看到这奥秘的底……

在写作《收获与播种》的过程中,浮现出两个形象,分别代表人类这场冒险的两个方面。那就是孩子(亦名“工匠”),以及老板。在即将开始的巡游里,要谈的几乎只是“孩子”。副题“孩子与母亲”里的,也是他。这名字,我希望,会在巡游途中渐渐明朗起来。

\origpage{A4} 而在反思其余的全部篇幅里,相反,占据前台的主要是老板。他可不是白当老板的! % ---------------- p.29 ----------------
更确切地说,不是一个老板,而是些相互竞争的企业的老板们。不过也确实:所有老板在本质上都彼此相像。而一旦开始谈老板,也就意味着要出现“坏蛋”了。在反思的第(I)部分(“愚妄与更新”,紧接本篇导论性的部分——即“四个乐章的前奏”——之后),那个“坏蛋”主要是我。在其后三个部分里,则主要是“别人”。轮着来!

这就是说,除了深刻的哲学思考与(毫无悔意的)“忏悔”,还会有“硫酸肖像”(借用我一位同行兼朋友的说法——他自己也挨了几下……)。更不必说那些大手笔、毫不含糊的“手术”。罗贝尔·若兰(Robert Jaulin)\fnote{1}{}{罗贝尔·若兰是老朋友了。我似乎明白,面对民族学界的建制,他处在一个“白狼”般的地位,与我在数学“上流社会”面前的处境略有相似。}向我保证(半开玩笑地):我在《收获与播种》里做的是“数学界的民族志”(也许是社会学,我也说不准了)。当然,听说自己(连自己都不知道)做着高深学问,总是受用的!事实是:在反思的“调查”部分里(而且违背我的意愿……),我看见我正在写下的那些纸页上,数学界建制的一大半人列队走过,更不必说许多地位更卑微的同行与朋友。而在最近这几个月里,自从去年十月我把《收获与播种》的试印本寄出去之后,这事又“来了一遍”。我这份见证,真是像一块石头扔进了死水塘。回响各色各样,几乎什么调子都有(只除了无聊那一种……)。几乎每一次,都完全出乎我的预料。此外还有许多沉默,意味深长。显然,关于我那些从前的学生、以及处境或高或低的同行们脑子里在转什么,我过去(如今也还)有许多要学,而且是形形色色的功课——请原谅,我是说“数学界的社会学”!对所有已经前来为我晚年这部社会学巨著添砖加瓦的人,我要在此表达我的感激之情。

当然,对那些暖色调的回响,我格外敏感。也有少数几位同行,向我道出了他们的感动,或一种(此前未表达过的)危机感、衰败感——对他们自认所属的这个数学界内部的危机与衰败。

\origpage{A5} 在这个圈子之外,最早对我的见证报以热烈、甚至激动回应的人里,我想在此举出西尔维·谢瓦莱与卡特琳·谢瓦莱(Sylvie et Catherine Chevalley)\fnote{2}{}{西尔维与卡特琳·谢瓦莱是克劳德·谢瓦莱的遗孀与女儿;这部《收获与播种》的核心部分(ReS III,“阴阳之钥”)正是献给他的。反思中多处谈到他,以及他在我的历程中所扮演的角色。}、罗贝尔·若兰、斯泰凡·德利若尔日(Stéphane Deligeorge)、克里斯蒂安·布尔茹瓦(Christian Bourgois)。如果说《收获与播种》将获得比最初那份(只给极狭小圈子的)试印本更广的流传,那主要归功于他们。尤其归功于他们那种有感染力的确信:我竭力去把握、去说出的东西,是应当被说出的;而且它能够被听到的范围,比我那些同行(往往阴沉,甚至乖戾,毫不愿意自我质疑……)的圈子更大。于是克里斯蒂安·布尔茹瓦毫不犹豫地承担了出版这“无法出版之物”的风险,而斯泰凡·德利若尔日则给我这份难以下咽的见证以殊荣,把它收进“Epistémé”丛书,(眼下)与牛顿、居维叶、阿拉戈并列。(我想不出更好的同伴了!)对每一位在我格外“敏感”的时刻一再表示同情与信任的人,我很高兴在此说出我全部的谢意。

于是我们这就出发,开始一趟巡游一部作品的旅程,作为穿越一生的旅行的开场。是的,一趟漫长的旅行,一千多页,每一页都塞得满满。这旅行我花了一辈子去做,也没做完,又花了一年多一页一页地重新发现它。有时词语迟迟不肯来,来表达一份仍在躲闪、尚未被把握的经验的全部汁液——就像压榨机里堆得满满的、熟透而饱满的葡萄,有时似乎想躲开紧攥着它的那股力量……但即便在词语仿佛互相推挤、汩汩涌出的时刻,它们的推挤与涌出也绝不是信手拈来。每一个词都在途中被称过,或者事后被称过,若发现太轻或太重,就仔细调准。所以这份反思—见证—旅行,不是为着让人一天读一点、一个月读完的,也不是为着那个急于读到“最后一句话”的读者。

% ---------------- p.30 ----------------
《收获与播种》里没有“最后一句话”,没有“结论”,正如我的一生里、你的一生里也没有。有的是一坛酒,在我这个人的木桶里陈了一辈子。你喝下的最后一杯不会比第一杯或第一百杯更好。它们全是“同一杯”,又各不相同。而倘若第一杯坏了,整桶都坏了;那还不如喝好水(若找得到),别喝坏酒。

\origpage{A6} 但好酒不宜匆匆喝,也不宜站着喝。

% ================================================================
%  第 2 章 巡游一部作品——或孩子与母亲 (原书 p.31-53, 边码 P1-P37)
%  源文: _work/src1/task-ch02a.md
%  说明: 章内 Contents 目录不排印
% ================================================================
\chapter{巡游一部作品——或孩子与母亲}

\noindent 1986年1月\par

\section{(1) 事物的魔法}

% ---------------- p.31 ----------------
\origpage{P1} 小时候,我很喜欢上学。读和写、算术、唱歌(他拉一把小提琴给我们伴奏),还有史前人类和火的发现,都是同一位老师教我们。我不记得那时候我们在学校里曾经觉得无聊。他掌握着数字的魔法,也掌握着词语、符号和声音的魔法。还有韵脚的魔法,在歌谣里,在小小的诗里。韵脚里似乎藏着某种超出词语之外的奥秘。这情形一直持续到有一天
% ---------------- p.32 ----------------
有人向我解释,那不过是个再简单不过的“窍门”:所谓押韵,无非是让两个前后相接的说话片段以同一个音节收尾,于是它们就像中了魔法一样,变成了诗句。这可是一场启示!在家里,我身边有人接应,一连几个星期、几个月,我都玩着作诗。有一阵子,我说话只押韵了。幸好这毛病过去了。可即便今天,偶尔我还会写写诗——只是不大再去找韵脚了,除非它自己跑来。

另有一回,一个比我大的伙伴,已经在读中学了,教了我负数。那是另一种很好玩的游戏,但很快玩腻了。还有填字游戏——我花了好几天、好几个星期来编,越编越交错复杂。这游戏里结合着形式的魔法,以及符号和词语的魔法。但这份热情也离开了我,似乎没留下什么痕迹。

在中学里,先是在德国读第一年,后来在法国,我是个不错的学生,但算不上“出色的学生”。凡是我最感兴趣的,我不计工本地投入;不太感兴趣的,就往往被撂在一边,也不太在乎那位“老师”怎么评价。在法国上中学的第一年,1940年,我和母亲一起被关在芒德附近的里厄克罗集中营。那是战争年代,我们是外国人——用当时的话说,是“不受欢迎的人”。但集中营的管理当局对营里的孩子睁一只眼闭一只眼,尽管他们同样是“不受欢迎的人”。我们进出多少随自己的便。我年纪最大,也是唯一去上中学的,学校在四五公里外,不管下雪还是刮风,脚上穿着总是进水的破鞋。\origpage{P2} 我还记得第一次“数学考试”,老师给了我一个坏分数,为的是证明“三角形全等的三种情形”之一。我的证明不是书上那个,而他是虔诚地照书办的。可我心里很清楚,我的证明并不比书上那个更有说服力,也不更没说服力——我遵循的正是书里的精神,用的还是那套老掉牙的“把某个图形这样挪到另一个图形上去”。显然,教我的这个人觉得自己没有能力凭自己的眼光来判断(在这里,就是判断一个推理是否成立)。他必须求助于一种权威,此处即一本书的权威。这些做派一定深深触动了我,我才会记得这桩小事。此后直到今天,我有很多机会看到:这样的做派绝不是例外,而是几乎普遍的规则。这个话题有很多可说的——在《收获与播种》里,我不止一次以这种或那种形式触及它。可直到今天,不管我愿不愿意,每次再撞上它,我都感到手足无措……

战争最后几年,母亲仍被关在集中营里,我则待在“瑞士救援会”办的一所儿童之家,收容难民儿童,在利尼翁河畔勒尚邦。我们大多是犹太人;每当接到通知(由当地警察来通知)说盖世太保要来搜捕,我们就到树林里躲上一两夜,三三两两一小群,并不太意识到这真的关系到我们的性命。那一带塞文山区藏满了犹太人,许多人靠当地居民的互助活了下来。

在“塞文中学”(我在那里受教育)最让我吃惊的,是我的同学们对所学的东西兴趣之少。至于我,学年一开始就把课本狼吞虎咽地读完,心想这一回总算要学到真正有趣的东西了;而这一年余下的时间,我就尽量自己打发,与此同时,预定的教学大纲一个学期接一个学期,无情地照本宣科。不过我们的老师都很好。博物学老师弗里德尔先生,人品与学识都极出色。可他不会“镇压”,被人起哄到死,以致到了学年末,课已经没法再上下去了,他那无力的声音
% ---------------- p.33 ----------------
被一片喧闹淹没。说不定正因为这个,我才没成为生物学家!

\origpage{P3} 我相当多的时间——甚至上课的时候(嘘……)——都在做数学题。很快,书上的题就不够我做了。也许是因为它们久而久之彼此太相像了;但我相信,尤其是因为它们来得太像从天而降,就这么一个接一个排着队,不说自己从哪儿来、往哪儿去。那是书上的问题,不是我的问题。其实真正自然的问题并不缺。比如,已知三角形三边之长 $a$、$b$、$c$,这个三角形就确定了(不计其位置),那么就该有一个显式的“公式”,例如把三角形的面积表示成 $a$、$b$、$c$ 的函数。四面体也一样:已知六条棱的长,体积是多少?这一回我想我大概费了不少劲,但硬啃下去,终究还是做出来了。反正,一件事“抓住”了我,我就不数在上面花了多少小时、多少天,哪怕把别的全都忘了!(如今也还是这样……)

我们的数学书里最让我不满意的,是对长度(一条曲线的)、面积(一个曲面的)、体积(一个立体的)这些概念,完全没有认真的定义。我向自己保证,一旦有闲,就要补上这个空缺。1945到1948年间,我在蒙彼利埃大学念书,把最精壮的精力都花在了这上面。学院的课不合我的意。虽然从未对自己明说,我大概觉得老师们不过是在重复他们的书,就像芒德中学那位第一位数学老师一样。所以我也只是隔很久才去一趟学院,为的是了解那个老掉牙的“教学大纲”讲到了哪里。就那大纲而言,书是够用的;但同样清楚的是,它们根本不回答我给自己提的问题。说实话,它们连看都没看见那些问题,就像我的中学课本也没看见一样。只要它们随手给出计算长度、面积、体积的配方,靠单重、双重、三重积分(超过三维的情形则谨慎地绕开……),那么为它们给出一个内蕴的定义这回事,似乎就根本不成为问题——对我的老师们如此,对教科书的作者们也如此。

\origpage{P4} 以我当时有限的经验看,完全可能显得我是世界上唯一对数学问题怀有好奇心的人。无论如何,在那些年里,这就是我未曾说出口的信念;那些年我处于完全的智识孤独之中,而这份孤独并不使我难受\fnote{1}{}{1945到1948年间,我和母亲住在蒙彼利埃十来公里外的一个小村子里,迈拉尔格(属旺达尔格乡),四野是葡萄园中的一处僻静地。(父亲1942年在奥斯维辛失踪。)我们靠我那份微薄的学生助学金过活,过得很拮据。为了收支相抵,我每年去摘葡萄,摘完葡萄再做些“拾穗酒”,勉强脱手(据说这违反了当时的法律……)。此外还有一处园子,从不用侍弄,却大量供给我们无花果、菠菜,甚至(到最后)还有番茄——那是一位好心的邻居种下的,种在一片绚丽的罂粟正中间。日子过得美——只是有时紧得刚好差那么一点,比如要换一副眼镜架,或者一双穿到露出麻线的鞋。幸而母亲在集中营里待了那么久,身体虚弱多病,可以享受免费医疗。否则我们永远付不起看医生的钱……}。说实话,我想那时候我从未深究过:自己是否真是世上唯一可能对我所做的事感兴趣的人。我的精力已充分投入那场自己给自己出的挑战:发展一套能让我完全满意的理论。

我丝毫没有怀疑自己终能做到、终能找到事物的最终答案——只要我肯花力气去细究它们,把它们对我说的随说随记到纸上。比如说,体积的直觉是无可置疑的。它只能是一个现实的映照,这现实眼下还难以捉摸,但完全可靠。要把握的就是这个现实,如此而已——也许有点像有一天“韵脚”那个神奇的现实被把握住、“被理解”那样。

% ---------------- p.34 ----------------
动手的时候,我十七岁,刚出中学,我以为这不过是几个星期的事。结果我在上面待了三年。我还想方设法,硬是在大学二年级期末考砸了一门——球面三角(在“深入天文学”选修方向里,原文如此),因为一个愚蠢的数值计算错误。(得承认,一离开中学,我的计算就一直不怎么行……)正因如此,我才不得不在蒙彼利埃再待第三年,在那里读完学士学位,而不能马上去巴黎——据人向我保证,那是唯一有机会遇到懂行之人的地方,懂数学里什么才算重要。给我通风报信的苏拉先生还向我保证:数学里最后剩下的一些问题,二三十年前已由一位名叫勒贝格的人解决了。他恰好发展了一套测度与积分理论,(说真的,巧得可笑!)给数学画上了句号。

\origpage{P5} 苏拉先生教我们“微积分”,是个和善的人,对我一片好意。我并不认为他因此说服了我。我心里大概早已预感到:数学在广度与深度上都是无边无际的。大海有“句号”吗?总之,我从未动过念头去把苏拉先生提到的勒贝格那本书找来,而他自己大概也从未把这本书拿在手里过。在我心里,一本书里可能装的东西,和我为了满足自己对那些令我不解之事的求知欲、照自己的方式所做的工作,毫无共同之处。

\section{(2) 独处的重要}

一两年后,我终于在巴黎与数学界有了接触,最终在那里得知(在许许多多别的之外):我在自己角落里用凑合的手段做出来的工作,(几乎)就是“人人”都熟知的东西,叫做勒贝格测度与积分理论。在我向其谈起这项工作(甚至给他们看一份手稿)的两三位年长者眼里,这有点像我白白浪费了时间,重做了一遍“已经知道的东西”。我也不记得自己有过失望。那时候,为自己所做的工作去获取什么“信誉”、哪怕只是一句赞许或别人一点兴趣,这个念头大概还离我的头脑很远。何况我的精力已经够忙的了:要熟悉一个完全不同的圈子,尤其是要学会巴黎人眼中数学家的 ABC\fnote{2}{}{我在《收获与播种》第一部分(ReS I)“受欢迎的异乡人”一节(第9号)里,对这段有点粗粝的过渡时期作了简短的叙述。}。

然而,现在回想这三年,我意识到它们绝非白费。在不知不觉中,我在孤独里学到了数学家这一行当最本质的东西——那是任何老师都无法真正教给人的。我从未需要对自己说明,也从未遇到一个可以分享我求知之渴的人,可我却知道,用俗话说,“打骨子里”知道:我是一个数学家——一个在完全意义上“做”数学的人,就像人们“做”爱那样。数学对我来说,成了一位永远迎候我欲望的情人。这些孤独的岁月奠定了一份从未动摇过的信赖的基础——无论是二十岁到巴黎时发现自己的无知有多么深、要学的东西有多么浩瀚;还是二十多年后,我永远离开数学界时那些动荡的插曲;还是最近这些年,由我昔日最亲近的伙伴们操办的那场针对我本人和我的作品的、相当疯狂的“葬礼”
% ---------------- p.35 ----------------
(提前举行,而且办得干净利落)……\origpage{P6}

换一种说法:在这些关键的年月里,我学会了独处\fnote{3}{}{这个说法有些不确切。我从来不需要“学会独处”,原因很简单:在我的童年里,我从未丢掉出生时就在我身上、也就在每个人身上的那种天生能力。但这三年孤独的工作——在其中我得以按自己那套自发的严格要求的标准,给自己量了尺寸——却在我身上确认并重新奠下了一份信赖与安然的确据,这一次是在我与数学工作的关系上;这份确据与那些立法的共识和风尚毫无关系。我在《根与孤独》一注(ReS IV,第1713号,尤其第1080页)里有机会再次提到它。}。我的意思是:凭自己的眼光去接近我想认识的事物,而不是信赖那些明说的或默认的观念与共识——它们来自我自认为属于的某个或大或小的群体,或由于任何别的理由而被赋予对我的权威。中学也好,大学也好,那些无声的共识曾告诉我:对“体积”这个概念本身,没有必要提出问题,它被说成“人所共知”“显而易见”“没有问题”。我越过了它,当作理所当然——正如几十年前勒贝格想必也曾越过它。正是在这种“越过”的行动里,在总之做自己、而不只是做那些立法共识的表达式里,在不被它们为我们划定的那个强制性圈子关在里面——首先是在这种孤独的行动里,才有“创造”。其余的一切都是额外来的。

后来,在接纳我的这个数学界里,我有机会遇到许多人,无论是年长者,还是与我年纪相仿的年轻人,他们显然比我聪明得多、也“有天赋”得多。我欣赏他们学新概念时那种轻松,像玩一样,拿来耍弄,仿佛从摇篮里就认识它们——而我却觉得自己笨重迟钝,像一只鼹鼠,在由那些(人们向我保证)我理当学会的东西堆成的一座不成形的山里面艰难地开路,而那些东西的来龙去脉我根本抓不住。事实上,我一点也没有那种出色学生的样子:轻松通过名校会考,三下两下就消化掉那些吓人的教学大纲。

\origpage{P7} 何况,我那些更出色的同学大多成了有本事、有名望的数学家。然而,隔了三十或三十五年回头看,我发现他们并没有在我们这个时代的数学上留下真正深刻的印记。他们在已经现成的环境里做了些事情,有时是漂亮的事情,而那环境他们从没想过要去触动。他们不知不觉地成了那些看不见而专横的圈子的囚徒——那些圈子在某个圈子、某个时代里界定出一个“宇宙”。要跨越它们,他们本应在自己身上重新找到出生时即有的那种能力,那种我也有的能力:独处的能力。

小孩子呢,独处毫无困难。他天生就是孤独的,虽然偶尔有人做伴他也不讨厌,而且知道到了吃奶的钟点就吵着要妈妈的“咪咪”。他也清楚地知道,不用谁来告诉他,那“咪咪”是给他的,而他会吃。可我们常常已经和心里的那个孩子失去了联系。我们不断地从最好的东西旁边走过,不屑于看它一眼……

如果说在《收获与播种》里我除了对自己说话,还对别的什么人说,那绝不是对一个“公众”说。我在那里是对你、对读我的你说话,把你当作一个人,一个孤独的人。我想说话的,是你身上那个懂得独处的人,那个孩子,而不是别人。那孩子常常很远,我很清楚。他什么没见过,而且早就见惯了。他不知躲到哪里去了,要走到他跟前,常常不容易。人家会赌咒说,他早就死了,或者干脆从未存在过——可我确信他就在某处,活得好好的。

我也知道我被听懂的标志是什么。那就是:超越一切文化与命运的差异
% ---------------- p.36 ----------------
,我所说的关于我本人、我的生活的话,在你那里得到回声与共鸣;当你在其中也重新找到你自己的生活、你自己对自己的经验,也许还是你此前未曾留意的一面。这不是什么“认同”,不是去认同与你相距遥远的什么东西或什么人。而也许是,你稍微地,通过我对自身生活的重新发现,也重新发现了你自己的生活,那离你最近的东西——顺着《收获与播种》一页一页,直到我今天此刻正在写下的这些页面。\origpage{P8}

\section{(3) 内在的冒险——或神话与见证}

首先,《收获与播种》是一场对我自己、对我的生活的反思。因此它也是一份见证,而且是从两方面来说的见证。它是对我过去的见证,反思的主要重量正落在这过去上。但同时它也是对最切近的当下的见证——对我此刻正在写作的见证,对《收获与播种》这些页面在一个个钟点、一个个夜晚和白天里诞生的此刻的见证。这些页面是一场关于我的生活的长久沉思的忠实见证,那沉思如它实际所经历的那样延续着(并且此刻仍在延续……)。

这些页面没有文学上的野心。它们是一份关于我自己的文献。我只在极狭窄的限度内才允许自己动它(比如偶尔作些文体上的修饰)\fnote{4}{}{因此,对错误(事实上的,或视角上的,等等)的可能的更正,并不是对初稿加以修补的机会,而是放在页底注里,或在日后对所述情境的“回访”中进行。}。若说有什么野心,那也只是:是真的。而这已经很多了。

此外,这份文献一点也不是“自传”。你在里面不会知道我的出生日期(它除了用来画一张星象图,恐怕没什么意思),也不会知道我母亲和父亲的名字、他们以什么为生,不会知道那位曾是我的妻子的人和其他在我生命中有过重要位置的女人的名字,也不会知道这些爱情中生下的孩子们的名字,以及他们各自把自己的生活过成了什么样。这并不是因为这些事情在我生命里不重要,也不是因为它们今天已不再重要。而是因为,这场对我的反思既已如此开始并延续下去,我从未在某一刻感到被促使去稍稍着手描述这些我在这里那里掠过的事情,更不会去认真地罗列名字和数字。从未有一刻,我会觉得这能给当时我所追索的“话题”增添什么。(而在前面那几页里,我却像不由自主似的,写进了也许比接下来一千页里还要多的、关于我生活的物质细节……)

如果你问我,这千页一路追索的“话题”究竟是什么,我会回答:是讲述、并由此发现我一生所是、且现在仍是的那场内在冒险。这份关于一场冒险的叙述—见证,同时在我刚才说的两个层面上展开。一个是对过去那场冒险的勘探,从它的根与起源一直追溯到我的童年。另一个是这“同一场”冒险的延续与更新,在我写《收获与播种》的一个个瞬间和一天天里,作为对来自外部世界的一声猛烈呼唤\fnote{5}{}{关于这声“猛烈呼唤”的细节,见《信》,特别是第3至8节。}的自发回应。\origpage{P9}

外部的事实为反思提供养料,但只在这样的限度内:它们激起并引发内在冒险的又一次跃起,或有助于照亮它。我的数学作品被埋葬、被洗劫这件事(后面会详谈),就是这样一次挑衅。它在我心里激起了强大的自我反应的全员出动,同时又向我揭示了那些深的、此前未被察觉的纽带,它们仍把我与从我这里生出的作品连在一起。

% ---------------- p.37 ----------------
\origpage{P10} 确实,我属于“数学尖子”这件事,并不必然成为你对我这段特殊冒险感兴趣的理由(更不是好理由)——我改了圈子、改了生活方式之后跟同行们闹出的麻烦,也不是理由。而且并不缺少同行,甚至朋友,他们觉得当众摊开自己的“心境”(用他们的话说)是最可笑不过的。要紧的是“成果”。“灵魂”呢,也就是说我们身上那个经历着这些“成果”的“生产”、以及它的各种余波(既在“生产者”的生活里,也在他的同类的生活里)的东西,却遭人轻视,甚至被公然嘲弄。这种态度自命为“谦虚”的表现;我在其中看到的却是一种逃避的征兆,一种古怪的失序,而且是我们呼吸的空气本身所助长的。我当然不是为那种被对自己的暗暗蔑视所击中的人写作的——那蔑视使他看不上我能给他的最好的东西。一种对真正构成他自己生活之物的蔑视,对构成我生活之物的蔑视:那些驱动着心灵的浅层与深层运动,粗糙的或细微的,那正是经历着经验、并对之作出反应的“灵魂”,它僵住或舒展,它退缩或在学习……

一场内在冒险的叙述,只能由经历它的人来做,别人不行。但即便这叙述只给自己看,它也很少不滑进建造神话的车辙里,而那神话的主角就是叙述者本人。这样的神话,并非诞生于一个民族、一种文化的创造想象,而是诞生于那个不敢承担一份谦卑现实的人的虚荣,他乐于用一件构造物、一件他心智的作品去替换它。但一份真实的叙述(如果有的话),关于一场真正被那样经历过的冒险,是珍贵之物。而这并不因为围绕叙述者的某种威望(不论有理没理),而只因为它存在,带着它的真之品质。这样的见证是珍贵的,无论它来自一位有名望甚至显赫的人物,还是来自一个没有前途、拖家带口的小职员,还是来自一个普通刑事犯。

这样一份叙述若对别人有什么功效,那首先是通过对另一个人经验的不加粉饰的见证,让他重新面对自己。或者(换个说法)也许是抹去他心中(哪怕只在读一次书那么长的时间里)那份蔑视——他借它来看待自己的冒险,以及那个既是乘客又是船长的“灵魂”……

\section{(4) 风俗画}

说到我作为数学家的过去,后来又(像违心地)发现我那场巨大“葬礼”的种种波折与内幕,我不由自主地——并非刻意——画出某个圈子、某个时代的画卷:一个某些价值正在瓦解的时代,那些价值曾赋予人的劳动以意义。这就是“风俗画”的一面,围绕一桩在“科学”编年史上大概是独一无二的“社会新闻”画成。我前面说的话,我想已经足够清楚:你在《收获与播种》里找不到一份关于某桩不寻常“事件”的“卷宗”,好让你很快了解情况。可是我有一位朋友,一直在找那份卷宗,他闭着眼睛走过,什么也没看见,几乎错过了构成《收获与播种》的实质与血肉的一切。

正如我在《信》里详细得多地解释的,“调查”(或“风俗画”)主要在第二部和第四部展开:《葬礼(一)——或中国皇帝的新衣》和《葬礼(三)——或四种运算》。一页一页地,我在那里执拗地把一大批(至少可以说)有汁水的事实一件一件挖出来,又尽量随挖随“安顿”。渐渐地,这些事实聚成一幅总体画面,它逐步从雾中显出
% ---------------- p.38 ----------------
,颜色越来越鲜,轮廓越来越清晰。在这些逐日写下的注里,刚出现的“原始事实”与个人的回忆、与心理的、哲学的、乃至(偶尔)数学的评论和思考,不可分割地混在一起。就是这样,我也无能为力!

以我做过的那项工作(它让我绷紧了一年多)为基础,用“调查结论”的体例编一份卷宗,应该是几个钟头或几天的额外工作,视有兴趣的读者的好奇与要求而定。我确实试过一回,要编那份著名的卷宗。那是我开始写一个本该叫《四种运算》的注的时候\fnote{6}{}{那条原计划的注最终炸开成了《收获与播种》第四部(同名《四种运算》)里的大约70条注,足足铺开四百来页。}。可后来,不行,一点办法也没有。我做不到!这确实不是我的表达方式,到了晚年更是如此。而且我现在认为,有《收获与播种》,我为“数学共同体”做的已经够多了,可以毫无悔意地把编那份该编的“卷宗”的事,留给我以外的人(如果我的同行中有谁觉得自己与此有关)。\origpage{P11}

\section{(5) 继承人与建造者}

现在该在这里说几句我的数学作品了——它在我生命里占据了一个重要的位置,而且(令我自己惊讶)至今仍占着。《收获与播种》里我不止一次回到这部作品——有时用每个人都能清楚明白的方式,有时则用多少带点技术性的词语\fnote{7}{}{除了关于我过去作品的数学概览,这里那里还有一些含新数学发展的段落。最长的一篇是 ReS IV 第171(ix)号注里的《五张照片(晶体与 D-模)》。}。后一类段落,在很大程度上会“越过”不止“门外汉”的头顶,连对其中所谈的数学已经不怎么“在行”的同行数学家也一样。你当然可以干脆跳过那些你觉得太“专业”的段落。你也可以把它们浏览一遍,也许顺带捕捉到数学事物世界那“神秘的美”(这是一位非数学家朋友来信中的说法)的一线反光——它们像一座座“不可企及的奇异小岛”,在反思的浩瀚流动的水面上浮现……

大多数数学家,我刚才说过,倾向于把自己圈在一个概念框架里、一个一劳永逸地定下来的“宇宙”里——基本上就是他们求学时发现已经“现成”的那个。他们像一座宽敞漂亮的、一切都安置好了的房子的继承人:有客厅、厨房、作坊,有成套的锅碗瓢盆和随手可用的工具,凭这些,天哪,足够做饭和做点手工了。这房子在几代人手里是怎样一点点盖起来的,某些工具(而不是别的……)是怎样、为什么被设计和打造出来的,为什么这间屋子这样布置、那间那样安排——这些问题是这些继承人从不会想到去问的。这就是“宇宙”,就是必须生活在其中的“给定之物”,如此而已!它显得很大(而且多数时候,离把它的所有房间都走遍还远着呢),同时又很熟悉,尤其是:不可变易。他们忙活,是为了维护和装点一份家产:修一张摇晃的柜子,给外墙抹灰,磨快一件工具,甚至有时候,对比较有事业心的来说,在作坊里从头做一件新家具。而他们全心投入时,这件家具也可能美极了,整座房子也显得更美了。\origpage{P12}

更罕见的是,他们中有人会想到去改动一件备用的工具,或者甚至在需求的反复催促下,去设想并造出一件新的。这样做的时候
% ---------------- p.39 ----------------
,他几乎要连声道歉,因为他觉得自己这样有点像冒犯了家族传统应有的虔敬,觉得自己用一桩古怪的创新把它冲撞了。

房子里大多数房间的窗和百叶窗都仔细关着——大概是怕有一股别处来的风灌进来。而当漂亮的新家具这里一件、那里一件,还不算子嗣,开始把变得狭窄的房间塞满,一直侵占到走廊时,这些继承人里没有一个愿意意识到:他那熟悉而舒适的宇宙开始有点紧得转不开身了。与其承认这一点,他们各人宁愿勉勉强强地侧身挤过去、卡在那里:有人挤在路易十五式碗柜和藤编摇椅之间,有人挤在流鼻涕的娃娃和埃及石棺之间,最后还有一位,在无计可施之下,尽其所能地爬上一堆杂七杂八、摇摇欲坠的椅子和长凳……

我刚画的这幅小画,并非数学界所特有。它说明的是一些根深蒂固、由来已久的制约,在所有圈子和人类活动的一切领域里都能遇到,而且(据我所知)在所有社会、所有时代都是如此。我此前已有机会提及,我也绝不声称自己可以免于它。我的见证将表明,实情恰恰相反。只不过,在智力创造活动这个相对有限的层面上,我受这种制约的触动相当少\fnote{8}{}{我认为主要原因在于某种有利的氛围,它环绕着我的童年,直到我五岁。参见《无辜》一注(ReS III,第107号)。}——它可以叫做“文化盲”:看不见(也无法活动于)周围文化所固定的那个“宇宙”之外。

至于我,我觉得自己属于那样一脉数学家:他们自发的天职与快乐,是不断地建造新房子\fnote{9}{}{这个“待建之屋”的原型意象,第一次浮现并被表述出来,是在《阴,仆役,与新主子们》一注中(ReS III,第135号)。}。一路之上\origpage{P13},他们也忍不住随时发明并打造一切必需的工具、器具、家具和仪器,既为了从地基到屋脊把房子建起来,也为了给未来的厨房和作坊备足东西,把房子布置得能住、能舒服。然而,一旦连最后一条檐沟、最后一只凳子都安放好了,这位工匠却很少在这些地方久留——这里的每一块石头、每一根椽子,都带着加工并安放它的那只手的痕迹。他的位置不在那些现成宇宙的安宁里,无论它们多么宜人、多么和谐——无论是他亲手布置的,还是他的前人们布置的。别的任务已经在召唤他去新的工地:在某种需要的专横推动下,那需要也许只有他一个人清楚地感到,或者(更常见的是)他预感到某种只有他一个人预感到的需要。他的位置在露天里。他是风的朋友,不怕独自干活,一干几个月、几年,必要时干一辈子,除非有可喜的接替者来搭把手。他当然和所有人一样只有两只手——但那是两只随时能猜到自己该做什么的手,既不嫌最粗笨的活,也不嫌最精细的活,而且从不厌倦去认识、反复认识那些无数次召唤它们去认识的事物。两只手也许太少,因为世界是无穷的。它们永远也穷尽不了它!然而,两只手又很多……

我这个人历史不行,但若要我举出这一脉数学家的名字,我自然而然想到伽罗瓦和黎曼(上世纪),以及希尔伯特(本世纪初)。若要在当初接纳我进入数学界的年长者中找一位代表\fnote{10}{}{我在《受欢迎的异乡人》一节(ReS I,第9号)里谈到这些开端。},我最先想到的是让·勒雷(Jean Leray)的名字,尽管我和他的接触
% ---------------- p.40 ----------------
始终极为零星\fnote{11}{}{不过这并不妨碍我(继 H. 嘉当与 J.-P. 塞尔之后)成为勒雷引入的某个重大革新概念——层——的主要使用者和推广者之一,而层贯穿我整个几何学家的工作,是其中一件本质的工具。也正是它给了我钥匙,把(拓扑)空间概念扩展为 topos 概念,后面会谈到。\par 此外,在我看来,勒雷与我勾画的“建造者”肖像有所不同:他似乎并不倾向于“从地基到屋脊建造房子”。不如说,他忍不住在无人想得到的地方开出了广阔的地基,却把完成它们、在上面建造、以及房子建成后住进去(哪怕只住一阵)的事留给别人……}。

我刚才大致勾画了两幅肖像:一位是“恋家”的数学家,满足于维护和装点一份遗产;一位是建造者—拓荒者\fnote{12}{}{我刚刚偷偷地、“从侧面”把两个带有雄性回响的修饰语(即“建造者”和“拓荒者”)并置在一起,而它们表达的却是发现冲动相当不同的方面,其本性比这些名字所能唤起的更精细。这将在这次巡游—反思的后文、“发现母亲——或两个山坡”(第17号)那一步里显现。},他忍不住不断越过那些界定一个宇宙的“看不见而专横的圈子”\fnote{13}{}{而且,他也就这样在无意中给这个旧宇宙(不是为他自己,至少是为那些不如他好动的同类)指定了新的界限,划在新的圈子里——这些圈子当然更广,但和它们所取代的那些一样看不见、一样专横。}。也可以用两个有点一刀切、但颇有意味的名字称他们:“守成者”与“革新者”。两者在同一个延续于世代、世纪、千年的集体冒险中,都有存在的理由和要扮演的角色。在一门科学或一门艺术的繁盛期,这两种气质之间没有对立,也没有敌意\fnote{14}{}{尤其在数学界,在我亲身见证的时期(1948–1969),情况就是如此,当时我自己也是这个世界的一部分。1970年我离开之后,似乎出现了一种大范围的反应,一种对“想法”总体的“蔑视的共识”,尤其针对我引入的那些重大革新想法。}。它们不同,而且互相补充,就像面团与酵母互相补充。\origpage{P14}

在这两个极端类型之间(它们在本性上丝毫不相反对),当然有一整套中间气质。某位“恋家者”从不想离开一处熟悉的住所,更不会去揽下到天知道什么地方另建一座的活计;然而当住处真的开始变窄时,他会毫不犹豫地拿起抹刀,收拾一间地窖或阁楼,加高一层,必要时还在墙边接出一间规模不大的新附属房\fnote{15}{}{我的大多数“前辈”(例如《一笔情愿的债》导论第10节里谈到的那些)都属于这种中间气质。我特别想到亨利·嘉当、克劳德·谢瓦莱、安德烈·韦伊、让–皮埃尔·塞尔、洛朗·施瓦茨。除也许韦伊之外,他们对我眼看要登上的那些孤独冒险,都投以“善意的目光”,没有“暗暗的担忧或责备”。}。他骨子里不是建造者,但常常会带着善意的眼光——至少是不带暗暗的担忧与责备——看着另一位曾与他同住一座房子的人,如今
% ---------------- p.41 ----------------
在某个穷乡僻壤辛苦地聚拢梁木和石头,一副仿佛已经看见宫殿的神气……\origpage{P15}

\section{(6) 视角与眼光}

但还是回到我自己和我的作品上来。

如果说我在数学家的技艺上有所出众,那与其说靠解题的灵巧与坚持(解前人传下的问题),不如说靠我身上这种天然的倾向:它驱使我看见一些显然关键、却无人看见的问题,或者提取出那些缺失的“好概念”(往往在新概念出现之前,谁也没意识到它们缺失),以及那些谁也没想到的“好陈述”。常常,概念与陈述配合得如此完美,使我心里不可能怀疑它们是对的(顶多需要些修补)——于是往往,当这不是要发表的“部件活”时,我就省得再往前走,省得花时间去把证明弄妥;而那证明,一旦陈述和它的语境看清了,多半只是个“手艺”问题,甚至可以说是例行公事。值得注意的事情数不清,不可能把每一个的召唤都跟到底!尽管如此,在我已写出并发表的作品里,正式证明过的命题和定理数以千计,我想我可以说,除了极少的例外,它们全都进入了共同的遗产,被普遍当作“已知”,在数学各处被日常使用。

但我的独特天赋把我带向的,比发现新问题、新概念、新陈述更进一步的,是发现富有成果的视角——它不断引导我引入并或多或少地发展一些全新的主题。我觉得,这就是我对我那个时代的数学所贡献的最本质的东西。说实话,我刚才说的那无数问题、概念、陈述,在我看来只有在这样一种“视角”的光照下才有意义——更确切地说,它们是从中自发地、带着显然性之力生出来的;就像黑夜里亮起的一道光(哪怕是散漫的光),仿佛从虚无中生出它突然向我们揭示的那些或模糊或清晰的轮廓。没有这道把它们束成一束的光,那十个、一百个、一千个问题、概念、陈述,就会显得像一堆杂七杂八、没有形状的“心智小玩意儿”,彼此孤立——而不是一个整体的各个部分;那整体也许仍不可见,仍躲在夜的褶皱里,却已被清楚地预感到。

富有成果的视角,就是那样一种视角:它把那些无人感到的灼热问题,揭示为同一个包容它们、赋予它们意义的整体的活的各部分;并且(也许是对这些问题的回答)把那些如此自然、却无人想到去提取的概念,以及那些仿佛顺理成章、却在激发它们的问题和使之得以表述的概念尚未出现之前谁也不敢提出的陈述,一并这样揭示出来。在我们这门技艺里\fnote{16}{}{毫无疑问,情况不只在我们这门“技艺”里如此,而且(在我看来)在一切发现的工作里都是如此,至少当这工作处于智识认识的层面时。},比所谓“关键定理”更有力的发现工具,是富有成果的视角——更确切地说,它们不是工具,而是研究者本人的眼睛,那双热切地想要认识数学事物之本性的眼睛。

这样,富有成果的视角不外是那只“眼睛”:它同时让我们发现,又让我们在被发现者的多样性中认出统一。而这统一,真正是把这些多样之物联系起来并赋予生命的生活本身与气息。

但正如其名所示,一个“视角”本身仍是局部的。它向我们揭示一处风景或全景的一个侧面,而此外还有许多同样有效、同样“真实”的侧面。只有当对同一现实的互补视角彼此结合起来,当我们的“眼睛”增多,目光才更深入地进入对事物的认识。我们想认识的现实越丰富、越复杂,拥有好几只“眼睛”\fnote{17}{}{每个视角都会带来发展出一套表达它、并专属于它的语言。拥有好几只“眼睛”或好几个“视角”来把握一个情境,也等于(至少在数学里)拥有好几种不同的语言来圈定它。}以便在它的全部广度与全部精微中把握它,也就越重要。

而有时候,一束汇聚在同一片广阔风景上的视角,凭着
% ---------------- p.42 ----------------
在我们身上那种能在多中把握一的能力,使一件新事物有了身体;这件新事物超出每一个局部透视,正如一个活物超出它的各个肢体和器官。这件新事物,可以叫做一种眼光。眼光把体现它的那些已知视角统一起来,并向我们揭示另一些此前被忽略的视角,正如富有成果的视角让人发现并把握众多新问题、新概念、新陈述,把它们当作同一整体的部分。\origpage{P16}

换个说法:眼光对于它仿佛从中生出、并由它统一起来的那些视角,就像白昼明亮温暖的光对于太阳光谱的不同成分。一种广阔而深沉的眼光,像一处取之不竭的源泉,注定要去启发和照亮的工作,不仅是在某一天生出它、并使自己成为它的仆人的那个人的工作,也是世世代代人的工作——他们也许(像他本人一样)被它让我们隐约看见的那些遥远边界所迷住……\origpage{P17}

\section{(7) “大想法”——或树木与森林}

我数学活动所谓“多产”的时期,也就是有正式发表物作证的时期,在1950到1969年之间,共二十年。而在二十五年里,从1945年(我十七岁)到1969年(我将满四十二岁),我几乎把全部精力投进了数学研究。这投入无疑过分。我为此付出的是长期的精神停滞、逐步的“变厚”,在《收获与播种》的篇幅里我不止一次有机会谈到。然而,在一项纯智力活动的有限领域内,凭一种局限于数学事物世界的眼光的绽放与成熟,那些年仍是创造力旺盛的年份。

在我生命的这段漫长时期,几乎全部时间和精力都用在所谓“部件活”上:那种细致的打造、装配和磨合工作,是从头建造一座座房子所必需的,而一个内心的声音(或一个魔鬼……)命令我建造它们,并按工作进展随时向我口授一位总施工的方案。我被“手艺”的活计占住:时而当石匠,时而当砌工、木匠,甚至当铅管工、细木工和家具匠——我很少偷闲把那份除我之外无人可见的总图(后来才显出它的样子……)用白纸黑字记下来,那总图在日、月、年之中以梦游者般的稳当引导着我的手\fnote{18}{}{“梦游者”这个意象,是柯斯特勒(Koestler)那本出色的书《梦游者》的标题给我的(Calmann-Lévy 版),那本书是一篇“关于宇宙观念史的研究”,从科学思想的开端直到牛顿。这段历史中令柯斯特勒深受触动、并由他特别指出的一点是:我们的世界认识常常从某一点走到另一点,那另一点(从逻辑上、事后看来)显得近在咫尺,而道路却要经过最离奇的迂回,仿佛在挑战健全的理性;然而,尽管有这千百道似乎注定要让人永远迷路的迂回,出发去寻找宇宙之“钥匙”的人们,却带着“梦游者般的稳当”,不由自主地、往往甚至连自己也没有察觉,碰上了另一些他们远未预料到的“钥匙”,而那些竟然正是“对的”。\par 据我在周围所能观察到的,在数学发现的层面上,这些发现之路上的惊人迂回,是某些大格局研究者的事,但绝不是所有研究者的事。这也许是因为,两三个世纪以来,自然科学的研究,尤其是数学研究,已从与某一文化、某一时代相联系的宗教或形而上学强制前提中解放出来——那些前提曾是对宇宙的“科学”理解之展开(无论好坏)特别强大的制动。然而确实,数学中某些最根本、最显然的观念与概念(如位移、群、零这个数、文字运算、空间中一点的坐标、集合的概念,或拓扑“形式”的概念,更不必说负数与复数),都经过几千年才出现。这些都是那种根深蒂固、深深植入心灵之中的“阻塞”的明白迹象——它抵制全新的观念,即便这些观念天真地简单,似乎以显然之力自行强加于人,而且几代、乃至几千年都是如此……\par{}回到我自己的工作,我觉得其中的“失手”(也许比多数同行更多)只限于细节,通常很快就由我自己发现。那不过是纯属“局部”的“途中意外”,并不严重影响关于所考察情境的基本直觉的可靠性。相反,在观念与重大指导性直觉的层面上,我的作品似乎没有任何“失误”,尽管这听起来难以置信。每一刻,即便不能把握一条进路的最终归宿(它多半仍藏在目光之外),也至少能把握那些最富成果、引我直达本质事物的方向——正是这种从未失灵的可靠性,使柯斯特勒的“梦游者”形象重新浮上我的心头。}。应该说,我乐于在“部件活”里倾注一份充满爱意的细致,这活并不令我讨厌。此外,前辈们所宣讲和实行的数学表达方式,(至少可以说)给了工作的技术面以优先,而不大鼓励那些停留在“动机”上的“离题话”;甚至对那些仿佛要从雾里唤出某个也许富有启发性的形象或眼光的离题话也不鼓励——因为它们还没有化身为木头的、石头的或纯硬水泥的、可触摸的建筑,更像是梦的碎片,而不是工匠专注而尽责的活计。\origpage{P18}

从量上说,我在这些高强度多产的年份里的工作,具体表现为大约一万二千页的发表物——论文、专著或讨论班讲义\fnote{19}{}{从1960年代起,其中一部分发表物是与同行(尤其是 J. 迪厄多内)和学生合作写成的。}——以及数以百计(如果不是数以千计)的新概念,它们连同我提取它们时给的名字一起,进入了共同的遗产\fnote{20}{}{这些概念中最重要的那些,《主题概述》以及随附的《历史评注》都作了回顾,两者将收入《反思》第4卷。有些名字是朋友或学生向我建议的,比如“光滑态射”这个说法(J. 迪厄多内),或“site、champ、gerbe、lien”这一整套,那是让·吉罗(Jean Giraud)的博士论文里发展出来的。}。在数学史上,我想我大概是在我们这门科学里引入新概念最多的人,同时也正因如此,不得不发明最多新名字的人——为了把这些概念表达得细致,并尽我所能地富有提示性。\origpage{P19}

这些纯“量”的指标,当然只给出对我作品的一种极其粗略的把握,错过了真正构成其灵魂、生命与力量的东西。正如我刚才写的,我在数学里贡献的最好的东西,是我先得以窥见、随后耐心提取并或多或少加以发展的新“视角”。像我刚谈到的那些概念一样,这些新视角进入极其多样的情境,它们本身也几乎数不清。

然而有些视角比另一些更广阔,单凭它们就引发并包揽大量局部视角,涉及大量各不相同的具体情境。这样一个视角,也可以恰当地称为一个“大想法”。凭它自身的丰产性,这样一个想法生出成群的后代,这些想法都继承它的丰产性,但大多数(如果不是全部)比母想法范围更小。

至于表达一个大想法,也就是“说出”它,这往往是几乎和它被构想出来、和在构想者心中缓慢孕育一样精细的事——或者更确切地说,这费力的孕育与成形工作,恰恰就是“表达”该想法的那项工作:耐心地、日复一日地把它从诞生时环绕它的雾幕中提取出来,逐渐给它一个可触的形态,成为一幅在星期、月份和年份中不断丰富、稳固、精细起来的图画\origpage{P20}。仅仅给这想法命名,用某个醒目的说法,或用多少带点技术性的关键词,也许只要几行、甚至几页——但很少有人在尚未相当熟悉它的情况下,能听懂这个“名字”,并在其中认出一张面孔。而当这想法到了完全成熟的时候,也许一百页就足以表达它,让孕育它的那位工匠完全满意——也可能一万页经过长期打磨和掂量的文字仍不够\fnote{21}{}{1970年离开数学舞台时,我关于概形这一中心主题的全部发表物(其中不少是合作),加起来该有一万来页。然而这在我眼前看到的那个关于概形的宏大纲领里,只占不大的一部分。这纲领在我离开后便被无限期搁置了,尽管此前已经发展并发表、可供所有人使用的东西,除极少数例外,一开头就进入了共同的遗产,成为大家日常使用、视为“人所共知”的概念与结果。\par 我离开时已完成的那部分纲领,关于概形主题及其延伸与分枝,单是它本身就代表了数学史上所完成的最庞大的基础工作,而且肯定也是科学史上最庞大的之一。}。

而无论在哪种情形下,在那些为了把想法据为己有而了解了这项终于呈现出想法的蓬勃生长的工作的人当中——它像一片在林间荒地上长起来的宽阔乔木林——很有理由打赌:许多人会清楚看见所有这些挺拔有力的树,并且各有所用(有人为了爬上去,有人为了取梁木和木板,还有人为了在自己壁炉里烧旺火……),但很少有人懂得看见森林……

\section{(8) 眼光——或十二个主题的和谐}

也许可以说,“大想法”就是那样一种视角:它不仅显得新而丰产,而且把一个新的、广阔的主题引入科学,而这主题正是它的化身。而任何一门科学,当我们不把它当作权力与支配的工具,而当作我们这个物种跨越时代的认识冒险来听时,它不外就是这种和谐——在不同时代或广或狭、或贫或富——它在世代与世纪中展开,靠的是所有先后出现的主题之间那种精细的对位,那些主题像被从虚无中唤来,加入其中并在其中交织。\origpage{P21}

在我于数学中提取出的众多新视角里,回头看去,有十二个我会称之为“大想法”\fnote{22}{}{以下是这份十二个主导想法、即我作品的“主主题”的清单(按出现的时间顺序),供好奇的数学读者一阅。\par 1. 拓扑张量积与核空间。\par 2. “连续”与“离散”的对偶(导出范畴、“六种运算”)。\par 3. 黎曼–罗赫–格罗滕迪克瑜伽($K$-理论,与相交理论的关系)。\par 4. 概形。\par 5. Topos。\par 6. 平展上同调与 $\ell$-进上同调。\par 7. Motif 与 Motivic 伽罗瓦群(Grothendieck 的 $\otimes$-范畴)。\par 8. 晶体与晶体上同调,“德拉姆系数”“霍奇系数”的瑜伽……\par 9. “拓扑代数”:$\infty$-栈、derivator;topos 的上同调形式主义,作为新的同伦代数的灵感。\par 10. 温和拓扑学。\par 11. 远阿贝尔代数几何的瑜伽,伽罗瓦–泰希米勒理论。\par 12. 正多面体与各种正则构型的“概形”或“算术”视角。\par 除这些主题中的第一个(其中重要的一支属于我1953年的博士论文,并在1950至1955年间我的泛函分析时期得到发展)之外,其余十一个都是在我1955年以后的几何学家时期提取出来的。}。看我的数学家作品,“感受”它,至少是稍稍看见并“感受”其中某些想法,以及它们引入的那些构成作品的经纬与灵魂的大主题。

% ---------------- p.45 ----------------
事所必然,这些想法中有些比另一些“更大”(而那些也就因此“更小”!)。换句话说,在这些新主题里,有些比另一些更广,有些更深地潜入数学事物之谜的核心\fnote{23}{}{在这些主题中,就意义之广而言,我觉得最广的是 topos,它提供了把代数几何、拓扑学与算术综合起来的想法。就已产生的发展之规模而言,最广的是概形主题。(关于这一点,见第20页页底注(*)。)正是它为所论主题中的另外八个(即除第1、5、10之外的所有主题)提供了“典范的”框架,同时为代数几何以及代数—几何语言的自下而上的更新提供了中心概念。\par 在另一端,十二个主题中的第一个和最后一个,在我看来比其余的规模更有限。不过,就最后一个而言,它把一种新视野引入正多面体与正则构型这个极古老的主题,我怀疑一个全身心投入其中的数学家,一生也不足以把它穷尽。至于所有这些主题中的第一个,即拓扑张量积,它扮演的更多是一件现成新工具的角色,而不是后人发展的灵感来源。不过,直到最近这些年,我仍偶尔收到多少较近的工作的零星回响,那些工作(二三十年后)解决了我留下的某些问题。\par 在我看来,这十二个主题中最深的是 Motif,以及与之紧密相连的远阿贝尔代数几何和伽罗瓦–泰希米勒瑜伽。\par 从工具经过我亲手调校、臻于完善、并在近二十年研究的各个“前沿地带”日常使用的力量看,我觉得最值得注目的是“概形”与“平展上同调、$\ell$-进上同调”这两块。对一个消息灵通的数学家来说,我想从现在起几乎不可能怀疑:概形工具,以及由它产生的 $\ell$-进上同调工具,属于本世纪少数几项重大成就之列,它们在最近几代人之中滋养并更新了我们的科学。}。有三个(在我看来并非最不重要的)只是在我离开数学界之后才出现的,至今仍处胚胎状态;“正式地说”,它们甚至并不存在,因为没有任何正式发表物可以充当它们的出生证明\fnote{24}{}{这三者得到哪怕最简略勾勒的唯一“半正式”文本,是1984年1月为申请 CNRS 脱产而写的《纲领概要》。这一文本(导论第3节“罗盘与行囊”里也谈到它)原则上将收入《反思》第4卷。}。在我离开之前出现的九个主题里,最后三个我留在蓬勃生长之中,今天仍处在童年状态,因为(我走后)没有充满爱的手去照料这些“孤儿”的所需,它们在一个敌对的世界里被弃之不顾\fnote{25}{}{在我不久离开之后,这三个孤儿便被无声无息地埋葬了;其中两个后来又被大吹大擂地掘出,却不提工匠的名字——一个在1981年,另一个(鉴于此举毫无瑕疵的成功)就在次年。}。至于另外六个主题,在我离开前的二十年里已完全成熟,可以说(有一两点保留\fnote{26}{}{这“除了极少例外”,主要指格罗滕迪克式的对偶瑜伽(导出范畴与六种运算),以及 topos 的瑜伽。这一点将在《收获与播种》第二部和第四部(《葬礼(一)》与《葬礼(三)》)里详加讨论。})它们那时已经进入共同的遗产:尤其在几何人中间,今天“人人”都在整天、随时吟诵它们,甚至不再知道自己是在吟诵(就像茹尔丹先生说着散文)。它们是我们呼吸的空气的一部分,当我们“做几何”,或当我们做多少带点“几何性”的算术、代数或分析时。\origpage{P22}

我作品这十二个大主题绝不是彼此孤立的。在我眼里,它们属于一种精神与话题的统一,这统一像一个共同而持久的低音,贯穿我全部“写出的”和“未写出的”作品。写这些行的时候,我仿佛又听见了同一个低音——像一个召唤!——贯穿那三年“无偿的”、顽强的、孤独的工作,那时候我还不曾操心世上除我之外是否还有别的数学家,因为我当时完全被那召唤我之物的魅力攫住了……\origpage{P23}

这种统一,并不只是同一个工匠在他手下出来的作品上留下的印记。这些主题由无数既精细又显然的纽带彼此相连,就像那些各各清晰可辨的不同主题之间的相连:它们
% ---------------- p.46 ----------------
在同一片广阔的对位中展开、交缠——在一个把它们聚合、推向前去、并赋予每一个以意义的和谐之中,每一个的意义、运动与圆满都有所有其他者参与。每一个局部主题似乎都从这更广的和谐中生出,又随时从中重生,远不是这和谐显得像那些先于它存在的构成性主题的“总和”或“结果”。说实话,我无法摆脱这种感觉(大概很可笑……):在某种意义上,正是这尚未出现、却无疑早已确确实实“存在”于尚未诞生之物的幽暗怀抱中某处的和谐——正是它先后唤起了这些只有通过它才获得全部意义的主题,也正是它在我刚出少年、那些炽烈孤独的年月里,已经用低而急迫的声音呼唤我……

无论如何,我作品的这十二个主主题确实都像出于一种秘密的预定,共同汇成同一部交响曲——或者换一个比喻,它们体现为同样多的不同“视角”,全都汇入同一个广阔的眼光。

这眼光只是到1957、58年前后才开始从雾中浮现、显出可辨认的轮廓——那是紧张孕育的年份\fnote{27}{}{1957年正是我被引向提取“黎曼–罗赫”主题(格罗滕迪克版)的那一年——这个主题一夜之间把我捧成了“大明星”。这一年也是我母亲去世之年,因此也是我生命中一道重要分界线之年。那是我一生中创造力最旺盛的年份之一,而且不限于数学层面。那时,全部精力投入数学工作已有十二年。就在那一年,一种感觉浮现出来:关于什么是数学工作,我大致已经“转完了一圈”,也许现在是时候投入别的事情了。那显然是一种内在更新的需要,一生中第一次浮出水面。那时我想过要当作家,并且有好几个月停止了全部数学活动。最后我决定,至少还要把手上已有的数学工作写成白纸黑字,也就是几个月的事,顶多一年……\par 无疑,那一跃的时机还没有成熟。总之,一旦重新拾起数学工作,便是它重新攫住了我。此后又十二年,它再没有放开我!\par 这段间歇之后的一年(1958),也许是我数学生涯中最丰产的一年。新几何的两个中心主题正是在这一年绽放:概形理论的强力启动(这是我当年夏天在爱丁堡国际数学家大会上报告的主题),以及“site”概念的出现——它是关键的 topos 概念的技术性临时版本。隔了将近三十年回头看,我现在可以说:新几何的眼光真正诞生,就在这一年;它紧随这门几何的两件主要工具之后:概形(它代表旧“代数簇”概念的一次变形),以及 topos(它代表空间概念的一次更深层的变形)。}。也许有点奇怪:这眼光对我如此切近,如此“显然”,以至于直到一年前\fnote{28}{}{我头一次想到给这眼光一个名字,是在1984年12月4日的反思里,在《阴,仆役(二)——或慷慨》一注(ReS III,第637页)的第1361号子注中。},我都没想过给它一个名字。(而我的一大癖好,恰恰是不停地给向我显现的事物命名,作为把握它们的第一步……)确实,我说不出某一个特定时刻,可以被经历为这眼光出现的时刻,或事后可以认作那样的时刻。一种新眼光是那样广阔,它的出现无疑不能定在某一个特定时刻,而必须在漫长的年月里(如果不是在世代之中)逐步渗入并占有那个或那些细看与凝视它的人;仿佛在那些注定要被逐步取代的熟悉的眼睛后面,必须辛苦地长出新眼睛来。而这眼光也太大,谈不上像把握路上拐角处出现的第一个概念那样去“把握”它。所以,最终不必惊讶:给这样广阔、又如此切近而弥散的东西命名的念头,只是事后、只是当这东西达到完全成熟之后才出现。\origpage{P24}

说实话,直到两年前,我与数学的关系(除了教书的任务之外)仅限于做数学——跟从一股不断把我向前拉的冲动,进入一个不断吸引我的“未知”。我从未想过在这股冲劲中停下来,哪怕只停一瞬\origpage{P25}
% ---------------- p.47 ----------------
,回头看看也许已画出的走过之路,甚至为一部已过去的作品定位。(无论是把它定位在我的生命里,当作一件仍有深的、长期未被察觉的纽带把我与它相连的东西;还是把它定位在“数学”这场集体冒险里。)

还有一件怪事:要让我终于“放下”并重新认识这部半被遗忘的作品,或者仅仅想到给那曾是它灵魂的眼光一个名字,竟需要我突然面对一场规模巨大的“葬礼”的现实:那场以沉默和嘲弄进行的葬礼,埋葬的是那眼光,以及生出那眼光的那位工匠……

\section{(9) 形式与结构——或事物的道路}

不知不觉,这篇“前言”一针引一线,最后成了对我作品的一次正式介绍,主要是给非数学读者看的。已经陷得太深,退不回去了,我只剩下把“介绍”讲完!我想尽力至少说几句,谈谈那些我在前面篇幅里晃动过的奇妙“大想法”(或“主主题”)的实质,以及那道著名的“眼光”的本性——那些主导想法据说要汇入其中。由于无法求助于多少带点技术性的语言,我大概只能传递一个极其模糊的形象(如果真有什么东西能“传过去”的话……)\fnote{29}{}{这个形象不得不“模糊”,丝毫也不妨碍它是忠实的,也不妨碍它确实再现了所看之物的本质的某种东西(在此即我的作品)。反过来,一个形象无论多么清晰,也可能完全走样,而且只包含次要的东西,完全错过本质。所以,如果你“接住”了我关于自己作品所要说的东西(那肯定我心中形象的一部分会真的“传过去”),你就可以得意地说,你对我的作品中最本质的东西,把握得比任何一位博学的同行也许都更好!}。

传统上,人们区分宇宙之物的三类“性质”或“方面”作为数学思考的对象:数\fnote{30}{}{这里指的是所谓“自然数”0、1、2、3等等,或者(至多)指那些借它们、通过初等性质的操作表达出来的数(如分数)。这些数不像“实数”那样,可以用来度量一个可连续变化的量,比如直线上、平面上或空间中两个变动点之间的距离。}、量、形。也可以把它们叫做事物的“算术”方面、“度量”(或“分析”)方面、“几何”方面。在数学所研究的大多数情境里,这三个方面同时在场并紧密互动。然而,最常见的是其中一个明显占主导。在我看来,对大多数数学家来说(对认识他们或了解他们作品的人来说),他们的基本气质相当清楚:是“算术家”、“分析家”还是“几何家”——即便他们本领很多,在各种音区和调门里都工作过。\origpage{P26}

我最初那些孤独的思考,关于测度与积分理论,毫无歧义地属于“量”或“分析”一类。我引入数学的第一个新主题也是如此(它在我看来比另外十一个范围要小)。我之从“分析”这个“侧门”进入数学,在我看来并非由于我的特殊气质,而是由于可以叫做“偶然境遇”的东西:对于我这颗爱好普遍性与严格性的头脑来说,中学和大学向我提供的教学里最大的空缺,恰好在事物的“度量”或“分析”方面。

1955年是我数学工作的一个关键转折:从“分析”
% ---------------- p.48 ----------------
转向“几何”。我还记得那种震撼的印象(当然完全是主观的):仿佛我离开了干瘦而硌人的草原,忽然来到某种“应许之地”,那里丰饶的财富无穷无尽地增殖,凡是手愿意落下去采摘或挖掘的地方都是如此……而这种压倒性的丰饶印象,超出了任何尺度\fnote{31}{}{我用“压倒性的,超出任何尺度”这个词语组合,勉强译出德文的“überwältigend”及其英文对应词“overwhelming”。上一句里那个(并不恰当的)“震撼的印象”,也应当带着这层意味来理解:当面对一种非凡的壮丽、伟大或美时,我们心中激起的印象与情感突然把我们淹没,以致任何表达我们所感的意愿都仿佛预先被摧毁了。},在此后这些年里只得到确认和加深,直到今天。\origpage{P27}

这就是说,如果在数学里有一件东西(大概从来如此)比任何别的更令我着迷,那不是“数”,也不是“量”,而始终是形。而在形为了向我们显现而选取的千零一张面孔中,比任何别的更令我着迷、并继续令我着迷的,是数学事物中隐藏的结构。

一件事物的结构绝不是我们可以“发明”的东西。我们只能耐心地把它揭示出来,谦卑地去认识它,“发现”它。这项工作里若有创造性,我们若有时做着铁匠或不倦的建造者的活,那绝不是为了“塑造”或“建造”什么“结构”。那些结构根本不曾等我们才存在,才恰好是它们所是的样子!而是为了尽可能忠实地表达我们正在发现和探测的这些事物,以及这个不肯轻易交付的结构——我们摸索着,用也许还结结巴巴的语言,试着把它圈定。于是我们不断地“发明”适合的语言,以越来越精细地表达数学事物的内在结构,并借助这语言,“建造”那些理应说明已被把握和看见之物的“理论”,一点一点、从头到尾。这里有一个持续不断的往返运动:在把握事物与表达被把握者之间,那语言在工作中、在当下需要的不断压力下变得精细并重新被创造。

读者大概已经猜到,这些“从头建造”的“理论”,也正是前面谈到的那些“漂亮的房子”:我们从先人那里继承的房子,以及我们在事物的召唤与倾听中亲手建造的房子。而我刚才说到建造者或铁匠的“创造性”(或想象力),还得补充一句:构成其灵魂与秘密神经的,绝不是那个说“我要这个,不要那个!”、并乐于凭自己心意决断的人的傲气;那就像一个糟糕的建筑师,还没看见、感受过一块地,还没探明它的可能与要求,脑子里就已经准备好了图纸。研究者的创造性与想象力的品质,在于他注意力的品质,在于他倾听事物之声。因为宇宙之物从不厌倦谈论自己、向那个有心去听的人显现。而最美的房子,显出工匠之爱的那一座,并不是比别人更大或更高的那座。美的房子,是忠实映照事物隐藏的结构与美的那座。\origpage{P28}

\section{(10) 新几何——或数与量的联姻}

可我又跑题了——我本打算谈主主题,它们在同一位母眼光中结合起来,像许多条河流回归它们为其子嗣的大海……

% ---------------- p.49 ----------------
这广阔的统一性眼光可以描述为一种新几何。据说,这就是上个世纪克罗内克所梦想的那一种\fnote{32}{}{关于这个“克罗内克的梦”,我只是听人说起:有人(也许正是约翰·泰特)对我说,我正在实现那个梦。在我从前辈那里接受的教育里,历史性的参考极少见;哺育我的并不是阅读多少古代的、甚至同时代的作者,而主要是与其他数学家——首先是我的前辈们——的口头或通信交流。1958年概形理论突然而强劲的启动,其主要的、也许是唯一的外来灵感,是塞尔那篇以 FAC(《代数凝聚层》)这个缩写广为人知的文章,它发表于此前几年。除此之外,理论后来发展中的主要灵感都是从自身产生的,并在年复一年中自我更新,靠的只是内部简洁性与一致性的要求,以及一种努力:在新语境中说明代数几何里“人所共知”的东西(那些东西在我手中变形的同时被我逐步消化),以及这个“已知”让我预感到的东西。}。但现实(一个大胆的梦有时让人预感到或窥见它,并鼓励我们去发现它……)每次都在丰富与共鸣上超过最大胆、最深的梦。可以肯定,对这新几何的不止一个部分(如果不是全部),在它出现的前一天,没有人会想到——工匠本人也不比别人更早想到。

可以说,“数”适合把握“不连续”或“离散”的聚集体的结构:那些往往有限的系统,由可以说是彼此孤立的“元素”或“对象”组成,没有某种从一个到另一个的“连续过渡”原则。“量”则相反,是典型地能够“连续变化”的性质;因此它适合把握连续的结构与现象:运动、空间、各种“簇”、力场等等。于是,算术显得(大致上)是离散结构的科学,分析则是连续结构的科学。

至于几何,可以说,它以现代意义上的科学的形式存在了两千多年,一直“骑”在这两类结构上,即“离散的”与“连续的”\fnote{33}{}{说实话,传统上处在几何家注意力中心的是“连续”方面,而“离散”性质,尤其是数值的和组合的性质,则被略而不谈,或草草打发。十年前我惊叹地发现二十面体组合理论的丰富,而这个主题在克莱因论二十面体那本经典著作里连碰都没有碰到(大概根本没看到)。我看到几何家对自发进入几何的离散结构这种(长达两千年的)忽视的另一个显著迹象:群(尤其对称群)的概念直到上世纪才出现,而且它最初还是(由埃瓦里斯特·伽罗瓦)在一个当时并不被视为属于“几何”的语境中引入的。确实,直到今天,仍有许多代数学家始终不明白:伽罗瓦理论本质上正是一种“几何”眼光,它更新了我们对所谓“算术”现象的理解……}。而且长期以来,两种本属不同种类的几何——一种离散、一种连续——之间并没有真正“离婚”。不如说,那是对同样的几何图形作研究时的两种不同视角:一个强调“离散”性质(特别是数值的和组合的性质),另一个强调“连续”性质(比如在周围空间中的位置,或以点的相互距离来度量的“量”等等)。\origpage{P29}

直到上世纪末,才出现了一场离婚,随着人们有时称为“抽象(代数)几何”的东西的出现与发展。大致说来,它在于为每个素数 $p$ 引入一种“特征 $p$ 的(代数)几何”,照着前几个世纪传下来的(代数)几何的(连续)模型来仿制,然而所处的语境却显得不可还原地“不连续”“离散”。这些新的几何对象自本世纪初以来地位日益重要,尤其是由于它们与算术——离散结构的典型科学——的紧密关系。这似乎是安德烈·韦伊(André Weil)\fnote{34}{}{安德烈·韦伊,旅居美国的法国数学家,是“布尔巴基小组”的“创始成员”之一;《收获与播种》第一部里会不少谈到这个小组(偶尔也会谈到韦伊本人)。}作品中的一条指导性想法,也许还是主要的力量线(在他的书面作品中,它或多或少停留在默认状态,这也是应该的):即“代数”几何,尤其是与不同素数相联系的“离散”几何,应当提供钥匙,让算术得到大范围更新。他就是本着这种精神,在1949年提取出著名的“韦伊猜想”的。

% ---------------- p.50 ----------------
说实话,这些猜想令人目瞪口呆:它们让人们对这些离散本性的新“簇”(或“空间”)隐约看见某些类型的构造与论证\fnote{35}{}{(供数学读者参考。)这里指的是与微分簇或复簇的上同调理论相联系的“构造与论证”,尤其是涉及莱夫谢茨不动点公式与霍奇理论的那些。}的可能性——而在此前,这些似乎只在分析家认为配得上这个名号的“空间”,即所谓“拓扑”空间(连续变化的概念在其中通行)的框架里才可设想。\origpage{P30}

可以认为,新几何首先是这两个世界之间的综合——它们此前相邻、紧密相依,却又彼此分隔:“算术”世界,其中住着那些没有连续性原则的(所谓)“空间”;以及连续量的世界,其中住着本来意义上的“空间”,它们对分析家的手段敞开,(正因如此)也被他接受为配得上在数学之城里有住所。在新眼光里,这两个昔日分隔的世界只成了一个。

这种“算术几何”(我提议这样称呼这种新几何)眼光的第一个胚胎,就在韦伊猜想里。在我某些主要主题\fnote{36}{}{指四个“中间”主题(第5至8号),即 topos、平展上同调与 $\ell$-进上同调、Motif,以及(较小程度上)晶体。这些主题是我在1958到1966年间先后提取出来的。}的发展中,从1958到1969这些年,这些猜想始终是我主要的灵感来源。而且,在我之前,一边是奥斯卡·扎里斯基(Oscar Zariski),另一边是让–皮埃尔·塞尔,已经为“抽象”代数几何那些无法无天的空间发展出某些“拓扑”方法,其灵感来自此前用于大家那些“正经”空间的方法\fnote{37}{}{(供数学读者参考。)在我看来,扎里斯基在这方面的主要贡献是引入“扎里斯基拓扑”(它后来成了塞尔在 FAC 中的一件本质工具),以及他的“连通性原则”和他所谓“全纯函数论”——后者在他手中变成了形式概形理论,以及形式与代数之间的“比较定理”(第二个灵感来源是塞尔那篇奠基性的 GAGA)。至于我在正文中提到的塞尔的贡献,当然首先是他把层(由让·勒雷在十几年前、在一个完全不同的语境中引入)的视角引入抽象代数几何,即在那篇已经提到的另一篇奠基性文章 FAC(《代数凝聚层》)中。\par 借这些“回顾”之光,如果我要举出新几何眼光最近的“祖先”,立刻想到的就是奥斯卡·扎里斯基、安德烈·韦伊、让·勒雷和让–皮埃尔·塞尔这几个名字。其中塞尔起了独特的作用,因为主要正是通过他,我才不仅知道了他自己的想法,也知道了扎里斯基、韦伊和勒雷那些在新几何的绽放与发展中起过作用的想法的。}。

当然,他们的想法在我构筑算术几何的最初几步中起了重要作用;不过说实话,更多是作为出发点和工具(为了一个广阔得多的语境的需要,我不得不或多或少把它们从头重做),而不是作为一路继续滋养我的梦与计划的灵感来源。无论如何,一开始就很清楚:即便重做过,这些工具也远远达不到所需——连朝着那些奇妙猜想走出最初几步都不够。\origpage{P31}

\section{(11) 魔法扇——或天真}

在新几何的启动与发展中,两个关键的力量线是概形与 topos。它们几乎同时出现,紧密共生
% ---------------- p.51 ----------------
\fnote{38}{}{这次启动发生在1958年,第23页的页底注里谈到它。“site”或“格罗滕迪克拓扑”的概念(topos 概念的临时版本)紧随着概形概念出现。反过来,正是它提供了“局部化”或“下降”的新语言,在概形主题与工具的发展中每一步都要用到。更为内蕴、更具几何性的 topos 概念,起初在随后几年里仍是隐性的,主要从1963年起随着平展上同调的发展而浮现,并逐步在我这里确立为最根本的概念。},在新几何引人注目的兴起中像是同一根驱动神经,而且从它们出现的那一年起就如此。为了结束这趟对我作品的巡览,我至少还得就这两个想法说几句。

概形这个概念,是把此前人们摆弄的那无穷系列的(代数)“簇”概念(每个素数一个这样的概念\fnote{39}{}{这一系列里还应当包括 $p=\infty$ 的情形,它对应“特征零”的代数簇。}……)统摄进唯一概念的想得到的最自然、最“显然”的办法。此外,同一个“概形”(或新版“簇”)对每个素数 $p$ 都产生一个完全确定的“特征 $p$ 的(代数)簇”。这些不同特征的簇的集合,于是可以被看成某种“簇的(无穷)扇”(每个特征一把)。而“概形”就是这把魔法扇,它把自己的各种可能特征下的“化身”或“转世”彼此连起来,像许多不同的“扇骨”。由此,它提供了一个有效的“过渡原则”,把那些此前显得多少彼此孤立、互不相通的几何所属的“簇”彼此连起来。如今它们被统摄在一门共同的“几何”里,并由它相连。可以把它叫做概形几何,它是那种“算术几何”的最初草图,而后者将在随后几年里在其中充分展开。\origpage{P32}

概形这个想法本身天真地简单——如此简单、如此卑微,以致在我之前没人想到要弯得这么低。说实话,甚至“蠢”得可以:以致好些年里,尽管事情明摆着,对我许多博学的同行来说,这实在“不严肃”!而我也确实需要好几个月紧凑而孤独的工作,才能在自己的角落里说服自己:“它行得通”——我那股不可救药的幼稚使我执意要试验的新语言,那个蠢得可以的语言,确实适合以新的光亮和新的精细,并且从此在一个共同的框架里,把握那些与前此“特征 $p$ 的几何”相连的最早的几何直觉中的某些。这类练习,任何“消息灵通”的人都会预先判定为愚蠢而无望;而在我的所有同行和朋友中,大概只有我一个人会想到把它放进脑子,甚至(被一个秘密的魔鬼驱使……)不顾所有人反对把它做到底!

我并没有让自己被周围那些关于什么“严肃”、什么不严肃的立法共识分散注意力,而是像过去一样,单纯地信任事物卑微的声音,信任我身上那个懂得倾听的东西。报偿是即时的,而且超出一切期待。在这几个月里,我甚至没有“特意”,就摸到了强大而出人意料的工具。它们让我不仅能以更透彻的光重新找到(像玩一样)那些公认难懂的旧结果并超越它们,而且终于能够着手并解决一些“特征 $p$ 的几何”问题——此前用当时一切已知手段都显得无法企及\fnote{40}{}{概形理论这次“强力启动”的报道,就是我1958年在爱丁堡国际数学家大会上的报告。在我看来,这篇报告是概形视角最好的导论之一,(也许)能让一位几何读者有动力去勉强啃一啃那部(后来的)厚重论著《代数几何原理》——它详详细细地(一个技术细节都不放过)阐述了代数几何的新基础与新技巧。}。\origpage{P33}

% ---------------- p.52 ----------------
在我们对宇宙之物的认识里(无论数学的或别的),我们身上那革新的力量不外是天真。那是我们出生时人人分得的原初天真,它安息在每个人心里,常常成为我们蔑视的对象,也成为我们最隐秘的恐惧的对象。唯有它把谦卑与大胆结合起来,让我们深入事物之心,也让我们允许事物深入我们、浸透我们。

这种力量绝不是“天赋”异禀的特权——比如说有一种超常的脑力,能灵巧轻便地消化和摆弄大量已知的事实、想法和技术。这样的天赋当然可贵,对于像我这样出生时没有被“超乎一切尺度地”赋予它的人,当然值得羡慕。

然而,不是这些天赋,也不是最炽烈的野心配上无懈可击的意志,能让人越过那些把我们的宇宙关起来的“看不见而专横的圈子”。唯有天真能越过它们,不知不觉也不在意,在那些我们独自倾听事物、全神贯注于一场儿童游戏的时刻……

\section{(12) 拓扑学——或雾中的丈量}

我们刚看到,“概形”这个革新性的想法,让人能把与不同素数(或不同“特征”)相联系的不同“几何”彼此连起来。然而这些几何各自本质上仍是“离散”或“不连续”的,与过去几个世纪传下来的(可上溯到欧几里得的)传统几何相对照。扎里斯基和塞尔引入的新想法,在一定程度上为这些几何恢复了某种连续性的“维度”,它随即被刚出现、旨在把它们统一起来的“概形几何”继承。但就那些(韦伊的)“奇妙猜想”而言,还差得远。从这个角度看,这些“扎里斯基拓扑”粗疏到这种地步,几乎就像还停留在“离散聚集体”的阶段。显然缺少的是某种新原则,能把那些几何对象(或“簇”,或“概形”)与通常的、或“正经”的(拓扑)“空间”连起来;所谓正经的空间,姑且说,是那些“点”显得彼此清楚分开的空间,而在扎里斯基引入的那些无法无天的空间里,点却有一种讨厌的倾向,彼此黏成一团……\origpage{P34}

确实,只有这样一个“新原则”的出现,而且非它不可,才能让那场“数与量的联姻”,或者说“不连续几何”与“连续几何”的联姻得以完成——韦伊猜想里已经有了对它的第一次预感。

“空间”概念大概是数学里最古老的概念之一。它在我们对世界的“几何”把握中如此基本,以致两千多年来它或多或少一直是默认的。只是在上一个世纪里,这个概念才逐步摆脱直接知觉(对我们周围那唯一同一个“空间”的知觉)及其传统理论化(“欧几里得式”)的专制控制,获得自身的自主与动力。今天,它属于数学中最普遍、最常用的少数概念之列,无疑为无一例外的所有数学家所熟悉。它也是再千变万化不过的概念,有千百张面孔,取决于人们给这些空间加上什么类型的结构:从最丰富的(如那些可敬的“欧几里得”结构,或“仿射”与“射影”结构,或同名“簇”的“代数”结构,后者推广并柔化了前者),直到最贫乏的:那些任何“定量”信息要素似乎都一去不返地消失了
% ---------------- p.53 ----------------
,只剩“邻近”或“极限”概念之定性精髓\fnote{41}{}{说到“极限”概念,我在这里想的尤其是“取极限”那个概念,而不是(对非数学人士更熟悉的)“边界”那个概念。},以及形式直觉最难以捉摸的版本(所谓“拓扑的”)。在这些概念中,最贫乏的一个,也是过去半个世纪里一直充当某种广阔的公共概念怀抱、把其余一切统摄于其中的那个,就是拓扑空间。对这些空间的研究构成几何中最迷人、最有生命力的分支之一:拓扑学。\origpage{P35}

由一个(所谓“拓扑”)“空间”所体现的这种“纯品质”结构,乍看之下无论多么难以捉摸——因为一切定量性质的数据(特别是两点之间的距离)都缺席,使我们无法挂靠任何关于“大”“小”的熟悉直觉——上个世纪里人们毕竟还是把这些空间精细地圈在了一套仔细“量身剪裁”的语言的紧密而柔韧的网眼里。更妙的是,人们还从头发明并制造出各种“米尺”或“量绳”,无论如何也要给那些触手般伸展的“空间”挂上各种“度量”(叫做“拓扑不变量”)——那些空间像抓不住的雾,似乎躲开一切测量的尝试。确实,这些不变量中大多数、而且是最本质的那些,其本性比单纯的“数”或“量”更精细——它们本身就是多少有些精巧的数学结构,以多少复杂些的构造挂到所考虑的空间上。这些不变量中最古老、最关键之一,上个世纪已经(由意大利数学家贝蒂)引入,是由与该空间相联系的各个所谓“上同调”“群”(或“空间”)构成的\fnote{42}{}{说实话,贝蒂引入的不变量是同调不变量。上同调是它的一种大致等价、“对偶”的版本,出现得晚得多。这个方面压过了最初的“同调”方面,主要(无疑)是由于让·勒雷引入了层的视角,下面会谈到。从技术角度看,可以说我大部分几何学家的工作,都在于提取并或多或少地发展那些缺失的上同调理论,为各种空间与簇、尤其是为“代数簇”和概形而设。一路之上,我还被引向用上同调的语言重新解释传统的同调不变量,并由此让它们在一种全新的光下显现。\par 拓扑学家还引入了许多别的“拓扑不变量”,用来把握拓扑空间的这类或那类性质。除了空间的“维数”和(上)同调不变量,最早的其他不变量是“同伦群”。1957年我又引入了另一个,即所谓“格罗滕迪克”群 $K(X)$,它随即大行其道,其重要性(无论在拓扑学还是在算术中)至今不断得到确认。\par 一大批新不变量,本性比目前已知和使用的更精细,但我觉得是根本性的,都写进了我的“温和拓扑学”纲领(其极简的概述见《纲领概要》,将刊于《反思》第4卷)。这一纲领基于“温和理论”或“温和空间”的概念,它有点像 topos 那样,构成空间概念的(第二次)“变形”。它比后者显然得多(在我看来),也没有那么深。我预计它对“严格意义上的”拓扑学的直接影响会明显更有冲击力,并且会彻底改变几何拓扑学家的“手艺”,办法是深刻改造他所工作的概念语境。(正如代数几何里引入概形视角时所发生的那样。)此外,我把自己的《纲领概要》寄给了好几位旧友和著名的拓扑学家,但看来它并没有引起他们中任何一位的兴趣……}。\origpage{P36}

正是它们(确实,主要是“在字里行间”)出现在韦伊猜想里,构成其深刻的“存在理由”,并使它们(至少对我来说,被塞尔的解释“带进了门”)具有全部意义。但要把这样的不变量挂到这些猜想里出现的“抽象”代数簇上,以回应为那一事业的需要而提出的非常精确的要求——这还只是一个希望。我怀疑,除塞尔和我自己之外,没有别人(甚至,尤其是,安德烈·韦伊本人!\fnote{43}{}{说来矛盾,韦伊对上同调形式主义有一种顽固的、看来是发自内脏的“阻塞”——而这恰恰主要是他那些著名的猜想,从1955年起启发了代数几何中那些宏大上同调理论的发展(塞尔以其已在前面的页底注中提到过的奠基性文章 FAC 开了第一炮)。\par 在我看来,这种“阻塞”在韦伊那里属于一种普遍的厌恶:厌恶一切“大堆头”,厌恶一切近似于形式主义(当它无法用几页纸概括时)或多少有点复杂交错的“构造”的东西。他确实一点也不像“建造者”;显然,他在三十年代是被迫地动手发展代数几何最初的“抽象”基础,而那些基础(鉴于他这种性情)对使用者来说实在成了一张普罗克鲁斯忒斯之床。\par 我不知道他是否怪我得太过分,怪我投入去建造那些宽敞的住宅,使克罗内克的梦和他的梦得以化身为一套精细而有效的语言与工具。总之,他从未就他看见我从事的工作、或已经完成的工作,对我说过一句话的评论。我把《收获与播种》寄给他三个多月,还附上我亲手写的热情题献,也没有得到任何回音。})真的相信……\origpage{P37}

% ================================================================
%  第 2 章 巡游一部作品——或孩子与母亲 (后半, 原书 p.54-71, 边码 P38-P65)
%  源文: _work/src1/task-ch02b.md
%  说明: 原书注 49、注 60 的注文在源文中缺失, 正文保留裸上标
% ================================================================

% ---------------- p.54 ----------------
\origpage{P38} 还在此前不久,我们关于这些上同调不变量的构想,已因让·勒雷的工作而大大丰富、深刻更新——那工作是战争期间他在德国当俘虏时继续做的,时在四十年代前半。其决定性的创新想法,是空间上的(阿贝尔)层,勒雷给每个层配上一列相应的“上同调群”(所谓“以该层为系数的”上同调群)。这就好像那根一直用来“丈量”空间的、可靠的老式“上同调标尺”,忽然被增殖为多得难以想象的众多新“标尺”,各种尺寸、形状、质地都有,每一根都贴合于所论的那个空间,每一根都就它向我们交出精确无比的信息,而且是唯有它才能给出的信息。这正是我们对各类空间的研究方式发生深刻转变的主导想法,也肯定是本世纪出现过的最关键的想法之一。多亏让–皮埃尔·塞尔后来的工作,勒雷的想法很快结出最初的果实——在其出现后的头一个十年里,拓扑空间(尤其是与上同调紧密相关的所谓“同伦”不变量)理论有了一次令人瞩目的重新启动;而所谓“抽象”代数几何也有另一次同样根本的重新启动(靠塞尔那篇奠基性文章 FAC,1955年发表)。我自1955年起在几何方面的工作,正与塞尔这些工作相接续,因而也与勒雷那些创新想法相接续。

\section{(13) topos——或双人床}

勒雷引入的层的视角与语言,让我们以一种新的光照去看各类“空间”与“簇”。然而它们并未触及空间概念本身,只是让我们以新的眼睛更精细地把握这些人人早已熟悉的传统“空间”。可后来表明:这个空间概念,并不足以承载那些最本质的“拓扑不变量”——正是它们表达着“抽象”代数簇(例如韦伊猜想所适用的那些)的“形状”,乃至一般“概形”(旧日簇的推广)的形状。对于那场期待中的“数与量的联姻”来说,这就像一张分明过于狭窄的床:充其量只有未来夫妇中的一个(即新娘)能勉强挤进去安身,却从来容不下两个!而那尚待发现的“新原则”——好让仁慈的仙女们许下的婚约得以圆房——不是别的,也正是未来夫妇所缺的那张宽敞的“床”;而在此之前,竟没有人哪怕察觉到这一点。

这张“双人床”,随着 topos(广义空间)这个想法出现了(像被魔杖一点……)。这个想法在同一种拓扑直觉里,既容纳了体现着连续量世界的传统(拓扑)空间,也容纳了那些不知悔改的抽象代数几何学家所谓的“空间”(或“簇”),还容纳了无数其他类型的结构——它们此前似乎无可挽回地被铆在“不连续”或“离散”的聚集体的“算术世界”上。

正是层的视角做了那沉默而可靠的向导,那把有效(且绝非秘密)的钥匙,
% ---------------- p.55 ----------------
\origpage{P39} 毫不迟疑、毫不绕弯地把我领进洞房,领到那张宽大的婚床前。这床确实宽得(像一条宽广、宁静、深不见底的河……),宽到——
\begin{quote}
“国王的全部马匹\\
都能一齐来饮……”
\end{quote}
——正如一支老歌里唱的;那歌你肯定也唱过,至少听人唱过。而头一个唱这歌的人,对 topos 那种隐秘的美、那种宁静的力量,感受得比我那些博学的学生、比我旧日的任何朋友都要真切……

钥匙是同一把,无论在最初那套临时性的进路里(经由“site”这个十分方便、却并非内蕴的概念),还是在 topos 的进路里。现在我想试着描述一下 topos 这个想法。

试想一个给定(拓扑)空间上所有层所构成的“集合”,或者,愿意的话,由所有这些用来“丈量”它的“标尺”所构成的那套了不起的装备\fnote{44}{}{(给数学家看的)严格说来,这里说的是集合的层,而不是勒雷为了构造“上同调群”而引入、作为最一般系数的阿贝尔层。顺便说,我相信自己是头一个系统地使用集合的层的人(从1955年起,在堪萨斯大学写的那篇《具有结构层的纤维空间的一般理论》里)。}。我们把这一“集合”或这套“装备”看作带有它最显然的那种结构——如果可以说的话,那结构“一眼就看得见”;也就是所谓“范畴”结构。(非数学的读者不必因为不知道这个术语的技术含义而发慌。下文他完全用不着它。)正是这种可称为“(所论空间上的)层范畴”的“丈量上层建筑”,从此将被看作“体现”着空间最本质的东西。这在(“数学上的健全常识”看来)完全正当,因为事实是:人们可以凭一个拓扑空间所关联的这个“层范畴”(或这套“丈量装备”),把它整个地“重新构造”出来\fnote{45}{}{(给数学家看的)严格说来,这只对所谓“索伯(sobre)空间”成立。不过这类空间几乎包括了我们平常遇到的全部空间,尤其包括分析学家们所钟爱的一切“分离”空间。}。要核实这一点只是个简单的练习——当然,得先把问题提出来……有了这一条,就足以确信:只要我们出于某种理由觉得合适,从此便可以“忘掉”最初那个空间,只保留、只使用与之关联的那个“范畴”(或那套“装备”);它将被看作所要表达的“拓扑(或空间)结构”最恰当的化身。

正如数学中常有的事,我们在这里(靠“层”或“上同调标尺”这个关键想法)成功地把某个概念(这里是“空间”)用另一个概念(这里是“范畴”)表达了出来。每一次,当人们发现可以把一个概念(表达某一类情境)译成另一个概念(对应另一类情境)时,我们对这两个概念的理解都会同时加深——因为分别附着于这一个或那一个的特定直觉,会出人意料地汇流到一起。于是,一个“拓扑”性质的情境(由某个给定空间体现),在这里被译成了一个“代数”性质的情境(由某个“范畴”体现);或者,愿意的话也可以说:由空间体现的“连续”,被这个“代数”性质的范畴结构“译”出、“表达”出来了——甚至直到那“不连续”或“离散”性质的外包装。

但这里还有更多的东西。这些概念中的头一个,即空间概念,在我们看来多少是个“极大”的概念——它已经如此一般,以至很难想象还能找到一个仍然“合理”的扩张。然而,事实是:在镜子的另一面\fnote{46}{}{这里说的“镜子”,像《爱丽丝镜中奇遇》里那面一样,是把一个空间放在它前面、便给出该空间之“像”的那面镜子;像就是与之关联的“范畴”,被看作空间的一种“替身”,“在镜子的另一面”……},我们从拓扑空间出发所遇到的这些“范畴”(或“装备”),性质非常特别。它们确实享有一组特征极强的性质\fnote{47}{}{(给数学家看的)这里说的主要是我在范畴论中以“正合性”为名引入的那些性质(同时还引入了现代范畴论中一般的归纳极限与投射极限概念)。见《论同调代数中的若干要点》,《东北数学杂志》,1957年(第119–221页)。},使它们看上去像是其中最简单那一个的各式“仿制品”——那最简单的,就是从缩成单独一点的空间出发所得到的那个。话虽如此,一个推广了传统拓扑空间的“新型空间(或 topos)”,将被径直描述为一个“范畴”:它未必来自某个寻常空间,却仍然具备这样一个“层范畴”所应有的全部好性质(当然,这些性质是一次性地明文规定好的)。

% ---------------- p.56 ----------------
\origpage{P41} 这就是那个新想法。它的出现,可以看作下面这个其实近乎孩子气的观察的结果:在一个拓扑空间里,真正要紧的绝不是它的“点”或点的子集\fnote{48}{}{例如,人们可以造出非常“大”的 topos,它只有一个“点”,甚至根本没有“点”!},也不是点与点之间的邻近关系之类,而是这个空间上的层,以及它们所构成的范畴。总而言之,我不过是把勒雷最初的想法推到了它最终的结论——推完之后,再跨出那一步。

如同层这个想法本身(源自勒雷),如同概形这个想法,如同一切要来撞翻某种积习已深的看法的“大想法”,topos 这个想法也自有一种令人不知所措之处:它那么自然,那么“显然”,那么单纯(可以说,到了天真乃至简陋的极限,甚至到了“犯傻”的地步)——就是那种特别的性质,常让我们脱口而出:“哦,原来就这么点东西!”语气里一半失望,一半嫉妒;也许还外加一层“古怪”“不正经”的弦外之音,凡是因为出人意料的过度简单而叫人摸不着头脑的东西,往往就落得这评语。这也许是在提醒我们,想起童年那些久已埋没、久已否认的日子……

\section{(14) 空间概念的变异——或气息与信仰}

概形概念构成了“代数簇”概念的一次大扩张,因此它把我前人所传的代数几何从里到外更新了一遍。topos 概念则构成了空间概念的一次意想不到的外延,更确切地说,一次蜕变。由此,它承载着拓扑学、并在此之外还有几何学获得类似更新的许诺。而且,从现在起它已经在新几何的兴起中起了关键作用(主要通过由它而生的 $\ell$-进上同调与晶体上同调这两个主题,并通过它们,在韦伊猜想的证明中起作用)。像它的姐姐(而且几乎是孪生姐姐)一样,它具备任何一次丰产的推广所必备的两个互补特征,如下。

% ---------------- p.57 ----------------
\origpage{P42} 第一,新概念不太宽:也就是说,在新的“空间”里(人们更愿意叫它们“topos”,免得冒犯一些敏感的耳朵\fnote{49}{}{“topos”这个名字是(连同“topologie”或“topologique”之名一起)选定的,意在暗示:它正是拓扑直觉所适用的“那个对象”。凭这个名字唤起的丰富心象之云,应当把它看作大致相当于“(拓扑)空间”一词,只是更强调这个概念“拓扑”上的特殊性。(于是,有“向量空间”,但在另行通知之前,并没有“向量 topos”!)看来应当把两个说法一并保留,各存其特殊性。}),那些对旧日好空间来说司空见惯的、最本质的“几何”直觉与构造\fnote{50}{}{这些“构造”里,尤其包括一切熟悉的“拓扑不变量”的构造,上同调不变量也在其中。至于后者,我在已经提到的那篇文章(“东北”1955)里把该做的都做了,好让它们对任何“topos”都有意义。},大体上都能搬过来。换句话说,对这些新对象,我们手头有整套丰富的意象与联想、概念以及至少某些技术,而这些东西以前只限于老式对象。

第二,新概念同时又足够宽,能把许多情境包进来——这些情境在此之前并不被认为会产生“拓扑–几何”性质的直觉,而那正是过去只留给寻常拓扑空间的直觉(而且自有其道理……)。

从韦伊猜想的眼光看,这里的关键在于:新概念确实足够宽,使我们能给每个“概形”配上这样一个“广义空间”或“topos”(称为所论概形的“平展 topos”)。这个 topos 的某些“上同调不变量”(全是些“傻乎乎”的东西!)当时看来颇有希望提供“所需要的东西”,使这些猜想获得全部意义,而且(谁知道呢!)也许还能提供证明它们的办法。

正是在我此刻正在写下的这些纸页里,我有生以来第一次(作为数学家)有闲暇去述说(哪怕只是对自己述说)我数学作品中全部的主旋律与重大指导性想法。这使我更好地领会这些主题、以及它们所体现的种种“视角”,在统一它们、并生出它们的那幅宏大几何图景中各自的位置与分量。正是通过这项工作,新几何最初那强劲兴起中两个要害的创新想法才在充分的光照下显现出来:概形的想法,与 topos 的想法。

这两个想法里,如今在我看来更深的是第二个,即 topos 的想法。假如五十年代末我没有挽起袖子,日复一日固执地干下去,一干十二个长年,做出一件精细与力量都臻完美的“概形工具”——那么我几乎无法想象,在此后的一二十年里,别人最终会一直忍住不去引入那个明明非引入不可的概念(哪怕是被迫的),也不会好歹搭起几间“预制件”的破棚屋,聊胜于我这番有心一石一石拼合、亲手竖起的那一座座宽敞舒适的宅邸。然而反过来,在近三十年的数学舞台上,我看不出还有谁会有这份天真,或这份单纯,去(替我)跨出那另一步——那一步比什么都关键,它引入了 topos 这个如此孩子气的想法(哪怕只是“site”那个想法)。就算这个概念已经有人殷勤地奉上,连同它似乎暗藏的那点怯生生的许诺——我也看不出还有谁,无论旧友还是学生,会有那份气息,尤其是那份信仰,把这卑微的想法\fnote{51}{}{(给数学家读者)我说“把这卑微的想法做到底”,指的是把平展上同调当作通往韦伊猜想的一条进路这个想法。正是在这句话的启发下,我于1958年发现了 site 概念;而这个概念(或与之极相近的 topos 概念)以及平展上同调的形式体系,则是在我的推动下于1962至1966年间发展起来的(有几位合作者协助,他们将在适当的地方谈到)。\par{}我说“气息”和“信仰”,指的是“非技术”性质的品质,而在这里,它们在我看来正是本质的品质。在另一个层面上,我还可以加上我所谓的“上同调的嗅觉”,也就是当初在我身上发展起来的、用于建造上同调理论的那种嗅觉。我原以为把它传给了我的上同调学生们。离开数学界十七年之后回头一看,我发现他们当中没有一个人把它保留下来。}一直做到底(它表面上如此微不足道,而目标看来无限遥远……):从最初结结巴巴的起步,直到“平展上同调之掌握”完全成熟——在随后那些年里,它最终在我手中化成了这个。

% ---------------- p.58 ----------------
\section{(15) 国王的全部马匹}

\origpage{P43}\origpage{P44} 是的,河很深;我童年的水又宽又静,在那片我以为早已离开的国土上。国王的全部马匹都能在那里从容地、喝个痛快地一齐来饮,也饮不尽它!那水来自冰川,炽热如那些遥远的积雪,又有着平原黏土的柔和。我刚说到其中一匹马,一个孩子把它牵来饮水,它久久地喝了个够。我还看见另一匹来饮了一会儿,说不定是循着同一个孩子的足迹来的——可那一回没耽搁多久。一定是有人把它赶走了。也就这样了,可以说,再没有别的。然而我看见平原上流浪着无数干渴的马群——就在今天早上,它们的嘶鸣还把我从床上拽起来,时辰还早得不像话,而我已是快六十的人,喜欢清静。没办法,我只好起床。看着它们瘦成一副副饿马的骨架,我心里难受,而好水明明不缺,青草地也不缺。可是看来,有人对这片我记忆中殷勤好客的国土施了恶意的魔法,封住了通往那些丰沛之水的路。又或许这是本地的马贩子设的局,好把价钱压下来,谁知道呢?又或许这是个再也没有孩子牵着马去饮水的国度,马儿们只好渴着,因为再没有一个孩子找得到通往那条河的路……

\section{(16) Motif——或心中的心}

topos 这个主题出自概形主题,就在概形出现的那一年——但在广度上它远远超出了母题。正是 topos 主题(而非概形主题)才是那张“床”,或那条“深河”:几何与代数、拓扑学与算术、数理逻辑与范畴论、连续的世界与“不连续”或“离散”结构的世界,都在那里结亲。如果说概形主题好比新几何的心脏,那么 topos 主题就是它的外壳,或它的居所。它是我所构想的最宽广的东西,为的是用同一种充满几何回响的语言,精细地把握那些彼此相距最远的情境所共有的某种“本质”——那些情境来自数学事物这个辽阔宇宙的这一片或那一片。

然而 topos 这个主题远没有像概形主题那样走运。我在《收获与播种》里多处谈到此事,这里不是细说这个概念所遭遇的种种离奇变故的地方。不过新几何的两个主旋律主题都出自 topos 主题,那是两个互补的“上同调理论”,它们各自都是为着给韦伊猜想提供一条进路而构想出来的:平展(或“$\ell$-进”)主题,与晶体主题。前者在我手中化成了 $\ell$-进上同调工具,如今它已显出是本世纪最强有力的数学工具之一。至于晶体主题,在我离开之后被压缩到一种近乎隐秘的存在,最终(在需要的压力下)于1981年6月被发掘出来,在聚光灯下,用一个化名,而情形比围绕 topos 的那些还要离奇。

正如所预料的,$\ell$-进上同调工具成了建立韦伊猜想的基本工具。我自己证明了其中一大摞,而最后一步是在我离开三年后由皮埃尔·德利涅出色地完成的,他是我那些“上同调”学生中最出色的一位。

另外,约在1968年,我曾提炼出韦伊猜想的一个更强的版本,而且尤其是一个更“几何”的版本。这些猜想一直“带着”(如果可以说的话!)一层看来不可约化的“算术”色彩,而这些猜想的本意恰恰是要通过“几何的”(或“连续的”)中介,来表达和把握“算术”(或“离散的”)\fnote{52}{}{(给数学家看的)韦伊猜想从属于一些“算术”性质的假设,尤其因为所论的簇必须定义在有限域上。从上同调形式主义的观点看,这导致要给与这种情形相关联的 Frobenius 自同态一个特殊位置。在我的进路里,那些关键性质(“广义指标定理”一类)涉及任意的代数对应,并且对事先给定的基域不作任何算术性质的假设。}。
% ---------------- p.59 ----------------
\origpage{P45}\origpage{P46} 在这个意义上,我所提炼的那个版本,比韦伊本人的版本更“忠于”“韦伊的哲学”——那种从未写出、也很少被说出的哲学;在过去四十年几何的非凡兴起中,它也许是主要的、未明言的动机\fnote{53}{}{不过,在我1970年离开之后,出现了一股十分明显的反挫,具体表现为一种相对停滞的局面;我在《收获与播种》的字里行间不止一次有机会谈到它。}。我的改写,要点在于提炼出某种“精华”:在被称作“抽象”代数簇的框架内,那仍然应当有效的东西——出自对“寻常”代数簇有效的经典“霍奇理论”\fnote{54}{}{这里“寻常”的意思是:“定义在复数域上”。霍奇理论(所谓“调和积分”理论)是复数代数簇背景下已知上同调理论中最强的一个。}。我把这些著名猜想的这个全新版本(就代数闭链而言)称为“标准猜想”。

在我心里,这是在发展出 $\ell$-进上同调工具之后,朝这些猜想迈出的新的一步。但同时,而且尤其是:它也是通往那个在我看来仍是我在数学中引入的最深主题的可能进路之一\fnote{55}{}{这是最深的主题,至少在我数学家活动“公开”的时期(1950至1969年)里是这样,也就是说直到我离开数学舞台为止。我认为自1977年起发展起来的远阿贝尔代数几何与伽罗瓦–泰希米勒理论这一主题,深度与之相当。}:Motif 的主题(它本身也生于“$\ell$-进上同调主题”)。这个主题好比概形主题的心脏或灵魂,是它最隐蔽、最躲开目光的部分;而概形主题本身又处于新图景的核心。而标准猜想\fnote{56}{}{(给代数几何读者)必要时,这些猜想需要重新表述。更详细的评论见《工地巡览》(ReS IV 第178号注,第1215–1216页)以及《信念与知识》第769页的页底注(ReS III 第162号注)。}里提炼出的少数关键现象,可以看作构成了 Motif 主题的终极精华,看作这个在一切主题中最精微者的“气息”,看作新几何这颗“心中的心”的“气息”。

大致是这么回事。我们已经看到,对给定的素数 $p$,(尤其为着韦伊猜想)知道如何为“特征 $p$ 的(代数)簇”构造“上同调理论”,这有多重要。而那条著名的“$\ell$-进上同调工具”恰恰提供了这样一种理论,甚至提供了无穷多个不同的上同调理论——对每个异于特征 $p$ 的素数 $\ell$,各有一个。这里显然又有一件“缺失的理论”,它对应的正是 $\ell$ 等于 $p$ 的情形。为补这个缺,我专门设想了另一种上同调理论(前面已经提到过),叫作“晶体上同调”。另外,在 $p$ 为无穷这个重要的情形里,还另有三种上同调理论\fnote{57}{}{(给数学家读者)这些理论分别对应于 Betti 上同调(用超越方法定义,借助把基域嵌入复数域)、霍奇上同调(塞尔定义)和 de Rham 上同调(我定义);后两者已可上溯到五十年代(而 Betti 上同调则上溯到上个世纪)。}——而且谁也不能保证人们迟早不会再去引入新的上同调理论,其形式性质全都类似。于是,与寻常拓扑学里的情形相反,这里摆在我们面前的是一大堆多得令人惶惑的不同上同调理论。人们有一种很清楚的印象:在某种起初还相当模糊的意义上,所有这些理论都该“归于同一”,它们“给出同样的结果”\fnote{58}{}{(给数学家读者)例如,若 $f$ 是代数簇 $X$ 的自同态,它诱导出上同调空间 $H^{i}(X)$ 的一个自同态,那么后者的“特征多项式”应当整系数,并且不依赖于所选取的那个特定上同调理论(例如 $\ell$-进理论,$\ell$ 可变)。对于一般的代数对应也一样,只要 $X$ 假定是真且光滑的。可悲的真相(它也让人看出,在我离开之后,特征 $p>0$ 的代数簇的上同调理论处于多么可悲的荒废状态)是:这件事至今仍未得到证明,甚至在 $X$ 是射影光滑曲面、$i=2$ 这个特殊情形里也没有。事实上,据我所知,在我离开之后还没有人肯费心关注这个关键问题——它正是那种看来从属于标准猜想的典型问题。时尚的法令是:唯一值得注意的自同态是 Frobenius 自同态(德利涅用现成的家伙什把它单独处理掉了……)。}。

% ---------------- p.60 ----------------
\origpage{P47} 正是为了表达不同上同调理论之间这种“亲缘”直觉,我提炼出了与一个代数簇相关联的“Motif”概念。用这个词,我想暗示:在借助先天地可能的全部上同调理论、与簇相关联的这一大堆不同上同调不变量之下,有一个“共同的动机”(或“共同的理由”)。这些不同的上同调理论,就好比同一个“基本动机”(称为“Motif 上同调理论”)的若干不同的主题展开,各自带着自己的“速度”、“调性”和“调式”(“大调”或“小调”);而这个基本动机,同时又是所有这些不同的主题化身(也就是说,所有这些可能的上同调理论)中最根本、最“精细”的一个。这样,与一个代数簇相关联的 Motif,便构成了“终极的”、“头等的”上同调不变量;其他一切(与各种可能的上同调理论相关联者)都从它推出,如同种种音乐化身,或种种“不同的实现”。簇的“上同调”的一切本质性质,都早已能在相应的 Motif 上“读出”(或“听出”),于是,那些特殊化的上同调不变量(例如 $\ell$-进或晶体不变量)上我们所熟悉的性质与结构,便只是 Motif 内部性质与结构的忠实映照\fnote{59}{}{(给数学家读者)另一种看域 $k$ 上 Motif 范畴的方式,是把它想象成 $k$ 上有限型分离概形范畴的一种“阿贝尔包络范畴”。于是,与这样一个概形 $X$ 相关联的 Motif(或“$X$ 的 Motif 上同调”,我记作 $H_{mot}^{*}(X)$)便显出是 $X$ 的一种阿贝尔化的“化身”。这里的关键在于:正如一个代数簇 $X$ 可以有“连续变化”(它的同构类因而依赖于连续的“参数”,或“模”),与 $X$ 相关联的 Motif、或更一般地一个“变化的”Motif,也同样可以有连续变化。这是 Motif 上同调的一个侧面,它与一切经典上同调不变量的情形形成鲜明对照,包括 $\ell$-进不变量——唯一的例外是复数代数簇的霍奇上同调。\par{}这就让人想见,“Motif 上同调”是一个精细得多的不变量,比传统上纯粹拓扑的不变量更紧地圈定了 $X$ 的“算术形状”(如果我敢冒险用这个说法)。在我的 Motif 图景里,这些 Motif 构成一条极其隐蔽、极其精致的“绳索”,把一个代数簇的代数–几何性质,与由它的 Motif 所体现的“算术”性质联结起来。后者可以被看作一个就其精神本身而言“几何”性质的对象,但在其中,那些从属于几何的“算术”性质可以说被“裸露”出来了。\par{}这样,在我眼中,Motif 便是人们至今能够与一个代数簇相关联的最深刻的“形状不变量”,除了它的“Motif 基本群”之外。这两个不变量在我看来,都像是一种尚待描述的“Motif 同伦型”的“影子”(我在《工地巡览——或工具与眼光》那条注里顺便说了几句(ReS IV 第178号,见工地5(Motif),尤其是第1214页))。我觉得,正是后面这个对象,才该是任意代数簇那种难以捉摸的“算术”(或“Motif”)“形状”直觉最完美的化身。}。

以上用音乐比喻的非技术语言说出的,是一个仍然具有孩子般单纯的想法之精华,它既精细又大胆。我是在那些我认为更紧迫的基础工作之余发展这个想法的,名之曰“Motif 理论”或“Motif 的哲学(或‘瑜伽’)”,时为1963–1969年。这是一个结构丰富得令人着迷的理论,其中很大一部分至今仍是猜想性的\fnote{60}{}{这些年里,但凡有人愿意听,我就把自己的 Motif 图景讲给他听,却始终不肯费事把这件事白纸黑字地发表出来(服务于众人的其他任务,我并不缺)。这就使后来我的某些学生能更从容地劫掠,而在我那些对情形一清二楚的旧友们的温情注视之下。(参见下页页底的注。)}。

% ---------------- p.61 ----------------
\origpage{P48}\origpage{P49} 我在《收获与播种》里多次谈到这套“Motif 的瑜伽”,它格外让我牵挂。这里不是回头重述我在别处所说的地方。只需说:那些“标准猜想”再自然不过地从这套 Motif 瑜伽中流出。同时,它们也为 Motif 概念的一种可能的成形构造提供了一个进路。

这些猜想,在我看来——今天在我看来仍然如此——是代数几何中最基本的两个问题之一。这个问题,以及另一个同样关键的问题(所谓“奇点消解”),至今都还没有解决。然而,第二个问题在今天仍像一百年前一样,显得是一个显赫而令人生畏的问题;而我曾有幸提炼出的那一个,却被时尚的专断法令(就在我离开数学舞台之后的那些年里,而且正如 Motif 主题本身一样\fnote{61}{}{事实上,这个主题在1982年(晶体主题之后一年)被发掘出来,这一回用的是它原来的名字(而且形式狭窄,只在基域特征为零这一种情形里),却没有说出工匠的名字。这只是一个例子,还有若干个:某个概念或主题在我离开之后随即被当作格罗滕迪克式的幻影埋掉,随后十年或十五年里又被我的某些学生一个个发掘出来,带着谦逊的骄傲,而且(还用得着再说明吗)不提工匠……})归为格罗滕迪克式的可爱胡闹。不过我又在预先说后面的事了……

\section{(17) 发现母亲——或两个坡面}

说实话,我为了建立韦伊猜想而对它们本身所作的思考,一直是零星的。那片已开始在我面前展开、我努力去细看、去把握的全景,在广度和深度上都远远超出一次证明的假定所需,甚至远远超出这些著名猜想最初所能让人瞥见的一切。随着概形主题与 topos 主题的出现,一个意想不到的新世界骤然打开了。“那些猜想”在其中确实占据中心位置,有点像一片幅员辽阔、行省无数的帝国或大陆的首都;而那些行省大多数与这个辉煌显赫的地方只有最遥远的关联。我从未需要对自己明说,但我知道自己从此是一项大任务的仆役:去探索这个巨大而未知的世界,把握它的轮廓直到最远的边界;也去四处走遍,以顽强而有条理的心思清点那些最近、最易到达的行省,为它们绘制忠实而精确得一丝不苟的地图,连最小的村落、最小的茅屋也要有它的位置……

主要是后面这项工作吸走了我大部分精力——那是一项耐心而庞大的基础工作,只有我一人看得清楚,尤其是只有我一人“用五脏六腑感觉得到”。从1958年(概形主题与 topos 主题接连出现的那一年)到1970年(我离开数学舞台的那一年),占去我时间远为最多的正是它。

而且我常常咬着嚼子忍耐:被这样拖住,像被一块黏人的重物坠着,没完没了地做那些(一旦看清了要害之后)在我看来更近乎“后勤”而非向未知冲刺的活儿。我得不断地按住那股向前冲的冲动——那是拓荒者或探险者的冲动,他出发去发现、去探索一个个未知而无名的世界,那些世界不停地呼唤我,要我去认识它们、为它们命名。那股冲动,以及我投在它上面的精力(几乎像是偷来的!),始终少得可怜。

然而我心底很清楚:正是这份从“任务”那里(可以说)偷来的精力,才是最稀有、最精微的本质——我作为数学家的“创造”,首先正落在那里:落在那份专注的凝神上,去把握那尚未诞生之物
% ---------------- p.62 ----------------
\origpage{P50}\origpage{P51} 在一个温暖而不竭的孕育母体中幽暗、无形、潮湿的褶皱里初现的形状与轮廓的痕迹;它似乎在呼唤我,要成形、要化身、要诞生……在发现的工作里,这份专注的凝神、这份炽热的关切是一种本质的力量,正如太阳的温暖之于埋在沃土里的种子的幽暗孕育,之于它们卑微而奇迹般地迎光绽放。

在我作为数学家的工作中,我看得最清楚的是这两股力量或冲动,它们同样深,而(我觉得)性质不同。为了说它们,我用了“建造者”的形象,以及“拓荒者”或“探险者”的形象。把两者并排放着,我忽然觉得它们都太“阳”,太“男性”,甚至“大男子气”!它们带着神话般高傲的回响,或“重大时刻”的回响。它们肯定是被我身上旧日那种“英雄式”的创造工作观(那超级阳刚的观点)的残迹所激发的。就这样,它们给出一种浓重着色的、甚至僵硬的、“立正”式的看法,去看一个远比这流动、远比这谦卑、远比这“简单”的现实——一个活生生的现实。

然而在这股“建造者”的男性冲动里(它似乎不断把我推向新工地),我同时也清楚辨认出恋家者的冲动:深深依恋“那个”家的人的冲动。首先,那是“他的”家,是“亲近者”的家——是一个亲密的活体之所在,他觉得自己是它的一部分。只是到了后来,随着被感受为“亲近”的圈子逐渐扩大,它才也成为一座“众人的家”。而在这股“造房子”的冲动里(就像人们“做”爱那样……),也有,而且首先有,一份柔情。还有与那些材料接触的冲动:一件一件地塑造它们,带着爱怜的用心,而只有通过这种充满爱的接触才能真正认识它们。而当墙立起来、梁和屋顶架好之后,还有一种深深的满足:一间一间地安置,看着在这些厅堂、房间和小角落里,渐渐建立起这座活着的房子的和谐秩序——美、好客、宜于居住。因为房子,首先并且隐秘地在每个人心里,也是母亲——是环绕我们、庇护我们的东西,既是庇护所又是安慰;而且也许(更深一层,甚至就在我们正把它整个造出来的时候)我们自身也正是从那里来的,是它在我们的出生之前那些永远被遗忘的时光里庇护了我们、哺育了我们。那也是“怀抱”。

而刚才为了越过“拓荒者”这个显赫名号、为了把握它所遮盖的那个更隐蔽的现实,自然而然浮现出来的形象,也同样剥去了一切“英雄”腔调。那里出现的,仍然是母性的原型意象——哺育的“母体”及其无形而幽暗的劳作……

这两股在我看来“性质不同”的冲动,最终比我以为的更接近。两者都属于一种“接触的冲动”,把我们带向“母亲”的相遇:带向那位既体现着切近的、“已知”之物,又体现着“未知”之物的她。委身于这一股或那一股冲动,都是“重新找到母亲”。都是重新接上与切近者、与“或多或少已知者”的接触,同时也接上与“遥远者”、与那“未知”却同时被预感着、即将被认识者的接触。

这里的差别是色调上的、配比上的,不是性质上的。当我“造房子”时,占上风的是“已知”;当我“探索”时,是未知。这两种发现的“方式”,或者更确切地说,同一个过程或同一项工作的这两个方面,不可分割地连在一起。两者都本质,而且互补。在我的数学工作里,我辨认出在这两种进路之间、更确切地说在其中一个占上风的时刻(或时期)与另一个占上风的时刻(或时期)之间,有一种不断的往返运动\fnote{62}{}{我在这里就数学工作所说的话,对“静观”的工作也同样成立(《收获与播种》里到处都会谈到它)。我几乎不怀疑:这是出现在一切发现工作中的东西,包括艺术家(譬如说,作家或诗人)的工作。我在这里描述的两个“坡面”,也可以看作:一个是表达及其“技术”要求的坡面,另一个是接受(种种知觉与印象)的坡面,后者靠一份专注的凝神而变成灵感。两者在工作的任何时刻都在场,而且在其中一个占上风的“时节”与另一个占上风的“时节”之间,总有这种不断的“往返”运动。}。但同样清楚的是:在任何时刻,两种方式都在场。

% ---------------- p.63 ----------------
\origpage{P52} 当我建造、布置,或者清理、打扫、归置时,定调的是工作中“阳”的或“男性”的“方式”或“坡面”。当我摸索着探索那不可把握者、那无形者、那无名者时,我便是自己存在中“阴”的或“女性”的那一面。

对我而言,问题绝不是要贬低或否弃我天性中的这一面或那一面,两者都本质——那建造并生育的“男性”,和那孕育、并庇护缓慢幽暗的怀胎的“女性”。我“是”两者——“阳”与“阴”,“男人”与“女人”。但我也知道,创造过程里最精细、最精微的本质,在“阴”的、“女性”的那一坡面一边——那谦卑、幽暗、往往其貌不扬的坡面。

工作中始终(我相信)对我施以最强魅惑的,正是这一坡面。可当时的共识却鼓励我把最明澈的那部分精力投到另一坡面,投到那在可触的“产品”中化身并自我确证的坡面——且不说那些产品还是完成、收尾了的,轮廓分明,以凿过的石头般的明证来担保自己的实在……

拉开距离看,我很清楚这些共识是怎样压在我身上的,也很清楚我是怎样“承受这重量”的——柔顺地承受!我工作中“构思”或“探索”的那一部分,直到我离开的时候都还被压缩在最低限度,好吧。然而,在这样回望我作为数学家的作品时,有一点以惊人的明显浮现出来:使这部作品有其本质与力量的,恰恰是这个如今被轻视、甚至被嘲笑或居高临下地鄙夷的坡面:那是“想法”的坡面,甚至是“梦”的坡面,绝不是“结果”的坡面。在这些纸页里,我试着界定我对我那个时代的数学所贡献的最本质的东西,用的是一种拥抱森林、而不在几棵树上流连的眼光——我看到的不是一串“大定理”的荣誉榜,而是一把活生生的、丰产想法的扇子\fnote{63}{}{并非我的作品里缺少人们所谓的“大定理”,其中还包括解决了别人提出、而我之前无人能解的问题的定理。(我在“上涨的海……”那条注(ReS III 第122号)第554页的页底注(***)里回顾了其中几个。)但正如我在这趟“巡游”开头(“视角与眼光”那一站,第6节)已经强调过的,这些定理对我来说,只有在一个大主题的滋养性语境里才获得全部意义,而那主题是由某个“丰产想法”开启的。于是它们的证明就像从源头、毫不费力地从承载它们的那个主题的本性、“深度”中流出——正如河水之浪仿佛从它水域的深处温柔地生出,没有断裂,也不费力。我在已经引过的那条注“上涨的海……”里,用别的意象表达了完全类似的意思。},它们全都汇聚到同一个广阔的图景里。

\section{(18) 孩子与母亲}

当这篇“前言”开始变成一趟穿越我数学作品的巡游,还加上我那段关于“(血统纯正的)继承人”和“(无可救药的)建造者”的小小品时,这个没写成的前言也开始有了一个名字:它该叫“孩子与建造者”。接下来几天里越来越清楚:“孩子”和“建造者”是同一个人物。于是这名字变得更简单了:“孩子—建造者”。天哪,这名字倒不缺少气派,而且完全对我的胃口!

可只见反思表明:这位高傲的“建造者”,或者(更谦逊地说)那个玩造房子游戏的孩子,只是那个著名的“玩耍的孩子”的一副面孔,而他有两副。还有另一个孩子:他爱探索东西,爱到处乱掏,爱钻进沙堆,钻进泥泞无名的淤滩,
% ---------------- p.64 ----------------
\origpage{P53}\origpage{P54} 钻进最不可能、最荒唐的角落……大概是为了掩人耳目(哪怕只是瞒过我自己……),我起初给他起了个火亮的名字“拓荒者”,后来又换成“探险者”——这名字更接地气,却仍罩着一圈威望的光晕。简直要问:在“建造者”和“拓荒者–探险者”之间,哪一个更男性、更诱人!正面还是反面?

可再凑近一看,我们这位无畏的“拓荒者”原来是个女孩(我却乐意把她打扮成男孩)——她是水洼、雨、细雨和夜的姐妹,沉默,几乎看不见,因为她总退到阴影里——那个总被遗忘的人(人们不对她做出一副嘲弄的样子就已经不错了……)。而我也真有办法,一天又一天地把她忘掉——可以说,是双倍地忘掉:我起初只愿意看见那个男孩(那个玩造房子的孩子……)——即便后来我终于忍不住,到底还是看见了另一个,我看到的她,也仍然是个男孩……

至于我那趟巡游的好名字,这下就完全站不住了。那是个全“阳”的、全“大男子”的名字,一个瘸腿的名字。要让它站得不歪,就得把另一个也写进去。可怪的是,“另一个”并没有真正的名字。勉强贴得上的只有“探险者”,但那还是个男孩的名字,没办法。这里的语言是个婊子,它甚至不让你察觉就把你套住,显然跟祖传的偏见串通一气。

也许可以用“造房子的孩子与探索的孩子”敷衍过去。把这一点留在不言之中:一个是“男孩”,另一个是“女孩”,而这是同一个男孩–女孩的孩子,他在建造时探索,在探索时建造……可是昨天,除了这观照与探索、命名与建造的阴阳两坡面之外,事情的又一面还显现了出来。

宇宙、世界,乃至万有,归根到底是陌生而遥远的东西。它们并不真正与我们相干。在我们最深处,求知的冲动并不是把我们带向它们。吸引我们的,是它们可触而切近的化身,最近、最“肉感”的那个,充满深沉的回响与奥秘——那化身与我们血肉之躯的起源相混同,也与我们族类的起源相混同——也是那自古以来就等待着我们、沉默着准备接纳我们的她,“在路的另一端”。冲动从她那里涌出,欲望的道路从她那里伸展——从她,母亲,从那位生我们如同生这世界的她;而那些道路把我们带向与她的相遇,朝她奔去,为的是一次次折返,消融在她之中。

% ---------------- p.65 ----------------
\origpage{P55} 就这样,在一趟不期而至的“巡游”的路上拐个弯,我意外地重又遇到一个曾让我觉得亲切、又被我有些遗忘了的寓言——孩子与母亲的寓言。可以把它看作一个关于“生命寻求自身”的寓言。或者在更谦卑的个体存在层面上,看作一个关于“存在者寻求事物”的寓言。

这是一个寓言,也是一种祖辈经验的表现,那经验深植于心灵之中——是滋养深层创造层的原初象征中最有力的一个。我相信在其中认出了以原型意象的远古语言表达出来的东西:人身上创造力的气息本身,它鼓动着人的血肉与精神,无论在它最卑微、最短暂的表现里,还是在最辉煌、最持久的表现里。

这“气息”,正如使它化身的血肉形象一样,是世上最卑微的东西。它也是最脆弱的、最被人忽视、最遭人轻贱的东西……

而这气息在你一生中的种种遭际,不是别的,正是你的冒险,你一生中“认识”的冒险。表达它的那个无言的寓言,就是孩子与母亲的寓言。

你是孩子,出自母亲,被庇护在她之中,以她的力量为食。孩子从母亲那里奔出——那至近的、熟识的母亲——去与母亲相遇:那无边的、永远未知而满是奥秘的母亲……

\begin{center}
“巡游一部作品”终
\end{center}

\subsection*{尾声:看不见的圈子}
\tocsubsub{尾声:看不见的圈子}

\section{(19) 死亡是我的摇篮(或三个娃娃给一个垂死者)}

直到五十年代末 topos 视角出现为止,空间概念的演变在我看来都是一种本质上“连续”的演变。它似乎平顺地、没有跳跃地延续着,从欧几里得对我们周围空间的理论化开始,从希腊人传下的几何学开始——那几何学致力于研究生活在这个空间里的某些“图形”(直线、平面、圆、三角形等)。诚然,数学家或“自然哲学家”构想“空间”\fnote{64}{}{我写尾声时最初的打算,是附一份极简的梗概,略述其中某些“深刻变化”,并显出我在其中看到的这种“本质连续性”。我放弃了,免得把这趟巡游拉得太长——它已经比预定的长多了!我想在《反思》第四卷拟议的“历史评注”里再谈,那一回是为数学家读者写的(这完全改变了阐述的任务)。}的方式发生过深刻的变化。但这些变化在我看来全都属于一种本质的“连续性”——它们从未把(像所有人一样)依恋着熟悉意象的数学家,突然抛进一片陌生的地方。它们就像一个人身上那些也许深刻、却是渐进的变化:我们从这人的童年就认识他,从他会走路一直看到他成年、完全成熟。有些漫长平淡的时期里,变化觉察不到;另一些时期里,也许又很剧烈。但即便在生长或成熟最强烈的时期,即便我们有几个月、甚至几年没再见到他,任何时刻都不可能有丝毫怀疑、丝毫犹豫:我们重又见到的,仍然是他,一个熟识的、亲近的存在,哪怕容貌已经变了。

另外我想可以说:到本世纪中叶,这个亲近的存在已经大大地老了——像一个终于筋疲力尽、磨损殆尽的人,被一大堆他毫无准备的新任务压垮。也许他甚至已经安然死去,却没人费心去记下一笔、验明此事。“大家”还都做出一副在活人家里忙活的样子,弄得几乎就像他确实还活得好好的。

那么,请想想这对家里的常客是多么难堪:那位可敬的老人原本僵直地端坐在扶手椅里,忽然却看见一个身量不过三个苹果高的壮实小崽子在那儿撒欢,还顺口宣称——毫不带笑,仿佛理所当然——空间先生(从现在起,您尽可随意把那“先生”二字去掉……)就是他!要是他至少有点像家里的相貌也好啊,说不定是个私生子,谁知道……可一点也没有!一眼看去,没有任何地方让人想起那位大家那么熟悉(或者以为熟悉……)的老空间父亲,那位大家无论如何都确信(这总该是最起码的……)是永恒的父亲……

% ---------------- p.66 ----------------
\origpage{P56}\origpage{P57} 这就是那著名的“空间概念的变异”。这就是我至少从六十年代初起就不得不“看见”、把它当作明摆着的事的东西,却一直没机会把它对自己明确说出来,直到我写下这些行字的此刻。而忽然,只凭这个形象的说法、只凭它立刻唤起的联想之云,我以新的清晰看见:传统的“空间”概念,以及与之紧密相关的“簇”的概念(各类簇,尤其是“代数簇”),在我来到这一带的时候,已经老得那么厉害,简直就像已经死了……\fnote{65}{}{这个断言(有人会觉得它武断)得“加点盐”来听。它的有效程度,不多不少,正如下面那个断言(我稍后也把它算在自己账上):力学的“牛顿模型”(地上的或天上的)在本世纪初已是“奄奄一息”,那时爱因斯坦赶来救驾。事实是,即便在今天,物理学里大多数“寻常”情形下,牛顿模型都完全够用,而去找相对论模型才是发疯(鉴于测量所允许的误差幅度)。同样,数学里许多情形下,旧的熟悉的“空间”与“簇”概念仍然完全够用,不必去找幂零元、topos 或“温和结构”。但在两种情形里,对于尖端研究中越来越多出现的语境,旧的概念框架都已变得无法表达哪怕最“寻常”的情境。}。我还可以说:正是随着概形视角(连同它的子嗣\fnote{66}{}{(给数学家看的)在这个“子嗣”里,我尤其算上形式概形、各类“重数”(特别是概形重数或形式重数),最后还有所谓“刚性解析”空间(由泰特引入,依循我提供的一个“总施工”,灵感来自 topos 这个新概念,同时也来自形式概形的概念)。这份清单当然远不完备……},外加一万页基础工作为其代价)、随后是 topos 视角的接连出现,一个不声不响的危机局面终于被解开了。

在刚才那幅图景里,其实不该说是一个小崽子,作为一次突然变异的产品,而该说是两个。而且这两个孩子之间有一种无可否认的“家族相似”,尽管他们一点也不像那位故去的老人。再者,凑近看还可以说:概形这小家伙,好比是已故的空间父亲(别名“各类簇”)与 topos 这小家伙之间的一个“亲缘环节”\fnote{67}{}{另外,对这两个娃娃,还该再加上第三个更小的,它出现在不那么宽厚的年代:就是“温和空间”这小崽子。正如我在别处指出过的,它没拿到出生证明,而我还是在完全不合法的状态下把它算进了我有幸引入数学的十二个“主旋律主题”之中。}。

\section{(20) 瞥一眼对门的邻居}

这局面,我觉得与本世纪初爱因斯坦的相对论理论出现时的局面非常接近。那时有一个更触目的概念死胡同,表现为一个突然的、看来无法解决的矛盾。理所当然,那个要来重整混沌的新想法,是一个孩子般单纯的想法。值得注意的是(而且符合一个再重复不过的剧本……):在所有那些忽然紧张起来、拼命想“保住家什”的聪明人、显要人、大人物当中,没有一个人想到这个想法。非得是一个无名的年轻人,刚从(说不定)学生的阶梯教室板凳上下来,跑来(也许还为自己那份大胆有点不好意思……)向他那些显赫的前辈解释该怎么做才能“拯救现象”:只要不再把空间与时间分开就行了\fnote{68}{}{这样描述爱因斯坦的想法,当然有点简略。在技术层面上,需要弄清的是在新时空上放什么结构(不过这一点其实已经“飘在空中”了,有麦克斯韦的理论和洛伦兹的想法)。这里本质性的一步不是技术性质的,而是“哲学”性质的:意识到对于相距遥远的事件而言,同时性概念没有任何实验上的实在性。就是这个“孩子气的观察”,这句“可皇帝没穿衣服!”,让人跨过了那道著名的“限定一个宇宙的、专横而看不见的圈子”……}!从技术上说,当时一切都已具备,足以让这个想法绽开并被接受。而爱因斯坦那些前辈的体面也正在于此:他们确实懂得接纳新想法,没有过多训斥。这是一个迹象,说明那还是个伟大的时代……

% ---------------- p.67 ----------------
\origpage{P58} 从数学角度看,爱因斯坦的新想法平淡无奇。但从我们对物理空间的构想来看,那却是一次深刻的变异,一次突然的“背井离乡”。这是自2400年前欧几里得提炼出物理空间的数学模型以来的第一次此类变异;那模型自古至今被所有物理学家和天文学家(包括牛顿)原封不动地拿来供力学之用,用以描述地上与天上的力学现象。

爱因斯坦这个最初的想法后来大大深化了,化成一个更精微、更丰富、更柔韧的数学模型,借助了已有数学概念这座丰富的武库\fnote{69}{}{这里主要是指“黎曼簇”的概念,以及这种簇上的张量计算。}。随着“广义相对论”,这个想法扩展为一幅物理世界的宏大图景,以同一道目光拥抱亚原子的极小微观世界、太阳系、银河与遥远的星系,以及电磁波在时空中行进——那时空每一点都被其中的物质所弯曲\fnote{70}{}{这个模型区别于空间与时间的欧几里得(或牛顿)模型、也区别于爱因斯坦最初那个模型(“狭义相对论”)的最触目的特征之一,是时空的整体拓扑形状始终不确定,而不是由模型本身的性质强制规定。这种整体形状究竟是什么,在我看来(作为数学家)是宇宙学中最迷人的问题之一。}。在宇宙学与物理学的历史上,这是第二次、也是最后一次(继三个世纪前牛顿那第一次大综合之后),出现一幅用数学模型的语言表达的、统一的宏大图景,涵盖宇宙中全部物理现象。

不过,爱因斯坦这幅物理宇宙的图景,后来也被事态淹没了。需要解释的“全部物理现象”,从本世纪初起就有时间不断扩充!出现了大量物理理论,各自以或多或少的成功,去解释那一大堆“被观察的事实”这座巨大杂物堆里有限的一捆。而人们至今还在等那个大胆的孩子,他会在玩耍中找到那把新钥匙(如果真有那么一把的话……),那个梦想中的“蛋糕模型”,肯于一下子“走通”,把所有现象一齐拯救……\fnote{71}{}{人们把这样一种假想的理论称为“统一理论”,它终将把前面谈到的众多局部理论“统一”起来、调和起来。我觉得,那项尚待着手的基本思考,将要在两个不同的层面上展开。\par 1°)一个“哲学”性质的思考,关于“数学模型”这个概念本身——为一个现实片段而设的数学模型。自牛顿理论成功以来,物理学家的一条默许公理便是:存在一个数学模型(甚至是唯一一个模型,或“那个”模型),能完美地表达物理实在,没有“脱节”,没有含糊。这个统治了两个多世纪的共识,像是毕达哥拉斯“万物皆数”那活生生的洞见留下的一块化石残迹。也许这就是新的“看不见的圈子”,取代了从前的形而上学圈子,用来限定物理学家的宇宙(而“自然哲学家”这一族类看来已彻底灭绝,被计算机那一族高奏凯歌地取代了……)。只要肯停下来哪怕一瞬,就很清楚:这个共识的有效性绝非显而易见。甚至有非常严肃的哲学理由,导致先天地怀疑它,或至少为它的有效性划定极严格的界限。现在正是时候,把这个公理交给一番严格的批判,甚至“证明”——超出一切可能的怀疑——它没有根据:不存在唯一严格的数学模型,能说明至今编目的全部所谓“物理”现象。\par 一旦“数学模型”这个概念本身、以及这样一个模型的“有效性”概念(在测量所允许的某些“误差幅度”之内)被令人满意地界定清楚,“统一理论”的问题、或至少“最优模型”(意义尚待明确)的问题,才终于会被清楚地提出来。同时,对于选择这样一个模型所带有的(也许是必然带有的)任意性程度,人们大概也会有更清楚的想法。\par 2°)在我看来,只有在这样的思考之后,提炼一个比前人的模型更令人满意的明确模型这一“技术”问题,才获得全部意义。那时也许正是时候,摆脱物理学家的第二条默许公理;这条公理可上溯到古代,深植于我们感知空间的方式本身:那就是空间与时间(或时空)的连续本性,也就是“物理现象”在其中展开的那个“场所”的连续本性。\par 想必已是十五、二十年前,我翻着构成黎曼全集的那本薄薄的书,被他“顺带”的一句评论打动了。他在那里指出:空间的终极结构完全可能是“离散的”,而我们关于它所形成的“连续”表象,也许只是对一个更复杂的实在的一种(久而久之也许过度的)简化;对人类精神来说,“连续”比“不连续”更容易把握,因此它便充当了我们把握不连续的一种“近似”。这话出自一位数学家之口,洞察力令人吃惊,而且当时物理空间的欧几里得模型还从未受到质疑;从严格的逻辑意义上说,传统上反倒是“不连续”被用作通向连续的技术进路。\par 何况近几十年来数学的发展表明,连续结构与不连续结构之间的共生,远比本世纪前半人们所想象的更为紧密。无论如何,要找到一个“令人满意的”模型(或者必要时,一组这样的模型,尽可能令人满意地“彼此衔接”),无论它是“连续”的、“离散”的还是“混合”性质的——这样的工作肯定要动用极大的概念想象力,以及一种炉火纯青的嗅觉,去把握并揭示新型的数学结构。这种想象力或“嗅觉”,我觉得是稀罕之物,不仅在物理学家中如此(那里爱因斯坦和薛定谔似乎是难得的例外),甚至在数学家中也如此(而这话我是完全有把握说的)。\par 总而言之,我预计,那期待中的更新(如果它还会来的话……)更可能出自一位骨子里的数学家、同时又熟知物理学的重大问题,而不是出自一位物理学家。但尤其是:那需要一个人具备“哲学上的开放”,才能抓住问题的症结。这个症结绝不是技术性质的,而是一个“自然哲学”的根本问题。})。

% ---------------- p.68 ----------------
\origpage{P59}\origpage{P60} 把我对我那个时代数学的贡献,与爱因斯坦对物理学的贡献相比较,这一念头是出于两个理由逼到我面前的:这两份作品都得益于我们对“空间”的构想发生的一次变异(一者取数学意义,另一者取物理意义);而且两者都采取一种统一性图景的形式,拥抱极大一批此前看来彼此分离的现象与情境。我在此看到他的作品\fnote{72}{}{我绝不声称自己熟悉爱因斯坦的作品。事实上,他的著作我一本也没读过,他的想法我只是通过道听途说、而且极为粗略地知道。不过我觉得自己还是看得出“森林”,尽管我从未费劲去细看他的任何一棵树……}与我的作品之间一种明显的精神亲缘。

这种亲缘,我觉得丝毫不被一种明显的“实质”差异所否定。正如我前面已经暗示的,爱因斯坦式的变异涉及物理空间概念,而爱因斯坦是从已知数学概念的武库里取用,从不需要扩充它,更不必颠覆它。他的贡献在于:在当时的已知数学结构中,挑出那些最适合\fnote{73}{}{关于“奄奄一息”这个形容词的评论,见前面的一条页底注(第55页的注)。}充当物理现象世界之“模型”的,用来取代前人传下的那个垂死的模型。在这个意义上,他的作品确实是一位物理学家的作品,并且在此之外,还是一位“自然哲学家”的作品——取牛顿及其同代人所理解的意义。这个“哲学”维度在我的数学作品里是没有的:我从未被引去追问那些在数学事物的宇宙中进行的“理想”概念构造,与物理宇宙中发生的现象(甚至心灵中经历的事件)之间可能有什么关系。我的作品是一位数学家的作品,他有意避开了“应用”(于其他科学)的问题,也避开了我工作的“动机”与心理根源的问题。而且是一位被其十分独特的天赋所驱使、不断扩充其技艺根基处概念武库的数学家的作品。就这样,我甚至没有察觉、仿佛在玩耍之间,被引去颠覆了对几何学家来说最根本的那个概念:空间概念(以及“簇”概念),也就是我们对几何对象所居住的那个“场所”本身的构想。

新的空间概念(像一种“广义空间”,但那些本该构成“空间”的点已经或多或少消失了)在实质上一点也不像爱因斯坦引入物理学的那概念(后者对数学家来说毫不令人迷失)。相反,比较倒该与薛定谔发现的量子力学去做\fnote{74}{}{据我理解(通过从各方传回的回声),人们一般认为本世纪物理学有过三次“革命”或大动荡:爱因斯坦的理论、居里夫妇对放射性的发现、以及薛定谔引入量子力学。}。在这种新力学里,传统的“质点”消失了,代之以一种“概率云”,它在此处与彼处的周围空间中疏密不同,取决于质点
% ---------------- p.69 ----------------
\origpage{P61}\origpage{P62} 位于该区域的概率。在这种新眼光里,人们清楚感到我们构想力学现象的方式发生了比爱因斯坦模型所体现的更深的一次“变异”——这次变异并不是简单地把一个在肩宽上略嫌狭窄的数学模型,换成一个类似但裁得更宽、更合身的模型。这一回,新模型与那些好的旧传统模型像得如此之少,以至连力学的大专家数学家也必定突然感到背井离乡,甚至迷失(或者愤慨……)。从牛顿力学走到爱因斯坦力学,对数学家来说,想必有点像从地道的普罗旺斯土话走到最时新的巴黎黑话。而走到量子力学,我想,那是从法语走到汉语。

而这些“概率云”——取代了昔日那些令人安心的物质粒子——奇异地让我想起充满 topos 的那些飘忽不定的“开邻域”:它们像一个个消散的幽灵,环绕着一些想象出来的“点”;而不肯驯服的想象力仍旧不顾一切地攀附着那些点……

\section{(21) “唯一者”——或孤独的馈赠}

这趟到“对门邻居”物理学家家里的小小远足,对一个(像大多数人那样)对数学家的世界一无所知、却肯定听人说起过爱因斯坦和他那著名的“第四维”、甚至听说过量子力学的读者,也许可以充当一个参照点。说到底,即便发明者当初并未料到他们的发现会具体化为广岛,后来又化为军备与(所谓)“和平”用途的原子竞赛,事实是:物理学上的发现有形地、几乎立刻地影响着一般人的世界。数学发现的影响,尤其在所谓“纯”数学(也就是说,不以“应用”为动机)里,则不那么直接,肯定也更难界定。譬如我就不知道,我对数学的贡献曾“服务”于什么,比如说造出个什么机器。事情如此,我当然无功可居,但这不妨碍它让我安心。一旦有了应用,就可以肯定:头一个扑上去的是军方(其后是警察)——至于工业(哪怕是所谓“和平”工业),也不见得总好多少……

当然,无论为我自己导航,还是为一位数学家读者导航,更该做的其实是在数学史本身之中找“参照点”来定位我的作品,而不是到别处去找类比。这些天我想过这事,以我对那段历史相当模糊的了解为限\fnote{75}{}{我还是个小鬼的时候,就对历史不太上心(地理也一样)。(《收获与播种》第五部分(只写了一部分)里,我曾“顺带”探到我觉得这种对历史的局部“阻塞”的深层原因——我相信近些年来这阻塞正在消解。)何况我的前辈们在“布尔巴基圈子”里所接受的数学教育,也没让事情好转——那里偶尔出现的历史引用少而又少。}。在“巡游”途中,我已经有机会提到一条数学家的“谱系”,他们的性情是我所认同的:埃瓦里斯特·伽罗瓦(Evariste Galois)、黎曼、希尔伯特。如果我对本门技艺的历史知道得更多,说不定我能把这条谱系往过去延伸得更远,或者在其中插入几个我只有耳闻的名字。让我吃惊的是:我不记得曾听说过(哪怕只是通过比我更熟悉历史的朋友或同行的暗示)除我之外还有哪个数学家,带来过一大串创新想法,而这些想法并不是彼此多少脱节的,而是像一个宏大统一图景的各个部分(物理与宇宙学中的牛顿和爱因斯坦、生物学中的达尔文和巴斯德就是这样的情形)。在数学史上,我只知道两个“时刻”诞生过一幅宏大格局的新图景。其一是2500年前古希腊数学作为我们今天所理解的科学诞生之时。另一个首先是十七世纪微积分诞生之时,那个时代以牛顿、莱布尼茨、笛卡尔等人的名字为标志。据我所知,在这两个时刻之一诞生的图景,都不是一个人的作品,而是一个时代的集体作品。

% ---------------- p.70 ----------------
\origpage{P63} 当然,从毕达哥拉斯和欧几里得的时代到十七世纪初,数学有时间改变面貌;同样,从十七世纪数学家创造的“无穷小计算”到本世纪中叶,也是如此。但据我所知,在这两个时期(一个两千多年,一个三百年)里发生的深刻变化,从未具体化、凝聚为一幅表达在某一部给定作品里的新图景\fnote{76}{}{写完这些行字几个小时后,我忽然想到:我在这里竟没有想到布尔巴基先生那部(集体)论著所力图呈现的当代数学的大综合。(在《收获与播种》第一部分里还会大量谈到布尔巴基小组。)我觉得这出于两个理由。\par 一方面,这一综合仅限于把一大批已经知道的想法与结果来一番“整理”,并不从中贡献它自己独创的新想法。若说其中有新想法,那大概是“结构”概念的一个精确数学定义;它确实成了贯穿整部论著的一条宝贵线索。但这个想法在我看来,更像一位聪明而富有想象力的词典编纂者的想法,而不像一种语言获得更新的要素——那种更新会带来对实在(这里是数学事物的实在)的新把握。\par 另一方面,从五十年代起,结构这一想法就被事态超越了:随着“范畴”方法骤然涌入数学中某些最有活力的部分,例如拓扑学或代数几何。(这样,“topos”概念就拒绝进入布尔巴基那只结构的“口袋”,那口袋确实窄得在肩宽上不够用!)布尔巴基固然是在完全知情的情况下决定不去趟这“浑水”,但这样一来,它也就放弃了最初的雄心:为整个当代数学提供基础与基本语言。\par 不过,它倒是固定下了一套语言,同时也固定下某种写作风格与处理数学的方式。这种风格起初是某种精神的(很不完整的)反映,那精神是希尔伯特活泼而直接的遗产。五十、六十年代,这种风格终于占了上风——好的方面如此,(更主要地)坏的方面也如此。近二十年来,它终于变成一种僵硬的“教规”,其所标榜的“严格”纯是门面,而当年鼓舞它的那种精神,看来已一去不返。},不像物理学和宇宙学里随大综合而发生的情形:牛顿的,然后是爱因斯坦的,在它们历史中的两个关键时刻。

看来,作为我心中诞生的那幅统一性宏大图景的仆役,我在数学史上从开端到如今都算“独一无二”。抱歉,我这样子像是要把自己抬举得超过了允许的分寸!不过,让我自己松一口气的是,我似乎辨认出一个潜在的(而且是天赐的!)兄弟。

刚才我已经有机会提到他,作为我那些“性情上的兄弟”中的第一位:埃瓦里斯特·伽罗瓦。在他短暂而闪电般的一生中\fnote{77}{}{埃瓦里斯特·伽罗瓦(1811–1832)死于一场决斗,年二十一岁。我记得他有好几部传记。我年轻时读过物理学家英费尔德写的一部传记体小说,当时给我很深的印象。},我相信辨认出了一幅宏大图景的开端——正是那“数与量的联姻”的图景,在一幅新的几何图景之中。我在《收获与播种》别处提到\fnote{78}{}{见《伽罗瓦的遗产》(ReS I,第7节)。},两年前我心中如何突然出现这个直觉:在当时最强烈地魅惑我的那项数学工作里,我正在“接续伽罗瓦的遗产”。这个直觉此后很少被提起,却有时间在沉默中成熟。我三周来对作品的回顾性思考,肯定又对此有所助益。如今我相信自己与一位过去数学家认出的最直接的谱系关系,正是把我与埃瓦里斯特·伽罗瓦连起来的那一条。不管对错,我觉得我一生中十五年发展起来、并在我离开数学舞台之后十六年里仍在我心中继续成熟、继续丰富的那幅图景——这图景也是伽罗瓦不可能不去发展的那一幅\fnote{79}{}{另外我深信,一个伽罗瓦会走得比我远得多。一方面因为他那份完全卓异的天赋(就我而言,并未分得这份天赋)。另一方面因为他大概不会像我这样,让大部分精力被那些没完没了的活儿分去,去把已经多少到手的东西逐一细致地整理成文……},
% ---------------- p.71 ----------------
\origpage{P64}\origpage{P65} 假如他处在我这个位置待在这一带,而没有早逝粗暴地截断一段壮丽的冲势。

肯定还有一个理由,促成我这种“本质亲缘”的感觉——这种亲缘不能归结为单单“数学性情”,也不能归结为一部作品的显著方面。在他的生活与我的生活之间,我还感到一种命运上的亲缘。诚然,伽罗瓦死得愚蠢,才二十一岁;而我呢,快六十了,还决心活得长长久久。可这并不妨碍:埃瓦里斯特·伽罗瓦生前,正如一个半世纪之后的我,在官方数学界里都是一个“边缘人”。就伽罗瓦的情形,浅看之下会觉得这种边缘性是“偶然的”,他不过是还没来得及靠自己的创新想法和作品“站住脚”。就我的情形,在我数学家的最初三年里,我的边缘性是由于我对数学界本身的存在一无所知(也许是故意的……)——那个我本该去面对的圈子;而自我离开数学舞台以来的十六年里,它是深思熟虑的选择的结果。肯定正是这个选择,招来了报复性的“无懈可击的集体意志”:要把我的名字的一切痕迹从数学中抹去,连同我所侍奉的那幅图景一起抹去。

但越过这些偶然的差异,我相信在这种“边缘性”里辨认出一个共同的原因,我觉得它是本质的。这个原因,我不把它看作历史境况,也不看作“性情”或“性格”上的特点(这些在我与他之间,大概就像人与人之间可能有的那么不同),当然更不在“天赋”的层面(伽罗瓦显然天赋惊人,而我比较起来只是平平)。如果真有一种“本质亲缘”,我看见它在一个远远更谦卑、远远更基本的层面上。

在我一生中,我曾在少数难得的场合感到这样的亲缘。也正因它,我觉得自己与另一位数学家“相近”,他是我的前辈:克劳德·谢瓦莱\fnote{80}{}{我在《收获与播种》里多处谈到克劳德·谢瓦莱,尤其是“与克劳德·谢瓦莱相遇——或自由与善意”那一节(ReS I 第11节),以及“向克劳德·谢瓦莱告别”那条注(ReS III 第100号注)。}。我要说的联系,是某种“天真”或“单纯”的联系,我前面有机会谈到过它。它表现为一种倾向(常不为周围人所赏识):用自己的眼睛去看事物,而不是透过某个人类群体殷勤奉送的、有专利的眼镜去看——那群体或大或小,出于这样那样的理由被赋予了权威。

这种“倾向”,或这种内心的态度,不是成熟的特权,而是童年的特权。它与生命一同,是出生时得到的馈赠——一份谦卑而可畏的馈赠。一份常常被深埋的馈赠,有些人还能多少保住它,或者也许能重新找回它……

也可以把它叫作孤独的馈赠。

% ================================================================
%  第 3 章 一封信 (前半: 原书 p.73-91, 边码 L1-L30)
%  源文: _work/src1/task-ch03a.md
% ================================================================
\chapter{一封信}

\noindent 1985年5月\par

% ---------------- p.73 ----------------
\section{(1) 千页之信}

\origpage{L1}我在这里寄给你的这份文本,是我所在大学张罗着打字、并只印了有限几册的,可它既不是抽印本,也不是预印本。它的名字——《收获与播种》——已经够清楚地预告了这一点。我把它当作一封长信寄给你——而且是一封再私人不过的信。我之所以寄给你,而不是满足于让你哪天(如果有这份好奇心)在书店里卖的某本书里看到它(倘若有哪个出版商疯到肯冒这个险……),是因为在它里面我更多是对你说话,而不是对别人。写它的时候,我不止一次想到你——得说,这封信我写了一年多,是把自己整个儿放进去写的。这是我给你的一份馈赠;写的时候我尽心竭力,把我(每一刻)能拿得出的最好的东西给出来。这份馈赠会不会被接受,我不知道——你的回答(或者你的不回答……)会让我知道……

在寄给你的同时,我也把《收获与播种》寄给数学界里所有那些同事、朋友或(旧日的)学生:我曾在某个时候与他们密切相连,或者他们以某种方式出现在我的反思里,点了名也好,没点名也好。你很有可能就在其中;而只要你用心去读,不只用眼睛和脑子去读,你必定会在那些没有点你名字的地方认出自己。我也把《收获与播种》寄给另外几位朋友,是科学家也好,不是也好。

% ---------------- p.74 ----------------
\origpage{L2}这封你正在读的“导引之信”,它向你预告并呈上一封“千页之信”(这还只是开头……),也将权作前言。写这几行的时候,那篇前言还没有写出来。此外,《收获与播种》由五个部分组成(还不算一篇“抽屉式”的导论)。我这里寄给你的是第一部分(《愚妄与更新》)、第二部分(《葬礼(1)——或中国皇帝的袍子》)和第四部分(《葬礼(3)——或四种运算》)\fnote{1}{}{我把那些以某种名义出现在我的反思里、但我并不亲自认识的同行另作别论。我只给他们寄去《四种运算》(那一部分与他们关系更近),同时寄去由这封信构成的“O 分册”,以及《收获与播种》的导论(外加前四部分全体的详细目录)。}。我觉得这几部分与你关系更近。第三部分(《葬礼(2)——或阴阳之钥》)无疑是我这份见证里最私人的部分,同时也是在我看来——比别的部分更甚——具有“普遍”价值的部分,超越它诞生时周遭的那些特殊境况。第四部分(《四种运算》)里,我这儿那儿提到这一部分;不过第四部分可以独立来读,甚至在很大程度上可以独立于前面三个部分来读\fnote{2}{(*)}{总的说来你会看到,每一个“节”(在《愚妄与更新》里)或每一则“按语”(在《收获与播种》其后三个部分的任何一个里),都有它自己的统一与自治。它可以脱离其余部分单独读,正如人看着一只手、一只脚、一根手指或一根脚趾、或整个身体上任何一大块一小块,也能看出兴味、得到乐趣,同时却不会忘记:那只是一个整体的一部分,而唯有这个整体(它始终留在未言说之中)才赋予它全部的意义。}。如果你读了这里寄去的东西,想给我回信(这也是我的愿望),或者它让你想读读缺的那一部分,那就告诉我。只要你回信让我感到你的兴趣超出了纯粹表面的好奇,我很乐意把它寄给你。

\section{(2) 《收获与播种》的诞生(一份闪电式的回顾)}

\origpage{L3}在这封前信里,我现在想用几页篇幅(如果办得到的话)告诉你,《收获与播种》里讲的是什么——说得比那个副标题本身更详细些:“一位数学家对往昔的反思与见证”(我的往昔,你猜得到的……)。《收获与播种》里东西很多,各色人等在其中看到的想必也很不同:一趟发现一段往昔的旅行;一场关于存在的冥想;一幅某个圈子、某个时代的风俗画(或者说,是从一个时代悄悄而不可抗拒地滑向另一个时代的图画……);一场调查(有时近乎侦探小说,另一些时候则近乎在数学这座大都市的底层上演的斗篷与长剑小说……);一次大规模的数学漫游(它会播下不止一颗……);一部应用精神分析的实用论著(或者,随你挑,一本“精神分析小说”);一篇关于自我认识的颂词;“我的忏悔”;一本私人日记;一种关于发现与创造的心理学;一份起诉书(照例是毫不留情的……),甚至是在“数学的美好世界”里算的一笔总账(而且不给人留半点情面……)。可以肯定的是,写它的时候我一刻也没有感到无聊,而与此同时我见识了、也看尽了各色各样的东西。要是你那些要紧的事务让你有空读它,那我倒奇怪你会读得无聊。除非是逼着你读,谁知道呢……

显然,它并不只是对数学家说话。不过也确有一些时候,它更多是对数学家说的,而不是对别人。在这封写给“《收获与播种》之信”的前信里,我却正想把那些可能与你、作为数学家,关系较近的东西概括出来、并着重凸显出来。
% ---------------- p.75 ----------------
\origpage{L4}要这么做,最自然的方式就是照直对你讲讲:我怎样一针引一线,接二连三写出了刚才提到的那四五块“砖头”。

如你所知,1970年我离开了数学的“大世界”,起因是我所属的机构(IHÉS)里一桩军费的事。此后几年,我搞反军备、搞生态,那种“文化大革命”式的行动,你大概在这儿那儿也听到过一些回声;然后我几乎从流通中消失了,天知道丢到了外省的哪所大学里。有传闻说,我把时间都花在放羊和打井上。实情是:除了许多别的营生,我还像所有人一样,硬着头皮去系里上课(这就是我那点毫无新意的糊口活儿,今天还是)。偶尔,甚至能连着几天、几星期、几个月,断断续续地重做一些数学——我手头满是装满我鬼画符的纸盒,能看懂的想必只有我自己。但那都是些与我从前做的事,至少乍看之下,极为不同的东西。1955年到1970年之间,我心爱的题目是上同调,尤其是各色簇的上同调(特别是代数簇)。我判断自己在这个方向上已经做得够多,尽可以让别人不用我也能应付;既然要做数学,那也该换一张唱片了……

1976年,一种新的激情进入我的生命,其强烈不亚于当年的数学激情,而且说来还是这激情的近亲。这就是对被我称作“冥想”的东西的激情(事物总得有个名字)。这个名字——正如任何别的名字在这里也免不了——必然会招来无数误解。与数学中一样,这是一件发现的工作。关于它,我在《收获与播种》里这儿那儿有所表述。总之,显然这里有的是活儿,够我干到死。而且,不止一次,我当真以为数学已成往事,今后我只会去管那些更严肃的事——以为我要去“冥想”了。

然而我最终还是不得不承认(那是四年前),数学的激情并未就此熄灭。甚至,连我自己也不太明白是怎么回事,并且暗自吃惊:我这个人(近十五年来)早已不再想在有生之年发表一行数学,却忽然发现自己被卷进了写一部数学著作的活儿里,那部书显然没完没了,还要一卷一卷地出;既然已经上了路,我索性要把自己以为在数学上有话要说的东西,统统倒进一套叫《数学反思》的(无穷?)系列书里,从此再没人提它。

那是两年前,1983年春。当时我已经太忙于写《追寻栈》(第一卷)了——那本该也构成《数学反思》(数学的)第一卷——顾不上追问自己身上发生了什么事。九个月后,理所当然地,这第一卷算是写完了,差不多也就是说,只剩导论要写、全稿要重读、注要加——然后就可以付印……

% ---------------- p.76 ----------------
\origpage{L5}直到此刻,那部书仍未完成——一年半以来它一动没动。那篇还欠着的导论,已经越过了(打字稿)一千二百页的大关;等它真正干完,少说也有一千四百页。你已经猜到了,这篇所谓“导论”,不是别的,正是《收获与播种》。据最新的消息,它该构成那套预定“系列”的第一、二卷,再加上第三卷的一部分。于是那套系列就此改名,干脆叫《反思》(不一定非得是数学的)。第三卷剩下的部分,主要由数学文本构成,如今这些对我比对《追寻栈》还要烫手。《追寻栈》就先等着吧,等明年再弄它的注、索引,当然还要一篇导论……

第一幕终!

\section{(3) 老板之死——弃置的工地}

\origpage{L6}我觉得,现在是时候做几点说明了:为什么我那么突然地离开了一个我在其中显然过得挺自在、过了二十多年的世界;为什么我起了那个古怪的念头,要“回来”(像个还魂的人……),而在我缺席的这十五年里,人家本来过得挺好;最后,为什么一部六七百页的数学著作的导论,竟写到一千二(或者一千四)百页。而在这里,一进入正题,我大概就要让你难受了(抱歉!),甚至惹你恼火。因为毫无疑问,你跟我从前一样,喜欢把你所属的那个圈子看得很“玫瑰色”——那个圈子里有你的位置、你的名头,诸如此类。我知道那是什么滋味……而这一回,恐怕要有点刺耳了……

我在《收获与播种》里这儿那儿谈到我离去的那段插曲,没有太在其中停留。那次“离去”在那里更像是我作为数学家的生命里一个重要的断口——我作为数学家的种种事件,都不断地相对于这个“点”来定位,作为“之前”与“之后”。要把我从一个我深深扎根的圈子、从一条早已深深划定的“轨迹”上拽开,需要一记极猛的冲击。这冲击来自碰撞:在一个我深深与之认同的圈子里,撞上了某种形式的腐败\fnote{3}{}{这里说的是各国科学家整体——以“建制”为首——与军事机器之间毫无保留的合作,把它当作经费、声望与权力的方便来源。这个问题在《收获与播种》里只顺带掠过一两次,例如去年4月2日那则《敬重》(第179号,第1221—1223页)。};对这腐败,我直到那时都选择闭上眼睛(仅仅做到自己不参与)。事后回看,我明白:在那事件之外,还有一股更深的力量在我身上起作用。那是一种强烈的、内在更新的需要。这样的更新,在一个地位显赫的科研机构那温吞的学术氛围里,是无法完成、也无法持续的。我身后是二十年紧张的数学创造、是二十年过度的数学投入——同时也是二十年漫长的、在“密闭罐”里的精神停滞……不知不觉间,我窒息了——我需要的是开阔处的空气!我那“天赐的”离去,标志着一场漫长停滞的骤然终结,也是我身上那些深刻力量走向平衡的第一步——那些力量曾被弯折、被拧紧在一种极度失衡、僵固的状态里……那次离去,真真正正是重新出发——是新旅程的第一步……

如我所说,我的数学激情并未就此熄灭。它在一阵阵零星的思考中找到了表达,走的路与我“从前”为自己划定的全然不同。至于我留在身后的那部作品,那部“从前”的作品——无论是白纸黑字印出来的,还是那部也许更为本质、却还没能找到书写之途、未能成书的作品——它大可以显得,而且在我眼里也确实显得,已经与我脱了钩。在去年之前(那是在《收获与播种》里),我从没起过念头,要在那些零零星星传回来的回声上稍作停留。我很清楚:我在数学里所做的一切,尤其在我1955到1970年那个“几何的”时期所做的,都是些本该被做出来的事;而我看见或瞥见的那些东西,都是本该出现、本该被拉到光天化日之下的东西。我也清楚:我做的工作,以及我让人做的工作,都是做得好的工作,是我整个人都投进去的工作。我把全部力气与全部爱都放了进去;而(在我看来)它从此是自治的了——一个活生生的、强健的东西——不再需要我来当它的母亲。在这方面,我离去时心里完全平静。我毫不怀疑:我留下的那些写出来的与没写出来的东西,我是把它们留在了可靠的手里;那些手会照看着它们展开、生长、繁衍,依照它们作为活物、作为强健之物的本性。

在这十五年紧张的数学工作中,在我身上绽放、成熟、长大的,是一种宏大的、统一的眼光,
% ---------------- p.77 ----------------
\origpage{L7}它化入少数几个极简单的有力观念。那眼光是一种“算术几何”的眼光,是拓扑学、几何学(代数的与解析的)与算术的综合,我在韦伊猜想里找到了它的第一个胚胎。正是它,成了那些年里我主要的灵感之源——在我看来,那些年尤其是这样一段时期:我从中提炼出这门新几何的主导观念,并打制出它几件主要的工具。这眼光、这些有力观念,对我来说已成了某种第二本性。(而在与它们断绝一切接触近十五年之后,我如今发现,这“第二本性”在我身上依然活着!)它们对我来说如此简单、如此自明,以至于理所当然地以为“所有人”都随时间与我一道把它们吸收、据为己有了。只是最近,就在这几个月里,我才明白:无论是那眼光,还是那几条一直是我恒常向导的“有力观念”,在任何一个已出版的文本里都找不到逐字的书写,至多不过是字里行间。而尤其要紧的是:我自以为已经传达出去的那眼光,以及承载它的那些有力观念,在达到完全成熟二十年后的今天,仍然不为任何人所知。我这个工匠,这些我有幸发现之物的仆役,同时也是它们在其中仍然活着的唯一一人。

\origpage{L8}我打制的这样或那样的工具,被这儿那儿用来“撬开”一个公认困难的问题,就像撬开一只保险箱。那工具显然结实。可我知道它还有另一种“力”,不只是万能钳子的力。它是一个整体的一部分,正如肢体是身体的一部分——那整体是它由之而来的,赋予它意义,它从那里取得力与生命。你尽可以用一根骨头(如果够粗)去砸开一个脑壳,这是明摆着的事。但那不是它真正的功能,不是它存在的理由。而我看见这些工具散落各处,被这个人那个人抓在手里,有点像被仔细剔净、剥光的骨头,是他们从一具身体上扯下来的——一具活的身体,而他们装作不认识它……

我在这里用字斟句酌的话说出的东西,是经过长久反思之后的结论;但它想必早已被我一点一点、含含糊糊地感知到了,年复一年,在那个尚未成形、还在寻求进入自觉的思想与形象、进入清晰分节之言说的未言说层面。我曾判定:这段过去,说到底与我无关了。那些远远近近传到我这里来的回声,尽管都被过滤过,却仍然意味深长,只消我稍作停留。我曾以为自己不过是众多工匠中的一个,在五六个正开足马力干着的“工地”上忙活\fnote{4}{}{关于这些被弃置的“工地”,我最后在三个月前那组《荒废的工地》(第176$'$至178号)里做了梳理。一年前,在发现“葬礼”之前,这事已经谈过一次,就在我重新与自己的作品接上联系、谈到它后来的遭遇的第一则按语《我的孤儿们》(第46号)里。}——也许是个更有经验的工匠,是那个长辈:当年他曾独自在这些地方干了许多年,直到一拨称心的接班人来了;是长辈,不错,但说到底与别人并无不同。可是你看,这一位一走,那就像一家石作行在老板突然去世之后宣告破产:几乎可以说一夜之间,工地就空了。“工匠们”走了,每人腋下夹走一点小零碎,觉得自家用得着。钱箱也没了,于是从此再没有任何理由让他继续费力干活了……

这同样是一种说法,是从一年多的反思与调查里澄清出来的。但毫无疑问,早在我离去后的最初几年,这已经是“在什么地方”被感知到的东西了。除了德利涅关于弗罗贝尼乌斯特征值绝对值的那些工作(就是那个“威望问题”,我最近才弄明白……)——每当我在远远近近碰上一位旧日的亲近同伴,曾与我同在工地上干过的,我问他“那么……呢?”,得到的总是同一个意味深长的手势:双臂往上一抬,像是在讨饶……显然,所有人都在忙着比那些曾放在我心头的事更重要的事;也显然,就在所有人摆出一副忙碌而要紧的神气忙活的时候,其实没做出多少东西。最要紧的
% ---------------- p.78 ----------------
\origpage{L9}那一块已经不见了——那种赋予局部任务以意义的统一,还有,我想,一种温度。剩下的是一堆从整体上脱落下来的零散任务,各人在自己的角落里孵着自己的小宝贝,或者好歹让它生点利息。

即便我想抗拒,眼看一切都戛然而止,我当然还是难受:再也听不到有人谈起 Motif,谈起 topos,谈起六种运算,谈起德拉姆系数,也听不到有人谈起霍奇系数,谈起那个“神秘函子”——它本该把各个素数对应的 $\ell$-进系数,围绕德拉姆系数,在同一把扇面上彼此联系起来;听不到有人谈起晶体(除非是得知它们始终原地不动),也听不到有人谈起“标准猜想”以及其他那些我提炼出来、显然代表着关键问题的东西。就连随《代数几何原理》(靠着迪厄多内不倦的协助)开始的那项宏大的基础工作——那几乎只消顺着已经取得的势头继续下去——也被撇在一边:所有人都满足于搬进别人耐心拼装、架设、擦亮了的墙壁与家具里。工匠一走,竟没有一个人想到轮到自己挽起袖子、动手拿起抹刀,去盖那些还等着盖的许多房子——那些好房子,能住人的,为自己,也为所有人……

我又一次忍不住,用完全自觉的形象把话接了下去,这些形象是从反思的工作里析出、浮上来的。但我毫不怀疑:这些形象早已以某种形式存在于我存在的深处。我一定早已感到,对我作品、同时也对我本人的那场“葬礼”,有着某种悄悄逼近的真实性;而这真实在去年4月19日突然强加于我,带着不可辩驳的力量,并且就叫那个名字:“葬礼”。相反,在自觉的层面上,我几乎不会想到要为此愤慨,甚至不会为此悲伤。说到底,不管是不是旧日的“亲近者”,那只是当事人自己的事:他选择拿自己的时间干什么。如果当年似乎推动过他、启发过他的东西不再启发他了,那是他的事,不是我的事。如果同样的事,以完美的一致,无一例外地落在我所有旧日的学生头上,那仍然是他们各自的事;我还有别的猫要抽,没工夫去追究这能有什么意思,就这样完事!至于我留下的那些东西,那些与我仍有某种深刻而不被察觉的联系的东西——即便它们显然被弃置在那片荒废的工地上,我自己却很清楚:它们不是那种会怕“时间的凌辱”、会怕风尚起伏的东西。如果它们还没有进入公共的遗产(当年我本以为已经进入了),它们迟早会在那里扎根,十年或一百年,说到底无关紧要……

\section{(4) 一阵葬礼的风}

可是,尽管这些年来我一直乐于回避那场大规模“葬礼”的含混感知,它却不依不饶地提醒我记起它,以别的面貌,而且不那么无害——不只是对一部作品的冷淡而已。我渐渐得知——我也说不清是怎么得知的——那被遗忘的眼光里的好几个概念,不仅已经废弃不用,而且在一个体面的圈子里,成了居高临下鄙夷的对象。情形尤其如此的是那个关键的统一性概念:topos(广义空间)——它正处在新几何的核心,正是它提供了拓扑学、代数几何与算术共有的几何直觉,也正是它让我既能提炼出平展的、$\ell$-进的层论工具,又能提炼出晶体上同调的那些主导观念(诚然,此后它们多多少少被遗忘了……)。说实话,
% ---------------- p.79 ----------------
\origpage{L10}连我自己的名字,也年复一年地、悄悄地、莫名其妙地成了取笑的对象——成了没完没了、稀里糊涂的“轰炸”的同义词(正是针对那些著名的 topos,或者那些“Motif”——他拿这些玩意儿把你耳朵都磨出茧子来了,可谁也没见过……),成了长达千页地把一根头发劈成四瓣的同义词,成了关于那些人人都一向知道、根本不用等他说的事滔滔不绝、臃肿巨大的唠叨的同义词……大致就是这类调子,不过是压低了声音,靠言外之意,带着“在高蹈之士与高雅之交之间”应有的那份体贴。

在《收获与播种》里进行的那场反思中,我想我已摸到了在各人身上起作用的那股深刻力量——它就藏在那副取笑与俯就的神气背后,面对一部其分量、其生命、其气息都非他们所能企及的作品。我还发现(撇开我这个人身上那些给我的作品与命运打下印记的特殊之处)那个秘密的“催化剂”,它促使这些力量以这种形式显现出来:面对一个完好无损的创造力所给出的意味深长的迹象,露出那种漫不经心的轻蔑。总之,在这场被取笑和轻蔑捂得悄无声息的葬礼上,他就是那位大主祭。说来奇怪,他同时也是所有人当中离我最近的那一个——而且是唯一一个曾经领会过某种眼光、并把它据为己有的人,那眼光里满是生命与强烈的力量。不过我又抢先了……

说实话,那些这儿那儿传回我这里的“隐约取笑的小阵风”,并没有特别伤到我。直到三四年前,它们在某种意义上还是无名的。我当然从中看到一个不大令人愉快的时代征兆,但它们并没有真正质疑我,也没有在我心里引起焦虑或不安。相反,更直接触动我的,是那些与我本人拉开距离的迹象,它们这儿那儿来自我昔日数学界里的不少朋友——尽管我离开了那个曾与我们共有的世界,我仍觉得自己与他们由同情的纽带相连——除了共同的激情与一段共同的过去所造就的那些纽带之外。在这方面,每一次我固然都难受,可也并没有太往心里去,而且(就我记忆所及)从没想过要把这三串迹象联系起来:被弃置的工地(以及被遗忘的眼光)、“取笑之风”,以及不少旧日朋友拉开距离。我给他们每个人都写过信,没有一个人回信。再说,从那以后,我写给旧友或学生、谈那些搁在我心头的事的信,常常没有回音,这已经不稀奇了。新时代,新风尚——我又能怎样?我只好不再给他们写信。可是(如果你是其中一位)我现在正在写的这封信,将是那个例外——一番再次向你递上的话——这一次接不接受,你自己看吧……

\origpage{L11}某些旧友与我本人拉开距离的最初迹象,如果我没有弄错,要追溯到1976年。也正是那一年,另一“串”迹象开始出现,那是我在回到《收获与播种》之前还要谈的。说得更准确些,这后两串迹象是同时出现的。就在我写下这些的此刻,我觉得它们其实不可分割,说到底,它们是同一现实的两个不同侧面或“面孔”,那一年闯进了我亲身经历的场域。至于我方才正要谈的那个侧面,那是一种系统性的“拒不受理”,\dmark 不声不响、无可辩驳,由一个“毫无裂缝的共识”\fnote{5}{}{这个“毫无裂缝的共识”,在《愚妄与更新》里这儿那儿偶尔被提到,到下一部分《葬礼(1)》里才终于成为一份详尽的见证与一番反思的对象,其中有“第X支队列”即“灵车”,由那些“棺材按语”(第93—96号)以及按语《掘墓人——或整个教团》组成。后者为《收获与播种》的这一部分收尾,同时也构成反思这“第二次吐息”的第一个成果。}专门留给1970年以后那少数几个学生及其同类——他们凭着各自的工作、工作方式与灵感,清清楚楚带着我的影响的印记。
% ---------------- p.80 ----------------
\origpage{L12}也许正是借这个机缘,我第一次察觉到那阵“隐约取笑的气息”,它通过他们,针对的是某种数学风格与某种数学进路——一种(按照当时在数学建制里显然已成普遍的共识)本不该存在的风格与眼光。

这里同样,是某种在无意识层面上被清楚感知到的东西。就在同一年,同一个荒诞的剧本(演示的是:一篇显然出色的博士论文竟无法发表)连续重复了五次,带着马戏团噱头般执拗的滑稽,它最终还是硬闯进了我自觉的注意。如今回想,我明白:当时某种现实正以一种善意的执拗在“向我示意”,而我在装聋:“喂,睁眼看看吧,大个子,留神点你眼皮底下正发生的事,这事跟你有关,真的……!!”我稍稍振作了一下,看了一眼(就那么一小会儿),半是发懵,半是心不在焉:“哦,是啊,好吧,有点怪,看来这儿是有人跟谁过不去,有什么东西横竖过不去,而且还这么齐整,说实话简直难以置信!”

这事简直难以置信到了这个地步,以至于我赶紧把那个噱头、把那场马戏都忘了。确实,我并不缺别的有趣的事可做。可这并没挡住马戏在随后几年里继续提醒我记起它——不再是噱头的调子,而是暗中以羞辱为乐的调子,或者照准脸上抡一拳的调子;只不过这里的人物都是些高雅人士,那一拳在这里也就必然采取更高雅的形式,但一样有效,全凭那些高雅人士的创意……

我感受到如同“照准脸上抡来的一拳”(是别人挨的)的那个事件,发生在1981年10月\fnote{6}{}{这段插曲见按语《棺材之三——或有点过于相对的雅可比簇》(第95号),尤其第404—406页。}。那一次,而且是从种种新精神的迹象不断向我袭来以来的第一次,我被打中了——大概比拳头落在我自己身上还要更重;那一拳是别人挨的,而那个人是我所疼爱的。他多少有点像个学生,而且是一位天分非凡的数学家,刚做出了一些漂亮的东西——不过这是个细节,说到底。不是细节的,是当时我有三个“从前的”学生,直接与那件事站在一边,而当事人(并非没有道理地)把那件事当作一种羞辱与侮辱。另外两个旧日的学生,早已有机会居高临下地对待过他,像有钱人打发一个穷酸书生那样\fnote{7}{}{在上一则脚注所引的那则按语里顺带谈到过。}。还有一个学生,三年后也会跟上同样的步子(而且依旧是“照脸一拳”的风格)——不过这一点我当时当然还不知道。当时叫住我的东西,已经绰绰有余。仿佛我那段从未被审视过的数学家的过去,忽然借着我五个旧日学生中某些人的嘴脸,朝我做了一个狰狞的鬼脸——这些人如今都成了要人,有权有势,而且倨傲……

那本该是再也合适不过的时机,去停下来,去探问那突然以如此暴力向我招呼的东西究竟是什么意思。可在我身上的某个地方早已决定(而这事从来不必说出口……):那段“从前”的过去,说到底与我无关了,我没有必要在那里停留;如果它现在似乎用一种我认得再清楚不过的声音——轻蔑时代的声音——在招呼我,那一定是认错了人。然而,我连着几天、也许几周,被焦虑捆得紧紧的,却连承认它都不肯。(只是去年,借着写《收获与播种》使我重回到这段插曲,我才终于得知有这份焦虑,它一出现就被控制住了。)我没有去确认它、探问它的意义,而是四处活动,我写信给
% ---------------- p.81 ----------------
\origpage{L13}左右各方,写“那些非写不可的信”。当事人甚至还费心给我回了信,自然是些含糊其辞的信,根本不触及任何实质。浪头终于平息,一切恢复秩序。直到去年,我几乎没再想过这事。可这一次,它留下了一道伤,或者更像一根扎人的刺,让人避着不去碰;一根养着那道伤口的刺,而那伤口只等着愈合……

这大概是我作为数学家的生涯里经历过的最痛苦、最难熬的经验——当我有幸看到(却又不肯真正去认清我眼睛所看见的):“某个我疼爱过的旧日学生或同伴,以悄悄压碎另一个我疼爱的人为乐,并且在那人身上认出了我”。它当时给我的打击,肯定比去年我那些相当疯狂的发见还要重——那些发见(在浮浅的目光看来)也许显得更加不可思议……确实,这经验在同样的调子上引起了另外几次共鸣,只是没那么暴烈,当时也就有点“漏过去了”。

\origpage{L14}这也让我想起:同样是1981年,也是我与旧日学生中唯一一个在我离去后仍保持经常联系的人的关系发生急剧转折的一年;而这个人,十五年来在数学层面上一直是我“优先的对话者”。正是那一年,已经出现几年的“那些表示轻蔑的迹象”\fnote{8}{}{这段插曲见按语《两个转折》(第66号)。}“突然变得如此粗暴”,于是我中断了与他的一切数学交流。那是在前面那件“照脸一拳”的事之前几个月。事后回看,这巧合让我觉得触目;但我不认为当时自己做过任何联系。那些事被归进一个个“分开的抽屉”;何况有人已经宣布过,这些抽屉其实无关紧要——案子已经了结!

这也让我想起:还是这1981年6月,已经开过一场颇为精彩的研讨会,在不止一个意义上值得记住——那场研讨会着实配得上以“反常讨论会”这个抹不掉的名字进入历史(或者说进入历史剩下的那点东西……)。我是在去年5月2日认识它的(更确切地说,是它朝我砸下来的!),就在发现那有血有肉的“葬礼”(4月19日)两周之后——我立刻明白,我撞上了“封神”。当然,是一场葬礼的封神,但同时也是一种轻蔑的封神——轻蔑的对象,是我们这门科学存在两千多年来数学家伦理那条心照不宣、不可动摇的根基:那条最基本的规矩,即不得把从别人那里拿来的想法与结果当作自己的东西拿出来。而我此刻注意到这两个事件在时间上那引人注目的巧合——它们看上去性质与分量都大不相同——不禁愕然:这里显露出的,正是对一个人的敬重与对一门技艺、一门科学的基本伦理规则的敬重之间那种深刻而明显的联系;正是这些规则,使这门技艺的从事不成为一场“混战”,使那些以擅长此道著称、并在此定调的人整体上不成为一个毫无顾忌的“黑手党”。不过我又抢先了……

\section{(5) 旅行}

我想,关于我的“重返数学”、并一针引一线地关于《收获与播种》的写作所处的来龙去脉,我大致已经绕了一圈。是去年3月底,在《愚妄与更新》的最后一节(《一段过去的重量》,第50号)里,我才终于想到要追问这次意外重返的缘由与意义。说到“缘由”,最强的一条想必是那种既含混又迫人的感觉:那些强健有力的东西,我曾
% ---------------- p.82 ----------------
\origpage{L15}当年托付给一些充满爱意的手——“它们却在坟墓里,与风、雨、阳光的恩泽隔绝,在我看不见它们的这十五年里萎坏着”\fnote{9}{}{引自按语《坟前的旋律——或自得》(第167号),第826页。}。我一定是渐渐地明白了,而且此前从未想到要对自己说出来:最终撬开那些虫蛀的木板、放出那些被囚禁的活物的,不是别人,正是我自己;那些活物生来不是为了让它们在封闭的棺材里烂掉,而是为了在开阔的空气中绽放。而那些围绕这些包了衬垫、臃肿不堪的棺材(无疑是照着那位可敬的死者的模样……)的假作哀戚与暗含取笑的神气,想必也“终于在我心里唤醒了一根近十年来有点睡过去的斗志之弦”,以及投身混战的欲望……\fnote{10}{}{见《一段过去的重量》(第50节),尤其第137页。(**)。}。

就这样,两年前,一件本来只打算作几天、顶多几周的快速踏勘——对某一处被撇下的“工地”的探看——变成了一部N卷本的数学长篇连载,插进了那套著名的新丛书《反思》(“数学的”,等哪天把那个多余的限定词删掉再说)。而且,从我得知自己正在写的是一部要出版的数学著作的那一刻起,我也就知道:除了或多或少合乎惯例的“数学”导论之外,我还要给它加另一篇“导论”,性质更私人。我觉得重要的是,我要解释我的“重返”——那根本不是重返一个圈子,而仅仅是“重返”紧张的数学投入,重返发表我笔下写出的数学文本,为期不定。同样,我也想解释清楚我现在写数学时所怀的精神,它在某些方面与我离去前的写作精神大为不同——那是一种发现之旅的“航海日志”精神。更何况,还有些别的事搁在我心上,它们无疑与这些相关,但我觉得更为本质。我早已认定,我要慢慢来,把该说的话说出来。那些事还含混不清,对我来说,却与我将要写的这些卷书、以及它们将插入其中的《反思》的意义不可分割。绝不该把它们草草塞进去,像为占用了匆忙读者宝贵的时间而道歉似的。如果说《追寻栈》里有些东西值得他、也值得所有人知道,那恰恰是我留待在这篇导论里说的。要是二三十页不够把它们说完,我就写上四十页,甚至五十页,那又何妨——何况我并没有强迫谁来读我……

\origpage{L16}《收获与播种》就是这样诞生的。我1983年6月写下了预定那篇导论的最初几页,那是在写《追寻栈》第一卷中间的一段空档。然后去年2月我又接着干,那时我那卷书已经基本完成好几个月了\fnote{11}{}{其间我花了足足一个月,思考一个伪直线系统的“结构曲面”,那是用一条伪直线相对于这样一个系统的所有可能“相对位置”的全体来表示的。我还写了《纲领概要》,它将收入《反思》第三卷。}。我本以为这篇导论会是个机会,让我把心里还有两三处略微模糊的东西弄清楚。可我丝毫没料到,它会像我刚写完的那卷书一样,成为一趟发现之旅;而且这趟旅行所进入的世界,比我准备在那卷已成的书和随后几卷里踏勘的世界要丰富得多、广阔得多。日子一天天、一周周、一月月过去,我也不太清楚发生了什么,这趟新的旅行就这样继续下去:去发现一段过去(三十多年来被执意回避的……),发现我自己,发现我与这段过去的种种联系;也去发现数学界里某些曾与我亲近、而我了解得如此之差的人;最后甚至,在
% ---------------- p.83 ----------------
\origpage{L17}这股势头之中,还成了一趟数学发现的旅行——因为十五二十年来,我第一次有闲情回到某些我离去时留得滚烫的问题上去\fnote{12}{}{五六十年代,我常常压抑自己想投身这类油水十足、滚烫的问题的欲望,因为被没完没了的基础工作占住了;那些工作,没有人能、也没有人愿意替我接着做,而我离去之后,也没有人有心把它继续下去……}。总之我可以说,我在《收获与播种》的篇幅里进行的,是三趟彼此紧紧交织的发现之旅。而这三趟,没有一趟随着第一千二百来页那个句号而告终。我这份见证将要收拢的回声(甚至包括那由沉默构成的回声……)都将成为这趟旅行的“续篇”。至于它的终点,这趟旅行肯定属于那种从未被走到头的旅行——也许,甚至到我们死的那天也走不到头……

于是我终于又回到了出发点:如果办得到,事先告诉你,《收获与播种》里“讲的是什么”。可也真是,前头那些篇幅,即使我没去找,也已经多多少少告诉你了。也许更有意思的做法是顺着势头往下讲,而不是去“预告”。

\noindent 1985年6月\par

\section{(6) 阴影的一面——或创造与轻蔑}

前面那些篇幅,是上个月借着一段短促的“空档”写下的。其间,我终于把《四种运算》(《收获与播种》的第四部分)最后收拾停当——现在只差把眼前这封信、这封“前信”也写完(它同样摆出一副要长到无法无天的架势……),一切就终于可以付排、可以复印了。我都不再相信这事了,快一年半了,我一直“就要写完”那批出了名的按语!去年2月(还有前一年的6月),当我动手写这篇对一部数学著作来说性质有点不寻常的“导论”时,(我想)主要有三类东西是我当时想说的。首先,我想说明我重返数学活动的意图,说明我写《追寻栈》这第一卷(我刚宣布它写完)时所怀的精神,也说明我打算带着《反思》去进行一趟更广阔的数学踏勘与发现之旅时所怀的精神。从此对我来说,不再是为某个正在孕育中的数学新世界呈上一套细致入微、一丝不苟的基础。那更该是些“航海日志”,工作在其中日复一日地继续,什么也不隐瞒,就跟它真实进行的样子一样,带着它的卡壳与失手,带着它执拗的回头,也带着它突然的跃进——一件日复一日被不可抗拒地拉着往前的工作(尽管有无数意外与不测),仿佛被一根看不见的线牵着——被某种飘忽、执着而可靠的眼光牵着。一件常常在摸索中前行的工作,尤其在那些“敏感的时刻”:某种还没有名字、还没有面孔的直觉,几乎无法察觉地浮上来;或者在某个新旅程的起点,应着最初几个想法与直觉的召唤、追着它们走——它们往往飘忽,不肯被语言的网眼捉住,而恰好是那种能够轻柔地捉住它们的语言,常常还欠着。首先要做的,正是让这样一种语言从一片显然空无的、摸不着的雾霭中凝结出来。那还只是被预感、尚未被瞥见、更谈不上被“看见”和被手指触到的东西,一点点地从不可秤量之物中澄清出来,从它的阴影与雾霭的外套里脱出来,取得形状、血肉与重量……
% ---------------- p.84 ----------------
\origpage{L18}工作的这一部分,样子不起眼,甚至(许多时候)可以说干得糟糕,却同时也是最微妙、最本质的部分——真正说来,正是在那里,某种新东西出现了,靠着一种紧张的注意、一种关切、一种对这脆弱得无限微妙、即将诞生之物的敬重。这是创造的部分,是受孕与在养育之母体温暖的幽暗中缓慢孕育的部分:从那看不见的双重始基细胞起,变成无形的胚胎,又日复一日、月复一月地,在一种幽暗而紧张、看不见也无形状的劳动中,变成一个血肉俱全的新生命。

这也正是发现工作中“幽暗的”部分,那“阴”的或“女性”的部分。与之互补的一面,那“明亮”的、“阳”的或“男性”的部分,则更像抡起锤子或大槌干活,干在磨得锋利的凿子上,或干在钢火淬好的楔子上。(那是些现成可用的工具,效力早已得到验证……)两个方面各有其存在的理由与功能,彼此不可分割地共生——或者说得更好些,它们正是那对不可分离的原始宇宙力量夫妇中的妻与夫;它们不断更新的拥抱,不断重新唤起那幽暗的创造劳作:受孕、孕育与诞生——孩子、新生之物的诞生。

\origpage{L19}在那篇出了名的、私人的、“哲学的”导论里,我觉出需要说的第二件事,恰恰是关于创造工作的本性。多年来我已明白:这本性一般都被忽视,被随手拈来的套话、被祖祖辈辈的压抑与恐惧掩盖了。它究竟在多大程度上确实如此,我此后才渐渐发现,日复一日、月复一月,在《收获与播种》里进行的那场反思与“调查”的全过程中。正是在这场反思的“发令”之初,在1983年6月写下的那几页里,我第一次被一个外表平淡、实则令人目瞪口呆的事实攫住——只消稍微停留一下就会看到:我刚才说到的发现工作中那“最富创造性”的部分,在那些本该呈现这样一种工作(至少呈现它最摸得着的果实)的文本或言谈里,几乎哪儿都不露面;无论是教科书和其他教学文本,还是原创论文与论著,还是口头课程与讨论班报告等等。看来几千年来,自从数学与其他技艺、科学诞生之初起,就存在着某种“沉默的合谋”,围绕着这些“不可告人的劳作”——它们是一切新观念(或大或小)绽放的前奏,那些观念更新着我们对这个世界某一部分的认识,而我们就在这个不断创造着的世界里活着。

说白了,看来对那个侧面、那个阶段——一切发现工作(以及一般创造工作)中最关键的阶段——的认识所受到的压抑,竟如此有效,如此被那些明明亲身了解这种工作的人内化,以至于人们常常会发誓说:即便这些人,也已经把关于它的一切自觉记忆的痕迹都铲除了。有点像在一个清教到极点的社会里,一个女人会从自己的记忆里铲除掉:与每一个她自认为有责任去擤鼻涕、擦屁股的孩子相关的,那使她受孕的一刻(她勉强承受的拥抱),那漫长的几个月妊娠(被当作一件不成体统的事),以及那漫长的几个钟点分娩(被当作一场有点恶心的苦难,最后总算解脱)。

这个比喻也许显得夸张,而且,如果我把它用到我今天还记得的、我二十年前自己所属的那个数学圈子的精神上,也许确实夸张。可在我于《收获与播种》里的反思过程中,我能够认识到——尤其在这最后几个月(随着《四种运算》的写作)以触目的方式认识到——自我离开数学舞台以来,在今天那些我曾熟悉的圈子里、并且(在我看来,至少在很大程度上)在整个数学界里成为律令的那种精神,发生了令人目瞪口呆的败坏\fnote{13}{}{何况这种败坏绝不限于“数学世界”。在整个科学生活里都能看到它,而在科学生活之外,在当代世界的全球尺度上也能看到。朝这个方向所作的初步确认与反思,见开启阴阳之思的那则按语《肌肉与内脏》(第106号)。}。甚至可能,无论凭我极为特殊的数学个性,还是凭围绕我离去时的那些境况,我那次离去都在一场本已在进行的演变中起了催化剂的作用\fnote{14}{}{这就是上页脚注所引那则按语里考察的演变。它与(我本人和我作品的)“葬礼”之间的联系,在《阴的葬仪(阳葬阴之四)》《天赐的时机——或封神》《不认(一)——或召回》《不认(二)——或变形》(第124、151、152、153号)诸按语里出现并得到考察。另见较晚的(在 RS IV 里的)按语《无用的细节》(第171(v)号,(c)部分“那些什么也不像的东西——或干枯”)与《家庭相册》(第173号,c 部分“那唯一者——或默许”)。}——那场演变,我当时丝毫未察觉(我那些同事与朋友里也没有人察觉,也许只有克劳德·谢瓦莱(Claude Chevalley)是唯一的例外)。我在这里尤其想到的这种败坏的那个侧面(它只是众多侧面中的一个\fnote{15}{}{在《收获与播种》里最常处在注意中心的侧面,尤其是在两个“调查”部分(RS II,即“中国皇帝的袍子”;RS IV,即“四种运算”)里,而且也许正是最让我“目瞪口呆”的那个侧面,是这一行当伦理的败坏:它表现为肆无忌惮的掠夺、贬损与舞弊,行之于当下某些最负盛名、最出色的数学家之间,而且(在很大程度上)是在众目睽睽之下。至于另外一些更微妙的侧面——它们其实与这一侧面直接相关——我请读者去看已经引过的那则按语(第173号 c 部分)《那些什么也不像的东西——或干枯》。})是对以下这一切心照不宣的轻蔑,乃至毫不含糊的取笑:在数学里,凡是不像锤子砸在铁砧上或凿子上那种纯粹劳作的——那些最微妙的创造过程(而且往往样子最不起眼)所受的轻蔑;一切灵感、梦、眼光(无论它们多么有力、多么丰饶)所受的轻蔑;乃至(推到极处)一切观念所受的轻蔑,无论它被构想、被表述得多么清楚:凡是没有白纸黑字写出来、没有以纯粹而硬邦邦的命题形式被写出、被登记、被编目、可以进入那些灌满我们巨型计算机无穷记忆的“数据库”的东西。
% ---------------- p.85 ----------------
\origpage{L20}(借用 C. L. 西格尔(C. L. Siegel)的一个说法\fnote{16}{}{这个说法在上一页脚注所引的那则按语里被引用并加以评论。})曾经有过一次非同寻常的“压平”,数学思想的一次“收缩”:它被剥去一个本质的维度,剥去它全部的“阴影的一面”、那“女性”的一面。诚然,由于一种祖传的传统,发现工作的这一面在很大程度上一直隐而不显,可以说从没有人谈过它——但与梦的深处源头(那些滋养伟大眼光与伟大构想的源头)的活的接触,此前(据我所知)从未丧失过。看来我们从今起已经进入了一个干枯的时代:在那里,这一源头不是被从中心处切断,而是通往它的路被堵死了,凭的是普遍轻蔑那不容上诉的判决,以及嘲弄的那些代表们。

我们似乎正在接近这样一个时刻:在每个人身上,不仅将被铲除一切与源头相近的工作、那“女性气质”的工作(被讥为“稀里糊涂”、“软塌塌”、“不牢靠”——或者在另一极端,被讥为“琐屑”、“小儿科”、“瞎轰炸”……)的记忆,而且这工作本身连同它的果实也将被连根拔除:那构想、酝酿并生出新概念与新眼光的所在。那也将是我们的技艺之实践被缩减为干巴巴、空洞无物的脑力“举重与哑铃”表演的时代,缩减为在竞赛中“攻破”(“难度有口皆碑”的)问题的种种壮举不断加码的时代——那是一个“超猛男”式的、亢奋而贫瘠的肥大的时代,接在三百多年创造性的更新之后。

\section{(7) 敬重与坚毅}

可是我又离题了,预先说出了反思教给我的东西。我原本是从双重的意图出发的,这双重的意图在反思开始之前就已清清楚楚在我心里:一是作一份“意图声
% ---------------- p.86 ----------------
\origpage{L21}明”,二是(与前者紧密相连,正如刚才显出的那样)就创造工作的本性说点什么。然而还有第三个意图,它在自觉层面上想必没那么清楚,却回应着一种更深、更本质的需要。它是由那些有时令人不知所措的“招呼”引起的,那些招呼从我的数学家的过去向我传来,借着那些曾是我的学生或朋友的人的声音(或者至少是其中不少人的声音)。在表皮层面上,这需要表现为一种想“把肚子里的东西倒空”、想说几句“不中听的实话”的欲望。但更深处,肯定是一种终于要去认识一段过去的需要,而那段过去我此前一直选择回避。首先正是从这种需要里,生出了《收获与播种》。这场漫长的反思,就是我对自身这种求知冲动、对那从外部世界、从我头也不回地离开的“数学世界”不断重新传来的招呼,日复一日作出的“回答”。除了《愚妄与更新》最开头那几页——构成其最初两章(《劳动与发现》与《梦与做梦者》)的那几页——之外,而且从接下《恐惧的诞生》(第18页)的那一章起,伴着一份根本不在计划之内的“见证”,我认为,正是这种认识我的过去并完全承担它的需要,成了写作《收获与播种》时起作用的主要力量。

\origpage{L22}从数学家的世界传到我这里、并在《收获与播种》全过程中(尤其是在第二、第四部分所进行的“调查”中)以新的力量反复落回我身上的那种招呼,一上来就戴上了自得的面具,要不就是(经过“精细调配”的)鄙夷、取笑或轻蔑的面具:有时是针对我,而(更主要地)是针对那些竟敢受我启发的人(他们当然没有料到等着他们的是什么),那些被某道心照不宣、不可更改的判决“归了类”,算作与我有一手的人。于是我又一次看见那条“明显”而深刻的联系在此显现:对他人这个人的敬重(或缺乏敬重);对创造之举、对其某些最微妙最本质的果实的敬重;最后还有对科学伦理那些最明显的规则的敬重——那些规则扎根于对自己和他人最基本的敬重,我忍不住要称之为我们这一行当里的“体面规则”。这些无疑都是某种基本的、本质的“自我敬重”的各个侧面。如果我试着用一句斩钉截铁的话,来总结《收获与播种》就某个曾是我的世界、我曾与之认同二十多年的圈子所教给我的一切,我会说:那是一个失去了敬重的世界\fnote{17}{}{这里同样,这个说法不只适用于某个有限的圈子——在那里我有充分机会就近看到这事——我觉得它概括了当今整个世界某种普遍的败坏。(可与第 L19 页脚注比较。)在《收获与播种》里所进行的“调查”之总结这一更有限的框架内,这个说法出现在去年4月2日那则《敬重》(第179号)里。}。

这件事在此前那些年里已经强烈地被感觉到,即便没有被说出来。它只是不断得到确认、变得明确,总是出人意料,有时令人目瞪口呆,贯穿《收获与播种》的全过程。它清楚显现在这样一个时刻:一种“哲学的”、一般的反思,忽然变成了个人的见证(就在开启前面提到的《恐惧的诞生》那一章的《受欢迎的异乡人》一节(第9号,第18页)里)。

然而这种感知并不是以尖锐或苦涩的怨责的调子出现的,而是(凭写作的内在逻辑、凭写作所唤起的另一种态度)以一种追问的调子出现的:在这场败坏里,在今天我所确认的这种敬重的丧失里,我自己那一份是什么?正是这个追问,贯穿并支撑着《收获与播种》这第一部分,直到它最终化解为一个清楚而毫不含糊的确认\fnote{18}{}{见《竞技数学》与《旋转木马收摊了》两节(第40、41号)。}。在此以前,这场败坏对我来
% ---------------- p.87 ----------------
\origpage{L23}说,像是忽然“从天而降”,无法解释,因而更加令人愤慨、难以容忍。而在反思的过程中,我发现它一直悄悄地进行着,想必没有人能在自己周围或自己身上察觉它,贯穿整个五十与六十年代,也包括我本人。

对这个卑微事实的确认——它想必明显得很,平淡无奇——标志着见证中第一个关键转折,以及一种立竿见影的质变\fnote{19}{}{次日,见证便深化为一场对我的沉思,并在随后的几周里保持这种特殊的品质,直到《收获与播种》这“第一次吐息”结束(以《一段过去的重量》一节,第50号)。}。这是我关于自己数学家的过去、关于我自己所要学的第一件本质的事。这种对我在普遍败坏中所负责任的认知(其敏锐程度随反思的时刻而或深或浅),贯穿《收获与播种》的全过程,成了一条底色、一个提醒。尤其当我的反思显出是在调查一个时代的失宠与不公时,情形更是如此。与理解的欲望、与那种驱动并推动一切真正的发现工作的好奇心一道,正是这种卑微的认知(它一再在途中被忘掉,又不顾一切地浮上来,而且总在最料不到的地方……)使我的见证免于(我相信)沦为对世界忘恩负义的干瘪怨责,甚至免于沦为与某些曾是我的学生或朋友(或两者兼是)的人“算总账”。这种对自己的不徇情,也给了我那种内心的平静,或者说那种坚毅,使我免于徇情于人的陷阱,哪怕只是免于一种虚假“谨慎”的陷阱。凡是我认为在反思的某一刻该说的话——无论关于我自己,还是关于我某位同事、某个旧日学生或朋友,还是关于一个圈子、一个时代——我都说了,从不需要去硬压自己的迟疑。至于那些迟疑,每次我只消仔细审视它们,它们便消失得无影无踪。

\section{(8) “我的亲近者”——或默契}

\origpage{L24}这封信里,我不打算把《收获与播种》写作中(或它某一阶段中)的所有“要紧时刻”(或所有“敏感时刻”)逐一巡览\fnote{20}{}{关于《收获与播种》前三部分全体的一个简短回顾与总结,你会在这两组按语里找到:《暮熟的果实》(第179—182号)与《发现一段过去》(第183—186号)。}。我只想说:这件工作里有四个界限分明的大阶段,或者说四口“气”——像呼吸的几口气,又像一道涌起的浪涛里接连而来的几道浪;那道浪涛我也不知是怎么从一片广袤、无声、静止而又流动、无界无名的大海里涌出来的,那海便是“我”;或者更确切地说,是一片无限辽阔、无限深邃的大海,远过于它所承载、所养育的那个“我”。这几口“气”、这几道“浪”,化成了《收获与播种》目前已经写成的四个部分。每一道浪来了,都不是我呼唤来的,也丝毫未曾预料;而任何时候我都说不清它要把我带到哪里,又何时结束。而当它结束了、新的一道浪已经接了上来,我还总以为自己在某一段冲势的收尾(那也终将,在最终的最终,是《收获与播种》的收尾!),其实我早已被同一个宏大运动的另一口气抬起来、托着往前走了。只有事后回看,这运动才清楚地显出,才毫不含糊地显出:在那被当作行动、被当作涌动而经历的东西里,有一个结构。

而且,那运动当然没有随着我在《收获与播种》上打的那个句号(完全是暂时的!)而结束,
% ---------------- p.88 ----------------
\origpage{L25}也不会随着我写给你的这封信的句号而结束;这封信也是那运动的一个“时段”。它也不是在某一个1983年6月或1984年2月的某一天诞生的——那天我坐到打字机前,要写(或接着写)某部数学著作的某篇导论。它诞生了(或者更确切地说,它重又诞生了……),就在冥想进入我生命的那一天……

可是我又离题了,任由即时的形象与联想把我带走(甚至卷走……),而不是乖乖守着一条“话头”、一条预定的线索。我今天本来是要接着讲,尽可能简短地讲去年4月那场“发现葬礼”的故事:当时我已有两个星期相信《收获与播种》已经写完——而就在那三四个星期的工夫里,种种发现怎样一个接一个、一个比一个更大更不可思议地朝我劈头砸下来——大得、疯得,以至于此后好几个月,我都极难“相信我健全官能的证词”,极难摆脱面对明证时那种暗暗的难以置信\fnote{21}{}{我试着用那个故事《中国皇帝的袍子》来表达这种困难,见同名按语(第77号);后来又在那则《尽到的义务——或真相的时刻》(第163号)里回到这一点。}。这种隐秘而顽固的不信,直到去年10月(在发现“光彩夺目的葬礼”六个月之后)才终于消散,那是在我的朋友、旧日学生(不错,是隐形的)皮埃尔·德利涅(Pierre Deligne)来我家拜访之后\fnote{22}{}{这次拜访的叙述,见我刚引的那则按语(在上页脚注中)。}。我这才第一次面对那场葬礼,不再是借文本的中介——那些文本固然用意味深长的言辞对我说着贬损、掠夺与对一部作品的屠杀,说着某种数学风格与某种数学进路(以缺席的大师为其人身)的埋葬——而是这一次,以一种直接而可触的方式,带着熟悉的面貌,借着一个熟悉的声音,音调殷勤而天真。葬礼终于就在我眼前,“有血有肉”,带着那副我如今已认得的殷勤而平淡的面貌;但这是我第一次用新的眼睛、新的注意去看它。于是,那个在我前几个月的反思中已显出是我隆重葬仪上的大主祭、是那位“披着祭袍的祭司”、同时又是一场史无前例的“手术”的主要匠人与主要“受益者”、一部被交付给取笑与掠夺的作品的隐形继承人的人,就这样在我面前展开……

这次会面正落在《收获与播种》里“第三道浪”的开端,当时我刚投身那场关于阴与阳的漫长冥想,追着一串飘忽而执着的联想。当时,这段短短的插曲只留下几行回声的痕迹,一笔带过。然而它标志着一个重要的时刻,它的果实要几个月后才清楚地显现。

还有第二次这样面对“有血有肉的葬礼”的时刻。那不过是十天前的事,它又一次“在最后一刻”重新启动了一场总也不断重新开张的调查。这一次,只是一个打给让–皮埃尔·塞尔(Jean-Pierre Serre)的电话\fnote{23}{}{这几乎是逐字引自按语《掘墓人——或整个教团》(第97号,第417页)。}。这场“漫无边际”的谈话,以触目的方式、甚至超出一切预料地证实了我(不过几天前)刚刚长篇大论地对自己讲明白的东西\fnote{24}{}{在同一则按语(第173号)的 c 部分(“那唯一者——或默许”)。},而且几乎是违我本意地讲明白的:关于塞尔在我的葬礼里所扮演的角色,关于他内心里对那件“就在他眼皮底下”发生的事的一种“秘密默许”,而他从不装作看见或感觉到什么。

这里同样,理所当然,谈话是再“酷”不过、再友好不过的;而且塞尔对我这种友好的态度显然也再真诚不过、再真实不过。这并不妨碍这一次我真正看见,或者——我真想写——“摸到”了我终于不得不承认的那种“默许”;它无疑是“秘密的”(如我先前所写),但尤其
% ---------------- p.89 ----------------
\origpage{L26}是殷勤的,正如我当时所能毫无疑义地看到的那样。一种殷勤而毫无保留的默许:让该被埋掉的被埋掉;并且,凡在看来合宜之处,无论用什么手段,都要让一种真实的(塞尔对此有第一手的了解)、不受欢迎的生父,被一个伪造而受欢迎的生父替换掉……\fnote{25}{}{这几乎是逐字引自按语《掘墓人——或整个教团》(第97号,第417页)。}。

\origpage{L27}\pmark 这触目地证实了一年前已经出现的一个直觉,当时我写道\fnote{26}{}{这段引文出自同一则按语(见上页脚注),同在第417页。}:

\begin{quote}
“从这一光照下看\fnote{27}{}{所谓“这一光照”,指的是刚才谈到的那种蓄意,即不惜一切代价除掉“不受欢迎的生父”(甚至“不可容忍的”,这里沿用所引按语中的说法)。},主祭德利涅便不再显得是那个按照决定他自身生命与行动的那些深刻力量、塑造出一种风尚的人,而更像是那“合法的继承人”\fnote{28}{}{德利涅这个“继承人”的角色,是一个既隐形(德利涅已发表的文字里,没有一行能让人猜到他可能从我的嘴里学到过什么),同时又清楚地被所有人感觉到、承认的角色。这正是德利涅两面手法及其“风格”的典型侧面之一:他精熟地在这含混之上做戏,收取这种心照不宣的继承人角色所带来的好处,同时又不认那位已故的师父,并主持着规模宏大的埋葬作业。}的指定工具,服务于一种具有毫无裂缝之一致性的集体意志——那意志执着于一项不可能的任务:把我和我个人的风格从当代数学里抹掉。”
\end{quote}

如果说德利涅当时在我眼里是那种“具有毫无裂缝之一致性的集体意志”的“指定工具”(同时也是头一个、主要的“受益者”),那么塞尔现在在我眼里,就是这同一种集体意志的化身,是它毫无保留之默许的担保人;那默许给予一切数不清的舞弊与欺诈,直至那些大规模的集体蒙骗与肆无忌惮的侵占“作业”,只要这些都有助于那项“不可能的任务”——针对我这点微不足道的、已故的人身,或者针对另一位\fnote{29}{}{我在这里想到的是左格曼·梅布胡特(Zoghman Mebkhout),第一次谈到他是在导论第6节(“葬礼”),然后是在按语《我的孤儿们》(第46号)里,以及在(较晚写成的、在发现葬礼之后的)按语《一场教学的失败(二)——或创造与自得》和《一种不义与无力的感觉》(第44$'$、44$''$号)里。我是在构成葬礼“第七支队列”的那十一则按语《讨论会——或梅布胡特的层与反常》(第75—80号)中,一点一点发现那桩对梅布胡特先驱性工作的不义抹除与侵占的。关于这(第四次、也是最后一次)“运算”,一份更详尽的调查与叙述,构成了“四种运算”这场调查里最充实的部分,用的正是那个理所当然的名字:《封神》(第171(i)至171号诸按语)。}竟敢自称出于我、并且不顾所有人反对而摆出“格罗滕迪克的继承者”架势的人。

\pmark 葬礼之中有许多悖谬而令人不知所措的侧面,其中之一便是:它首先是——甚至可以说是唯独是——由那些曾是我的朋友或学生的人做成的,而在那个世界里,我从不曾知道自己有敌人。我想,首先正是由于这一点,《收获与播种》才与你关系格外近,而我这封正在写的信也才要反过来成为一次招呼。因为如果你是数学家,如果你是他们中的一个——曾是我的学生,或曾是我的朋友——那么你与这场葬礼大概并非无关:或者以行动,或者以默契,哪怕只是以你对我的沉默——对一件就在你家门口台阶前进行的事保持沉默。而如果(说来稀奇)你接受我这些卑微的话、接受它们带给你的见证,而不是闭门不出、把送信人当不受欢迎的人打发走,那么你或许就会明白:被所有人、并且有你参与(积极的参与,或心照不宣的默许)埋掉的,不只是另一个人的作品——那是他与数学相爱的果实与活生生的见证;而且,在这个埋葬(它从不叫出自己的名字……)之下更深、更隐秘的层面上,被埋掉的是你自己存在中活着而本质的一部分,是你原初的认识、爱与创造的力量;你乐于亲手借着另一个人的身体把它埋掉。

% ---------------- p.90 ----------------
\origpage{L28}在我所有的学生里,德利涅占着一个很特殊的位置,关于这一点我在反思中谈得很多\fnote{30}{}{关于这一点,尤其见 RS II 里那十七则一组的按语《吾友皮埃尔》(第60—71号)。}。他曾经是——而且远远是——最“亲近”的一个,也是唯一一个(不论学生与否)把某种宏大的眼光亲密地吸收过来、并据为己有\fnote{31}{}{这种“宏大的眼光”,德利涅确确实实“亲密地吸收过来、并据为己有”了;它对他也曾有过强大的魅力,而且至今仍不由自主地魅惑着他,尽管同时有一股强横的力量推着他去摧毁它,去炸开它根本的统一,去把散落的碎片据为己有。于是,他对一位被否认的、“已故的”师父那种隐形的敌意,正是他自身内部分裂的表现;这分裂深深地印在他离去之后的作品上——那部作品远远停在我所了解的他那相当惊人的才能之下。}的人;那眼光在我身上诞生、长大,远在我俩相遇之前。而在我所有与我共有数学激情的朋友里,最亲近的是塞尔——他同时也多少算个长辈——(同样远远地)最亲近,尤其是因为他在十年间在我的工作里扮演了一个独一无二的角色:某些重大投入的“引爆器”,以及五十、六十年代启发我数学思考的多数有力观念(直到我离去之时)的“引爆器”。他们两人与我这个人的那种很特别的关系,当然与他们各自非凡的才能不无关联;这才能使他们对同代以及其后几代的数学家拥有同样非凡的威望。撇开这些共同点,塞尔与德利涅的性情与作风,在我看来却再相异不过,在许多方面彼此正相对立。

不管怎样,如果说有哪位数学家曾以某种名义“亲近”我这个人、我的作品(而且,更有甚者,是出了名地如此),那正是塞尔与德利涅:一位是长辈,是在某种眼光孕育的关键时期我作品中的灵感之源;另一位是我最有天分的学生,而他反过来把我当作(并且一直当作,不管有没有那场葬礼……)他主要的(也是秘密的……)灵感之源\fnote{32}{}{关于这一点,见上页脚注。}。如果在我离去之后不久,一场葬礼就开动起来(并且正儿八经地变成了“死亡”),并且落实为一支没完没了的、大大小小“作业”的队列,服务于同一个目的,那么这事只可能是在这两位——前长辈与前学生(甚至前“门徒”)——协同而紧密团结的配合下完成的:一位悄悄而有效地主持着作业,同时召集我某些学生归队\fnote{33}{}{这里说的,确确实实是另外那五个以簇的上同调为主题(与德利涅一样)的学生。},他们正巴不得屠杀父亲(以那个肥硕臃肿、嗡嗡饶舌的超级娘们的怪诞可笑形象为靶子);另一位则毫无保留、无条件、无限制地对(四项)作业的继续开“绿灯”(那些作业是:贬损、屠杀、宰割、分赃,分一具取之不尽的尸首……)。

\section{(9) 清空}

\origpage{L29}正如我刚才已经暗示过的,我必须克服相当大的内心抵抗,或者更确切地说,要靠一种耐心、细致而顽强的工作让它们自行消解,才能与某些熟悉的形象分开——那些形象牢牢盘踞、惰性极大,几十年来在我身上(正如在所有人身上,想必也在你身上)取代了对现实直接而细腻的感知,具体说来,是对某个数学世界的感知,而一段过去与一部作品仍然把我与那个世界连在一起。这些形象、这些现成观念里最根深蒂固的一条是:一个国际知名的学者,甚至一个看上去是大数学家的人,竟然会去干大大小小的欺诈,这从一开始就被排除在外(哪怕只是作为例外,更不必说作为一种亲爱的
% ---------------- p.91 ----------------
\origpage{L30}习惯……);或者,即便他(又出于老习惯)自己不伸手去沾,竟会张开双臂欢迎别人搞起来的那些“(有时简直不顾一切体面感)”的作业,而在其中,由于这样或那样的原因,他还能得利。

我身上这种精神的惰性竟到了这个地步:仅仅不到两个月前,在整整一年之久的漫长反思之后,我才终于怯生生地瞥见:塞尔在这葬礼里也许也有份——而这件事如今在我看来已是明证,甚至与最近那次意味深长的谈话无关。如同对那个曾在我起步时善意接纳我的“布尔巴基圈子”的所有成员一样,在他身上尤其如此,对我来说存在着某种心照不宣的“禁忌”,围绕着他这个人。他代表着某种“优雅”的化身——这种优雅绝不限于形式,它还包括一种严格、一种一丝不苟的正直。

在我去年4月19日发现那场葬礼之前,我做梦也不会想到:我某个旧日的学生竟会在执业中做出不诚实的事,无论是对我,还是对任何人;而如果这假设落在我最有天分的那一个、也最亲近的那一个身上,在我看来就最为荒诞不经!然而,早从我离去之时起,直到其后这些年、直到今天,我一直有充分机会看清他与我的关系是何等分裂。也不止一次,我见他(看来只为取乐)运用那种使人气馁并羞辱人的权力,只要时机合适。每一次我都深受触动(大概比我愿意向自己承认的更深……)。这些都是某种深刻失序的相当意味深长的迹象;而这失序(我有充分机会确认)绝不限于他一个人,即使在我旧日学生中最狭小的圈子里也是如此。这种失序,由于丧失对他人的敬重,其触目与深刻程度,并不亚于那种表现为所谓“职业上的不诚实”的失序。尽管如此,\dmark 发现这样一种不诚实,对我仍是一个十足的意外、一记重击。

在那场令人目瞪口呆的揭示之后——随后又是一整串同一路数的“瀑布”般的发现——的几周里,我渐渐明白:在我某些学生中间\fnote{34}{}{见上页脚注。},某种舞弊早在我离去之前的那些年里就已经开始了。恰恰在我最有天分的那一个身上,这一点格外触目——就是他,在我离去之后定了调子,并且(如我刚才所写)“悄悄而有效地主持着作业”。隔了将近二十年回看,这舞弊如今在我看来已是明摆着的事,简直“扎眼”。如果说当时我选择闭上眼不看正在发生的事,一心在那个“一切都只是秩序与美”的世界里(如我乐于想象的那样)追逐“白鲸”,那么我今天确认:我当时未能承担起我应负的责任——对那些在我身边学一门我所爱的行当的学生应负的责任;那行当还不只是一门手艺,也不只是养成某种“嗅觉”。由于对出色的学生徇情,由于我乐于(凭一道心照不宣的判决)把他们当作“特殊的存在”、当作无可怀疑的人,我当时就以我那一份\fnote{35}{}{那份“贡献”,尤其见于按语《特殊的存在》(第67$'$号),以及《攀升》与《含混》两则按语(第63$'$、63$''$号),后来又(在一种略有不同的光照下)见于按语《逐出》(第169号)末尾。另一类“贡献”见于《愚妄与更新》,那是对天分不那么出色的年轻数学家的自得态度。这种对普遍败坏中所负一份责任的认识,在《竞技数学》一节(第40号)里达到顶点。}促成了我今天看到一个世界、以及一些曾为我所爱的人中间蔓延开来的那种(在我看来史无前例的)腐败。

当然,鉴于那些东西惰性极大,要与之分离——与人们习惯称之为“幻觉”的东西分离(那语调里不无几分惋惜……),而我宁愿称之为
% 下接 part1-ch03b.tex

% ================================================================
%  第 3 章 一封信 (后半: 原书 p.92-108, 边码 L31-L56)
%  源文: _work/src1/task-ch03b.md
% ================================================================
%
% ---------------- p.92 ----------------
\origpage{L31}现成的观念;关于我自己,关于一个我从前与之认同的圈子,关于一些我爱过、也许如今依然爱着的人——与这些观念“分离”,或者更确切地说,让它们从我身上自行脱落。这是一件工作,确确实实是工作,但从来不是一场搏斗——这工作带给我的,除了许多别的可贵之物,还有时而出现的悲伤时刻,却从来没有一刻悔恨,也没有一刻苦涩。苦涩是回避一种认识的手段之一,是回避一段经历所传消息的手段之一;是以关于世界、关于他人的另一种“幻觉”(不妨说是负片式的)为代价,把自己维持在对自身的一种顽固幻觉之中。

\origpage{L32}我是没有苦涩、也没有悔恨地看着这些现成的观念一件一件从我身上脱落,像卸下一件件累赘、甚至压人的重负——它们曾为我所“珍爱”,出于积习,也因为它们“从来”就在那里。它们确已变成了某种第二本性。但这“第二本性”并不是“我”。把它们一块一块地剥离,并不是撕裂,甚至也不是那种被剥夺了珍视之物的人的失落。“剥离”一词,在我这里是作为工作的报偿与果实而来的。它的标志,是一种当下即至的、有益身心的轻松,一种受欢迎的解放。

\section{(10) 运动中的四道浪}

果不其然,这封信与我动笔时所设想的全然不同。我本来主要是想就“葬礼”做一份小小的“概述”:大致情形如此如此,你信也罢不信也罢(我自己当初也几乎难以相信……),可事情就是这样,无可置疑,不管你乐意不乐意,某年某月某日某页,白纸黑字印在某期刊某本书上,去看就是了——再说,《收获与播种》里一切都已逐条拧开来细讲;见“四种运算”诸注——爱信不信!你要是宁可不去读我,自有别人替你去读……

结果这些全都没有。可这封信已经写到三十页的坎上了,而我原计划统共五六页。也不是我存心如此:一页一页写下来,我被引着对你讲的,竟是那些最本质的东西;而我当初那么急不可耐要倒空的“口袋”(这一回它可真是明明白白摆在头几页上!),到现在还没解开呢!它甚至不再挠我的手指了,那点欲望在路上就消散了。我明白了:这里不是地方……

说实话,《收获与播种》的第四部分(也是最长的一部分),名为“葬礼(3)”或“四种运算”,原本出自一则“按语”,当初正是打算做成一份小小的“概述”,好把去年那场突袭式(而且风风火火)的调查所揭示的东西大致勾勒一下——那场调查续在第二部分里(“葬礼(1)”,或“中国皇帝的袍子”)。我以为那不过是一则五页或十页的“按语”,不会更多。结果一环扣一环,调查又被它带着重新跑了起来,写成了将近四百页——几乎是我本该为之做摘要、做总结的那部分的两倍!于是那份小小的概述至今仍然缺着,而《收获与播种》里已有六百页是专门调查“葬礼”的。这确实有点蠢。不过,在把稿子交给印刷厂之前,总还来得及把它补进《导论》的第三部分(那部分早就不是在乎多十页二十页的时候了)。

《收获与播种》的五个部分(最后一部分尚未写完,恐怕几个月之内也写不完),是(三道)“冥想”之浪与(两道)“调查”之浪的交替。这里仿佛以一种缩影的方式,映照出我最近九年的生命——那九年同样也是“浪”的交替,那些浪涌自如今支配我生命的两种激情:冥想的激情与数学的激情。而说实话,\origpage{L33}%
% ---------------- p.93 ----------------
我方才用“调查”这个信手拈来的名字所称的那两个部分(或“两道浪”),恰恰是直接从我作为数学家的过往的扎根处涌出来的,被我对数学的激情、以及扎根于这激情之中的种种自我执念所驱动。

第一道浪,“愚妄与更新”,是与我的数学家过往的第一次相遇,由此引出对我当下的一次冥想——我刚刚发现,这当下正扎根于那段过往之中。诚然,这一点绝无半点预谋;可这一部分为《收获与播种》此后全部篇幅定下了“基调”,它像是一种内在的、天赐而不可或缺的准备,好让我承受紧随其后的“光华四射的葬礼”之发现——那是在第二道浪里,“葬礼(1)——或中国皇帝的袍子”。

其实,与其说这是一场“调查”,不如说这就是那场发现逐日展开的故事:它对我是怎样的冲击,我如何竭力面对这样不打招呼就劈头砸下来的东西,如何设法用我的亲身经历去安放那难以置信之物,直到它终于变得为我所熟悉,好歹让它变得可以理解。这一运动引出一个最初的、暂定的终点,在按语《掘墓人——或整个教团》(第97号)里:这是第一次尝试,要在一个多年来、如今更比任何时候都更尖锐地显出“对常理的可怕挑战”之面貌的东西里,辨认出一种解释、一种意义!

这同一场第二运动,也引出一次“病中插曲”\fnote{36}{}{这次插曲是两则按语的主题:《事件——或身体与精神》和《陷阱——或轻易与耗竭》(第98、99号),它们开启了名为《亡者(始终没有故去)》的“第十一支队列”。},逼我绝对静养,三个多月里中止了一切智力活动。那正是我自以为又一次即将把《收获与播种》收尾的时候(只差最后几件“后勤”活儿……)。去年九月底前后恢复常态活动,我正准备终于给那些还陷在困境里的按语做最后的收尾,仍以为只消再添两三则收尾的按语,其中一则讲我刚经历的那次“健康事件”。结果,一周又一周、一月又一月,又来了上千页——比已经写好的还多出一倍——而这一回,显然我还是没写完\fnote{37}{}{“始终没有写完”——哪怕只因为后面还应有一道第五部分,而在写这几行的时候,它还没有完成。}!

事实上,在这段漫长的中断期间,我几乎与一种(在搁下它时)再热不过(甚至滚烫!)的质料失去了接触;这几乎强迫我以新的眼睛回到那质料上去——如果我不想愚蠢地“了结”一个“程序”的最后收尾,而那个程序我早已与它断了活的联系。

第三道浪就是这样诞生的,诞生在《收获与播种》这一片浩瀚的运动之中——那是一道漫长的“冥想之浪”,主题是阴与阳,是事物的动态与人存在中“暗”与“明”的两面。这道冥想出自一种愿望:更深入地理解在“葬礼”中起作用的那些深层力量;然而它从一开始就获得了自身的自治与统一,并且一上来就朝向最普遍的东西,同时也朝向最私密个人的东西。正是在这道冥想之中,我发现了这件事(说实话它显而易见,只要人肯把问题提出来):在我自发地发现事物的进程里,无论在数学还是在别处,“基调”都是“阴”的、“女性”的;而且尤其重要的是,与最常见的情形相反,我始终忠实于我身上这种原初本性\fnote{38}{}{不过,这种“对我原初本性的忠实”绝非彻底的。长久以来,它仅限于我的数学工作;而在别的一切地方,尤其在我与他人的关系里,我随大流,抬高并优先看待我身上那些被感受为“男性”的特征,压抑那些“女性”的特征。关于这一点,在《一个生命的故事:三乐章的循环》(第107—110号)那一组按语里有相当详细的讨论,那组按语差不多是《阴阳之钥》的开篇。},从未让它弯折或修正,以
%
% ---------------- p.94 ----------------
\origpage{L34}迁就周遭圈子里被尊奉的主导价值。这一发现起初在我眼里只是一种有趣的好奇。只是渐渐地,它才显出是理解“葬礼”的一把关键钥匙。而且——这一点在我看来意义还要大——如今我极其清楚、不留丝毫疑问地看到:如果说,以绝非出众的智力天赋,我竟然能始终在数学工作中使出自己的全部本事,创造出一部作品,孕育出一种广大、有力而丰饶的视野,那么这完全要归功于这种忠实,归功于这种全不理会去迁就规范的心态——靠了它,我才怀着完全的信任委身于原初的求知冲动,不去削它、砍它,不减去任何造就其力量、其精细、其不可分割之本性的东西。

\origpage{L35}然而,在这道名为“葬礼(2)——或阴阳之钥”的冥想里,处在注意力中心的并不是创造力及其源头,而毋宁是“冲突”,是创造力的阻塞状态,是创造的能量因心灵中敌对力量(多半是隐秘的力量)的对抗而涣散。暴力的那些面貌——那种(表面上)“无端的”、“为了好玩”的暴力——曾不止一次让我在“葬礼”中不知所措,并让一大堆相似的亲身处境重新浮现。对这种暴力的经验,在我一生中就像是“冲突经验那坚硬的、不可化约的核”。我此前从未直面过这样一个可怕的谜:这种暴力在一般人类存在中、尤其在我自己的存在中,其存在本身与普遍性究竟是怎么回事。整个下半部(即“阴”的一面,或“衰”的一面),注意力都集中在这个谜上。正是在冥想的这一部分里,对“葬礼”之意义、对其中表达出的那些力量,一种更深的看法逐渐浮现出来。这大概也是《收获与播种》在自我认识层面上最丰饶的一部分:它让我接触到一些要害的问题与处境,并且恰恰让我感觉到那种“要害”的性质——而直到去年,这一点还一直被回避着。

在这段关于阴阳的没完没了的“离题”走完之后,我发觉自己大体上仍旧剩着那“两三则按语”要写(顶多加一两则,其中一则早已有了现成的名字:“四种运算”……),写完了就与《收获与播种》两清。

后来发生的事大家都知道了:这“最后几则按语”最终成了《收获与播种》最长的一部分,将近五百页。这就是这场运动的“第四道浪”。它也是“葬礼”的第三即最后一部分,我给它取名“四种运算”——这也是那一组按语(“对一具尸体做的四种运算”)的名字,那组按语构成了这第四次思考吐息的核心。在《收获与播种》里,这一部分才是严格意义上的“调查”部分——不过还带点调味:这场调查并不局限于纯粹的“技术”方面、即“侦探”方面;这里的思考,如同《收获与播种》别处一样,首先是由求知与求理解的欲望驱动的。这里的调子当然比“葬礼”第一部分要“带劲”得多——那会儿我还有点揉着眼睛、自问是不是在做梦呢!话虽如此,一页页披露出来的事实却常常来得正是时候,让许多此前在这里那里只是一带而过的东西当场得到说明,落实到具体而触目的例子上。也正是在这一部分里,数学上的离题占了重要位置,那是被一次重新接上的接触激发起来的(出于调查的需要)——与一种我十五年未见的质料重新接触。在光谱的另一端,还有对吾友左格曼·梅布胡特(Zoghman Mebkhout)遭遇的现场叙述(这一部分即献给他):他落到一伙手段高强、毫无顾忌的“黑手党”手里,而当初他投身于那个(固然引人入胜、表面看来又无伤大雅的)课题——各类簇的上同调——时,根本没有梦到过这些。为了给
%
% ---------------- p.95 ----------------
\origpage{L36}这整个“调查”部分那纠缠错综的按语、副按语、副副按语……的迷宫提供一条简明的线索,我请你去看目录(第167$'$号至1767号),以及这一叠里的第一则:《侦探——或玫瑰色的人生》(第167$'$号)。不过我要提一句:这则按语注明4月22日,后来有点被“事态甩在了后面”——因为一波三折,这场我当时以为(基本上)已经查到底的调查,又拖拖拉拉地继续了两个月。

\origpage{L37}这第四次吐息一口气延续了四个多月,从二月中旬直到六月底。主要是在思考的这一部分里,靠一种细致而顽强的“档案工作”,日子一天天过去、页面一页页累积,“葬礼”的现实才逐渐建立起一种具体的、可触的接触;总之,我才好歹与它“熟识”起来——尽管它在我身上引起的(并且继续引起的)那种发自脏腑的拒绝反应,一直阻碍着真正的认识。这段漫长的思考,以对德利涅来访的回顾开篇(本信中已经提到过),而以关于我与塞尔的关系、关于塞尔在“葬礼”中的角色的一番“最后一刻”的思考收尾\fnote{39}{}{见按语《家庭相册》(第173号)的 c、d、e 三部分,其中最后一部分注明的日期是6月18日(正好是十天前)。只有一则按语或按语片段日期更晚(即《一场屠杀中的五篇论文——或孝道》,第1767号,注明次日,6月19日)。你会注意到,在《收获与播种》这第四部分、即“调查部分”里,与别处相反,按语往往按逻辑顺序而非时间顺序相随。例如“葬礼”最后两则按语(构成结尾的《自深处》)注明4月7日,比我方才引的那则按语早两个半月。我还要指出:除了《葬礼(3)》真正的“调查”部分(第167$'$—1767号按语)——那部分构成殡葬大典的“第五个时辰”(《阴阳之钥》是第二个时辰)——按语都按写作顺序相随,极少例外。}。此前那个空缺,就在于我暗中把塞尔“排除在外”,借着我已说过的那个“禁忌”——如今我觉得,这或许正是直到上个月为止我对“葬礼”的理解中留下的最严重的空缺;于是这番“最后一刻”的思考,反倒成了《收获与播种》这“第四次吐息”带给我的最重要的东西:让我对“葬礼”及其间表达出的力量,有一种不那么单薄、更加充实的领会。

\section{(11) 运动与结构}

我想,关于《收获与播种》我有什么最要紧的话想对你说,都已兜了一圈说完了;为了让你先知道“这是怎么回事”,我说的肯定已经足够,足够让你判断:接下去那一千多页的信,是否“与你有关”——因而,你是否还要继续读下去。若答案是“是”,那么我觉得,再就《收获与播种》的形式补几句说明(尤其是实用方面的)是有用的。

这形式是某种精神的映照与表达,我竭力想把那种精神在前面的篇幅里“传递”出来。与我过去的出版物相比,如果说《收获与播种》里(以及它所出自的《追寻栈》里)出现了一种新的品质,那大概就是自发性。当然,这里有贯穿的线索,有大的追问,它们赋予整个思考以连贯与统一。然而思考是逐日进行的,没有预先立定的“纲领”或“计划”,从来没有什么“应该证明的东西”事先定给我。我的意旨不是证明,而是发现,是往一种未知的质料里更深地走进去,是让那些还只是被预感、被猜疑、被瞥见的东西凝结起来。我可以毫不夸张地说,在这件工作里,没有一天、也没有一夜的思考是在“预料之中”的范围内展开的——无论就观念、意象、联想而言,那些东西都是我坐到
%
% ---------------- p.96 ----------------
\origpage{L38}白纸面前、要在上面顽强地追一条线索、或者接上另一条刚冒出来的线索时,已经在那里的。每一次,思考中出现的东西都与我事先能够预言的不同——假使我竟敢勉力去预先描述我以为在眼前看见的东西的话。最常见的是,思考走上了开头全然未料的道路,通向同样未料的新风景。然而即便它沿着一条大致预定的路线走,几小时里旅程向我显示的东西,也与我上路时对它抱有的形象相去之远,正如一处真实的风景——带着它清凉的阴影与温暖的光的嬉戏,随着散步者脚步而微妙变换的透视,还有被一阵吹得草舞、吹得芦苇歌唱的微风送来的无数声响与无名芬芳……——这样一处活的、不可捉摸的风景,与一张明信片相去之远,无论那明信片多么美、多么成功、多么“准确”。

是那种一口气追下去的思考——一天之内、或一夜之内——构成不可分割的单位,可以说是活的、个体的细胞;在整片思考(《收获与播种》在此即是)之中。整场反思之于每一个这样的单位(或那些共同构成旋律的“按语”\fnote{40}{}{最初写《愚妄与更新》时,“按语”这个名字对我就是“注解”的同义词,起的是脚注的作用。出于排版上的方便,我宁愿把这些注解放到正文末尾(第1至44号按语,第141页与第171页)。这样做的原因之一是:这些“按语”或“注解”中有一些要占上一页乃至好几页,甚至比它们本该评论的正文还要长。至于思考“初稿”那些不可分割的“单位”,我当时因为找不到更好的名字,就把它们叫作“节”(比“段落”不那么令人生畏!)。\par 这一情形和文本的结构,随着下一部分而改变;那一部分最初叫“葬礼”,后来成了“葬礼(1)”(或“中国皇帝的袍子”)。这段思考接在双重按语《我的孤儿们》与《拒受一份遗产——或一个矛盾的代价》(第46、47号,第177、192页)之后,那两则按语是给《收获与播种》最后那一“节”(更确切地说,是给它本来要有的第一部分、即《愚妄与更新》)《一段过去的重量》(第50号,第131页)作注解的。后来,又给同一节添了别的注解(第44$'$、50号按语),还有别的按语来给《我的孤儿们》作注解,这些又依次生出新的作注解的按语;这一回还不算上真正的脚注——当预定的注解(写成白纸黑字之后仍然是)篇幅不大的时候。于是,从理论上说,《收获与播种》那一部分(当时它本该构成第二即最后一部分)就表现为一组给《一段过去的重量》那一“节”作的“按语”。由于惯性,这种“按语”(而非“节”)的细分法在随后三个部分里仍旧保留着;在那里,我给思考“初稿”作注解的手段兼用两种:脚注(当篇幅允许时),以及正文中引到的、写在后面的按语。\par 从排版上说,“按语”与“节”(在 RS I 里用作思考“初稿”的基本单位)的区别在于一个记号,如(1)、(2)等(即把按语编号放在括号里、排成上标,这是引注的通行做法),它或者放在该按语的开头,或者作为一种引注放在正文中提及它的适当位置。而“节”则用1至50的阿拉伯数字标出(不用那些令人生畏的下标与上标——出于实际需要,我在按语里不得不使用它们)。话虽如此,可以说《收获与播种》第一部分里“节”的功能与后面各部分里“按语”的功能并无本质区别。我在本信这一部分(《自发与结构》)中关于这一功能所作的评论,同样适用于 RS I 的“节”,尽管我一律用“按语”这个通名。\par 至于别的细则与约定,尤其是关于如何读《葬礼(1)》的目录,我请参见《导论》第7节(《葬礼的编排》),特别是第 xiv–xv 页。}……),正如一个活机体的身体之于它每一个个体细胞——那些细胞千差万别,每一个都占着一个只属于它的位置、担着一个只属于它的功能。不过有时候,在同一段一口气追下来的思考里,事后会觉察到一些重要的断口,把它分成好几个这样的单位或消息,于是每一个都得到自己的名字,由此获得自己的身份与自治。另一些时候则相反,一段由于某种原因(多半是偶然原因)被打断的思考,第二天或第三天又自发地接了下去;或者一段横跨两三天写成的思考,事后看来却像是从头到尾一口气写成的;仿佛只有睡眠的需要,才迫使我们违心地在那里插入某种断口(可以说是“生理性的”),它只由一行干巴巴的日期
%
% ---------------- p.97 ----------------
\origpage{L39}(有时还由好几行)标出来,标在所谈那则“按语”的某些相邻段落之间——这时那则按语便以一个独一无二的名字把自己区分出来。

\origpage{L40}就这样,《收获与播种》的每一则按语都有它自己的个体性,有使它区别于任何别的按语的面孔与功能。对每一则,我都设法用一个名字表达它独有的特性;那名字应当还原或唤起它“要说”的东西之本质,至少是某种本质。每一则,我真正认它,首先就是凭它的名字;此后每当需要它来帮忙,我也是凭这个名字唤它。

名字往往在我还没想到它之前就自己来了。正是它不期而至的出现,向我示意:我还在写的这则按语已近完成——它已经把要说的话说完,只等我写完眼下这一段……同样常见的是,在继续思考之前重读昨天或前天的按语时,名字也照样自发地出现。有时,在刚来的那则按语出现之后的几天或几周里,名字会有所变动,或者添上第二个我起初没想到的名字。许多按语有两个名字,表达对它所传消息的两种不同的照明,有时是互为补充的照明。这些双名中第一个向我呈现的,是在“愚妄与更新”开头不久,即《与克劳德·谢瓦莱相遇——或自由与善意》(第11号)。

只有两次,我在开始写一则按语之前就先有了名字——而且两次都后来被事态的发展搅乱了!

只有拉开距离——几周、甚至几月的距离——之后,在逐日相继的按语整体里,才会显出某种整体的运动与结构。我设法靠各种分组与再分组来把握这两者,每一组都有它自己的名字,那名字赋予它独有的存在、它的功能或消息;有点像方才那个比喻里同一个身体的器官与四肢,以及四肢的某些部位。于是,在《收获与播种》这“整体”里,有我已说过的五个“部分”,每一部分各有自己的结构:“愚妄与更新”分组成八“章”,I 至 VIII\fnote{41}{}{在《愚妄与更新》里,我有时把这些章称作《收获与播种》的“部分”;当然不要把它们与已说过的那五个部分相混,那五个部分是后来才出现的。};而构成“葬礼”的那三个部分合起来(它们也是几个月里逐渐显出来的……),则由一支漫长而庄严的、十二支“队列”I 至 XII 的行列组成。其中最后一支,更确切地说是那场“殡葬大典”(这就是它的名字)——前面十一支“队列”都朝它走去(肯定没怎么料到什么……)——其规模真是巨大,与它所隆重安葬的那部作品的规模相称:它囊括了 RS III(《葬礼(2)》)的几乎全部和 RS IV(《葬礼(3)》)的全部,连同它近八百页、一百五十来则按语(而当初,这场大名鼎鼎的典礼只预定包含两则!)。这场典礼由大主祭亲自(以他众所周知的谦逊……)手法娴熟地主持,分为九个彼此分开的“时辰”或礼仪环节,以《悼词》开场(这还用说),并以《自深处》收尾(理当如此)。这些“时辰”之中另有两个,一名《阴阳之钥》,一名《四种运算》,每一个都(远远)构成它所嵌入的《收获与播种》那部分(III 或 IV)的最大份额,并且顺带把自己的名字给了那一部分。

在《收获与播种》全书中,我像护着眼珠一样照料目录,不断改动它,以顾及源源不断涌来的未预料的按语\fnote{42}{}{这些未预料的按语里,尤其有那些“出自一则尺寸变得无法接受的脚注”的。最常见的是,我把它紧放在它所隶属的那则按语之后,给它同一个编号,必要时加上一个 $'$ 或 $''$ 号,甚至加上 $'''$——这就免去了把此后所有已写好的按语一并重新编号这桩干不下来的活儿!这些由一则脚注生出另一则的按语,在目录里前面带一个 $!$ 号(至少在《葬礼(1)》里如此)。},并
%
% ---------------- p.98 ----------------
\origpage{L41}且尽我所能让它精细地映照出思考的整体运动与其中显现的微妙结构。正是在第三、尤其是第四部分(刚才说到的“钥匙”与“四种运算”)里,这结构最为复杂、最为交错。

\origpage{L42}为了给文本保住自发性的品格、保住思考如其实际进行与实际被经历时那种未预料的面貌,我不愿意把名字放在每则按语前面——何况名字每次都是事后才出现的。所以我建议你:每读完一则按语,就翻到目录去看看这则按语叫什么;顺便也可以一眼看出它如何嵌入已经进行的思考,甚至嵌入尚未来到的思考。否则,你就有望无可救药地迷失在一堆看来难以消化、五花八门的按语里,其编号有时古怪,甚至可说令人生畏\fnote{43}{}{至于这些有时看来也许古怪的编号何以必要,我请你参见这封滔滔不绝的信里前一条脚注。};就像一个旅人在一座陌生的城市里迷了路(那城市是按世代与世纪的任性古怪地长出来的……),既没有向导,连一张地图也没有,帮他辨认方向。在准备付印的稿子里,我打算随文加入“章”以及别的按语组与节组的名字,只把按语(或节)本身除外。但即便那样,偶尔查阅目录在我看来仍不可少,免得迷失在几百则按语的一团乱麻里——它们一串接一串,绵延一千多页……

\section{(12) 自发与严格}

自发与严格,是同一品质那“暗”与“明”的两面。唯有二者的婚配,才生出文本或存在者的那种特殊品质——人们也许可以用“真之品质”这样的说法去唤它。如果说在我过去的出版物里,自发性(即便不是缺席,也是)被压缩到极小的份额,我并不认为,随着它在晚年于我身上绽放,严格性因此就变小了。毋宁说,它那位阴性的伴侣的完全在场,给了严格性一种新的维度、新的丰饶。

这种严格也针对它自己施行:它留意使自己必须做的那桩微妙的“筛选”——在意识领域里流过的大量东西中,不断地把有意义的、本质的东西从偶然的、次要的东西里澄清出来——不至于变厚、凝固成审查与自满的自动机制。唯有我们心中求知的渴望,才能唤醒并激发这样一种不笨重的警觉、这样一种灵敏,去对抗那巨大而无处不在的惰性,即那些由现成观念凿成的“(所谓)天然的坡道”——它们是我们的恐惧与我们所受的塑造的表达。

而这同一种严格、同一种警觉的注意,也朝向自发性,也朝向那些带自发性面貌的东西,为的是在那里同样分辨出成分:分辨出那些固然再自然不过的“坡道”,把它们与真正从存在深处、从原初的求知与行动冲动中喷涌出来的东西区别开来——那冲动把我们带向与世界的相遇。

在写作层面上,严格表现为一种持久的用心:用语言尽可能精细、尽可能忠实地圈定要表达的思想、情感、知觉、意象、直觉……;在所要表达之物轮廓分明之处,不满足于一个含糊或近似的词;而在所要表达之物仍被“还只是被预感”的迷雾围绕之处,也不满足于一个虚假精确(因而同样走形)的词。当我们试图就地捕捉它如其所是的样子时——也只有在那时——未知之物才向我们显露其真正的本性,也许甚至直到大白天光之中,如果它生来属于白昼、而我们的欲望又促它卸下阴影与雾的帷幕。我们的角色不是妄称描述并固定我们所不知、从我们手中逃逸的东西,而是谦卑地、炽烈地去认识那从四面八方围绕着我们的未知与神秘。

% ---------------- p.99 ----------------
\origpage{L43}这就是说,写作的角色不是记录一项研究的结果,而正是记录研究的过程本身——记录爱的工作,记录我们与我们的母亲、那“世界”、那“未知者”的爱的作品;她不知疲倦地召唤我们进入她之中,去认识她那取之不尽的躯体里的她,去到她之中那神秘的欲望之路把我们带去的任何地方。

为了传达这一过程,那些回头的修正——它们给写作的“初稿”、甚至第二稿第三稿添上细微差别、加以精确化、深化,有时加以改正——本身就是发现进程的一部分。它们构成文本本质的一部分,赋予文本全部意义。所以,放在《愚妄与更新》末尾的那些“按语”(或“注解”)——构成文本“初稿”的那五十“节”里处处引到它们——是文本不可分割的、本质的一部分。我竭力劝你边读边翻过去看,至少每读完一节的末尾(那节里出现了一处或几处对这类“按语”的指引)就去看。对《收获与播种》其余部分里的脚注也是如此;或者,当某则“按语”(在此充当“正文”)里引到后面某一则按语,那后一则便起“回看”它或注解它的作用。这一点,连同我请你不要在阅读中途离开目录的劝告,就是我要给你的阅读建议中最主要的一条。

最后一个问题,是实用性的,它将(有点平淡地)结束这封该收尾的信。有些时候有点“慌乱”,为了把《收获与播种》的各个分册准备好,交给系里的复印部去印,好让印刷(如果可能)赶在大假之前。匆忙之中,有一整页最后一刻的脚注要补进第2分册(《葬礼(1)——或中国皇帝的袍子》),却“跳”掉了。那些主要是改正某些直到最近、在写《四种运算》的过程中才出现的材料性错误。这些脚注里有一条比其他几条更重要,我想在此指出。那是对按语《受害者——或两种沉默》(第78$'$号,第304页)的一则注解。在那则按语里,我竭力——除了别的——去圈定我的一些印象(当然全是主观的):关于吾友左格曼·梅布胡特当时如何“内化”他所遭受的那种令人不安的掠夺;而这则按语被他觉得对他不公,因为我似乎几乎把他与那些掠夺者“装进了同一个袋子”。可以肯定的是,在那则按语里——它并不自称给出别的东西,只给出与某个特定“时刻”相关的印象——我只敲响了一口钟,而把另一些同样真实(也许还不那么可争议)的钟声留在未言之中(而且大概还当作不言自明的事)。不管怎样,对这一个微妙题目的思考,在一年之后,在按语《根与孤独》(第171号)里大大深化了。这一则没有引起左格曼的保留。关于同一题目的另一些思考材料,也见于《三块界石——或无辜》与《死页》(第171(x)与(xii)号)这两则按语。这三则按语都属于《封神》,那是《四种运算》里专讲对左格曼·梅布胡特作品的据为己有与窃取那场运算的部分。

我只剩下祝你阅读愉快了——也盼着轮到我读你的信!

亚历山大·格罗滕迪克

\origpage{L44}
\subsection*{后记:或一场辩论的语境与前情}
\tocsubsub{后记:或一场辩论的语境与前情}

% ---------------- p.100 ----------------
\origpage{L45}
\section{(13) 瓶装光谱仪}

这封信写完之后,整整七个月过去了;寄出去也已近四个月,连同随信而去的那个“大砖头”。而且每一本上都有我手写的题赠\fnote{44}{}{只有少数例外,主要是那些我本人不认识的同行,他们只收到试印本的第0和第4分册,作为他们积极参与我的“葬礼”的奖赏。}。像一只“漂流瓶”,更确切地说,像一整群这样的漂流瓶子,我的消息一路漂去、流转,直到这个我曾熟悉的数学微观宇宙最偏远的角落。随着日子、星期、月份一天天过去,直接与间接的回声传回来,于我不期然地,面前仿佛出现了一张数学界的巨大射线照片——那是由一架触手四伸的“瓶装光谱仪”拍下的,而我那些天真的“瓶子”就是它一根根游走的天线。于是(责任随身份而来!),我这个本来并不缺事可做的人,如今面前又多了一件新任务:破译这张射线照片,并且尽我所能,把我从中读出的东西报告出来。这将作为《收获与播种》的第六部分(也是最后一部分,我保证!)。倘若上帝假我以岁月,这一部分便将为我“晚年那部社会学巨著”加冕。眼下,先说几句初步的评论。

用来迎接我这支可怜的手工小舰队的,占压倒优势的,是一种甜腻又带三分凶气的调子,唱的是“瞧,格罗滕迪克老了,犯起偏执狂来啦”,或者“这位可真把自己当回事”——如此这般,戏就演完了!然而这种调子的信我只收到一封\fnote{45}{}{这封信出自一位曾是我的学生的人,而且他还是我的一位“同葬者”。},另有两封属于那种含蓄的、自得其乐的嘲弄\fnote{46}{}{出自两位当年在布尔巴基共事的同行,其中一位还是我初出道时善意接待过我的年长者之一。}。我的数学家收信人,包括那些曾是我的学生的人,大多数以沉默作答\fnote{47}{}{寄给数学家的共一百三十一份,至今有五十三个收信人有了动静,哪怕只是告知收到。这其中有六个是我过去的学生——另外八个人里,我一个也没得到音信。}——这沉默对我说了很多。

不过我已经收到了大量通信。绝大多数信是客气而尴尬的调子,那尴尬常常还想显得友好,仿佛是出于礼数。有两三次,我在这尴尬背后、仿佛被它滤过之后,感到一种依然活着的情感的温度。最常见的情形是:尴尬若不是以对好情感的郑重声明来表达(为他自己,或为别人),就是以恭维来表达——我这一辈子也没收到过这么多恭维!唱的是“伟大的数学家”、“了不起的篇章”(关于创造力“以及诸如此类”……)、“无可争议的作家”,等等等等。为了凑足分量,我甚至还有幸得到一句情真意切(绝非反讽)的恭维,夸我内心生活之丰富。不用说,在所有这些信里,我的通信人都小心翼翼地不进入任何问题的实处,更不会亲自卷入;调子倒更像是那种“被请来发表意见”的人(借用其中一封信的说法)面对一桩有点不干不净的事——何况还是假设的、甚至想象的——总之尤其要紧的是:一桩与他本人无关的事。可是当他装出要碰一碰其中一个问题的样子时,那是指尖轻触,并且尽可能把它推得离自己远远的——或者借给我一番好劝告之机,或者用谨慎的条件式,或者在不知\origpage{L46}%
% ---------------- p.101 ----------------
该说什么时用那套现成的套话,或者用别的什么法子。\origpage{L47}不过还是有些人透露说,也许发生过一些不太正常的事——同时小心地把是什么事、是什么人留在最大的含糊之中……

我也收到过坦率热烈的回声,来自十五六位新老朋友。有些人表达了感动,并不想掩饰它、也不想让它噤声。这些回声,以及另一些来自数学界之外、同样热烈的回声,将是我对一项漫长而孤独的工作的报偿——那工作不仅为我自己,也是为所有人做的。

在收到我这封信的一百三十来位同行中,有三个人真正在完整意义上作了回应:他们把自己放了进去,而不是限于对时下事件作一番远距离的评论。我还从一位非数学界的通信者那里收到一个这样的回声。这些才是对我的消息的真正答复。这也正是我最好的报偿。

\section{(14) 三只脚踩一个盘子}

我的好几位同行和数学家朋友曾表示希望,《收获与播种》能在数学界引发一场广泛的辩论:关于这个圈子里的风气状况,关于数学家的伦理,关于他工作的意义与目的。目前,至少可以说,事情并没有往那个方向走。从现在起(说一句应景的双关语),关于一场“葬礼”的辩论,看样子已经径自被一场辩论的“葬礼”取代了!

不管人们愿不愿意,尽管大多数人的沉默与冷漠,一场辩论却确确实实已经开了场。它恐怕永远不会具有真正公开辩论那样的规模,更不会(上帝保佑别这样!)具有“官方”辩论那种排场与僵硬。无论如何,已经有很多人抢先一步、三下两下就把它了结了——在还没弄清它是什么之前,就在自己心里把它关掉了;他们仗着那条万古不易、不可动摇的共识:“在这最好的世界里,一切都再好不过”(在此即数学的世界)。不过,也许终究会有一次质疑从外面来,渐渐地从一些“证人”那里来——他们不属于同一个圈子,因而不是这个圈子的群体共识的囚徒,因而(即便在自己心里)也不觉得自己被牵连在内。

在收到的几乎全部回声里,我都看到同一种混淆,涉及两个前提性的问题:《收获与播种》所提出(至少是默契地提出)的“辩论”究竟针对什么;以及谁有资格了解它并就它表态,或者说:谁有资格在充分知情的基础上形成意见。关于这一点,我想在此标明三个“坐标点”。当然,这并不能阻止那些执意要混淆的人继续待在混淆里。至少,对于那些想知道究竟是怎么回事的人,这也许能帮他们不被四面八方来的响动(甚至包括那些最善意的响动……)分了神。

a)一些真诚的朋友向我保证“一切最终都会摆平”(而“一切”,我猜,指的是那些不幸被弄坏了的“东西”……);我只要重返舞台,“用新工作让自己站住脚”,做做演讲等等——其余的别人会做。人们会宽宏地说“对这个该死的格罗滕迪克,咱们也的确有点不公”,然后不声不响地、带着或多或少的诚意去拨正航向\fnote{48}{(*)}{我已经不止一次注意到这类微妙的迹象,表明人们已经好好记下:狮子醒了……};甚至慈父般地拍拍他的肩膀,说他是“伟大的数学家”,好让这么个毕竟还算可敬的家伙平静下来——可惜他偏偏作势要发火,要掀起不受欢迎的风浪。

正如这些朋友所暗示的,这里根本不是要“松松绳子”、也不是要让人松松绳子。就我而言,我全然不需要恭维,不需要真诚的仰慕者,也不需要“盟友”——为“我的”事业,或为任何别的事业。问题不在我身上——我过得好得很;也不在我的作品——它自己会说话,哪怕是对聋子说。如果说这场辩论也(在别的之外)涉及我这个人、我的作品,那只是因为它们是对别的东西的揭示——通过一场葬礼的现实(的确是最具揭示性的现实之一)。

如果真有“某个人”在我看来该唤起警觉、不安与紧迫之感,那绝不是我这个人物,甚至也不是我的任何一位“同葬者”。那是一个集体的存在者,既不可捉摸又十分具体,人们常常谈起它,却从来小心翼翼地不去审视它——它的名字叫“数学共同体”。

最近这几周,我终于把它看成一个有血有肉的人,而它的身体正遭受着深处的坏疽。最好的食物、最精美的菜肴,在它体内都变成了毒,使病害传播得更广、扎得更深。然而它又有一种不可遏制的贪食,要永远吃得更多,仿佛这肯定是它自欺的一种方式——关于一种它无论如何也不愿去认识的病。对它说什么都是白费——最最简单的词也失去了意义。它们不再是消息的载体,只剩下引发恐惧与拒绝的喀嗒声……

b)我的大多数同行,甚至那些心肠好的老朋友,在贸然发表意见时,都要围上一圈谨慎的条件式,诸如“如果事情真是……那确实是不能容许的”——好让他们心满意足地再躺回去睡大觉。可我原以为我说得够清楚了……

\origpage{L48}隔着七个月的距离,如今我可以说明:对于《收获与播种》里报告并评论的几乎全部事实,它们的真实性不成为任何争论的对象。我会在后面回到那少数几个例外(它们也会各自在相应的地方被标明是这样的);至于其余全部事实,在《收获与播种》初版写成之后,与几位主要当事人的仔细对证(即皮埃尔·德利涅、让–皮埃尔·塞尔与吕克·吕西耶)消除了细节上的错误,并就材料性事实本身达成了毫不含糊的一致\fnote{49}{(*)}{我很高兴向这三位表达我的感谢,感谢他们在此事上表现出的善意;并确认他们在一切涉及材料性事实的问题上完全诚实。}。

所以,辩论根本不涉及事实的真实性——那不成问题——而涉及这样一个问题:这些事实所描述的作派与态度,是否应当被视为被认可的、“正常的”,还是不是。

这里说的是这样一些作派:在我的见证里,我(也许错了……)称之为丑闻式的;如背信弃义或滥用权力,如明目张胆的不诚实,不止一次达到了不义与无耻的程度。在了解了这些事实(十五年前还不可想象的事实)之后,我还得学的一件事——相当难以想象的一件事——就是:我的同行中的大多数,甚至包括那些曾是我的学生或朋友的人,今天竟把这些作派视为正常,视为完全体面。

c)我的许多同行与老朋友维持混淆还有第二种办法。调子是这样的:“抱歉,可我们不是这一行的专家——别要求我们去了解那些(天佑如此……)从我们头顶上飞过去的事实……”

我则断言:要了解那些主要事实,根本不必是“专家”(这回轮到我抱歉了!),甚至不必会背乘法口诀表或知道毕达哥拉斯定理。

% ---------------- p.103 ----------------
\origpage{L49}连《熙德》或拉封丹《寓言》都没读过也不要紧。一个发育正常的十岁孩子,在这件事上跟最负盛名的专家一样有能力(甚至比他还强……)\fnote{50}{(**)}{当然,我写《收获与播种》并不是给十岁孩子看的;若要对他说话,我会选一种他熟悉的语言。}。

\origpage{L50}让我举一个例子来说明这一点,就从“葬礼”里随手拈来一个\fnote{51}{(***)}{这是我在1984年4月19日发现的第一场“葬礼大手术”,也正是那一天,“葬礼”这个名字强加给了我。关于此事,见同一天写的那两则按语:《一个梦的回忆——或母题的诞生》与《葬礼——或新父亲》(RS III,第51、52号)。那里也有即将谈到的这本书的完整出处。}。要了解下面这几件事并就它们作出判断,根本不必知道“Motif”这个多面而极微妙的数学概念的来龙去脉,也不必小学毕业。

1°) 1963年至1969年间,我引入了“Motif”这一概念;并围绕这一概念发展出一套“哲学”与一套“理论”,它们部分地仍停留在猜想阶段。无论对错(此处无关紧要),我把 Motif 理论看作我为我那个时代的数学带来的最深的东西。再说,如今“Motif 的瑜伽”的重要性与深度已无人质疑(在它被近乎完全沉默了十年之后——那沉默从我离开数学舞台之后随即开始)。

2°)在唯一一本(1981年出版)、基本上以 Motif 理论为专题的书里(这个由我引入的名字就出现在书名中),唯一一处能让读者疑心我这点微末之人与书中长篇大论发展的那种理论有远近关系的段落,在第261页。那段话(两行半)是在向读者解释:那里发展的理论,与一个名叫格罗滕迪克的人的理论毫无关系(该理论在那里被提到第一次也是最后一次,没有任何别的出处或说明)。

3°)有一个著名的猜想,人称“霍奇猜想”(它究竟讲什么都无所谓),它的成立将意味着那本光彩夺目的书里所谓“另一种”Motif 理论,与我将近二十年前在众目睽睽之下发展出的那种理论相同(是它的一个极特殊的特例)。

我还可以加上 4°):那本书四位署名者中最负盛名的一位曾是我的学生,他多年来那些精彩的想法,除了从我这里,还能从哪里学来?——而他在书里把它们呈现得像是刚刚才想出来似的\fnote{52}{(*)}{我并不是说这本书里没有一些想法,甚至一些漂亮的想法,出自这位作者或别的合著者。但整本书的问题框架、赋予它意义的概念语境,乃至那套微妙的 X-范畴理论(被误称为“Tannaka 式”)——那套技术构成了全书的核心——这些都是我的成果。};5°)这两件事在消息灵通的人中间是众所周知的,可是在文献里白费力气也找不到一处书面痕迹,能证明那位出色的作者也许从我口中听到过什么\fnote{53}{(*)}{不过有一个例外:1977年塞尔写的一份报告里有一行,将在相应之处谈到。};以及 6°)那个微妙的算术问题——据主要作者本人向我解释,它构成全书的核心问题(而我的名字始终未被提起)——是我在六十年代沿着“Motif 的瑜伽”的路子提出来并澄清的,而作者是通过我才知道它的;我还可以再堆上 7°、8° 等等(在相应之处我当然不会不这么做)。

前面这些,对我的用意来说已经够了,我的用意就是:要了解这类事实并就它们作出判断,根本不需要什么专门的“能力”——“事情不在那个层面上”。这里起作用的能力,除了健全的理性(原则上人人都有)之外,就是我愿意称之为体面感的东西。

这里说的那本书,如今已是数学文献中被引用最多的著作之一,而它的“主要作者”也是当代最负盛名的数学家之一。话虽如此,而且看得清楚,在这桩事情里,如今在我看来最最值得注意的是:在这本书数不清的读者当中,包括那些第一手知道究竟是怎么回事的人,包括那些曾是我的学生、我的朋友的人——竟没有一个人看出其中有什么不正常。无论如何,直到我写这几行的今天,还没有一个人来向我表示,对这本光彩夺目的书有丝毫保留\fnote{54}{(**)}{总共只有两位同行(包括左格曼·梅布胡特)向我表示过这样的“保留”。这两位都不能算是这本书的“读者”。他们是出于好奇翻了翻,想弄弄清楚……}。

% ---------------- p.104 ----------------
\origpage{L51}至于我的同行和老朋友中那些从未把这书拿在手里、却以此为由申辩自己无能的人,我要对他们说:要当“专家”,根本不必——只要到随便哪一家数学图书馆去把那卷书借来,翻一翻,亲眼看看那无人争议的东西就行了……

\section{(15) 坏疽——或时代精神(1)}

\origpage{L52}这场“Motif 运算”只是同一水准的四场“大运算”之一,此外还有一大批规模较小、精神相同的。它远不是用来充实我那份“一个时代的风俗画”的集体蒙骗中“最大”的一桩,更不是最独特的一桩。它所做的,不过是趁富人不在(或已故……)之机,只掠夺了他的羊群;而不是(在普遍的漠然中)为了取乐,当着他的面去掐死穷人的母羊。甚至直到数学语言里如今已进入日常用法的层面,那些看来无害的书名、概念名、命题名,时时被人引用,它们本身就已经是一种蒙骗或欺诈\fnote{55}{}{我这里想到的尤其是那个古怪的缩写“SGA 4 ½”(分数还是有用处的!),它一个就抵得上双重欺诈(而且是当代数学文献中被引用最多的缩写之一);还有“维尔迪耶对偶”或“维尔迪耶对偶对象”、“德利涅–格罗滕迪克猜想”之类的名称,最后还有“Tannaka 式范畴”(这一回 Tannaka 本人并无过错,因为从来没人问过他……)。这些将在相应之处更详细地谈到。},并以各自的方式见证一个时代的不体面。

如果我自认为曾为“数学共同体”做过什么有用的事,那就是把若干不光彩的事实带到了大白天光之下——那些事实原本在暗处发酵。这肯定是人人都天天、或差不多天天、远远近近接触到的那类事实。他们之中有多少人肯费点闲工夫停下来——哪怕只是一瞬——去闻一闻空气,看一看?

那个自己曾撞上一些人的傲慢、另一些人(或同一批人)的不诚实的人,也许还自以为那是一种专门分配给他的特别霉运。把自己的经验与我的见证相对照,也许他会感到,这“霉运”也是他给一种时代精神起的名字,那种时代精神压着他,一如压着所有人。而且(谁知道呢!)也许这会促使他投入一场辩论——那场辩论与他有关,正如与我有关。

可是,如果我“拿到公共广场上摊开”的这堆“脏衣服”,除了引来一些人没有笑意的冷笑、另一些人客气的尴尬,在众人的漠然中别无反响,那么一个本来浑浊的局面就会变得非常清楚。(至少对那个还愿意用自己眼睛的人是这样。)真诚与体面的传统共识\fnote{56}{}{我说这些“出于真诚与体面的共识”,并不是说它们从未被违反过。但即便在被违反的时候,那也确实是“违反”,而共识本身仍然是被接受的。},在数学家之间的关系里、在数学家与其技艺的关系里,从此都将成为过去之物,“过时”之物。无需有哪个国际数学家组织再隆重宣告,这从此也将是既成之事、几乎是官方之事:如今,一切手段都允许,不再有任何保留与限制——对于数学界里掌权的那伙“靠内部推举结成的兄弟会”而言。一切观念上的捣鬼
%
% ---------------- p.105 ----------------
\origpage{L53}把戏,只为把那个只等着相信的漠然读者牵着鼻子走;一切父权名分的黑市交易;同伙之间的假引注,以及对注定沉默者的沉默;各种拉帮结派与形形色色的伪造,乃至在众目睽睽之下最粗劣的剽窃——是的,对这一切都说“阿门”,并且带着数学公共广场上所有“大名头”、所有大大小小老板的祝福,无论是用言语还是用沉默(若不用积极而殷勤的参与的话)。是的,对那里风靡一时的“新风格”说“阿门”!那曾经是一门艺术的东西,如今经(几乎)一致的赞同,变成了一个搅浑水、抢地盘的集市,在老板们慈父般的目光下。

曾经有一个时候,在数学家的世界里,权力的行使受到一致的、不可触犯的共识的限制,那是集体体面感的表达。这些共识与这种感情,如今都成了过时之物、被超越之物,肯定配不上计算机、太空舱与中子弹的光荣时代。

从此这将是既成而封定之事:权力,对于那伙掌握它的人而言,是一种任意专断的权力。

\section{(16) 公开认错——或时代精神(2)}

在《信》里,我想我已充分说明了我写《收获与播种》时所怀的精神,足以让人清楚:我绝不声称在那里做史家的工作。那是一份真诚的见证,关乎第一手的亲身经历,以及对此经历的思考。见证与思考都供所有人取用,包括史家——他可以把它当作众多材料之一来用。那时,把这份材料交付批判性分析、依其技艺的严格准则来做,便是他的事。

当然,应当区分狭义上的事实(“生事实”或“材料性事实”)与对这些事实的“评估”或“解释”——后者赋予它们意义,而这意义对一位观察者(或共同行动者)和对另一个人并不相同。大体上可以说,《收获与播种》的“见证”方面关涉事实,而它的“思考”方面关涉对事实的解释,也就是我为赋予它们意义所做的工作。在构成见证的“事实”里,我也归入“心理事实”,尤其是我的见证所映照的各种情感、联想与意象,无论它们发生在或远或近的过去,还是就在写作的当下。

对于我在《收获与播种》里描述或提及的事实,我区分三类来源。有些事实是记忆还给我的,其清晰或模糊程度因场合而异,有时还被变形。关于这些,我可以为我在写作之时的真诚立意担保,却绝不为全无错误担保。相反,我曾有机会指出其中若干细节错误,并在相应之处用后来补写的脚注标明。另一方面,有书面文献,尤其是书信,特别是正式发表的科学出版物,我有时以所期望的全部精确性引到它们。最后,还有第三方的证词。有时它补充我自己的记忆,让我重新唤起、加以精确化,有时加以改正。在某些少见的场合(我马上会回过来谈),这证词给我带来了与我已知信息全然不同的新信息。当我转述这样一份证词时,这并不意味着我有能力自始至终核查它的准确与可靠,而只是说:它相当可信地嵌入了我第一手所知的那张丰富的事实之网,足以让我确信(无论对错……)这份证词在本质上确实合乎真相。

% ---------------- p.106 ----------------
\origpage{L54}对于一个细心的读者,我想在任何时候都不会有困难去“分清”事实的记述与对事实的解释,并且(在前一种情况下)辨明我方才描述的三类来源中,是哪一类在起作用。

\origpage{L55}\asterisms

方才我提到一份第三方证词,我转述了它,却未能“自始至终核查其可靠性”——那就是左格曼·梅布胡特的证词,关于围绕他著作的那场大规模掉包手术。在《收获与播种》里提到的“材料性事实”中,如今尚可争议、或按我现在的判断需要改正的,只有某些仅凭梅布胡特一人证词而立的事实。为结束这篇后记,我要在此对试印本《收获与播种》中呈现的“梅布胡特事件版本”提出一些批评性的评论。更详细的评论与更正将各自在相应之处收入印刷版(即《收获与播种》的定本)。

我曾想充当其转述者的那个“梅布胡特版本”,在我看来本质上由以下两个论点构成:

1. 1972年至1979年间,梅布胡特是孤身一人\fnote{57}{}{1975年柏原的可构造性定理是例外,它在理论中的重要性毫无争议。但按梅布胡特的说法,那该是柏原对这一正在诞生的理论的唯一贡献。这一(不准确的)说法,曾由柏原其他论文的缺席所佐证——在那些论文里,他本该至少暗示某些主要想法。},在普遍的漠然中,从我的著作汲取灵感,发展出“$D$-模哲学”,把它作为我意义上的“上同调系数”的新理论。

2. 无论在法国还是在国际上,都有过一种一致的共识,要把他的名字与角色从这一新理论中抹掉——一旦它的意义开始被承认。

这个版本有很强的文献支持:一方面是梅布胡特那些完全令人信服的出版物;另一方面是其他作者许多出版物的署名(尤其是1981年6月吕米尼讨论会论文集),其中掉包的蓄意用心不容任何怀疑。最后,梅布胡特另外提供给我的那些更详细的细节(我在“葬礼(3)——或四种运算”那部分里转述了它们),虽然无法直接核实,却与某种整体的氛围完全吻合,而那氛围的真实性在我看来已不容任何怀疑。

我刚得知若干新的事实\fnote{58}{}{我感谢皮埃尔·沙皮拉与克里斯蒂安·乌泽尔,他们愿意提醒我注意这些事实,以及我对梅布胡特–柏原之争的陈述中那种偏颇之处。},它们表明对上面第 1° 点须作大大的修正。梅布胡特所处的孤立\fnote{59}{}{这种孤立首先来自我过去的学生对梅布胡特的想法与工作的漠然——他固执地装作在追随一位被一致共识判给遗忘的“先祖”……}确实是真的,但那是一种相对的孤立。在法国,有 J.-P. 拉米在同一个题目上的工作(梅布胡特对我只字未提),而尤其是:看来梅布胡特所发展并完成到底、且被他归为自己名下的某些重要想法,可能应归于柏原\fnote{60}{}{这些想法中最重要的是所谓“黎曼–希尔伯特”对应(用新风格的术语来说)在 $D$-模上的版本。相关猜想已由梅布胡特证明,并且(据沙皮拉向我肯定)也由柏原证明(而梅布胡特向我保证他的证明是唯一发表的)。关于证明优先权的问题,对我仍是一片迷雾,我放弃把余生用来把它弄清楚……\par 至于以 $\nabla^{\infty}$-模表述的姊妹命题,看来毫无疑问:想法的归属与证明都属于梅布胡特。}。这一来,就使我(过于)忠实地转述的那个梅布胡特版本里关于柏原–梅布胡特之争的某些情节,显得不可信或可疑了。

% ---------------- p.107 ----------------
\origpage{L56}毫无疑问,无论就“档案工作”的层面而言,还是就他设法完成到底的某些想法的构想而言,梅布胡特都是 $D$-模新理论的主要拓荒者之一,也许甚至是首要的拓荒者;无论如何,他是唯一一个全身心投入这项任务的人,而这项任务的真正意义当时还没向他显现,正如它也没向所有人显现。也确实,围绕这部著作发生的那场掉包手术——以吕米尼讨论会为顶点——在我看来仍是数学界这个世纪的大不体面之一。但若声称(像我出于真诚所做的那样)梅布胡特是独自担此任务的人,那就错了。相反,他是唯一一个有诚实与勇气明确指出我的想法与著作在他的工作中、在新理论的萌发中所占分量的人。

这篇后记不是细谈此事的地方——我会在相应之处细谈,包括一些有助于说明“梅布胡特版本”心理语境的评论。如果说“梅布胡特–柏原的纠纷”对我有某种意义,那仅仅在于它照亮了一个时代的整体氛围。而对我来说,甚至它那些变形、甚至通过那些促其浮现而起作用的力量,“梅布胡特版本”也显得是我为“一个时代的卷宗”提供的众多不那么可争议的材料之中一份有力的“时代征兆”。

我还须为我的轻率作公开认错:我呈现梅布胡特–柏原之争时,只用了一幅仅仅依据梅布胡特所提供的证词与文献的图画,而且仿佛那个版本不容任何怀疑。那个版本把一位第三方置于可笑乃至可憎的光线下,这就更有理由谨慎行事。为我的轻率、为这种缺乏健全的谨慎,我很愿意在此向柏原先生表示我最诚挚的歉意。

% ================================================================
%  第 4 章 引言 (原书 p.109-122, 边码 P.I-P.XXII)
%  源文: _work/src1/task-ch04.md
% ================================================================
\chapter{引言}

% ---------------- p.109 ----------------
\section{(I) 五叶草}

\subsection{1. 梦与成就}

\origpage{I} 到今年七月,就是三年前了,我做了一个不大寻常的梦。我说“不大寻常”,这印象只是事后才有的,是醒来回想时才有的。梦本身来到我这里,却像世上最自然、最显然不过的事,不敲锣不打鼓——以至于醒来时我差点没留意它,随手把它扔进遗忘的角落,好去办“当天的正事”。从前一天起,我就已经动身去做一场关于我与数学之关系的反思了。这是我有生以来第一次肯费这个力气去看一看——而且,我之所以在那个时刻动手,实在是因为我几乎是被逼无奈、硬被赶上路的!在此前的那些月和那些年里,发生过一些如此奇特——若不是说暴烈——的事情,那是些数学激情的爆发,不打招呼就闯进我的生活,以至于实在再也没法继续不去看看到底发生了什么。

我说的这个梦没有任何情节,也没有任何动作。它只有一幅图像,静止不动,同时又极其鲜活。那是一个人的头,侧面像。看得出它正从右向左看。那是一个成年男子,无须,一头乱发在头的周围形成一圈力量的光轮。从这颗头上透出来的,尤其是那种年轻、欢快的力量感,仿佛从那柔韧有力的颈项弧线里涌出(那弧线是猜出来的,而不是看出来的)。脸上的表情,与其说是个成年人,或者一个已经沉稳下来的人(不管成熟与否),不如说是个顽皮的小淘气,正为某件他刚刚得手、或正在盘算着要干的勾当而洋洋得意。从那里透出来的,尤其是那种强烈、含而不露、迸发为嬉戏的生的欢乐……

% ---------------- p.110 ----------------
\origpage{II} 并没有第二个人在场,没有一个“我”在观看或凝视这另一个——我们只看得见他的头。但有的是对这头颅、对它透出的东西的一种强烈感知。也没有任何人去感受印象、去评论它、给它命名,或者给被感知的那个人贴上一个名字,称之为“某某”。有的只是这极其鲜活的东西,这颗男人的头,以及对这东西同样鲜活而强烈的感知。

至于醒来时,我并不是有意地,就记起了过去这一夜的梦;这个男人的头的影像,并没有以特别的强度从那堆梦里凸显出来,它并没有挤到前面来,冲我喊、向我低语:你该看的正是我!当这个梦出现在我对一夜之梦的匆匆一瞥里,在被窝里温暖的静谧中,我自然有了醒着的心智那个习惯性的反应:给看见的东西安上一个名字。何况我根本不必去找:只要我一问,我立刻就知道了,梦里那个在那儿出现的男人的头,不是别的,正是我自己的头。

这一手不赖,我当时想,到底还是得服气:在梦里这样看见自己,就像那是另一个人!这个梦的到来,有点像我在散步时,纯属最偶然的巧遇,撞上了一株四叶草,甚至五叶草,于是我按规矩惊异上片刻,然后继续赶路,就像什么也没发生过。

至少,事情差点就这样过去了。幸而——在这种情形下我碰过许多次——我还是为了问心无愧,把这个“不赖”的小插曲白纸黑字记了下来,并开始了一场反思,它本该顺着前一天的思路继续下去。可是,一针引一线,那一天的反思最终只是让我沉浸于这个不起眼的梦、这幅唯一的图像的含义,以及它带给我的关于我自己的那条消息。

这里不是细说我这一日的冥想教了我什么、带给我什么的地方。或者更确切地说,不是细说这个梦教了我什么、带给我什么的地方——那是在我已经进入一种专注与倾听的状态、从而能够接纳它要对我说的话之后。梦与这份倾听的第一枚直接的果实,是一股骤然涌来的新能量。正是这股能量支撑了随后几个月里那场长时间的冥想,与之相抗的是内心里顽强的阻力,我得靠耐心而执拗的功夫把它们一块一块拆掉。

我开始留意自己某些梦,已有五年;而这是第一个不以“信使之梦”如今已可辨认的那副面目出现的信使之梦——那种梦有惊人的场景手段,有异常强烈、有时令人翻肠倒肚的视象。这一个却“酷”到了家,没有任何强迫你注意的东西,简直是悄无声息——取亦可,舍亦可,毫无纠纷。

\origpage{III} 几周之前,倒来过一个老派风格的信使之梦,调门是戏剧性的,甚至是野性的;它突然而即刻地终结了一段漫长的数学狂热。两个梦之间唯一表面的亲缘关系是:两者之中都没有观察者。这个梦以一则简短有力的寓言,指出当时我生活里正在发生的一件事,而我却不肯费心去留意它——说实话,我还小心翼翼地不去理会它。正是这个梦让我当时明白了有一件反思的工作刻不容缓;几周后我便投入其中,它随后持续了将近六个月。\origpage{IV}在这部反思—见证《收获与播种》的最后一部分里,我多少有机会谈到它;这部作品开启本卷,并给了本卷它的名字\fnote{1}{}{尤其参见第43节《扫兴的老板——或压力锅》。}。

% ---------------- p.111 ----------------
\origpage{V} 我之所以用这另一个梦、用这幅我自己的影像—视象(我在德语笔记里称之为“Traumgesicht meiner selbst”)来开篇这则引言,是因为最近这几周里,这个梦的念头不止一次回到我这里,那时,那场“关于一位数学家的往昔”的冥想正走向它的终点。说实话,回头看,这个梦以来的三年在我眼里像是澄清与成熟的年月,走向对它那条单纯而清澈的消息的实现。梦把“我本来的样子”显示给我。同样清楚的是,在我醒着的生活里,我并不完全是梦所显示的那个人——来自久远的重负与僵硬(过去和现在都)常常妨碍我完全而单纯地做我自己。这些年里,尽管这个梦的念头只是偶尔回到我这里,它一定以某种方式起了作用。它绝不是某种范本或理想,让我努力去与之相像,而是对一种欢乐的单纯的悄声提醒——那单纯“就是我”,它以许多方式显现,并被召唤着去摆脱那些继续压着它的东西,去充分地绽放。这个梦是一条既纤细又有力的纽带,一头连着仍被来自过去的许多重负压着的现在,一头连着这现在已孕育在胚芽里的、近在眼前的“明天”——一个从此刻起就是我、而且肯定一直就在我之中的“明天”……

肯定,如果说最近这几周里这个难得提起的梦又一次鲜明地在场,那是因为在某个层面上——那不是探测与分析的思想的层面——我一定已经“知道”:我正在做、并正引向终点的这件工作,这件接续并深化了三年前那件工作的工作,是向着实现它带给我的、关于我自己的那条消息迈出的又一步。

如今在我看来,这就是《收获与播种》的主要意义所在,这件将近两个月的高强度工作的主要意义所在。直到它完成了,我才意识到我做这件事有多么重要。在这件工作的过程中,我经历过许多欢乐的时刻,那往往是一种促狭的、爱开玩笑的、兴高采烈的欢乐。也有过悲伤的时刻,有过一些时刻,我重新经历这些年里痛苦地击中过我的挫折与伤心——但没有过哪怕一刻的苦涩。我离开这件工作时,带着那种知道自己把一件工作做到了底的人的完全的满足。没有任何一件事,无论它多么“小”,是我在那里回避掉的,或者是我心里看重、本该说出而竟没有说出的,以至于此刻会在我心里留下一点不满、一点遗憾的残渣——无论那多么“小”。

在写这份见证时,我很清楚它不会讨所有人喜欢。甚至很可能,我找到了办法把所有人,一个不落,都得罪了。然而这根本不是我的本意,连得罪任何人也不是。我的本意不过是去看那些简单而重要的事,那些日常的事——我作为数学家的过去(有时也包括现在)里的事——好最终(晚知道总比不知道好!)毫无疑义、毫无保留地发现它们过去是什么、现在是什么;并且一路走一路用简单的话说出我所看见的。

\subsection{2. 一场旅行的精神}

这场反思最终成了《收获与播种》,它起初是作为《追寻栈》(À la Poursuite des Champs)第一卷(正在完成中)的“引言”开始的——《追寻栈》是我自1970年以来就打算发表的第一件数学工作。头几页是我去年六月在一个空闲的空当里写的;不到两个月前,我又从搁下的地方把这反思接了起来。我意识到有不少东西值得看、值得说,所以指望的是一篇比较充实的引言,三四十页。可是此后将近两个月里——直到此刻,我正在为这最初本是引言的文字写一篇新引言——我每天都以为今天就
% ---------------- p.112 ----------------
会把它做完,至不济明天或后天也就完了。等到几周之后我开始逼近一百页这个关口,这引言就被擢升为“导论性的一章”。又过了几周,当所谓“这一章”的规模远远超出正在准备的这卷里其他各章(写这几行时,除最后一章外都已写完),我终于明白了:它的位置不该在一本数学书里,这场反思与见证在那里实在太挤。它真正的位子是在单独的一卷里,那将是“数学反思”系列的第1卷——这个系列我打算今后几年顺着《追寻栈》的势头继续下去\fnote{2}{}{(5月28日)这里指的是《收获与播种》第一部分《愚妄与更新》的文本。写这几行时,第二部分尚未写出。}。

我不会说《收获与播种》——这部《数学反思》系列的第一卷(接下来先是《追寻栈》的两三卷)——是《反思》的“导论”卷。不如说,我把这第一卷看作将来之物的地基,或者更确切地说,看作定下基调的那一卷,给出我踏上这趟新旅程时所怀的那种精神的那一卷;这旅程我打算在今后几年里继续走下去,而它把我带向哪里,我说不上来。

关于本卷主干部分的这些说明,最后再交代几点实用性质的情况。读者在《收获与播种》的正文里偶见“本卷”的提法,不必奇怪——那暗指的是《追寻栈》的第一卷(《模型史》),我当时还以为自己仍在为它写引言。我不愿去“改正”这些地方,因为我最看重的是保住文本的自发性,保住它作为见证的真实性——不仅是对一个遥远过去的见证,也是对正在写作的当下这一刻的见证。

同样出于这个理由,我对文本初稿所作的修改,只限于纠正文笔上的笨拙,或某些有时含混、妨碍理解我要表达之意的说法。这些修改有时把我引向比写初稿时更清晰、更精细的领会。至于对初稿稍具实质性的改动——为给它添些层次、说得更精确、作些补充,或者(有时)加以更正——则放在大约五十条编号的“注”里,集中排在反思的末尾,它们占了全文四分之一还多。我用(1)之类的符号去指这些注……在这些注里,我挑出二十来条,觉得它们(以篇幅或实质而论)重要得可与这篇反思自发组织成的五十个“节”或“段”中的任何一节相比。这些较长的注也收进了目录,排在五十节的清单之后。不出所料,对某些长注来说,又需要给注再加一条或几条注。这些注便接在该注之后,用的是同一类指号;只有相当短的注例外,它们就排在同一页上作“脚注”,用的是诸如(46\textsubscript{1})、(46\textsubscript{2})那样的指号。

给正文的每一节、以及每一条最充实的注取一个名字,我从中得到很大的乐趣——何况后来这甚至成了我理清头绪的必需品。不用说,这些名字都是事后才找到的;当初开始写某一节或某条稍长的注时,我一条也说不出它的实质是什么。我用来称呼八个部分(I至VIII)的那些名字(如“劳动与发现”等等)就更是如此——这八部分是我事后把构成全文的五十节归并而成的。

至于这八个部分的内容,我只作极简短的评论。头两个部分——I(《劳动与发现》)与II(《梦与做梦者》)——含有一些关于数学工作、以及一般而论关于发现工作的反思要素。
% ---------------- p.113 ----------------
\origpage{VII}在其中,我本人牵涉得零散得多,也间接得多,远不如后面几个部分。后面这几个部分才尤其具有见证与冥想的品格。\origpage{VIII}III至VI部分,主要是关于我1948至1970年间“在数学世界里”那段数学家往昔的反思与见证。推动这场冥想的动机,首先是想理解这段过去,以此努力去理解并承担一个在某些方面有时令人失望或令人茫然的现在。VII(《孩子嬉戏》)与VIII(《孤独的历险》)两部分,更多涉及我从1970年至今、也就是从我离开“数学家的世界”不再回头以来,我与数学之关系的演变。那些动机、那些力量与境况,它们使我(连我自己都吃惊)在中断十三年多之后,重新开始一种“公开的”数学活动(即写作并发表《数学反思》)。

% ================================================================
%  第 4 章 引言(后半) — 原书 p.114-122, 边码 P.IX-P.XXII
%  源文: _work/src1/task-ch04b.md
%  承接 chapters/part1-ch04.tex(该文件止于 4.1.2 之末)
% ================================================================

% ---- [补译] 4.1.3 节余下部分(p.113; part1-ch04.tex 中缺此节正文) ----
\subsection{3. 罗盘与行囊}

我还该就与《收获与播种》一起构成本卷、并与它同名的另外两篇文字说几句话。

《纲领概要》勾画出我近十年来所追寻的数学思考的主要主题。我打算今后几年里,至少把其中若干稍加展开,写成一连串不拘形式的反思——我已有机会提起过它们,即“数学反思”。这份概要是去年一月我写的一份报告的逐字复制,当时是为了支持我申请国家科研中心(CNRS)的一个研究员职位。我把它收进本卷,是因为这个纲领显然远远超出我卑微个人之力所能及——哪怕上天让我再活一百年,并且我选择把这一百年用来尽可能远地推进这些主题。

《主题概要》写于1972年,是为另一次申请(法兰西学院的一个教授职位)。它按主题勾画了我当时视为自己主要数学贡献的东西。这篇文字带着它写作时心境的印记——那阵子我对数学的兴趣,说得客气些,也只能算是极边缘的。所以这份概要顶多不过是一份干巴巴、有条有理的罗列(幸好它并不求详尽……)。它看不出受到任何眼光的照耀,也看不出被任何欲望的气息吹动——仿佛我在其中像交差一样逐一过目的那些东西(而这确实就是我当时的心境),从来不曾被一种活的眼光触碰过,也从来不曾有过那种激情:要把它们带到光天化日之下,而它们当时还只是躲在雾与影的帷幕背后被隐约预感着……

然而,我之所以还是决定把这份恐怕没什么灵感的报告收在这里,主要是想堵住(假如这事办得到的话)某些高来高去的同行和某种风气的嘴:自从我离开那个我们曾共同所属的世界以来,他们总爱居高临下地看他们亲切地称之为“格罗滕迪克货色”(grothendieckeries)的东西。据说,这词的意思是死抠一些太琐碎的小事,以致一位严肃而有品味的数学家不肯为它们浪费自己当然宝贵的时间。也许这份“难以消化的文摘”在他们眼里会显得更“严肃”一些!至于我笔下那些由一种眼光、一种激情所驱动的文字,它们并不是写给那些人的——某种风气把他们维持、并且为他们开脱在一种自得自满之中,使他们对那些令我着迷的东西毫无感觉。如果说我写东西不只是为自己,那是为这样的人:他们不觉得自己的时间和自己本人珍贵到了不能无休无止地去追寻那些谁都不屑一顾的显而易见之物,不能为每一件被发现之物的内在之美而欢欣,并且把这美与我们所知的任何别的、各有其自身之美的东西区分开来。

% ---------------- p.114 ----------------
\origpage{IX} 如果我愿意把构成本卷的三篇文字彼此定位,并定出各自在这趟旅行中——我随《数学反思》一同登上了这趟旅行——所扮演的角色,我就可以说:反思—见证《收获与播种》反映并描述了我踏上这趟旅行时所怀的精神,正是这精神给了它意义。《纲领概要》描述了我的灵感之源,它们为这趟驶入未知的旅行定出一个方向——当然不是说定下一个目的地——有点像一只罗盘,或者一根结实的阿里阿德涅线。最后,《主题概要》匆匆过了一遍一份行囊,那是我1970年之前作为数学家的往昔里积攒下来的;其中至少有一部分,在旅行的这个或那个阶段会有用、会受欢迎(就像我的上同调代数与 topos 代数那套本能反应,如今在《追寻栈》里已是我不可或缺的东西)。而这三篇文字相随的次序,连同它们各自的篇幅,恰好(并非我刻意安排)反映出我在这一旅程中给予它们的分量与重视——这趟旅程的第一程,已近尾声。

\subsection{4. 一场追寻显而易见之物的旅行}

我还该就这趟一年多一点以前启程的旅行——即“数学反思”——再说几句更周详的话。关于我怀着怎样的精神踏上这趟旅行,我在《收获与播种》开头八节(即这场反思的第I、II部分)里已说得相当详细;我想,这种精神从现在起就在本卷里,也在接续它的那一卷(即《模型史》,它是《追寻栈》的第1卷)里显露出来了——后者正在完成中。因此,我觉得在这篇引言里再就此发挥,已是多余。

我当然无法预言这趟已经启程的旅行会是什么样子;这事我会随着它一路走下去而一点点发现。眼下我并没有一个哪怕只及大要的预定路线,而且我怀疑近期内会冒出这么一条来。正如我前面说过的,那些大概会启发我思考的主要主题,都多少勾画在《纲领概要》里了——那份“罗盘文本”。在这些主题当中,还有《追寻栈》的那个主要主题,即“栈”;我希望就在今年之内把它(并且就此打住)走完一遍,写成两卷,也许三卷。关于这个主题,我在《纲领概要》里写道:“……这有点像一笔债,我要向一段科学的过去偿还:在那十五年里(1955至1970年间),上同调工具的发展始终是我在代数几何基础方面工作中不变的主旋律”。因此,在预定要写的主题当中,这一个最深地扎根于我的科学“过去”。也正是这一个,在逝去的这十五年里始终作为一种遗憾在场——它也许是我离开数学舞台时留下的未竟工作中最显眼的一处空缺,而我旧日的学生或朋友里,没有一个人费心去填补它。关于这项进行中的工作,感兴趣的读者可以去查阅《纲领概要》中相应的那一节,或者去查阅《追寻栈》第一卷(正在完成中)的引言(这一回是真正的引言!)

作为我的科学过去留给我的另一份遗产,尤其还有那个我格外放在心上的 Motif 观念:它至今仍等着走出那一直被囚于其中的黑夜——尽管它现身已有足足十五年了。如果没有任何比我位置更好的人(更年轻,而且拥有他所掌握的工具与知识)决定在最近这几年里动手,那么最终由我自己来干这里所必需的基础工作,也并非不可能。

我趁此机会说明:Motif 观念的命运(或者不如说,厄运……),以及我发掘出来的另外几个观念——它们在我看来(潜在地)是其中最多产的——的命运,构成了一篇将近二十页的回顾性反思,是《收获与播种》诸“注”中
%
% ---------------- p.115 ----------------
\origpage{X} 最长的一篇(也是最后几篇之一)\fnote{3}{}{这条双注(第46、47号)及其子注已收入《收获与播种》第二部分“葬礼”,后者是它的直接续篇。}。事后我把这条注分成两部分(《我的孤儿们》与《拒绝一份遗产——或一个矛盾的代价》),此外还有紧随其后的三条“子注”\fnote{4}{}{这里指的是第48、49、50号子注(第48$'$号注是后来补加的)。}。这连续五条注合起来,是《收获与播种》里唯一不作泛泛一提、而是真正谈论数学概念的部分。这些概念在这里成了机会,用来例示数学家世界内部的某些矛盾——而这些矛盾本身又映照出人本身之中的矛盾。有一个时候我曾想把这株触手蔓生的注从它所属的文本里分出去,归入《主题概要》。那样做的好处,是能给后者一个透视,并给一篇实在太像目录的文字注入一点生气。然而我还是没有这样做,为的是保住一份见证的真实性——而这条巨型注,不管我乐意不乐意,确实就是这份见证的一部分。

\origpage{XI} 关于《收获与播种》里所说我以何种心境着手“反思”,我想在这里只补一句话——对此我已在其中一条注(《年轻人的势利——或纯洁的捍卫者们》)里表过态,那里我写道:“我毕生作为数学家的抱负,或者不如说我的欢乐与激情,始终是去发现那些显而易见之物;在眼下这部著作里,这也依然是我唯一的抱负”(《追寻栈》)。对于我一年来随“反思”进行的这趟新旅行,这也同样是我唯一的抱负。开启这趟旅行的这部《收获与播种》也不例外(至少对我的读者而言是如此,如果真有读者的话)。

\subsection{5. 一笔情愿的债}

我想以几句关于本卷《收获与播种》两则题献的话,来结束这篇引言。

那则题献——“献给曾是我的学生们的人,我曾把自己最好的给了他们——也给了最坏的”——至少去年夏天就已在我心里,尤其是在我写开头四节的时候——那四节当时还该是一本数学书的引言。这说明我很清楚(其实已经清楚好几年了):有一个“最坏的”东西需要审视——而现在正是时候,再没有比现在更好的时候了!(不过我万万没想到,这个“最坏的”最终会领着我走完一场将近两百页的冥想。)

相反,“献给曾是我的长辈们的人”这则题献,是半路上才出现的,就像这场反思的名字本身(它后来也成了一卷的名字)一样。这则题献向我揭示出他们在我作为数学家的生命中所扮演的重要角色——这角色的效力至今仍然活着。这一点,在后面的篇幅里大概会相当清楚地显现出来——所以在这里不必就此多说。这些“长辈”,按他们在我二十岁那年(大致)进入我生命的次序,是亨利·嘉当、克劳德·谢瓦莱、安德烈·韦伊、让-皮埃尔·塞尔、洛朗·施瓦茨、让·迪厄多内、罗歇·戈德芒、让·德尔萨特。当时那个初来乍到的我,受到了他们每一个人的善意接待;此后他们中许多人给了我持久的友谊与情谊。我还须在此提到让·勒雷:在我与“数学家世界”的第一次接触时(1948/49年),他的善意接待同样是一种珍贵的鼓励。我的反思显现出一笔感恩之债,欠他们每一个人——这些“来自另一个世界、另有另一份命运”的人。这笔债绝不是负担。发现它时,它像一份欢乐来临,让我变得更轻。

\noindent 1984年3月底\par

% ---------------- p.116 ----------------
\section{(II) 一桩敬重之举}

(5月4日—……6月)

\subsection{6. 葬礼}

\origpage{XII} 一件意料之外的事,把一场本已走到头的反思又重新发动起来。它开启了一连串大大小小的发现,在过去这几周里逐步揭开了一个一直模糊不清的局面,并使它的轮廓愈发分明。这尤其使我周详而深入地走进了一些事件与处境——此前它们只是被顺带一提,或隐约暗示过。于是,前面提到过的那篇关于一部作品之沉浮的“十五页回顾性反思”(引言,第4节),就获得了意想不到的规模,多出了大约两百页。

由于事势所迫,也由于一场反思的内在逻辑,我在路上被引到把别人牵涉进来,其程度不亚于牵涉我自己。被牵涉得比谁都深的(除了我自己),是一个与我有着将近二十年友谊的人。关于他,我曾写道(那是一种委婉说法\fnote{5}{}{关于这个“委婉说法”的含义,见注《与众不同者》,第57$'$号。}):在这份扎根于共同激情的亲切友谊的最初几年里,他“多少有点像个学生”;而在很长一段时间里,在我内心深处,我把他看作某种“合法继承人”——所继承的,是我相信自己还能在数学上带来的东西,是在那部始终零碎的已发表作品之外的东西。很多人早已认出他是谁了:他就是皮埃尔·德利涅。

我并不为随这些注一起公开一份对一段私人关系的个人反思、并且就这样未经征求他同意而把他牵涉进来而道歉。一个长期隐而不显、混乱不清的局面,终于被摊到光天化日之下加以审视,在我看来是重要的,对所有人都是健康的。这样做,我是在提供一份见证——无疑是主观的,既不自称穷尽了一个微妙而复杂的局面,也不自称没有错误。它的首要长处(正如我过去那些出版物、或眼下正在写的那些的首要长处一样)是它存在着,供可能对它感兴趣的人取用。我所操心的,既不是要说服人,也不是要躲在那些所谓“确凿”的东西背后,以免于错误或怀疑。我所操心的是要真实——把我在每一刻所看见、所感到的东西就那样说出来——以此作为深入它们、理解它们的手段。

“葬礼”这个名字——用于所有那些与“一段过去的重量”相关的注的总称——在反思过程中以一种越来越强的力量自显出来\fnote{6}{}{在这场反思接近尾声时,另一个名字浮现出来,表达了某幅图景的另一个触目惊心的方面——那幅图景在过去五周里逐步在我眼前展开。那是一个童话的名字,我将在相应之处回到它:《中国皇帝的袍子》……}。在其中,我扮演那位预先的死者的角色,与几位(年轻得多的)数学界同行组成凄凉的送葬队伍;他们的作品落在我1970年“离去”之后,并且凭着某种风格、某种对数学的进路,带着我影响的印记。这些人里排在头一个的,是我的朋友左格曼·梅布胡特;他有幸——这运气沉重得很——要承受一个被当作“1970年之后的格罗滕迪克的学生”的人所遭遇的一切不利,却又没享受到与我接触、得到我鼓励与指点的那种好处,因为他只是通过我的文字、从我的作品那里“受教”的。那阵子(在他所出没的那个世界里),我已经俨然一副“死者”的模样,以致长期以来,连见上一面的念头显然都不曾出现过;一段持续的往来(无论私人的还是数学的),直到去年才终于结上。

% ---------------- p.117 ----------------
\origpage{XIII} 这并没有妨碍梅布胡特逆着一种专横的风气、逆着他那些(曾是我的学生的)长辈们的轻蔑,在近乎完全的孤立中,做出独创而深刻的工作——靠的是一种意想不到的综合:把佐藤学派的想法与我的想法合到一起。这项工作提供了一种新的入手之处,可以把握解析簇与代数簇的上同调,并预示着我们对这种上同调的理解将有一场大规模的更新。毫无疑问,假如梅布胡特在那些最适合做这件事的人那里,得到了他们当年在我这里得到过的热情接待与毫无保留的支持,那么这场更新早已是既成之事,而且已经完成多年了。至少,自1980年10月起,他的想法与工作为代数簇上同调理论的一场引人注目的重启提供了灵感与技术手段,使这门理论终于(除了德利涅围绕韦伊猜想的那些结果之外)走出了一段漫长的停滞期。

\origpage{XIV} 难以置信却又千真万确:近四年来,他的想法与结果被“所有人”使用着(与我的想法同等待遇),而他的名字仍被小心地无视、被默不作声地略过——恰恰是被那些第一手了解他的工作、并在自己的著作里根本性地使用它的人。我不知道数学界在别的什么时候可曾见过这样的失格:其中某些最有影响、最有声望的成员,在普遍的漠然中,带头蔑视数学家这行伦理中最被普遍承认的规矩。

我看见四个男人,才能出众的数学家,他们和我一起享有——也有资格享有——这场以沉默与轻蔑举行的葬礼的哀荣。而我在每个人身上都看见:轻蔑如何啃咬着那曾激励过他的美好激情。

除了这几个人,我尤其看见两个人:他们一个接一个站在数学公共广场的聚光灯下,在众人的簇拥中主持葬礼,而同时(在一种更隐蔽的意义上)他们也被埋葬——被自己的双手埋葬,与他们刻意埋葬的那些人一同被埋葬。其中一个我已经点过名了。另一个同样是从前的学生、从前的朋友:让-路易·维尔迪耶。1970年我“离去”之后,他与我之间的接触没有维持下来,除了职业层面上几次匆促的会面。所以大概正因如此,他在这次反思里只通过他职业生涯中的某些举动出现;至于这些举动可能的动机——就他与我的关系这一层面而言——则未经审视,而且我对此也完全摸不着头脑。

如果说有一个追问在逝去的这些年里始终逼着我,它曾是《收获与播种》的一个深刻动机,也在这整场反思中一路跟着我,那就是:在某种精神、某种风气的到来之中,我占了怎样一份——正是这种精神与风气,使得像我说的那样的失格成为可能,而这一切发生在一个曾属于我的世界里,在我作为数学家的二十多年里,我曾与它认同为一。这场反思让我发现:由于我身上某些虚浮的态度——它们表现为对才具平平的同事一种无言的轻蔑,表现为对我自己、对某些才具出众的数学家的自我迁就——我对这种精神并不陌生;而今天,我看见这精神在我曾爱过的那些人当中肆意铺展,也在我曾把自己热爱的一门手艺教给的那些人当中铺展;那些我没有好好爱、也没有好好教的人,如今在这个曾为我所珍爱、而我已经离开的世界里定调子(当他们不是在立法的时候)。

我感到一阵风刮起来,那是自得、玩世不恭与轻蔑的风。“它刮着,不管什么‘功绩’什么‘过失’,用它的气息烧灼卑微的志向,也烧灼最美的激情……”。我明白了:这阵风,是我参与播下的盲目而不经心的种子那丰饶的收成。而如果它的气息回吹到我身上,回吹到我曾交托给别人手里的东西上,回吹到我今天所爱的人身上,而且
%
% ---------------- p.118 ----------------
\origpage{XV} 回吹到那些敢于自称属于我、或仅仅是受我启发的人身上,那么这是事物的回还,我没有理由抱怨,而且它有许多东西要教我。

\subsection{7. 葬仪的序次}

于是,在“葬礼”这个名目下,我在目录里把那些与那个表面平淡无奇的节“一段过去的重量”(第50节)相关的主要“注”的浩浩荡荡的队列归拢在一起,从而让这个名字充分显出了它的意义——当初写《收获与播种》“初稿”的这最后一节时,这名字一上来就落到了我头上。

\origpage{XVI} 在这支亲缘关系错综的长长的注的队列里,过去四周里才加入的那些(注(51)至注(97))\fnote{8}{}{还须加上1984年5月12日的第104号注。第98号及其后的各注(上一句提到的第104号注除外)构成反思的“第三口气”,自1984年9月22日起;它们也标了日期。}显得与众不同:它们是唯一标了日期的(4月19日至5月24日)\fnote{9}{}{在同一天写的一连串连续的注里,只有第一条标日期。其余未标日期的注是第44$'$号至第50号(构成第I、II、III支队伍)。第46、47、50号注写于3月30日或31日,第44$'$、48、48$'$、49号注写于四月上半月,最后,第44$''$号注是标了日期的(5月10日)。}。我觉得最自然的做法,是按它们在反思中相继出现的时序\fnote{10}{}{有时我在这个时间顺序里作过幅度不大的调换,换成一个“所谓逻辑的”次序,只要我觉得反思进程的整体印象不会因此失真。不过,作为仅有的例外,我要指出十一条注(其编号前有一个感叹号!)——它们源自某条注之后的页底脚注,篇幅涨到了让人却步的规模,于是我把它们各自放在与之相关的那条注之后(除了第98号注,它关联的是第47号)。}把它们排出来,而不是按某种所谓的“逻辑”次序,也不是按先前各注中提到它们的先后次序。为了能找回后一种(绝非线性的)承继关系——即各条参与其中的注之间的谱系——我在目录里让每条的编号后面跟着(在它之前的那些注中)第一次提到它的那一条的编号\fnote{11}{}{当对某条注的引用(如(45))就出现在“一段过去的重量”这一节本身之中时,则把这一节的编号(50)加括号,置于该注编号之后,如“46(50)”。},若没有这种情况,则跟着它所直接续写的那一条的编号\fnote{12}{}{在目录里,若某条注是前一条注的直接续写(此时两号相连),其编号前加一个 * 号。例如“*47, 46”表示第47号注是第46号注的直接续写(而这里第46号并非紧邻在它前面的那一条,紧邻的是第46\textsubscript{9}号注)。\par 最后,我在目录里给那些后面不跟别的编号的注加了着重号,也就是说,给那些代表反思“新起点”的注——它们并不嵌在已经做过的反思的某个确定位置上。}。(后一种关系在正文里另有标示:在前一条注的末尾放一个指号,如注(46)末行末尾的($\Rightarrow$ 47),指向续写它的注(47)。)最后,某一注中某些稍具技术性的细节,则集中在该注末尾,编成子注,用接在该注原编号之后的连续下标标示——如注(46)《我的孤儿们》的子注(46\textsubscript{1})至(46\textsubscript{9})。

为了多少给“葬礼”的整体安排立起结构,也为了让人在蜂拥而至的众多注里能认得出路,我觉得在此场合宜于在队列中插入几个庄重而富于暗示的小标题,每个小标题都走在并引领一支或长或短的送葬队——由主题相同、彼此相连的连续若干注组成。

于是我有幸看见,在一支前来为我送葬的庄严长队里,十支\fnote{13}{}{(9月29日)其实最终有十二支队伍,若把灵车(X)和“死者(仍未过世)”(XI)算进去的话——后者在最后关头又溜进了队列……}队伍一支接一支地聚拢起来——有的卑微,有的气派,有的悔恨,有的暗暗欢腾,在这样场合下本也免不了如此。于是走上前来的有:身后的学生(人人都把无视他当作本分)、孤儿们(为这场合才刚被掘出)、时风及其名流(这是我活该)、诸Motif(我所有孤儿中最晚出生、也最晚被掘出的),
%
% ---------------- p.119 ----------------
\origpage{XVII} ……我的朋友皮埃尔谦逊地领着最要紧的那一支队伍,紧随其后的是“一致同意”——那些(默默地)协同一致的注——以及全体出动的(所谓)“乖戾”研讨会(它与身后的学生、即“无名学生”划清界限的方式,是夹在中间的几支捧着鲜花与花环的送葬队);最后,为了给这气派的队列一个体面的收场,走上前来的还有“学生”(绝非身后的,更绝非无名的),又名“老板”,身后跟着我那一帮忙忙叨叨的学生(带着大量铁锹与绳索),最后是灵车(车上赫然四口漂亮的橡木棺材,用螺丝牢牢拧紧,还不算那个掘墓人)……十支队伍终于齐了(总算齐了),缓缓走向葬礼仪式。

\origpage{XVIII} 仪式的高潮是悼词,由我的朋友皮埃尔亲自出马,以完美的分寸献上——他应众人之请主持葬礼,人人满意。仪式在一曲终局而决定性的“自深处”(De Profundis)歌声中结束(至少人们这样希望)——这歌由可敬的死者本人唱出,作为真诚的感恩;他瞒着所有人,从自己那场声势骇人的葬礼上活了下来,并且按说从中得了教训,心满意足——这份满足,便构成了那场值得纪念的葬礼的最后一个音符、最后一个和弦。

\subsection{8. 一个秘密的终结}

在这场(但愿是)最后的反思阶段中,我看出有必要在《数学反思》第1卷的“附录”里,除前面谈到的三篇之外,再附上另外两篇性质属于数学的文字\fnote{14}{}{此外,我想给《主题概要》(见引言第3节“罗盘与行囊”)加一份“评注”,就我对其中粗略过目的那些“主题”的贡献作些说明,也就那些在我的数学作品中主要思想力的生成中起过作用的影响作些说明。最近六周的回顾已经显出(令我自己也吃惊):在大多数这些思想的启动上,以及在我1955至1970年间给自己定下的某些“大任务”上,塞尔都起了“雷管”的作用。\par 最后,还有一篇性质属于(通常意义上的)数学的文字,也是唯一(偶然)出现在非技术性文本《收获与播种》里的——我要指出注《屠杀》(第87号)下的子注第87\textsubscript{1}号:我在那里以它应有的仔细,阐明凝聚层语境中为人熟悉的黎曼–罗赫–格罗滕迪克定理的一个“谨慎的”(猜想性)变体。这个猜想(与另外若干一起)出现在1965/66年度 SGA~5 讨论班的收尾报告里;那份报告(如同许多别的报告一样)在十一年后以 SGA~5 之名出版的文集中没有留下任何痕迹。这个关键讨论班落在我某些学生手里之后的种种遭际,以及这些遭际与某种“SGA~4½ 行动”的关联,在注第63$''$、67、67$'$、68、68$'$、84、85、85$'$、86、87、88号所进行的反思过程中逐步显露出来。\par 又及:另一条给出相当充实的数学评注的注,讨论的是为那些已知的、拥有所谓“六种运算”对偶形式的情形(在可能的范围内)抽出一个共同“topos 式”框架是否合宜——我还要指出注《赊账的论文与全险》(第81号)下的子注第81\textsubscript{2}号。}。

第一篇是一份带评注、分两部分的报告的复制,是我1968与1969年就 P.\ 德利涅的工作(其中有些至今仍未发表)写的,对应于1965/67/68三年间在法国高等科学研究所(IHÉS)的数学活动。

另一篇是一份“六种变率公式表”的草稿,汇集了我在1956至1963年间抽出的某种对偶形式(灵感来自庞加莱对偶与塞尔对偶)的共同特征;这份公式表后来表明,对于迄今遇到的一切上同调对偶情形,它都具有“普适”的品格。随着我离开数学舞台,这套形式似乎就废弃不用了,以致据我所知,还没有人(除我之外)肯费神哪怕只写出那些基本运算的清单、它们所给出的基本典范同构的清单,以及这些同构之间那些本质的相容关系的清单。

这份融贯的公式表草稿,对我将是显然的第一步,通向那幅“Motif 之梦的宏伟全景图”——十五年来,它一直“等着那位敢于挥笔把它画出来的数学家”。

% ---------------- p.120 ----------------
\origpage{XIX} 看样子,这位数学家不是别人,正是我自己。确实,该是时候了:将近二十年前诞生、在私密中交托出去的东西——它本不是为了成为一个人的特权,而是为了供所有人取用——终于要走出秘密的黑夜,再一次诞生于白昼的充足光线下。

确实,除我之外,只有一个人对这套“Motif 的瑜伽”有贴身的了解——在我离去之前的那些日子里、那些年里,他一天天、一年年从我口中把它学了去。在我有幸发现并带到光天化日之下的一切数学事物当中,这个 Motif 的实在在我眼里仍是最迷人的、最充满奥秘的——它就在“几何”与“算术”之间那深层同一性的核心处。而这套“Motif 的瑜伽”——那长久被忽视的实在把我引向的瑜伽——也许是我在作为数学家的生命头一段里抽出的最有力的发现工具。

但同样真实的是:这个实在,以及这套努力把它圈得尽可能近的“瑜伽”,从来不曾被我守为秘密。我被撰写基础这一必须完成的任务占住了(此后人人都很乐意把它们原封不动地用在自己日常的工作里),没有拿出那必不可少的几个月,为这套 Motif 的瑜伽写一份宏阔的整体草稿,从而把它交到所有人手里。然而在我出乎意料地离去之前的那些年里,我并没有放过机会——每逢相遇,就向愿意听的人讲它,首先是我的学生们;而他们(其中一人除外)都把它忘了,就像所有人都把它忘了一样。我讲它,不是为了推销什么会挂上我名字的“发明”,而是为了引人注意一个每走一步都显现出来的实在——只要你关心代数簇的上同调,尤其是关心它们那些“算术”性质,以及迄今已知的各种上同调理论彼此之间的关系。这个实在,与当年“无穷小”的实在一样可触——那东西被感知到,远早于严格语言的出现,而正是那种语言使人能完满地把握它并把它“确立”下来。而为了把握 Motif 的实在,我们今天绝不缺一套灵活而合用的语言,也不缺在数学理论建构方面成熟老练的经验——这些,我们的前人是缺的。

如果我从前在屋顶上大声喊过的东西落进了聋耳朵,如果某一个人的沉默不语所收到的回声,是所有那些装作对上同调感兴趣的人(而他们明明和我一样有眼睛有手……)的沉默与麻木,那么我也不能把责任只归到那一个选择把“好处”留给自己的人身上——那本是我为所有人而交托给他的。我们不得不承认:我们这个时代,其科学产出的狂飙足以与投入军备或消费品的劲头相比,却离我们十七世纪前人那种“大胆的活力”相去甚远——他们为了发展无穷小的计算,“不走那么多弯路”,不会因为这东西是否“猜想性”的顾虑而停住,也不等他们中某位显赫人物肯来给他们开绿灯,才去抓住每个人自己亲眼看见、第一手感到的东西。

\subsection{9. 场景与角色}

凭它自身的内部结构,也凭它那个特殊的主题,“葬礼”(如今已占《收获与播种》全文一半以上)在很大程度上、并且从逻辑观点看,独立于它前面那场漫长的反思。然而这只是一种纯粹表面的独立。对我来说,这场围绕着一场“葬礼”的反思——那葬礼渐渐从“未言”与“预感”的雾中走出来——与在它之前、它由之而生并且由之获得全部意义的那场反思,是不可分的。它起头时只是“顺路”匆匆一瞥一部作品的沉浮——那部作品我有点(很大程度上)已经看不见
% ---------------- p.121 ----------------
\origpage{XX} 了;后来,既未预谋也未寻求,它变成了一场对一段重要关系的冥想,而这又把我引向对那部作品在“曾是我的学生们”手中的命运的反思。把这场反思从它自发从中生出的那场反思里割出去,在我看来就是把它降格为一幅简单的“风俗画”(甚至是在数学“上流社会”里算的一笔账)。

\origpage{XXI} 确实,若有人偏要如此,整个《收获与播种》也能作同样的降格,变成一幅“风俗画”。当然,某个时代、某个圈子里通行的风俗,那些参与塑造该圈中人生活的风俗,自有其重要性,也值得描写。然而《收获与播种》的细心读者会明白:我的打算不是描写风俗,也就是说,不是描写某一个场景——它随时代而变、因地而变,我们的行动就在它上面展开。这场景在很大程度上界定并划出了我们身上各种力量可用的手段,使它们得以表达。场景以及它提供的手段(还有它强加于人的“游戏规则”)千变万化,而我们身上那些深层力量的本性——它们在集体层面上塑造场景,在个人层面上又在场景上表达自己——看来从一个圈子、一种文化到另一个圈子、另一种文化,从一个时代到另一个时代,都是同一个。如果说我这一生里,在数学之外、在爱女人之外,还有什么东西让我感到奥秘与吸引(不错,是晚近的事),那正是这少数几种力量的隐秘本性——它们有能力驱使我们行动,为“最好”也为“最坏”,为埋藏也为创造。

\subsection{10. 一桩敬重之举}

这场最终取名为“葬礼”的反思,起初是一桩敬重之举。敬重那些我发现了的东西——我看着它们在一片虚无中凝聚、成形,我是第一个尝到它们的滋味与劲道的人,我给它们起了名字,既为表达我对它们的认识,也为表达我的敬重。对这些东西,我给出了自己身上最好的部分。它们汲取着我身上的力量,生长、绽放,像许多粗壮的枝条从同一根活着的树干上迸发出来,那树干有着粗壮而众多的根。这些是活着的、在场的东西,不是可有可无的发明——它们紧密相依,构成一个活生生的统一体,这统一体由它们每一个构成,又给每一个以位置与意义、一个起源与一个终点。我很久以前就把它们留下了,没有丝毫不安,也没有丝毫遗憾,因为我知道我留下的东西是健全而强壮的,它根本不需要我,就能依着它自己的本性继续生长、绽放、繁衍。我留下的不是一袋金币——那是可以被人偷走的;也不是一堆工具——那是会生锈、会腐烂的。

然而,年复一年,当我自以为离那个我离开的世界很远的时候,却有一阵阵阴险的轻蔑、隐隐的嘲弄,一直飘到我隐居之处,在这里那里回旋,指着这些东西中的某些——我明知它们强壮而美,有它们的位置,有它们独一无二、任何别的东西都永远无法代替的功能。我感到它们像是敌对世界里的孤儿——那个世界患着轻蔑的病,专门纠缠没有铠甲护身的东西。这场反思就是在这样的心境中开始的,作为对这些东西的一桩敬重之举,并因此也是对我自己的一桩敬重之举——作为对它们与我之间那条深沉纽带的提醒:谁若乐于对这样一件由我的爱滋养出来的东西摆出轻蔑的样子,他所乐于轻蔑的正是我,以及一切从我这里生出的东西。

对于另一种人,情形也一样:他第一手知道那条把我和某件东西连起来的纽带——那东西他不是从我这里、又是从哪里学来的呢——却摆出一副把这条纽带视为无足轻重、或者干脆无视它的样子,或者(哪怕只是默默地、以省略的方式)为自己、或为别人僭取一种虚假的“父权”。我在其中
% ---------------- p.122 ----------------
\origpage{XXII} 清清楚楚看到一桩轻蔑之举:轻蔑那件从工匠手里生出的东西,轻蔑那默默而精细、使这件东西得以生出并由此诞生的工作,轻蔑那工匠,而首先(以一种更隐蔽也更本质的方式)轻蔑他自己。

如果我的“回归数学”只消让我重新想起这条纽带,只消在我心里唤起这桩当着所有人——当着那些摆出轻蔑样子的人,也当着无动于衷的见证者——的敬重之举,那么这次回归也就不算白费。

确实,我与我留下的那部写出来的和没写出来的(至少是未发表的)作品,真的失去了接触。开始这场反思时——我看那些枝条还相当分明,却不怎么记得它们原是同一棵树的一部分。说来奇怪:非得那幅我留下之物遭劫掠的图景在我眼前一点点展开,我才在自己心里重新找回那种意义——那种关于被如此劫掠、如此拆散之物的活生生的统一性的意义。一个人拿走了一些金币,另一个拿走了一两件工具,用来炫耀,甚至用来使用——但是那赋予我所留之物以生命和真正力量的那种统一性,却从每一个人、也从所有人手里漏掉了。不过我清楚知道有一个人,深深感到过这种统一性和这种力量,在他内心深处今天仍然感到它;而他却乐于耗费自己身上的力量,要去摧毁他在别人身上、通过别人的作品感到过的那种统一性。作品的美与创造之德,恰恰就在这种活生生的统一性之中。尽管遭了劫掠,我仍把它们找回来,完好如初,仿佛我刚离开它们——只是我成熟了,今天用一双新的眼睛看它们。

然而,如果说有什么东西真被劫掠、被损毁、被卸掉了它原本的力量,那是在那些人身上:他们忘了自己身上安歇着的力量,自以为在劫掠一件任由他们摆布的东西,其实他们只是把自己与那创造之德割断——那东西就在他们手边,一如它在所有人手边,却绝不任他们摆布,也不在任何人的权力之下。

这样,这场反思,以及通过它而来的那次不期而至的“回归”,也让我重新接触到一种被遗忘的美。正是完完全全感受到了这种美,才使那桩敬重之举有了全部意义——它在注《我的孤儿们》\fnote{15}{}{这条注(第46号)在收入“葬礼”的所有注中,按时间是最早的一条。}里表达得笨拙,而我刚才在这里,在完全知情的情况下,又把它重述了一遍。


\end{document}
